% =========================================================================
%  DERECHO CIVIL: INTRODUCCIÓN Y PARTE GENERAL
%  Documento único e integrado de estudio (apuntes_integrados.tex)
%  Compilar con: pdflatex apuntes_integrados.tex (dos o tres pasadas)
% =========================================================================
\documentclass[12pt,oneside]{book}

% -------------------------------------------------------------------------
% METADATOS
% -------------------------------------------------------------------------
\newcommand{\universidad}{Universidad de Buenos Aires (UBA)}
\newcommand{\facultad}{Facultad de Derecho}
\newcommand{\materia}{Derecho Civil}
\newcommand{\titulo}{Derecho Civil: Introducción y Parte General}
\newcommand{\subtitulo}{Apuntes integrados de estudio}
\newcommand{\autor}{Isaac Edgar Camacho Ocampo}
\newcommand{\anio}{2026}

% -------------------------------------------------------------------------
% PAQUETES
% -------------------------------------------------------------------------
\usepackage[utf8]{inputenc}
\usepackage[T1]{fontenc}
\usepackage[spanish,es-nodecimaldot]{babel}
\usepackage[
    top=2.5cm,
    bottom=2.5cm,
    left=3cm,
    right=2.5cm,
    headheight=15pt
]{geometry}
\usepackage{amsmath}
\usepackage{amssymb}
\usepackage{mathtools}
\usepackage{graphicx}
\usepackage{float}
\usepackage{caption}
\usepackage{subcaption}
\usepackage{booktabs}
\usepackage{array}
\usepackage{longtable}
\usepackage{tabularx}
\usepackage{enumitem}
\usepackage[table]{xcolor}
\usepackage[most]{tcolorbox}
\usepackage{fancyhdr}
\usepackage{ragged2e}
\usepackage[
    colorlinks=true,
    linkcolor=blue!50!black,
    citecolor=green!40!black,
    urlcolor=blue!60!black,
    pdftitle={Derecho Civil: Introducción y Parte General},
    pdfauthor={Isaac Edgar Camacho Ocampo},
    pdfsubject={Derecho Civil --- Introducción y Parte General},
    bookmarks=true
]{hyperref}

% -------------------------------------------------------------------------
% CONFIGURACIÓN GENERAL
% -------------------------------------------------------------------------
\graphicspath{{./img/}{./graficos/}{./figuras/}}
\setcounter{tocdepth}{2}
\setcounter{secnumdepth}{2}

\setlength{\parindent}{0pt}
\setlength{\parskip}{0.5em}
\setlength{\emergencystretch}{2em}
\setlength{\LTpre}{6pt}
\setlength{\LTpost}{6pt}

\setlist[itemize]{topsep=0.3em, itemsep=0.2em}
\setlist[enumerate]{topsep=0.3em, itemsep=0.2em}

\clubpenalty=10000
\widowpenalty=10000

% -------------------------------------------------------------------------
% TÍTULOS DE PARTE Y CAPÍTULO (tamaños moderados) Y TEXTO DE PARTE
% -------------------------------------------------------------------------
\makeatletter
\renewcommand\@makechapterhead[1]{%
  \vspace*{0.5cm}%
  {\parindent\z@ \raggedright \normalfont
   \ifnum \c@secnumdepth >\m@ne
     \large\bfseries\color{black!65} \@chapapp\space\thechapter \par\nobreak\vskip 4\p@
   \fi
   \interlinepenalty\@M \Large\bfseries #1\par\nobreak
   \vskip 6\p@ \hrule height 0.6pt \vskip 22\p@}}
\renewcommand\@makeschapterhead[1]{%
  \vspace*{0.5cm}%
  {\parindent\z@ \raggedright \normalfont
   \interlinepenalty\@M \Large\bfseries #1\par\nobreak
   \vskip 6\p@ \hrule height 0.6pt \vskip 22\p@}}
\newcommand{\partintro}[1]{\gdef\@partintro@text{#1}}
\gdef\@partintro@text{}
\renewcommand\part{%
  \clearpage \thispagestyle{plain}%
  \null\vfil
  \secdef\@part\@spart}
\def\@part[#1]#2{%
  \ifnum \c@secnumdepth >-2\relax
    \refstepcounter{part}%
    \addcontentsline{toc}{part}{\thepart\hspace{1em}#1}%
  \else
    \addcontentsline{toc}{part}{#1}%
  \fi
  \markboth{}{}%
  {\centering \interlinepenalty \@M \normalfont
   \ifnum \c@secnumdepth >-2\relax
     \large\bfseries \partname\nobreakspace\thepart \par\vskip 12\p@
   \fi
   \LARGE\bfseries #2\par}%
  \@endpart}
\def\@endpart{%
  \vskip 2em
  {\centering\begin{minipage}{0.86\textwidth}\normalfont\normalsize\RaggedRight\@partintro@text\end{minipage}\par}%
  \gdef\@partintro@text{}%
  \vfil\newpage}
\makeatother

% -------------------------------------------------------------------------
% ENCABEZADOS Y PIES
% -------------------------------------------------------------------------
\pagestyle{fancy}
\fancyhf{}
\fancyhead[L]{\small\nouppercase{\leftmark}}
\fancyhead[R]{\small\materia\ --- Parte General}
\fancyfoot[C]{\thepage}
\renewcommand{\headrulewidth}{0.4pt}
\renewcommand{\footrulewidth}{0pt}
\fancypagestyle{plain}{%
  \fancyhf{}
  \fancyfoot[C]{\thepage}
  \renewcommand{\headrulewidth}{0pt}
}

% -------------------------------------------------------------------------
% COMANDOS MATEMÁTICOS
% -------------------------------------------------------------------------
\newcommand{\abs}[1]{\left|#1\right|}
\newcommand{\norm}[1]{\left\lVert#1\right\rVert}
\newcommand{\vect}[1]{\mathbf{#1}}
\newcommand{\deriv}[2]{\frac{d#1}{d#2}}
\newcommand{\pderiv}[2]{\frac{\partial#1}{\partial#2}}

% -------------------------------------------------------------------------
% TIPOS DE COLUMNA Y ENTORNO PARA TABLAS
% -------------------------------------------------------------------------
\newcolumntype{Y}{>{\RaggedRight\arraybackslash}X}
\newcolumntype{B}[1]{>{\bfseries\RaggedRight\arraybackslash}p{#1}}
\newcolumntype{P}[1]{>{\RaggedRight\arraybackslash}p{#1}}
\newenvironment{tabla}{%
    \par\medskip\begingroup\small\renewcommand{\arraystretch}{1.3}\noindent\ignorespaces
}{%
    \par\endgroup\medskip
}

% -------------------------------------------------------------------------
% CUADRO CORNELL (columna izquierda estrecha / columna derecha amplia)
% -------------------------------------------------------------------------
\newcommand{\cornellhead}{%
\toprule
\textbf{Preguntas / palabras clave} & \textbf{Notas / respuestas} \\
\midrule
}
\newcolumntype{Q}{>{\raggedright\arraybackslash}p{0.26\textwidth}}
\newcolumntype{Z}{>{\raggedright\arraybackslash}p{0.68\textwidth}}
\newcolumntype{V}{>{\raggedright\arraybackslash}p{\dimexpr0.94\textwidth+2\tabcolsep+\arrayrulewidth\relax}}
\newcommand{\cornellinicio}{\begingroup\small\renewcommand{\arraystretch}{1.4}}
\newcommand{\cornellfin}{\endgroup}
\newcommand{\cornellpregunta}[2]{\textbf{#1} & #2 \\ \hline}
\newcommand{\cornellresumen}[1]{\rowcolor{blue!5}\multicolumn{2}{|V|}{\textbf{Resumen:} #1} \\ \hline}

\newlength{\cornellL}
\newlength{\cornellR}
\setlength{\cornellL}{0.27\textwidth}
\setlength{\cornellR}{0.63\textwidth}

% -------------------------------------------------------------------------
% CAJAS ACADÉMICAS (uso semántico)
% -------------------------------------------------------------------------
\newtcolorbox{definition}[1]{
    enhanced, breakable,
    title={Definición: #1},
    fonttitle=\bfseries,
    colback=blue!3, colframe=blue!50!black,
    boxrule=0.8pt, arc=2mm
}
\newtcolorbox{important}{
    enhanced, breakable,
    title={Importante},
    fonttitle=\bfseries,
    colback=yellow!8, colframe=orange!70!black,
    boxrule=0.8pt, arc=2mm
}
\newtcolorbox{example}[1]{
    enhanced, breakable,
    title={Ejemplo: #1},
    fonttitle=\bfseries,
    colback=green!4, colframe=green!40!black,
    boxrule=0.8pt, arc=2mm
}
\newtcolorbox{formula}[1]{
    enhanced, breakable,
    title={Fórmula / Regla: #1},
    fonttitle=\bfseries,
    colback=purple!4, colframe=purple!50!black,
    boxrule=0.8pt, arc=2mm
}
\newtcolorbox{exercise}[1]{
    enhanced, breakable,
    title={Ejercicio: #1},
    fonttitle=\bfseries,
    colback=gray!5, colframe=gray!60!black,
    boxrule=0.8pt, arc=2mm
}
\newtcolorbox{note}{
    enhanced, breakable,
    title={Nota},
    fonttitle=\bfseries,
    colback=white, colframe=gray!50,
    boxrule=0.6pt, arc=2mm
}
% Transcripción de disposiciones normativas citadas
\newtcolorbox{norma}[1]{
    enhanced, breakable, frame hidden,
    borderline west={2.5pt}{0pt}{blue!50!black},
    colback=gray!6, boxrule=0pt, arc=0pt,
    left=8pt, right=6pt,
    before upper={\textbf{\upshape #1}\par\smallskip\itshape}
}
\newtcolorbox{articulo}[1]{
    enhanced, breakable, frame hidden,
    borderline west={2pt}{0pt}{gray!60},
    colback=gray!4, boxrule=0pt, sharp corners,
    left=10pt, top=3pt, bottom=3pt,
    before upper={\textbf{\upshape #1.}\ \itshape}
}


\begin{document}

\begin{titlepage}
\thispagestyle{empty}
\begin{center}
\vspace*{1cm}
{\large \textsc{\universidad}}\\[0.3cm]
{\large \textsc{\facultad}}\\[0.3cm]
{\large Cátedra de Derecho Civil, Parte General}

\vspace{3cm}
{\Huge \textbf{\materia}}\\[0.6cm]
{\LARGE Introducción y Parte General}

\vspace{1.2cm}
\rule{0.70\textwidth}{0.5pt}\\[0.5cm]
{\large \subtitulo}\\[0.5cm]
\rule{0.70\textwidth}{0.5pt}

\vfill
{\large \autor}\\[0.3cm]
{\large \anio}
\end{center}
\vfill
\begin{flushleft}
\textit{Este es un modesto aporte a los estudiantes de ``Derecho Civil --- Parte General''; no pretende sustituir el material oficial de la materia.}
\end{flushleft}
% TODO: verificar consistencia entre fuentes originales. Los archivos fuente indican universidad y año de forma no uniforme (por ejemplo, ``Universidad Argentina'' en un archivo y años entre 2020 y 2026); se adoptó la versión predominante.
\end{titlepage}

\frontmatter
\tableofcontents

\chapter{Introducción general}

Este libro reúne en un único recorrido el estudio de la introducción al derecho civil y de la parte general del derecho civil argentino, tomando como eje el Código Civil y Comercial de la Nación (CCyC). Su organización sigue los elementos de la relación jurídica: el \textbf{sujeto} (quién entabla la relación), el \textbf{objeto} (la materia sobre la que recae) y la \textbf{causa} (el origen del que emergen las relaciones jurídicas), y se completa con la \textbf{extinción} de esas relaciones, es decir, con el modo en que termina aquello que nació.

\begin{table}[H]
\centering
\small
\renewcommand{\arraystretch}{1.25}
\begin{tabularx}{\textwidth}{>{\raggedright\arraybackslash}p{0.17\textwidth} >{\raggedright\arraybackslash}p{0.33\textwidth} >{\raggedright\arraybackslash}X}
\toprule
\textbf{Parte} & \textbf{Cuestión que responde} & \textbf{Contenido} \\
\midrule
I. Fundamentos & ¿De dónde proviene el derecho aplicable y dónde están los límites de su ejercicio? & Fuentes del derecho; ley, costumbre, jurisprudencia y doctrina; abuso del derecho. \\
II. La persona humana & ¿Quién es sujeto de derechos y cómo se protege? & Capacidad y sus restricciones; derechos personalísimos; fin de la existencia; ausencia y presunción de fallecimiento. \\
III. Las personas jurídicas & ¿Qué otros sujetos pueden ser titulares de derechos? & Concepto, régimen, responsabilidad, disolución, fundaciones y asociaciones. \\
IV. Patrimonio, derechos y bienes & ¿Sobre qué recaen las relaciones y con qué responde el deudor? & Patrimonio; clasificación de los derechos; bienes y cosas; garantía patrimonial y acciones de los acreedores. \\
V. Hechos y actos jurídicos & ¿Cómo nacen y qué efectos producen las relaciones jurídicas? & Hechos jurídicos y acto voluntario; acto jurídico y sus elementos; forma e instrumentos; clasificación, modalidades y efectos. \\
VI. Vicios de los actos jurídicos & ¿Qué ocurre cuando el acto está afectado en su voluntad, en su equilibrio o en su sinceridad? & Error, dolo, violencia; lesión; simulación y fraude. \\
VII. Extinción & ¿Cómo terminan las relaciones jurídicas? & Hechos y actos extintivos; prescripción y caducidad. \\
\bottomrule
\end{tabularx}
\caption{Mapa general del libro}
\end{table}

Cada parte se despliega en capítulos que avanzan desde los conceptos fundamentales hacia sus aplicaciones, y la mayoría de ellos se cierra con un cuadro Cornell de repaso que permite recuperar las ideas centrales mediante preguntas. Las cajas de color se reservan para definiciones, reglas, ejemplos y advertencias, y el texto corrido sigue siendo el soporte principal de la exposición.

\mainmatter
\partintro{Antes de estudiar a las personas, los bienes y los actos, esta parte presenta el marco desde el cual se razona jurídicamente: de dónde proviene el derecho aplicable (las fuentes), con qué jerarquía y desde cuándo rige, cómo debe interpretarse y cuáles son los límites al modo de ejercer los derechos.}
\part{Fundamentos: las fuentes del derecho y el ejercicio de los derechos}
\chapter[Concepto y clasificación de las fuentes]{Concepto y clasificación de las fuentes del derecho}

Los seis primeros capítulos de esta parte abordan las fuentes del derecho en el derecho civil argentino, con especial énfasis en el Código Civil y Comercial de la Nación (Ley 26.994, en adelante CCyC). Se estudian la clasificación de las fuentes (formales y materiales), la ley como fuente principal (su concepto, clasificación, jerarquía y vigencia), la costumbre jurídica, la jurisprudencia (incluidos los fallos plenarios) y la doctrina como fuente material.

\begin{note}
El conocimiento de las fuentes del derecho constituye una competencia fundamental para abogados, jueces, escribanos, funcionarios públicos y asesores legales: todo el razonamiento jurídico descansa en saber dónde buscar la norma aplicable, cuál tiene mayor jerarquía y cómo interpretarla cuando el texto no es claro. Por ejemplo, quien redacta un contrato debe saber si una cláusula contradice una norma imperativa (y por tanto es nula) o si simplemente desplaza una norma supletoria (lo cual está permitido). Un juez que resuelve un litigio debe ordenar las fuentes: primero la Constitución y los tratados de derechos humanos, luego las leyes, luego los decretos, y sólo entonces la costumbre o la doctrina como orientación no vinculante. Dominar el tema evita también errores graves: citar jurisprudencia como si fuera obligatoria cuando no lo es, ignorar un fallo plenario que sí lo es, o desconocer que una costumbre mercantil puede integrar válidamente un contrato comercial.
\end{note}

\section{Concepto de fuentes del derecho}

Las fuentes son los \textbf{modos de expresión del derecho}, es decir, las maneras en que el derecho se exterioriza. Su estudio comprende el concepto y la clasificación de las fuentes, las disposiciones generales del Código Civil y Comercial y el desarrollo de la principal de ellas, que es la ley.

\section{Fuentes formales y fuentes materiales}

Es habitual y clásica la distinción entre fuentes formales y fuentes materiales.

\begin{definition}{Fuentes formales}
Son aquellas que son \textbf{obligatorias} y \textbf{vinculan al intérprete}; constituyen el modo de manifestarse externamente del derecho positivo.
\end{definition}

\begin{definition}{Fuentes materiales}
Son aquellas que \textbf{no son obligatorias}. Dan el contenido; persuaden por su razonabilidad; ofrecen razonamientos y argumentos en torno a lo justo, pero no son vinculantes.
\end{definition}

Quien tiene a su cargo interpretar la ley (por ejemplo, un juez) no está obligado a seguir una fuente material. Sin embargo, la fuente material muchas veces \textbf{nutre a las fuentes formales}: la jurisprudencia, que durante mucho tiempo es una fuente material, cuando se consolida puede dar lugar a que se sancione, por ejemplo, una ley que recepta algo que la jurisprudencia ya venía sosteniendo.

\subsection{Fuentes formales en el derecho argentino}

\begin{table}[H]
\centering
\begin{tabularx}{\textwidth}{>{\raggedright\arraybackslash}p{0.27\textwidth} >{\raggedright\arraybackslash}X}
\toprule
\textbf{Fuente formal} & \textbf{Definición} \\
\midrule
La ley en sentido material &
Norma escrita, obligatoria, de carácter general, dictada por quien tiene el cuidado de una sociedad en orden al bien común. Esta norma de carácter general, escrita, promulgada y conocida es una fuente en sentido formal. \\
\addlinespace
La costumbre jurídica &
Conducta repetida de manera generalizada, en forma uniforme, por una porción importante de una comunidad, con la convicción de que se trata de un deber jurídico. No es un mero repetir de un hecho como un uso puramente social: existe una convicción jurídica. \\
\addlinespace
Los fallos plenarios &
Sentencias emanadas de una cámara de apelaciones reunida en pleno, cuando resuelve una cuestión en torno a la cual existen sentencias contradictorias de salas de esa misma cámara. Establece una doctrina que se vuelve obligatoria por disposición del Código Procesal Civil y Comercial de la Nación. \\
\bottomrule
\end{tabularx}
\end{table}

\subsection{Fuentes materiales en el derecho argentino}

\begin{table}[H]
\centering
\begin{tabularx}{\textwidth}{>{\raggedright\arraybackslash}p{0.27\textwidth} >{\raggedright\arraybackslash}X}
\toprule
\textbf{Fuente material} & \textbf{Descripción} \\
\midrule
La jurisprudencia en general & El conjunto de sentencias que dictan los jueces. \\
\addlinespace
La doctrina & La opinión de los autores sobre los distintos temas jurídicos. \\
\addlinespace
El derecho comparado & El recurso al estudio y análisis del ordenamiento jurídico de otros países, de sus leyes y de las decisiones de sus tribunales. \\
\addlinespace
Los principios y valores jurídicos & Enunciados en el artículo 2 del CCyC. \\
\bottomrule
\end{tabularx}
\end{table}

\section{El desborde de las fuentes}

En el momento actual se asiste a un cierto \textbf{desborde de las fuentes del derecho}. Las razones que se señalan son las siguientes:

\begin{itemize}
    \item La relevancia otorgada a la Constitución, a los tratados de derechos humanos y a los principios jurídicos fundamentales hace que muchas veces se invoquen tratados y que la ley pierda relevancia o fuerza.
    \item Esta situación se acompaña de una cierta crisis del positivismo jurídico y de un avance de teorías interpretativas, como el neoconstitucionalismo en sus distintas variantes. Estas teorías sostienen que el juez debe realizar la mejor interpretación posible, mirar el derecho bajo su mejor luz y tratar de descubrir la mejor solución en cada caso, lo cual quita relevancia a la ley.
    \item La existencia de comités de derechos humanos y la abundante jurisprudencia de la Corte Interamericana de Derechos Humanos, que se pronuncia sobre muchísimos temas (opiniones consultivas, sentencias de homologación de acuerdos, de supervisión del cumplimiento de sentencias, de interpretación de sentencias), genera un enorme material y plantea la duda de si debe o no aplicarse.
    \item Aparece el denominado \textit{soft law}: derecho no tan rígido ni tan formalizado, que surge de recomendaciones de un comité o de algún considerando de la Corte Interamericana en un caso determinado, respecto del cual se plantea cómo debe ser aplicado.
\end{itemize}

Todo ello tiene como consecuencia una cierta \textbf{inseguridad jurídica}. Es el contexto en el que se ejerce el derecho en el siglo XXI: un contexto de creciente \textbf{constitucionalización del derecho privado}.

\begin{important}
Para el derecho civil, el Código Civil y Comercial sigue siendo la fuente principal (fuente formal especial), aunque exista una presencia muy fuerte de los tratados de derechos humanos. Esto marca nuevos desafíos para el intérprete.
\end{important}

\section{Cuadro Cornell: concepto y clasificación de las fuentes}

\input{/dev/null}
\begin{longtable}{@{}>{\raggedright\arraybackslash}p{0.25\textwidth} >{\raggedright\arraybackslash}p{0.68\textwidth}@{}}
\toprule
\textbf{Preguntas / palabras clave} & \textbf{Notas / respuestas} \\
\midrule
\endfirsthead
\toprule
\textbf{Preguntas / palabras clave} & \textbf{Notas / respuestas} \\
\midrule
\endhead
\bottomrule
\endfoot

\textbf{¿Qué son las fuentes del derecho?} &
Son los modos de expresión del derecho; las maneras en que el derecho se exterioriza. Se clasifican en fuentes formales (obligatorias, vinculantes para el intérprete, como la ley, la costumbre y los fallos plenarios) y fuentes materiales (no vinculantes, pero que orientan por su razonabilidad, como la jurisprudencia general, la doctrina y el derecho comparado). \\[4pt]

\textbf{¿Cuál es la diferencia entre fuente formal y fuente material?} &
La fuente formal obliga al intérprete: si existe una ley aplicable, el juez debe aplicarla. La fuente material persuade pero no obliga: un juez puede considerar la doctrina de un autor reconocido o un fallo de otro tribunal, pero no está obligado a seguirlos. La jurisprudencia, salvo los fallos plenarios, es fuente material en Argentina. \\[4pt]

\textbf{¿Cuáles son las fuentes formales y materiales en el derecho argentino?} &
Formales: la ley en sentido material, la costumbre jurídica y los fallos plenarios. Materiales: la jurisprudencia en general, la doctrina, el derecho comparado y los principios y valores jurídicos (art.~2 CCyC). \\[4pt]

\textbf{¿Qué se entiende por desborde de las fuentes?} &
La relevancia dada a la Constitución, a los tratados de derechos humanos y a los principios jurídicos, junto con la crisis del positivismo, el avance de teorías interpretativas como el neoconstitucionalismo y la abundante producción de comités y de la Corte Interamericana (incluido el \textit{soft law}), hace que la ley pierda relevancia y genera inseguridad jurídica. \\

\midrule
\multicolumn{2}{p{0.93\textwidth}}{%
\textbf{Resumen:} Las fuentes son los modos de expresión del derecho. Las formales (ley en sentido material, costumbre jurídica y fallos plenarios) obligan al intérprete; las materiales (jurisprudencia general, doctrina, derecho comparado y principios y valores jurídicos) persuaden por su razonabilidad pero no obligan, aunque pueden nutrir a las formales. En el contexto actual de desborde de las fuentes y constitucionalización del derecho privado, el CCyC sigue siendo la fuente principal del derecho civil.} \\
\end{longtable}


% ----------------------------------------------------------
\chapter[Arts. 1 y 2 del Código Civil y Comercial]{El Código Civil y Comercial: fuentes, aplicación e interpretación}

\section{Artículo 1: fuentes y aplicación}

El artículo 1, titulado ``Fuentes y aplicación'', marca el tono del tema de las fuentes.

\begin{quote}
\textbf{Art.~1 CCyC --- Fuentes y aplicación.} Los casos que este código rige deben ser resueltos según las leyes que resulten aplicables, conforme con la Constitución Nacional y los tratados de derechos humanos en los que la República sea parte. A tal efecto se tendrá en cuenta la finalidad de la norma. Los usos, prácticas y costumbres son vinculantes cuando las leyes o los interesados se refieren a ellos o en situaciones no regladas legalmente, siempre que no sean contrarios a derecho.
\end{quote}

De la primera parte surge la \textbf{ley como fuente formal}. La segunda parte refiere a la \textbf{costumbre}: los usos, prácticas y costumbres son vinculantes (fuente en sentido formal) cuando las leyes o los interesados se refieren a ellos, o en situaciones no regladas legalmente, siempre que no sean contrarios a derecho.

Respecto de la redacción, la palabra ``casos'' parecería remitir siempre a una conflictividad, cuando a veces puede tratarse de una situación o relación jurídica (que es lo que propiamente trata el derecho civil y comercial) sin que haya un caso en sentido jurídico-técnico propio. Se trata de una disquisición sin mayor relevancia.

\subsection{Tratados que no son de derechos humanos}

El artículo 1 tampoco menciona a los tratados y convenciones que no son de derechos humanos. La materia que rige el CCyC es mucho más amplia que los tratados de derechos humanos, y existen tratados y convenciones que, aunque no sean de derechos humanos, son aplicables a materias civiles y especialmente comerciales. Tienen un rango superior a la ley, por lo que son fuente formal aunque no tengan jerarquía constitucional.

\begin{quote}
\textbf{Art.~2594 CCyC --- Normas aplicables (Derecho Internacional Privado).} Las normas jurídicas aplicables a situaciones vinculadas con varios ordenamientos jurídicos nacionales se determinan por los tratados y las convenciones internacionales vigentes de aplicación en el caso y, en defecto de normas de fuente internacional, se aplican las normas del derecho internacional privado argentino de fuente interna.
\end{quote}

\section{La jurisprudencia en el artículo 1: el debate legislativo}

La comisión redactora del anteproyecto, integrada por quien era presidente de la Corte Suprema, por la magistrada que integró la Corte y por la autora del proyecto, terminó su tarea en marzo de 2012 y elevó el anteproyecto. En esa versión, el artículo 1 incluía a la jurisprudencia como fuente: establecía que los casos debían resolverse teniendo en cuenta la jurisprudencia en consonancia con las circunstancias del caso.
% TODO: contenido incompleto o ambiguo en las notas originales (redacción exacta de la frase suprimida y composición de la comisión)

Esto dio lugar a un debate en la comisión bicameral (el Ejecutivo envía el proyecto al Senado y se conforma una comisión bicameral). En el debate en el Congreso, los legisladores decidieron quitar esa frase.

El fundamento dado por la comisión bicameral en su dictamen final es que se optó por suprimir la referencia a la jurisprudencia porque esta importa un criterio vinculado al análisis del asunto judicial en cuestión, en procura de soluciones jurídicas más adecuadas que se fundamentan en el antecedente judicial: algo más que un criterio de interpretación directa. Básicamente se quiso señalar que \textbf{no se quería dar a la jurisprudencia el carácter de fuente formal}.

\begin{important}
Si la jurisprudencia hubiera figurado en el artículo 1, podría haber sido considerada fuente formal y podría haberse pensado en un pasaje hacia un sistema donde la jurisprudencia fuera obligatoria. La jurisprudencia es una fuente material y no tiene carácter de fuente formal, a diferencia de otros sistemas con un deber mucho más estricto de seguir el precedente, que no está presente en el sistema argentino.
\end{important}

\section{Artículo 2: interpretación de la ley}

El artículo 2 brinda criterios para interpretar la ley, tema importante porque la ley siempre está sujeta a interpretación. Se reconoce como antecedente el artículo 16 del Código Civil de Vélez.

\begin{quote}
\textbf{Art.~2 CCyC --- Interpretación.} La ley debe ser interpretada teniendo en cuenta sus palabras, sus finalidades, las leyes análogas, las disposiciones que surgen de los tratados sobre derechos humanos, los principios y los valores jurídicos, de modo coherente con todo el ordenamiento.
\end{quote}

\begin{table}[H]
\centering
\begin{tabularx}{\textwidth}{>{\raggedright\arraybackslash}p{0.24\textwidth} >{\raggedright\arraybackslash}X}
\toprule
\textbf{Criterio} & \textbf{Desarrollo} \\
\midrule
Palabras de la ley &
Se atiende a qué palabras usa la norma. Por ejemplo, el art.~765 dice ``puede liberarse'': el ``puede'' indica una opción o posibilidad del deudor; no dice que el acreedor deba aceptar. \\
\addlinespace
Finalidad de la norma &
Es un concepto difícil de determinar. Por ejemplo, la finalidad de la ley de alquileres podría ser favorecer una mayor seguridad jurídica y la protección de los vulnerables. \\
\addlinespace
Leyes análogas &
Cuando un tema no puede resolverse por el texto ni por la finalidad de la ley, se recurre a otras leyes. Durante años no hubo normas sobre consentimiento informado para actos médicos en general, pero sí una sobre consentimiento informado para trasplante de órganos; por analogía se aplicaba la ley de trasplante de órganos hasta que en 2009 llegó la Ley de Derechos del Paciente (26.529). \\
\addlinespace
Tratados de derechos humanos &
Criterio importante de interpretación, en base al proceso de constitucionalización y a la idea de que los tratados tienen rango constitucional. \\
\addlinespace
Principios y valores jurídicos &
Ha dado lugar a importantes discusiones. Los sectores positivistas señalan que esto abre demasiado la cuestión interpretativa. Una visión iusnaturalista diría que son principios que emanan de la razón humana y están por encima de la ley positiva; una visión neoconstitucionalista, que emanan del consenso internacional (como el principio de dignidad humana en la Declaración Universal de Derechos Humanos). Hay una coincidencia: al final hay una referencia a principios superiores a la ley positiva. \\
\addlinespace
Coherencia del ordenamiento &
Todo debe hacerse de forma coherente con todo el ordenamiento jurídico, tratando de salvar su coherencia interna. \\
\bottomrule
\end{tabularx}
\end{table}

\section{Cuadro Cornell: artículos 1 y 2 del Código}

\begin{longtable}{@{}>{\raggedright\arraybackslash}p{0.25\textwidth} >{\raggedright\arraybackslash}p{0.68\textwidth}@{}}
\toprule
\textbf{Preguntas / palabras clave} & \textbf{Notas / respuestas} \\
\midrule
\endfirsthead
\toprule
\textbf{Preguntas / palabras clave} & \textbf{Notas / respuestas} \\
\midrule
\endhead
\bottomrule
\endfoot

\textbf{¿Qué establece el art.~1 del CCyC sobre las fuentes?} &
Titulado ``Fuentes y aplicación'', dispone que los casos se resuelven según las leyes aplicables, conforme a la Constitución Nacional y los tratados de derechos humanos. También establece que los usos, prácticas y costumbres son vinculantes cuando las leyes o los interesados se refieren a ellos, o en situaciones no regladas, siempre que no sean contrarios a derecho. \\[4pt]

\textbf{¿Por qué se eliminó la jurisprudencia del art.~1?} &
La comisión bicameral, en su dictamen final, optó por suprimir la referencia a la jurisprudencia para no darle el carácter de fuente formal, evitando la posibilidad de pasar a un sistema de precedente obligatorio. La jurisprudencia es fuente material. \\[4pt]

\textbf{¿Qué ocurre con los tratados que no son de derechos humanos?} &
El art.~1 no los menciona, pero existen tratados y convenciones aplicables a materias civiles y comerciales que tienen rango superior a la ley; son fuente formal aunque no tengan jerarquía constitucional. \\[4pt]

\textbf{¿Cómo se interpreta la ley según el art.~2 CCyC?} &
Teniendo en cuenta: 1) las palabras de la ley; 2) su finalidad; 3) las leyes análogas; 4) las disposiciones de los tratados de derechos humanos; 5) los principios y valores jurídicos; 6) todo ello de modo coherente con el ordenamiento jurídico en su conjunto. \\

\midrule
\multicolumn{2}{p{0.93\textwidth}}{%
\textbf{Resumen:} El art.~1 del CCyC reconoce como fuentes formales la ley (conforme a la Constitución y a los tratados de derechos humanos) y la costumbre (cuando la ley o los interesados se refieren a ella o en situaciones no regladas, si no es contraria a derecho). La jurisprudencia fue suprimida del texto para no otorgarle carácter de fuente formal. El art.~2 establece los criterios de interpretación: palabras, finalidad, leyes análogas, tratados de derechos humanos, principios y valores jurídicos y coherencia con todo el ordenamiento.} \\
\end{longtable}


% ----------------------------------------------------------
\chapter[La ley como fuente formal]{La ley como fuente formal}

\section{Concepto de ley: definición iusnaturalista}

\begin{definition}{Ley (Tomás de Aquino, iusnaturalismo)}
La ley es una \textbf{ordenación de la razón} en orden al \textbf{bien común}, \textbf{promulgada} por quien tiene el cuidado de la comunidad.
\end{definition}

La ley marca y prescribe conducta, regula, establece qué se puede hacer, señala aquello que debe realizarse, prohíbe conductas e indica sanciones para quien quebranta una disposición.

\begin{itemize}
    \item \textbf{Ordenación de la razón:} se espera que la ley exprese fundamentos, se derive de principios lógicos y tenga una lógica interna que mueva a la voluntad. Las personas se mueven según fines y, para ver esos fines, cuentan con la racionalidad, que permite comprender qué es aquello que se debe hacer. La ley puede disponer algo que parezca irrazonable, lo que da lugar a la discusión sobre las leyes injustas; pero en principio se supone que toda ley prescribe algo para el bien común.
    \item \textbf{Bien común:} es la realización del conjunto de condiciones que permiten que las personas, las familias y las asociaciones puedan alcanzar su perfección. Para ello se necesita una autoridad que disponga qué debe hacer cada uno, señale los límites de las conductas y lo debido al otro.
    \item \textbf{Promulgada:} la ley necesita ser conocida; no puede quedar en la mente de quien gobierna o del Congreso, sino que debe ser dada a conocer.
    \item \textbf{Proviene de la autoridad:} quien no tiene autoridad no puede lograr, por ejemplo, que le transfieran un impuesto. La Constitución establece los mecanismos formales por los que se establecen las leyes, los decretos y las distintas normas.
\end{itemize}

\begin{example}{Alimentos en la familia}
En la familia se deben alimentos: los padres hacia los hijos. Es algo que hace al bien común porque resulta bueno que haya un deber de acompañar a los niños en su crianza, lo cual tiene implicancias económicas. La ley prescribe que se paguen alimentos, y eso apunta al bien común: si los niños crecen y se desarrollan, la sociedad progresa, con más personas libres, sanas y capaces de entablar relaciones y buscar el bien entre todos.
\end{example}

\section{Ley en sentido formal y ley en sentido material}

Esta clasificación suele generar muchas confusiones.

\begin{table}[H]
\centering
\begin{tabularx}{\textwidth}{>{\raggedright\arraybackslash}p{0.17\textwidth} >{\raggedright\arraybackslash}X >{\raggedright\arraybackslash}p{0.30\textwidth}}
\toprule
\textbf{Tipo} & \textbf{Concepto} & \textbf{Ejemplo} \\
\midrule
Ley en sentido \textbf{formal} &
Actos que emanan del Poder Legislativo y que han sido elaborados por el procedimiento establecido en la Constitución en el título de la formación y sanción de leyes, que requiere también la participación del Poder Ejecutivo, que promulga o veta. &
El Código Civil y Comercial (Ley N.° 26.994): el Congreso, bajo el procedimiento establecido, emite estos actos, los numera y los publica en el Boletín Oficial. \\
\addlinespace
Ley en sentido \textbf{material} &
Norma escrita, general y obligatoria que emana de la autoridad estatal competente. También es una fuente formal; la clave es que es norma escrita en general. Puede ser un decreto del Poder Ejecutivo o una resolución. El carácter de ley se lo dan la generalidad y la obligatoriedad. &
El Decreto 1089/2012, que reglamenta la Ley 26.529 de Derechos del Paciente: no es ley en sentido formal (no fue aprobada por el Congreso) pero sí en sentido material. \\
\bottomrule
\end{tabularx}
\end{table}

\begin{example}{Derechos del paciente}
Ante un problema como paciente (por ejemplo, en el contexto de la pandemia, por no haber recibido un trato adecuado), se recurre a la Ley de Derechos del Paciente (26.529, ley en sentido formal) y a su decreto reglamentario. Ambas son leyes aplicables al caso: leyes en sentido material.
\end{example}

Generalmente la ley en sentido formal es también ley en sentido material. Muy pocas veces hay leyes en sentido formal que no lo son: a veces se aprueba una pensión para alguien o alguna disposición muy específica que no tiene el alcance general de las leyes en sentido material. Por lo tanto, \textbf{no toda ley en sentido formal es ley en sentido material}.
% TODO: contenido incompleto o ambiguo en las notas originales (la frase "toda ley en sentido material formal puede ser también ley en sentido material" es confusa)

\section{Caracteres de la ley como fuente formal}

\begin{table}[H]
\centering
\begin{tabularx}{\textwidth}{>{\raggedright\arraybackslash}p{0.19\textwidth} >{\raggedright\arraybackslash}X}
\toprule
\textbf{Carácter} & \textbf{Explicación} \\
\midrule
Obligatoriedad &
Es imperativa; se impone a la voluntad porque ha sido establecida por la autoridad que tiene la competencia dada constitucionalmente. El art.~4 CCyC establece: ``Las leyes son obligatorias para todos los que habitan el territorio de la República, sean ciudadanos o extranjeros, residentes, domiciliados o transeúntes, sin perjuicio de lo que dispongan las leyes especiales''. \\
\addlinespace
Generalidad &
Rige a un número indeterminado de personas: a quienquiera que esté comprendido en el ámbito de su aplicación. Esto la distingue de los actos administrativos, donde el Estado decide respecto de un particular concreto (por ejemplo, al adjudicar una licitación). \\
\addlinespace
Origen público &
Emana de la autoridad con competencia en la materia. En sentido formal, del Congreso; en sentido material, de cada autoridad competente (por ejemplo, resoluciones de organismos como la Inspección General de Justicia), siempre dentro del marco de la ley y la Constitución. \\
\addlinespace
Coactividad &
Ante el incumplimiento existen maneras de forzar el cumplimiento. Por ejemplo, en el derecho civil, si alguien comete un ilícito civil y causa un daño, se puede lograr coactivamente que pague, embargando sus bienes y ejecutándolos. Las sanciones son resarcitorias o represivas. \\
\bottomrule
\end{tabularx}
\end{table}

\section{Procedimiento de formación de las leyes}

\begin{table}[H]
\centering
\begin{tabularx}{\textwidth}{>{\centering\arraybackslash}p{0.06\textwidth} >{\raggedright\arraybackslash}p{0.16\textwidth} >{\raggedright\arraybackslash}X}
\toprule
\textbf{N.°} & \textbf{Etapa} & \textbf{Descripción} \\
\midrule
1° & Iniciativa &
Presentar o proponer un proyecto de ley. Están facultados los miembros de cada cámara y el Poder Ejecutivo; también puede haber iniciativa popular si se dan ciertos requisitos que la Constitución establece con una ley reglamentaria. \\
\addlinespace
2° & Discusión &
Debate al interior de cada cámara. Hay cámara de origen y cámara revisora; ciertas materias deben comenzar por una cámara y no por la otra. Dentro de cada cámara hay comisiones de debate, se elaboran dictámenes, los dictámenes conforman el ``orden del día'', este se lleva al plenario y el plenario vota. \\
\addlinespace
3° & Sanción &
Acto por el cual el Poder Legislativo aprueba el proyecto. La cámara de origen aprueba; la otra revisa y, si hace cambios, vuelve a la cámara de origen. Si la cámara revisora aprobó sus cambios con dos terceras partes, la cámara de origen debe insistir con esa mayoría para vencer a la cámara revisora. \\
\addlinespace
4° & Promulgación &
El Poder Ejecutivo tiene la atribución de promulgar u observar (vetar) parcial o totalmente. La promulgación puede ser expresa (mediante un decreto) o tácita (si transcurren diez días sin que el Ejecutivo se pronuncie, se considera promulgada). \\
\addlinespace
5° & Publicación &
Acto por el cual se da a conocer la ley a través del Boletín Oficial, instancia en la que se puede conocer el texto completo. Ha ocurrido que alguien busca la ley en Google y obtiene un texto que no es el correcto (una versión previa o el proyecto originario); la versión vigente debe buscarse en el Boletín Oficial y en los sitios oficiales. \\
\bottomrule
\end{tabularx}
\end{table}

\section{Clasificación de las leyes}

\begin{table}[H]
\centering
\begin{tabularx}{\textwidth}{>{\raggedright\arraybackslash}p{0.17\textwidth} >{\raggedright\arraybackslash}p{0.18\textwidth} >{\raggedright\arraybackslash}X}
\toprule
\textbf{Criterio} & \textbf{Tipo} & \textbf{Descripción y ejemplo} \\
\midrule
Por su contenido dispositivo & Prohibitiva &
Ejemplo: art.~57 CCyC: ``Queda prohibida toda práctica destinada a producir una alteración genética del embrión que se transmita a su descendencia'' (normas de ingeniería genética sobre embriones, para preservar la integridad de la especie humana y evitar actos eugenésicos). \\
\addlinespace
& Dispositiva &
Regula un hecho. Ejemplo: el plazo mínimo de un contrato de locación según la nueva ley de alquileres del año 2020 debe ser de tres años. \\
\midrule
Por la materia & De fondo (sustantivas) &
Regulan materias que hacen a los derechos fundamentales. Según el art.~75 inc.~12 de la Constitución Nacional, dictarlas es competencia del Congreso de la Nación (como el CCyC). Las provincias no podrían dictar sus propios códigos civiles. \\
\addlinespace
& De forma (adjetivas) &
Códigos procesales: son competencia de las provincias. El sistema federal hace que las provincias conserven el poder que no han delegado en el gobierno central. \\
\midrule
Por su imperatividad & Imperativas (indisponibles) &
Obligatorias; se imponen a las partes, que no pueden dejarlas de lado en sus convenciones. Ejemplo: el plazo mínimo del contrato de locación (3 años); no puede pactarse por 1 año aunque las partes quieran. \\
\addlinespace
& Supletorias (disponibles) &
Las partes pueden dejarlas de lado. Rigen por defecto, cuando las partes no dispusieron nada. Ejemplo: si el código permite la sublocación salvo acuerdo en contrario y el contrato nada dice, rige supletoriamente la permisión de sublocar. \\
\bottomrule
\end{tabularx}
\end{table}

\begin{quote}
\textbf{Art.~962 CCyC --- Normas indisponibles y supletorias en contratos.} Las normas legales relativas a los contratos son supletorias por regla general, en materia de contratos. La excepción se da cuando del modo de expresión, de su contenido o de su contexto resulta su carácter indisponible.
\end{quote}

\subsection{Ejemplo problemático: el artículo 765 CCyC}

\begin{quote}
\textbf{Art.~765 CCyC --- Concepto de obligación dineraria (texto parcial).} Si por el acto por el que se ha constituido la obligación se estipuló dar moneda que no sea de curso legal en la República, la obligación debe considerarse como de dar cantidades de cosas y el deudor puede liberarse dando el equivalente en moneda de curso legal.
\end{quote}

La regla según la cual ``el deudor puede liberarse'' ha dado lugar a dos interpretaciones: una parte de la doctrina argentina considera que es una norma \textbf{supletoria} (porque dice ``puede liberarse''); otra parte sostiene que es \textbf{imperativa}.

\begin{example}{Consecuencias prácticas de la calificación del art.~765}
Supóngase que se prestaron dólares y el contrato dice que quien los recibe se obliga a devolver en dólares y renuncia a la potestad de entregar el equivalente en moneda de curso legal. Si se interpreta que la norma es imperativa, esa cláusula no tendría ningún valor y el deudor podría, de todos modos, entregar pesos. Si se la interpreta como supletoria, la cláusula es válida.
\end{example}

La mayoría en las Jornadas de Derecho Civil del año 2015 opinó que esta norma era de carácter supletorio, y esta es la interpretación que ha primado hasta el momento.

\section{Orden jerárquico de las leyes en el derecho argentino}

La norma suprema es la Constitución Nacional. Los tratados internacionales de derechos humanos con jerarquía constitucional están dotados de esa jerarquía por el artículo 75, inciso 22, y se interpretan como complementarios de las normas de la primera parte de la Constitución.

\begin{table}[H]
\centering
\begin{tabularx}{\textwidth}{>{\centering\arraybackslash}p{0.09\textwidth} >{\raggedright\arraybackslash}X >{\raggedright\arraybackslash}p{0.32\textwidth}}
\toprule
\textbf{Rango} & \textbf{Norma} & \textbf{Base constitucional} \\
\midrule
1° & Constitución Nacional y tratados de derechos humanos con jerarquía constitucional & Art.~75 inc.~22 CN: norma suprema. \\
2° & Tratados internacionales y concordatos (no de derechos humanos), con jerarquía supralegal & Art.~75 inc.~22 CN: superiores a las leyes, inferiores a la Constitución. \\
3° & Leyes nacionales & Art.~75 inc.~12 y otros. \\
4° & Decretos del Poder Ejecutivo & Art.~99 CN. \\
5° & Resoluciones y reglamentos de organismos de la administración pública & Marco dado por las leyes. \\
\bottomrule
\end{tabularx}
\end{table}

Esta misma estructura se puede repetir en cada provincia con sus propias constituciones provinciales, leyes y decretos.

\begin{example}{Resolución de la Inspección General de Justicia}
Al constituirse una asociación civil resulta aplicable una resolución de la Inspección General de Justicia (IGJ) en el ámbito de Buenos Aires, que debe ser obedecida. Pero esa resolución no puede contradecir ni el Código Civil ni la Constitución Nacional. Si, en su manera de interpretar este marco, la IGJ los contradijera, estaría ante algo inválido que se podría impugnar.
\end{example}

\section{Cuadro Cornell: la ley como fuente formal}

\begin{longtable}{@{}>{\raggedright\arraybackslash}p{0.25\textwidth} >{\raggedright\arraybackslash}p{0.68\textwidth}@{}}
\toprule
\textbf{Preguntas / palabras clave} & \textbf{Notas / respuestas} \\
\midrule
\endfirsthead
\toprule
\textbf{Preguntas / palabras clave} & \textbf{Notas / respuestas} \\
\midrule
\endhead
\bottomrule
\endfoot

\textbf{¿Cuál es la definición de ley según Tomás de Aquino?} &
Ordenación de la razón en orden al bien común, promulgada por quien tiene el cuidado de la comunidad. Esta definición iusnaturalista destaca que la ley debe ser racional (derivada de principios lógicos), orientada al bien común, proveniente de quien tiene autoridad y promulgada (conocida). \\[4pt]

\textbf{¿Cuál es la diferencia entre ley formal y ley material?} &
Ley en sentido formal: emana del Poder Legislativo siguiendo el procedimiento constitucional (con participación del Poder Ejecutivo, que promulga o veta). Ley en sentido material: cualquier norma escrita, general y obligatoria emanada de autoridad estatal competente, aunque no sea del Congreso (decretos, resoluciones, etc.). Generalmente la ley formal es también material, pero no toda ley formal lo es (por ejemplo, la que aprueba una pensión para una persona determinada). \\[4pt]

\textbf{¿Cuáles son los caracteres de la ley?} &
Obligatoriedad (art.~4 CCyC), generalidad (rige a un número indeterminado de personas), origen público (autoridad competente) y coactividad (existen medios para forzar el cumplimiento; las sanciones son resarcitorias o represivas). \\[4pt]

\textbf{¿Cómo se forma una ley?} &
Iniciativa, discusión, sanción, promulgación (expresa o tácita a los diez días) y publicación en el Boletín Oficial. \\[4pt]

\textbf{¿Qué son las leyes imperativas y supletorias?} &
Las imperativas (indisponibles) se imponen a las partes y no pueden ser dejadas de lado por convención privada (ej.: el plazo mínimo de 3 años en la locación). Las supletorias (disponibles) rigen por defecto cuando las partes no dispusieron nada y pueden ser desplazadas por acuerdo (ej.: la sublocación, permitida salvo pacto en contrario). El art.~962 CCyC establece que en materia de contratos la regla es que las normas son supletorias. \\[4pt]

\textbf{¿Cuál es el orden jerárquico normativo en Argentina?} &
1°) Constitución Nacional y tratados de derechos humanos con jerarquía constitucional (art.~75 inc.~22). 2°) Tratados internacionales sin jerarquía constitucional (supralegales, pero inferiores a la Constitución). 3°) Leyes nacionales. 4°) Decretos del Poder Ejecutivo. 5°) Resoluciones y reglamentos de organismos administrativos. Una resolución que contradiga las normas superiores es inválida y puede impugnarse. \\

\midrule
\multicolumn{2}{p{0.93\textwidth}}{%
\textbf{Resumen:} La ley, definida como ordenación de la razón al bien común promulgada por quien cuida de la comunidad, puede ser formal (emanada del Congreso por el procedimiento constitucional) o material (norma escrita, general y obligatoria de autoridad competente). Sus caracteres son obligatoriedad, generalidad, origen público y coactividad. Se clasifica por su contenido (prohibitiva, dispositiva), por la materia (de fondo, de forma) y por su imperatividad (imperativas y supletorias). La jerarquía normativa va de la Constitución y los tratados de derechos humanos con jerarquía constitucional hasta las resoluciones y reglamentos administrativos.} \\
\end{longtable}


% ----------------------------------------------------------
\chapter[Vigencia de la ley y cómputo de plazos]{Vigencia temporal de la ley y cómputo de plazos}

\section{Vigencia temporal de la ley}

\begin{quote}
\textbf{Art.~5 CCyC --- Vigencia.} Las leyes rigen después del octavo día de su publicación oficial, o desde el día que ellas determinen.
\end{quote}

La ley rige desde el octavo día de su publicación en el Boletín Oficial (desde el día de la publicación se cuentan ocho días y allí comienza su vigencia), salvo que la propia ley determine algo distinto. A veces la ley dice ``esta ley rige desde el día de su publicación'' o ``a los 90 días'' o ``a los 180 días''.

\begin{example}{Entrada en vigencia del Código Civil y Comercial}
El Código Civil y Comercial iba a regir inicialmente a partir del 1° de enero de 2016, pero otra ley dispuso que rigiera a partir del 1° de agosto de 2015. El Código fue sancionado en octubre de 2014 y publicado en el Boletín Oficial el 8 de octubre de 2014; desde octubre hasta agosto no estuvo vigente, porque la ley determinó que entrara en vigencia el 1° de agosto. El legislador puede establecer cuándo rige la ley.
\end{example}

\subsection{Modos de derogación}

\begin{table}[H]
\centering
\begin{tabularx}{\textwidth}{>{\raggedright\arraybackslash}p{0.27\textwidth} >{\raggedright\arraybackslash}X}
\toprule
\textbf{Tipo} & \textbf{Descripción} \\
\midrule
Expresa &
La ley posterior dispone explícitamente la derogación de la anterior. Ejemplo: el CCyC (Ley 26.994) derogó expresamente el Código Civil de Vélez Sársfield (Ley 340). \\
\addlinespace
Tácita por incompatibilidad &
Cuando la ley posterior y la anterior resultan incompatibles entre sí, la nueva ley deroga a la anterior. \\
\addlinespace
Ley especial deroga ley general &
Cuando hay una ley especial en la materia, se aplica la especial. Ejemplo: en materia de sociedades, la ley de sociedades es la ley especial respecto del CCyC (ley general). \\
\addlinespace
Ley general posterior no deroga ley especial &
Salvo que se pueda percibir la voluntad derogatoria del legislador. \\
\addlinespace
Caducidad &
A veces la propia ley establece su duración temporal; o pierde fuerza porque hay una costumbre contraria a ella o porque cambiaron las circunstancias y ya no tiene ese valor. Esto suele pasar muy raramente. \\
\bottomrule
\end{tabularx}
\end{table}

\section{Modo de contar los intervalos del derecho (art.~6 CCyC)}

\begin{quote}
\textbf{Art.~6 CCyC --- Modo de contar los intervalos del derecho.} El modo de contar los intervalos del derecho es el siguiente: día es el intervalo que corre de medianoche a medianoche. En los plazos fijados en días, a contar desde uno determinado, queda éste excluido del cómputo, el cual debe empezar al día siguiente. Los plazos de meses o años se computan de fecha a fecha. Cuando en el mes del vencimiento no hubiere día equivalente al inicial del cómputo, se entiende que el plazo expira el último día de ese mes. Los plazos vencen a la hora veinticuatro del día del vencimiento respectivo. El cómputo civil de los plazos es de días completos y continuos, y no se excluyen los días inhábiles o no laborables. En los plazos fijados en horas, a contar desde una hora determinada, queda ésta excluida del cómputo, el cual debe empezar desde la hora siguiente.
\end{quote}

\begin{table}[H]
\centering
\begin{tabularx}{\textwidth}{>{\raggedright\arraybackslash}p{0.21\textwidth} >{\raggedright\arraybackslash}X >{\raggedright\arraybackslash}p{0.28\textwidth}}
\toprule
\textbf{Tipo de plazo} & \textbf{Regla de cómputo} & \textbf{Ejemplo} \\
\midrule
En días &
Se excluye el día del cómputo; empieza a contarse desde el día siguiente. Un día se toma de medianoche a medianoche. &
Plazo de 15 días desde el 15/08: no cuenta el 15; comienza el 16; vence el 30/08 a las 24 h. \\
\addlinespace
En meses &
Se cuentan de fecha a fecha (mes a mes), no por días exactos. &
3 meses desde el 15/08: vence el 15/11, aunque tengan distintas cantidades de días. \\
\addlinespace
En años &
Se cuentan de fecha a fecha (año a año). &
1 año desde el 15/08/2020 vence el 15/08/2021. \\
\addlinespace
Mes de vencimiento sin día equivalente &
Si en el mes del vencimiento no existe esa fecha, se toma el último día del mes. &
Plazo que vence el 31/02 (inexistente): vence el 28/02 (o 29 en año bisiesto). \\
\addlinespace
Días inhábiles &
El cómputo civil es de días completos y continuos; no se excluyen los días inhábiles o no laborables (admite prueba en contrario por ley o acuerdo de partes). &
A diferencia del cómputo procesal, que sí excluye inhábiles. \\
\bottomrule
\end{tabularx}
\end{table}

\begin{note}
Estas normas son de carácter supletorio, porque las partes pueden disponer un cómputo distinto en sus contratos o una ley especial puede establecer otras reglas.
\end{note}

\section{Cuadro Cornell: vigencia de la ley y cómputo de plazos}

\begin{longtable}{@{}>{\raggedright\arraybackslash}p{0.25\textwidth} >{\raggedright\arraybackslash}p{0.68\textwidth}@{}}
\toprule
\textbf{Preguntas / palabras clave} & \textbf{Notas / respuestas} \\
\midrule
\endfirsthead
\toprule
\textbf{Preguntas / palabras clave} & \textbf{Notas / respuestas} \\
\midrule
\endhead
\bottomrule
\endfoot

\textbf{¿Cuándo entra en vigencia una ley?} &
Desde el octavo día de su publicación en el Boletín Oficial, salvo que la propia ley determine algo distinto (desde el día de su publicación, a los 90 días, a los 180 días, o una fecha determinada). \\[4pt]

\textbf{¿Cómo deja de regir una ley?} &
Por derogación expresa (la nueva ley lo indica), tácita (incompatibilidad entre leyes), o mediante la regla de que la ley especial desplaza a la general; una ley general posterior no deroga a la especial, salvo voluntad derogatoria del legislador. También por caducidad, situación que suele ocurrir muy raramente. \\[4pt]

\textbf{¿Cómo se cuentan los plazos según el art.~6 CCyC?} &
Un día va de medianoche a medianoche. Los plazos en días excluyen el día del cómputo (se empieza a contar desde el día siguiente). Los plazos en meses y años se cuentan de fecha a fecha. Si en el mes de vencimiento no existe el día equivalente, se toma el último día de ese mes. El cómputo civil es de días completos y continuos, sin excluir inhábiles. Estas reglas son supletorias: las partes o la ley especial pueden pactar otra cosa. \\

\midrule
\multicolumn{2}{p{0.93\textwidth}}{%
\textbf{Resumen:} Las leyes rigen desde el octavo día de su publicación oficial o desde la fecha que ellas determinen, y dejan de regir por derogación expresa o tácita, por la primacía de la ley especial o por caducidad. Los plazos civiles se computan de días completos y continuos, sin excluir inhábiles, salvo acuerdo o ley especial en contrario.} \\
\end{longtable}


% ----------------------------------------------------------
\chapter[La costumbre jurídica]{La costumbre como fuente formal}

\section{Concepto y elementos}

\begin{definition}{Costumbre}
Observancia constante y uniforme de un cierto comportamiento por parte de los miembros de una comunidad social, con la convicción de que ese comportamiento responde a una necesidad jurídica.
\end{definition}

La costumbre es una conducta que se repite, generalizada y uniforme, pero no cualquier conducta repetida: existe una \textbf{convicción de que es un deber jurídico}. Este elemento es muy importante, porque distingue la costumbre de un mero uso social.

\begin{table}[H]
\centering
\begin{tabularx}{\textwidth}{>{\raggedright\arraybackslash}p{0.16\textwidth} >{\raggedright\arraybackslash}p{0.20\textwidth} >{\raggedright\arraybackslash}X}
\toprule
\textbf{Elemento} & \textbf{Denominación} & \textbf{Descripción} \\
\midrule
Objetivo & \textit{Corpus} & Repetición de un comportamiento por cierto tiempo razonablemente prolongado, por parte de una comunidad. \\
\addlinespace
Subjetivo & \textit{Animus} (\textit{opinio iuris}) & Convicción de los miembros de la comunidad de la obligatoriedad de esa conducta. Permite distinguir la costumbre de cualquier otro uso social. \\
\bottomrule
\end{tabularx}
\end{table}

\begin{example}{Costumbre y uso social}
Un ejemplo trivial: una persona tiene la costumbre de ir siempre a un mismo café, pero ello no responde a un deber jurídico; no se obligó a consumir en ese café todos los días. Para demostrar que existe una costumbre jurídica hay que demostrar la observancia constante de un comportamiento, pero con la convicción de su obligatoriedad.
% TODO: contenido incompleto o ambiguo en las notas originales (el desarrollo del ejemplo del café y del cobro del doble es confuso)
\end{example}

\section{Importancia histórica de la costumbre}

La costumbre fue durante mucho tiempo la fuente más importante en la producción normativa. Ha influido en institutos típicos del derecho francés, que son la base del ordenamiento argentino:

\begin{table}[H]
\centering
\begin{tabularx}{\textwidth}{>{\raggedright\arraybackslash}p{0.34\textwidth} >{\raggedright\arraybackslash}X}
\toprule
\textbf{Instituto} & \textbf{Descripción} \\
\midrule
El patrimonio como garantía de los acreedores & La expectativa de poder exigir y ejecutar coactivamente el patrimonio (los bienes) del deudor. \\
\addlinespace
La fuerza vinculante de los contratos (\textit{pacta sunt servanda}) & Los pactos deben cumplirse y honrarse entre las partes. \\
\addlinespace
La revocabilidad de las donaciones & Instituto que surgió de la costumbre. \\
\addlinespace
La prescripción desde la exigibilidad & El cómputo de la prescripción comienza el día de la exigibilidad de una obligación. \\
\bottomrule
\end{tabularx}
\end{table}

Todas estas cuestiones surgieron del derecho consuetudinario en Francia, fueron recogidas por el Código Francés y pasaron así a los distintos ordenamientos. Además, la costumbre tiene importancia, y sigue teniéndola, en materia comercial y en muchas materias vinculadas con los contratos.

\section{Tipos de costumbre}

\begin{table}[H]
\centering
\begin{tabularx}{\textwidth}{>{\raggedright\arraybackslash}p{0.20\textwidth} >{\raggedright\arraybackslash}p{0.14\textwidth} >{\raggedright\arraybackslash}X}
\toprule
\textbf{Tipo} & \textbf{Latín} & \textbf{Descripción} \\
\midrule
Costumbre según la ley & \textit{Secundum legem} &
Rige cuando la ley delega en la costumbre la atribución de regular una cierta situación; su autoridad está de alguna manera delegada por la propia ley. Ejemplo: el art.~1255 CCyC sobre el precio de los contratos de obra y servicios, regulado por los usos y costumbres. \\
\addlinespace
Costumbre en ausencia de ley & \textit{Praeter legem} &
Rige en ausencia de disposición legal: en cierta materia nunca hubo una ley que regulara un tema, pero existió una costumbre. Ejemplo: antes de la Ley 18.248 (sobre nombre de personas físicas) existían costumbres que regulaban el nombre. \\
\addlinespace
Costumbre contra la ley & \textit{Contra legem} &
La que contradice una disposición expresa; en principio no es válida. Ejemplo histórico: una norma en Francia que prohibía a las mujeres usar pantalones, vigente hasta 2010; la costumbre ya la había dejado obsoleta aunque no hubiera sido derogada formalmente. \\
\bottomrule
\end{tabularx}
\end{table}

\section{La costumbre en el Código Civil y Comercial}

En el Código Civil de Vélez, el artículo 17 decía que los usos y costumbres no pueden crear derechos sino cuando las leyes se refieren a ellos o en situaciones no regladas legalmente: una fórmula negativa que el nuevo Código modifica. En el Código de Comercio, antes de la unificación, también había una gran importancia de la costumbre; su artículo 2 decía que en las materias en que las convenciones particulares pueden derogar la ley, la tradición (costumbre) autoriza al juez a indagar si la esencia del acto utiliza la costumbre.
% TODO: contenido incompleto o ambiguo en las notas originales (paráfrasis del art. 2 del Código de Comercio)

El CCyC amplía la noción respecto del código anterior: antes se hablaba de ``usos y costumbres''; ahora se agrega la expresión ``prácticas''. La costumbre es vinculante (fuente formal) cuando se dan los siguientes supuestos, conforme a la última parte del art.~1:

\begin{table}[H]
\centering
\begin{tabularx}{\textwidth}{>{\raggedright\arraybackslash}p{0.27\textwidth} >{\raggedright\arraybackslash}X}
\toprule
\textbf{Supuesto} & \textbf{Explicación} \\
\midrule
La ley se refiere a ellos & Costumbre \textit{secundum legem}. Igual que en el código anterior. \\
\addlinespace
Las partes se refieren a ellos & Novedad del CCyC: las partes pueden pactar que su relación jurídica se rija en ciertos aspectos por las costumbres. Tiene especial relevancia en materia comercial. \\
\addlinespace
Situaciones no regladas legalmente & Costumbre \textit{praeter legem}. Ya estaba en el código anterior. \\
\addlinespace
Que no sean contrarios a derecho & Condición agregada por el CCyC (no estaba en el código anterior). Es una fórmula de mayor amplitud que simplemente ``no oponerse a la ley'': la costumbre no puede ser contraria al derecho en su conjunto. \\
\bottomrule
\end{tabularx}
\end{table}

\subsection{Presencia de la costumbre en materia de contratos}

\begin{quote}
\textbf{Art.~964 CCyC --- Integración del contrato.} El contenido del contrato se integra con: a) las normas indisponibles, que se aplican en sustitución de las cláusulas incompatibles con ellas; b) las normas supletorias; c) los usos y prácticas del lugar de celebración, en cuanto sean aplicables porque hayan sido declarados obligatorios por las partes o porque sean ampliamente conocidos y regularmente observados en el ámbito en que se celebra el contrato, siempre que su aplicación sea razonable.
\end{quote}

\begin{quote}
\textbf{Art.~970 CCyC --- Contratos nominados e innominados (parcial).} Los contratos innominados están regidos, en el siguiente orden, por: a) la voluntad de las partes; b) las normas generales sobre contratos y obligaciones; c) los usos y prácticas del lugar de celebración [\ldots].
\end{quote}

\begin{quote}
\textbf{Art.~1063 CCyC --- Significado de las palabras (interpretación).} Las palabras empleadas en el contrato deben entenderse en el sentido que les da el uso general, excepto que tengan un significado específico que surja de la ley, del acuerdo de las partes o de los usos y prácticas del lugar de celebración.
\end{quote}

\begin{table}[H]
\centering
\begin{tabularx}{\textwidth}{>{\raggedright\arraybackslash}p{0.16\textwidth} >{\raggedright\arraybackslash}X >{\raggedright\arraybackslash}p{0.14\textwidth}}
\toprule
\textbf{Operación} & \textbf{En qué consiste} & \textbf{Art.} \\
\midrule
Integrar & Determinar qué cláusulas existen. Si un contrato no dice nada sobre un tema, se recurre a la ley imperativa, luego a la supletoria y luego a los usos y prácticas, que pasan a ser parte del contenido del contrato. & 964 \\
\addlinespace
Interpretar & Determinar el alcance y sentido de una cláusula existente. Uno de los criterios para interpretar las palabras de un contrato son los usos y prácticas del lugar. & 1063 \\
\bottomrule
\end{tabularx}
\end{table}

\section{Cuadro Cornell: la costumbre jurídica}

\begin{longtable}{@{}>{\raggedright\arraybackslash}p{0.25\textwidth} >{\raggedright\arraybackslash}p{0.68\textwidth}@{}}
\toprule
\textbf{Preguntas / palabras clave} & \textbf{Notas / respuestas} \\
\midrule
\endfirsthead
\toprule
\textbf{Preguntas / palabras clave} & \textbf{Notas / respuestas} \\
\midrule
\endhead
\bottomrule
\endfoot

\textbf{¿Cuáles son los elementos de la costumbre jurídica?} &
\textit{Corpus} (elemento objetivo): repetición constante y uniforme de un comportamiento por parte de una comunidad, durante un tiempo razonablemente prolongado. \textit{Animus} u \textit{opinio iuris} (elemento subjetivo): convicción de los miembros de la comunidad de que esa conducta responde a un deber jurídico. Sin el segundo elemento, es un mero uso social, no costumbre jurídica. \\[4pt]

\textbf{¿Cuáles son los tipos de costumbre?} &
\textit{Secundum legem}: la ley remite a la costumbre para regular cierta situación (ej.: precio de los contratos de obra). \textit{Praeter legem}: rige en ausencia de ley. \textit{Contra legem}: contradice una disposición legal expresa; en principio inválida. El CCyC exige que la costumbre no sea contraria a derecho. \\[4pt]

\textbf{¿Cuándo es vinculante la costumbre según el CCyC?} &
Cuando las leyes o los interesados se refieren a ella, o en situaciones no regladas legalmente, siempre que no sea contraria a derecho. Novedades respecto del código anterior: se agrega el término ``prácticas'', la posibilidad de que las partes se refieran a ellas y la condición de no ser contrarios a derecho. \\[4pt]

\textbf{¿Qué diferencia hay entre integrar e interpretar un contrato?} &
Integrar (art.~964) es determinar qué cláusulas existen: ley imperativa, luego supletoria y luego usos y prácticas. Interpretar (art.~1063) es determinar el alcance y sentido de una cláusula existente, para lo cual los usos y prácticas del lugar son un criterio. \\

\midrule
\multicolumn{2}{p{0.93\textwidth}}{%
\textbf{Resumen:} La costumbre jurídica requiere \textit{corpus} y \textit{animus}. Puede ser \textit{secundum}, \textit{praeter} o \textit{contra legem}. El art.~1 del CCyC la reconoce como vinculante cuando la ley o los interesados se refieren a ella o en situaciones no regladas, siempre que no sea contraria a derecho, y tiene presencia en materia de contratos (arts.~964, 970 y 1063).} \\
\end{longtable}


% ----------------------------------------------------------
\chapter[Jurisprudencia, plenarios y doctrina]{Jurisprudencia, fallos plenarios y doctrina}

\section{La jurisprudencia como fuente material}

La jurisprudencia es el conjunto de sentencias que dictan los tribunales de justicia. Es indudablemente una \textbf{fuente material}, porque no es vinculante.

La sentencia es la decisión judicial adoptada ante un caso concreto y, por lo tanto, tiene fuerza vinculante para las partes de ese caso. Por una aspiración de justicia, si ante una situación análoga se resolvió de una manera, existe la expectativa de obtener la misma resolución ante una situación parecida: el sentido de justicia lleva a preguntarse cómo puede ser que ante el mismo caso se resuelva distinto.

En el sistema argentino, sin embargo, la jurisprudencia no obliga: el hecho de que un juez haya dicho algo no obliga al siguiente, que puede resolver distinto. Tiene autoridad (un juez que ve que varios casos se resolvieron de una manera determinada puede considerar que no debe apartarse), pero no está obligado. Si hay una línea jurisprudencial muy uniforme o francamente mayoritaria, es difícil cambiarla, salvo que existan factores fundamentales superiores vinculados con principios fundamentales del ordenamiento jurídico o con la dignidad de la persona.

Como se vio en el capítulo sobre el CCyC, el artículo 1 no menciona la jurisprudencia. En el proyecto, elaborado por una comisión creada por decreto presidencial en febrero de 2011 y que entregó el anteproyecto en marzo de 2012, el artículo 1 la incluía como criterio para interpretar la ley y resolver un caso. Los legisladores de la comisión bicameral quitaron esa frase para evitar que la jurisprudencia pasara a ser fuente formal con obligatoriedad general.

\section{Los fallos plenarios como fuente formal}

Salvo los fallos plenarios, la jurisprudencia en principio no es fuente formal.

\begin{definition}{Fallo plenario}
Sentencia dictada por una cámara de apelaciones en pleno que resuelve una cuestión en torno a la cual existen sentencias contradictorias dictadas por salas del mismo tribunal, y que se convierte en \textbf{obligatoria} para todas las salas y tribunales inferiores de ese fuero, en esa jurisdicción y competencia.
\end{definition}

Para comprenderlo es necesario entender cómo funciona un tribunal. Los tribunales están organizados en instancias; se toma como ejemplo el Código Procesal Civil y Comercial de la Nación, en el ámbito de la Ciudad de Buenos Aires:

\begin{table}[H]
\centering
\begin{tabularx}{\textwidth}{>{\raggedright\arraybackslash}p{0.19\textwidth} >{\raggedright\arraybackslash}p{0.31\textwidth} >{\raggedright\arraybackslash}X}
\toprule
\textbf{Instancia} & \textbf{Órgano} & \textbf{Función} \\
\midrule
Primera instancia & Juzgados de primera instancia (unipersonales) & Tramita la demanda, contestación, producción de prueba, alegatos y sentencia. \\
\addlinespace
Segunda instancia & Cámara de Apelaciones (organizada en salas de 3 miembros cada una: Sala A, B, C\ldots hasta la M) & Revisa por vía de apelación las decisiones de primera instancia. Cada sala resuelve de manera independiente. \\
\addlinespace
Tribunal en pleno & Todas las salas de la Cámara reunidas & Resuelve cuando hay contradicción entre salas sobre una misma cuestión jurídica. El fallo resultante es obligatorio. \\
\bottomrule
\end{tabularx}
\end{table}

\begin{example}{Contradicción entre salas que motiva un fallo plenario}
Siguiendo con el tema del art.~765 (dólares): en una causa toca la Sala E, que sostiene que el art.~765 es supletorio, de modo que puede pactarse que la liberación sea sólo en dólares. La parte queda conforme porque pactó que el deudor no podría liberarse entregando pesos. En otro juicio se apela y la causa va a la Sala C, que entiende que la norma es imperativa y que las partes no pueden pactar que se entreguen siempre dólares: el deudor siempre tiene la potestad de entregar pesos. Hay contradicción: la Sala E dice una cosa y la Sala C otra.

Ante el fallo de la Sala C que perjudica a la parte, esta interpone el \textbf{recurso de inaplicabilidad de la ley}, que solicita a la Sala C que convoque el plenario de todas las salas. Todas las salas (A, B, C, \ldots, M) se reúnen en plenario y resuelven el tema, unificando la interpretación. Esa posición se convierte en obligatoria para todas las salas y para todos los tribunales inferiores de ese fuero, en esa jurisdicción y competencia.
\end{example}

\begin{quote}
\textbf{Arts.~300 y 303 del CPCCN --- Efectos del fallo plenario.} La interpretación de la ley establecida en una sentencia plenaria será obligatoria para la misma Cámara y para los jueces de primera instancia respecto de los cuales aquélla sea tribunal de alzada, sin perjuicio de que los jueces dejen a salvo su opinión personal. Sólo podrá modificarse dicha doctrina por medio de una nueva sentencia plenaria.
\end{quote}

El fundamento de los fallos plenarios es \textbf{unificar la jurisprudencia} y evitar criterios dispares, porque ofenden el sentido de justicia. El recurso se interpone dentro de los diez días de la sentencia definitiva, mencionando el precedente contradictorio y exponiendo los fundamentos; el tribunal resuelve si es procedente. Si se dicta un fallo plenario, este establece la doctrina legal aplicable y sólo puede ser cambiado por otro plenario.

\begin{important}
Un plenario de la Cámara Civil de Buenos Aires obliga a los jueces civiles de la Ciudad de Buenos Aires, pero no puede obligar a las cámaras de Córdoba o de Mendoza, que tendrán sus propios plenarios de acuerdo con sus códigos procesales.
\end{important}

\section{Difusión de la jurisprudencia}

\begin{table}[H]
\centering
\begin{tabularx}{\textwidth}{>{\raggedright\arraybackslash}p{0.32\textwidth} >{\raggedright\arraybackslash}X}
\toprule
\textbf{Fuente de difusión} & \textbf{Descripción} \\
\midrule
Publicaciones oficiales de la CSJN & La Corte Suprema de Justicia de la Nación tiene su sitio oficial (www.csjn.gov.ar) con sus fallos. \\
\addlinespace
Centro de Información Judicial (CIJ) & Publica a nivel nacional las sentencias que los tribunales tienen la obligación de remitir, con los datos debidamente analizados en los casos donde puede haber una cuestión de intimidad. \\
\addlinespace
Sistema Argentino de Información Jurídica (SAIJ) & También recoge mucha jurisprudencia, la publica y la difunde junto con distintos artículos y revistas. \\
\addlinespace
Revistas especializadas & ``La Ley'', ``El Derecho'', ``Jurisprudencia Argentina'' y muchas otras, varias de base online, que difunden la jurisprudencia. \\
\bottomrule
\end{tabularx}
\end{table}

\section{La doctrina como fuente material}

La doctrina es la opinión de los juristas. Es una fuente material y, obviamente, no es obligatoria.

Hubo un momento histórico, en el primer milenio romano, en que ciertos juristas tenían el \textit{ius respondendi ex auctoritate principis}: la autoridad para responder cuestiones y decir el derecho con la autoridad del príncipe. Fue un momento histórico determinado y ya no sería posible.

La doctrina tiene valor material: los autores escriben y publican en las revistas, esos textos llegan a los jueces, que los citan muchas veces al interpretar. Las partes también citan en sus escritos de demanda y contestación a autores que refuercen su postura. Por ejemplo, al interpretar el art.~765, se cita a un autor que dice que la norma es imperativa y a otro que dice que es supletoria; por la autoridad de ese jurista puede convenir citarlo, y así se convence a los jueces. El juez verá si los argumentos son convincentes, porque es una fuente material y por lo tanto no es obligatoria.

\section{Otras fuentes materiales: derecho comparado y principios jurídicos}

\begin{table}[H]
\centering
\begin{tabularx}{\textwidth}{>{\raggedright\arraybackslash}p{0.22\textwidth} >{\raggedright\arraybackslash}X}
\toprule
\textbf{Fuente material} & \textbf{Descripción y alcance} \\
\midrule
Derecho comparado &
Estudio y análisis de la legislación positiva existente en otros países, así como de la interpretación y aplicación que hacen sus propios tribunales. Tiene un valor relativo: que un juez de otro país haya dicho algo no implica que el juez argentino deba seguirlo. A veces tiene más valor si la ley argentina sigue a una ley extranjera (por ejemplo, el error reconocible del CCyC viene del Código italiano, caso en que el derecho comparado tiene más valor interpretativo). No es fuente formal: el juez no está obligado porque en España se resuelva de determinada manera. \\
\addlinespace
Principios y valores jurídicos &
El art.~2 señala que para interpretar la ley hay que tener en cuenta los principios y valores jurídicos. Es un criterio de interpretación que también de alguna manera es fuente material: informa la manera de determinar lo que es justo. Hay debate: los positivistas señalan que abre demasiado la cuestión interpretativa; el iusnaturalismo dice que son principios que emanan de la razón humana y están por encima de la ley positiva; el neoconstitucionalismo dice que emanan del consenso internacional (como la dignidad humana en la Declaración Universal de Derechos Humanos). \\
\bottomrule
\end{tabularx}
\end{table}

\section{Cuadro Cornell: jurisprudencia, plenarios y doctrina}

% TODO: contenido incompleto o ambiguo en las notas originales (el desarrollo cita los arts. 300 y 303 del CPCCN; el resumen original agrega "arts. 288 y ss.")
\begin{longtable}{@{}>{\raggedright\arraybackslash}p{0.25\textwidth} >{\raggedright\arraybackslash}p{0.68\textwidth}@{}}
\toprule
\textbf{Preguntas / palabras clave} & \textbf{Notas / respuestas} \\
\midrule
\endfirsthead
\toprule
\textbf{Preguntas / palabras clave} & \textbf{Notas / respuestas} \\
\midrule
\endhead
\bottomrule
\endfoot

\textbf{¿Qué significa que la jurisprudencia es fuente material?} &
Que los fallos de los jueces orientan pero no obligan a otros jueces. Cada juez puede resolver distinto del anterior, incluso en casos análogos. La sentencia es vinculante sólo para las partes del juicio. La excepción son los fallos plenarios, cuya doctrina sí es obligatoria. \\[4pt]

\textbf{¿Qué es un fallo plenario y por qué es fuente formal?} &
Es una sentencia dictada por una cámara de apelaciones reunida en pleno, convocada mediante el recurso de inaplicabilidad de la ley, cuando existen criterios contradictorios entre salas. Unifica la interpretación y su doctrina es obligatoria para todas las salas y jueces de primera instancia de ese fuero y jurisdicción. Sólo puede ser modificada por otro plenario. Se regula en los arts.~288 y ss. y 300 a 303 del CPCCN. \\[4pt]

\textbf{¿Qué valor tiene la doctrina como fuente del derecho?} &
Es la opinión de los juristas. Es fuente material: no obliga pero puede persuadir. Los jueces pueden citarla y las partes la utilizan en sus escritos para fundar sus posiciones. Su peso depende de la autoridad del jurista y de sus argumentos. El \textit{ius respondendi} fue un momento histórico que ya no sería posible. \\[4pt]

\textbf{¿Qué valor tienen el derecho comparado y los principios jurídicos?} &
El derecho comparado tiene un valor relativo y no es fuente formal; tiene más valor interpretativo cuando la ley argentina sigue a una ley extranjera. Los principios y valores jurídicos (art.~2) informan la manera de determinar lo que es justo y son objeto de debate entre positivistas, iusnaturalistas y neoconstitucionalistas. \\

\midrule
\multicolumn{2}{p{0.93\textwidth}}{%
\textbf{Resumen:} La jurisprudencia es fuente material: no obliga más allá de las partes del caso, salvo los fallos plenarios, que unifican criterios contradictorios entre salas y son obligatorios para las salas y los tribunales inferiores del fuero. La doctrina, el derecho comparado y los principios y valores jurídicos son fuentes materiales que orientan y persuaden, pero no obligan.} \\
\end{longtable}


% ----------------------------------------------------------

\section{Síntesis de las fuentes del derecho}

Las fuentes del derecho son el punto de partida de todo razonamiento jurídico: indican de dónde proviene la norma que debe aplicarse y qué jerarquía tiene frente a otras.

En el derecho argentino, el Código Civil y Comercial (art.~1) reconoce como fuentes formales (vinculantes) la ley (en sentido formal y material), la costumbre jurídica (cuando la ley o las partes se refieren a ella, o en situaciones no reguladas, siempre que no contravenga el derecho) y los fallos plenarios (sentencias de cámaras en pleno que unifican criterios contradictorios entre salas y resultan obligatorias para los tribunales inferiores del fuero). Son fuentes materiales (orientadoras pero no obligatorias) la jurisprudencia general, la doctrina, el derecho comparado y los principios y valores jurídicos (art.~2).

La ley, como fuente principal, puede ser formal (emanada del Congreso siguiendo el procedimiento constitucional) o material (cualquier norma escrita, general y obligatoria emanada de autoridad competente). Se clasifica, además, en imperativa (cuyas disposiciones las partes no pueden derogar) y supletoria (que rige por defecto cuando las partes no acordaron otra cosa). La interpretación de la ley exige integrar sus palabras con su finalidad, las leyes análogas, los tratados de derechos humanos y los principios jurídicos, todo de manera coherente con el ordenamiento.

Las leyes rigen desde el octavo día de su publicación oficial, salvo que establezcan otra fecha, y los plazos se computan de días completos y continuos, sin excluir inhábiles, salvo acuerdo en contrario. Comprender este sistema jerárquico y las reglas de vigencia e interpretación es indispensable para cualquier profesional del derecho: permite saber qué norma aplicar, con qué fuerza, desde cuándo y cómo.
\chapter{El problema conceptual del abuso del derecho}
% ============================================================

\section{Introducción}

El abuso del derecho es uno de los institutos más importantes del Derecho Civil argentino. Su estudio comprende tres cuestiones: qué significa ejercer un derecho de forma ``abusiva'', cuáles son los criterios teóricos para identificarlo y cómo fue incorporado progresivamente al ordenamiento jurídico argentino, desde el Código Civil de Vélez Sársfield hasta el Código Civil y Comercial de 2015.

\begin{important}
El artículo 10 del Código Civil y Comercial se aplica en todas las ramas del derecho: civil, comercial, laboral, administrativo y ambiental. En contratos, sociedades, propiedad horizontal o relaciones laborales, el abuso del derecho es un límite que los jueces aplican cotidianamente para evitar injusticias.
\end{important}

Desde el punto de vista profesional, el abuso del derecho constituye:
\begin{itemize}
    \item una \textbf{herramienta de defensa y ataque}: puede invocarse para proteger a un cliente cuando alguien ejerce un derecho de manera dañina, inequitativa o contraria a la buena fe;
    \item un \textbf{límite a las pretensiones desmedidas}, aplicado por los jueces en distintos ámbitos;
    \item un \textbf{principio transversal}, que rige también las relaciones de vecindad, de consumo, comerciales y familiares.
\end{itemize}

\begin{example}{hilo conductor del capítulo}
Un vecino construyó una parrilla en el espacio común del edificio. Otro copropietario tiene derecho a reclamarle que la retire. Pero si este último hizo algo similar y no le reclama al resto, cabe preguntarse si el derecho que ejerce sigue siendo justo. La pregunta que recorre el tema es: \textbf{¿cuándo el ejercicio de un derecho se vuelve abusivo?}
\end{example}

Los temas se encadenan lógicamente: primero se analiza el problema conceptual, luego los criterios teóricos y, finalmente, cómo el derecho argentino fue respondiendo a este desafío.

\section{El dilema del abuso del derecho}

La expresión \textbf{abuso del derecho} plantea, desde el comienzo, una aparente contradicción lógica: ¿cómo es posible que algo que de por sí es justo --- es decir, que es derecho, algo exigible, una potestad conforme a una cierta regla --- se torne abusivo? Este es el gran dilema de una noción que se presenta como contradictoria.

\subsection{La noción de derecho subjetivo como punto de partida}

La posición que se adopte frente al concepto de \textbf{derecho subjetivo} determina si se admite o no la posibilidad del abuso del derecho: según la noción de derecho subjetivo que se tenga, se tendrá una visión del abuso del derecho, admitiéndolo o no admitiéndolo.

\section{Caso práctico: propiedad horizontal}

Para ilustrar el concepto se presenta un caso concreto resuelto por los tribunales argentinos.

\begin{example}{reclamo selectivo en propiedad horizontal}
Un copropietario tiene la potestad de exigir que otro condómino que construyó algo en violación del reglamento de copropiedad lo derribe. Sin embargo, si todos los copropietarios también construyeron en infracción al reglamento, y quien reclama (a) solo le exige a uno y no al resto, y además (b) él mismo tampoco cumplió con el reglamento, los tribunales consideraron que ese ejercicio del derecho era abusivo. La pretensión fue declarada inadmisible.
\end{example}

El interrogante que plantea el caso es hasta qué punto se tiene derecho a exigir a otro que cumpla con el reglamento de copropiedad cuando quien reclama tampoco lo ha cumplido. Si bien existe un derecho a reclamar en una propiedad horizontal --- por ejemplo, contra el condómino que construyó en exceso, incluso invadiendo zonas comunes ---, en este caso se consideró que quien pedía también había realizado actos que violentaban el reglamento, por lo que su reclamo constituía un abuso del derecho.

\section{Críticas a la noción de abuso del derecho}

La noción de abuso del derecho ha recibido críticas doctrinarias, en particular desde la visión que concibe al derecho subjetivo como una potestad absoluta de la voluntad.

% TODO: contenido incompleto o ambiguo en las notas originales (la cita de Planiol figura como "confiada a la mitad de la voluntad")
\textbf{Planiol} sostuvo que, si el derecho es una potestad confiada a la voluntad, la expresión ``abuso del derecho'' choca contra sí misma: no podría haber abuso, porque la persona siempre estaría ejerciendo algo que de por sí es legítimo. Además, si es la voluntad la que determina el contenido de la noción de derecho, es la propia voluntad la que fija dicho contenido.

Esta posición individualista extrema tuvo su recepción en la redacción original del Código Civil argentino, aunque luego fue modificada con la reforma de 1968 y con el Código Civil y Comercial de 2015.

\section{Terminología técnica correcta}

\begin{important}
La terminología exacta es \textbf{``ejercicio abusivo de un derecho''}, y no ``abuso del derecho''. La pregunta correcta es \emph{cuándo el ejercicio de un derecho se torna abusivo}: el derecho en sí mismo es legítimo; lo abusivo es la manera en que se lo ejerce.
\end{important}

\section{Cuadro Cornell: el problema conceptual del abuso}

\begin{longtable}{
    >{\raggedright\arraybackslash}p{0.27\textwidth}
    >{\raggedright\arraybackslash}p{0.63\textwidth}
}
\toprule
\textbf{Preguntas / palabras clave} & \textbf{Notas / respuestas} \\
\midrule
\endfirsthead
\toprule
\textbf{Preguntas / palabras clave} & \textbf{Notas / respuestas} \\
\midrule
\endhead
\bottomrule
\endfoot

\textbf{¿Qué es el abuso del derecho?} &
Es el ejercicio de un derecho subjetivo legítimo de una manera que se torna abusiva. La terminología exacta es ``ejercicio abusivo de un derecho'', porque el derecho en sí mismo es legítimo; lo abusivo es la forma en que se lo ejerce. \\[0.8em]

\textbf{¿Por qué es una noción contradictoria?} &
Porque si un derecho es algo justo y exigible, resulta paradójico que pueda tornarse abusivo. Este es el gran dilema de la noción. \\[0.8em]

\textbf{¿De qué depende que se admita o no el abuso del derecho?} &
De la noción de derecho subjetivo que se adopte: según ella se tendrá una visión del abuso del derecho, admitiéndolo o no admitiéndolo. \\[0.8em]

\textbf{¿Cómo resolvieron los tribunales el caso del reclamo selectivo en propiedad horizontal?} &
Consideraron abusivo el reclamo de quien exigía a un solo condómino la demolición de lo construido en infracción al reglamento, cuando todos los copropietarios habían construido en infracción y quien reclamaba tampoco lo había cumplido. La pretensión fue declarada inadmisible. \\[0.8em]

\textbf{¿Cuál fue la crítica de Planiol?} &
Que si el derecho es una potestad confiada a la voluntad, la expresión ``abuso del derecho'' choca contra sí misma, porque la persona siempre estaría ejerciendo algo legítimo. \\

\midrule
\multicolumn{2}{p{\dimexpr0.90\textwidth+2\tabcolsep\relax}}{%
\textbf{Resumen:} el abuso del derecho plantea la paradoja de que un derecho legítimo pueda ejercerse de modo abusivo. La posibilidad de admitirlo depende de la noción de derecho subjetivo que se adopte, y fue criticado desde la visión de la voluntad como potestad absoluta. Por eso la expresión técnicamente exacta es ``ejercicio abusivo de un derecho''. Un caso de propiedad horizontal muestra cómo los tribunales rechazaron un reclamo selectivo por considerarlo abusivo.
} \\
\end{longtable}

% ============================================================
\chapter{Criterios teóricos para determinar el abuso}
% ============================================================

La doctrina elaboró criterios para identificar cuándo hay ejercicio abusivo de un derecho. Se distingue entre criterios \textbf{subjetivos} y criterios \textbf{objetivos}, y existe además una postura que los combina.

\section{Criterios subjetivos}

Los criterios subjetivos miran a las condiciones del agente que ejerce el derecho. Son tres:

\begin{table}[H]
\centering
\begin{tabularx}{\textwidth}{
    >{\centering\arraybackslash}p{1cm}
    >{\raggedright\arraybackslash}p{4.2cm}
    >{\raggedright\arraybackslash}X
}
\toprule
\textbf{N.°} & \textbf{Criterio} & \textbf{Descripción} \\
\midrule
1.° & \textbf{Intención de dañar} & El ejercicio es abusivo cuando la finalidad del agente es causar un daño. La intención dañosa es el elemento determinante. \\
\addlinespace
2.° & \textbf{Negligencia en el ejercicio} & El ejercicio es abusivo cuando el derecho se ejerce de forma negligente, sin tomar las diligencias que le son propias. \\
\addlinespace
3.° & \textbf{Sin interés alguno} & El ejercicio es abusivo cuando se realiza sin ningún interés detrás; es decir, sin un interés legítimo que lo justifique. \\
\bottomrule
\end{tabularx}
\caption{Criterios subjetivos de determinación del abuso.}
\end{table}

% TODO: contenido incompleto o ambiguo en las notas originales (en la transcripción oral, el segundo criterio subjetivo figura enunciado como "no es abusivo cuando se ejerce de forma negligente"; el cuadro del apunte y la conclusión indican que la negligencia torna abusivo el ejercicio)

\section{Criterios objetivos}

Los criterios objetivos atienden a circunstancias, elementos o factores objetivos, más allá de las disposiciones subjetivas del agente que ejerce el derecho.

\begin{table}[H]
\centering
\begin{tabularx}{\textwidth}{
    >{\centering\arraybackslash}p{1cm}
    >{\raggedright\arraybackslash}p{4.2cm}
    >{\raggedright\arraybackslash}X
}
\toprule
\textbf{N.°} & \textbf{Criterio} & \textbf{Descripción} \\
\midrule
1.° & \textbf{Fin económico y social} & Será abusivo el ejercicio que va en contra del fin económico y social que tiene ese derecho. \\
\addlinespace
2.° & \textbf{Fin de su institución} & Es abusivo el ejercicio contrario a los fines para los cuales fue instituido el derecho. Este criterio fue receptado por la Ley 17.711. \\
\addlinespace
3.° & \textbf{Buena fe y buenas costumbres} & Criterio de Borda: es abusivo el ejercicio que excede los límites de la moral, la buena fe y las buenas costumbres. Ejemplo: intereses usurarios. \\
\bottomrule
\end{tabularx}
\caption{Criterios objetivos de determinación del abuso.}
\end{table}

% TODO: contenido incompleto o ambiguo en las notas originales (en la transcripción oral, el criterio de Borda se enuncia con la expresión "buena fe y las malas costumbres")
\begin{example}{intereses usurarios}
El caso típico y habitual se da cuando se reclaman intereses que llegan a ser usurarios. La \textbf{usura} es una tasa de interés desmedida: se acepta una tasa de interés razonable, pero si esta se convierte en desmedida, se están excediendo los límites de la moral. Se trata de un exceso respecto de la moral; no de una contrariedad con la buena fe y las buenas costumbres, pues de ser así no sería derecho.
\end{example}

\section{Criterio mixto}

Existe también una postura que combina ambos tipos de criterios: será el juez, en cada caso, quien mediante una ponderación de los elementos objetivos y subjetivos evalúe si el ejercicio del derecho ha sido abusivo.

\section{Cuadro Cornell: criterios para determinar el abuso}

\begin{longtable}{
    >{\raggedright\arraybackslash}p{0.27\textwidth}
    >{\raggedright\arraybackslash}p{0.63\textwidth}
}
\toprule
\textbf{Preguntas / palabras clave} & \textbf{Notas / respuestas} \\
\midrule
\endfirsthead
\toprule
\textbf{Preguntas / palabras clave} & \textbf{Notas / respuestas} \\
\midrule
\endhead
\bottomrule
\endfoot

\textbf{¿Qué miran los criterios subjetivos y cuáles son?} &
Miran las condiciones del agente que ejerce el derecho. Son: 1.° intención de dañar (finalidad de causar un perjuicio); 2.° negligencia en el ejercicio (no se toman las diligencias debidas); 3.° ausencia de todo interés legítimo que justifique el ejercicio. \\[0.8em]

\textbf{¿Qué miran los criterios objetivos y cuáles son?} &
Miran circunstancias o factores objetivos, más allá de las disposiciones subjetivas del agente. Son: 1.° contrariedad con el fin económico y social del derecho; 2.° contrariedad con los fines de su institución (receptado por la Ley 17.711); 3.° exceso sobre los límites de la moral, la buena fe y las buenas costumbres (criterio de Borda; ejemplo: intereses usurarios). \\[0.8em]

\textbf{¿Qué es la usura?} &
Una tasa de interés desmedida. Se acepta una tasa razonable, pero cuando esta se vuelve desmedida se exceden los límites de la moral. \\[0.8em]

\textbf{¿Cuál es el criterio mixto?} &
El juez pondera en cada caso concreto tanto los elementos subjetivos como los objetivos para determinar si el ejercicio fue abusivo. \\

\midrule
\multicolumn{2}{p{\dimexpr0.90\textwidth+2\tabcolsep\relax}}{%
\textbf{Resumen:} los criterios subjetivos atienden a la situación del agente (intención de dañar, negligencia, ausencia de interés); los objetivos, a factores externos a él (fin económico y social, fin de la institución, límites de la moral, la buena fe y las buenas costumbres). El criterio mixto encomienda al juez la ponderación de ambos en cada caso. Los criterios objetivos son los que terminaron siendo receptados por el derecho positivo argentino.
} \\
\end{longtable}

% ============================================================
\chapter{El abuso del derecho en el derecho argentino}
% ============================================================

Este capítulo analiza la evolución histórica del instituto en el ordenamiento jurídico argentino.

\section{El Código Civil de Vélez Sársfield (1871)}

El Código Civil de 1871 está imbuido de una filosofía liberal e individualista propia de la época. Vélez Sársfield entiende el derecho subjetivo como una potestad que el orden jurídico confiere para alcanzar la voluntad y, por lo tanto, \textbf{rechaza la noción de abuso del derecho}.

% TODO: contenido incompleto o ambiguo en las notas originales (el apunte atribuye la redacción originaria a "Decreto-Ley 340/1869")
\begin{formula}{art. 1071, redacción original del Código Civil (1871)}
``El ejercicio de un derecho propio o el cumplimiento de una obligación legal no puede constituirse como ilícito.''

\smallskip
\textit{Redacción originaria según Decreto-Ley 340/1869, en vigencia desde 1871. Esta norma rechazaba expresamente la posibilidad del abuso del derecho: ejercer un derecho nunca podía configurar un acto ilícito.}
\end{formula}

Según esta norma no podía haber abuso del derecho, porque ejercer un derecho no es un ilícito. Vélez Sársfield se aparta así de toda teoría del abuso del derecho y no recepta esa idea en ningún artículo.

\section{La jurisprudencia anterior a la reforma de 1968}

Antes de la Ley 17.711, los tribunales comenzaron a reconocer el instituto por vía pretoriana. A través de distintos fallos --- entre ellos casos de intereses excesivos --- la jurisprudencia entendió que podía haber abuso de derechos. El instituto se mezclaba con otros, pero estaba presente la idea de limitar la visión del derecho subjetivo como algo absoluto.

\section{La Constitución de 1949}

La Constitución de 1949, que luego fue derogada en 1956, constituyó un hito importante en la evolución constitucional del concepto.

\begin{formula}{art. 35, Constitución Nacional de 1949 (derogada en 1956)}
``Los derechos y garantías reconocidos no podrán ser alterados por las leyes que reglamenten su ejercicio; tampoco para amparar a ningún habitante de la nación en perjuicio, detrimento o menoscabo de otros. Los abusos de esos derechos que perjudiquen a la comunidad o lleguen a cualquier forma de explotación del hombre por el hombre configuran delitos que serán castigados por las leyes.''
\end{formula}

Se trata de una norma del campo constitucional, no específicamente del derecho civil, pero con una resonancia indudable en este. Corresponde a un momento de impronta social y de fuerte intervención estatal, que coloca límites a las potestades individuales, sobre todo a la propiedad: esta queda sometida también a una finalidad de bien común y a ciertos límites, de modo que no puede hacerse un uso absoluto del derecho.

\section{El fallo Reina (Corte Suprema, 1956)}

Días antes de la derogación de la Constitución de 1949, la Corte Suprema dictó el fallo \textbf{Reina}, en un caso de abuso del derecho. La Corte sostuvo que, aunque no estuviera en ninguna norma --- previendo que la Constitución de 1949 sería derogada ---, la teoría que ve el ejercicio abusivo del derecho es \textbf{inherente a la noción del derecho}, porque todo derecho es relativo y tiene una finalidad que debe cumplir.

Los derechos encuentran, entonces, un límite, y la teoría del abuso del derecho --- aunque no esté registrada en una norma --- corresponde a la concepción misma del derecho. El fallo buscaba dejar a salvo la institución más allá de la inminente derogación del artículo 35 de la Constitución de 1949.

\section{La Ley 17.711 (1968): reforma del artículo 1071}

En 1968, la Ley 17.711 produjo una reforma integral del Código Civil de Vélez Sársfield, que tocó lugares centrales; el artículo 1071 es uno de ellos.

% TODO: contenido incompleto o ambiguo en las notas originales (la referencia al "artículo referido al 954" figura así en la transcripción)
Esta reforma es parangonada con la modificación, en aquel momento, del artículo referido al 954, que incorporó la imprevisión y la lesión, así como el factor moral y el factor social, como limitación de una visión excesivamente individualista.

\begin{formula}{art. 1071, redacción según Ley 17.711 (1968)}
``El ejercicio regular de un derecho propio o el cumplimiento de una obligación legal no puede constituirse como ilícito. La ley no ampara el ejercicio abusivo de los derechos. Se considerará tal al que contraríe los fines que aquélla tuvo en mira al reconocerlos o al que exceda los límites impuestos por la buena fe, la moral y las buenas costumbres.''
\end{formula}

\subsection{Tres cambios fundamentales}

\begin{table}[H]
\centering
\begin{tabularx}{\textwidth}{
    >{\centering\arraybackslash}p{1cm}
    >{\raggedright\arraybackslash}p{4.2cm}
    >{\raggedright\arraybackslash}X
}
\toprule
\textbf{N.°} & \textbf{Cambio} & \textbf{Explicación} \\
\midrule
1.° & \textbf{Se agrega ``regular''} & El ejercicio de un derecho propio pasa a ser ``el ejercicio regular'' de un derecho propio. Lo regular es conforme a una cierta regla, medida o título. Si el ejercicio es irregular, se acerca al abuso del derecho. \\
\addlinespace
2.° & \textbf{Reconocimiento expreso} & Se agrega la frase ``La ley no ampara el ejercicio abusivo de los derechos''. Es una definición expresa de que no se admitirá el abuso del derecho. \\
\addlinespace
3.° & \textbf{Dos criterios objetivos} & (a) Contrario a los fines que tuvo en mira la ley al reconocer ese derecho (criterio de Josserand). (b) Exceso de los límites impuestos por la buena fe, la moral y las buenas costumbres (criterio de Borda). Se usa ``exceder'' y no ``contrario'' porque, si fuera contrario, habría un problema de compatibilidad con el objeto del acto. \\
\bottomrule
\end{tabularx}
\caption{Cambios introducidos por la Ley 17.711 al artículo 1071.}
\end{table}

\begin{definition}{ejercicio regular}
Es el ejercicio conforme a una cierta regla, medida o título. Si el ejercicio es irregular, se acerca al abuso del derecho.
\end{definition}

La palabra ``regular'' es muy importante: en el nuevo código se la encuentra en muchos otros artículos, como una manera de calificar el ejercicio de los derechos.

\section{El Código Civil y Comercial de la Nación (2015)}

El Código Civil y Comercial del año 2014, con vigencia desde 2015, ratifica la vigencia de la institución del abuso del derecho en el derecho argentino y la amplía.

\begin{formula}{art. 10, Código Civil y Comercial de la Nación (2015)}
``El ejercicio regular de un derecho propio o el cumplimiento de una obligación legal no puede constituirse como ilícito. La ley no ampara el ejercicio abusivo de los derechos. Se considera tal el que contraríe los fines del ordenamiento jurídico o el que exceda los límites impuestos por la buena fe, la moral y las buenas costumbres. El juez debe ordenar lo necesario para evitar los efectos del ejercicio abusivo o de la situación jurídica abusiva y, si correspondiere, procurar la reposición al estado de hecho anterior o fijar una indemnización.''
\end{formula}

\subsection{Tres cambios respecto de la Ley 17.711}

\begin{table}[H]
\centering
\begin{tabularx}{\textwidth}{
    >{\centering\arraybackslash}p{1cm}
    >{\raggedright\arraybackslash}p{4.2cm}
    >{\raggedright\arraybackslash}X
}
\toprule
\textbf{N.°} & \textbf{Cambio} & \textbf{Explicación} \\
\midrule
1.° & \textbf{Ubicación metodológica: Título Preliminar} & Ya no está en la parte de daños, sino en el Título Preliminar (art. 10). Esto le da fuerza expansiva a todo el derecho civil y comercial. \\
\addlinespace
2.° & \textbf{``Fines del ordenamiento jurídico''} & El art. 1071 hablaba de ``los fines que tuvo en mira la ley al reconocer ese derecho''; el art. 10 habla de ``los fines del ordenamiento jurídico''. Es más amplio. \\
\addlinespace
3.° & \textbf{Atribuciones expresas para el juez} & Se explicitan los poderes del juez: debe ordenar lo necesario para evitar los efectos del ejercicio abusivo; procurar la reposición al estado anterior; y fijar una indemnización por los daños si ya ocurrieron. \\
\bottomrule
\end{tabularx}
\caption{Cambios introducidos por el art. 10 del Código Civil y Comercial.}
\end{table}

\section{Otros artículos relacionados del Código Civil y Comercial}

El ejercicio regular y el abuso del derecho aparecen en múltiples artículos del Código, formando un sistema coherente.

\begin{table}[H]
\centering
\begin{tabularx}{\textwidth}{
    >{\raggedright\arraybackslash}p{2.3cm}
    >{\raggedright\arraybackslash}X
}
\toprule
\textbf{Artículo} & \textbf{Contenido / relación con el abuso del derecho} \\
\midrule
Art. 9 & Principio de buena fe. Base axiológica del artículo 10 sobre abuso del derecho. \\
\addlinespace
Art. 10 & Abuso del derecho. Norma principal del instituto, ubicada en el Título Preliminar. \\
\addlinespace
Art. 11 & Abuso de posición dominante. Se aplican los principios de los arts. 9 y 10 a quien abusa de su posición dominante en el mercado. \\
\addlinespace
Art. 14 & Derechos individuales y colectivos. La ley no ampara el ejercicio abusivo de los derechos individuales cuando pueda afectar al ambiente y a los derechos de incidencia colectiva. \\
\addlinespace
Art. 53 & Uso de la expresión ``ejercicio regular'' en materia de derechos de la personalidad. \\
\addlinespace
Art. 235 & Ejercicio regular de extraer aguas y otras disposiciones relativas al dominio público. \\
\addlinespace
Art. 240 & Límites al ejercicio de los derechos individuales sobre bienes en función de los derechos de incidencia colectiva. \\
\addlinespace
Art. 1.718 & Responsabilidad civil. El daño causado por el ejercicio regular de un derecho propio está justificado y no es antijurídico (inciso a). \\
\bottomrule
\end{tabularx}
\caption{Artículos del Código Civil y Comercial relacionados con el ejercicio regular y el abuso del derecho.}
\end{table}

La expresión ``ejercicio regular'' encontró mucho arraigo en el derecho civil y tiene su origen en la reforma de 1968.

\section{Cuadro Cornell: el abuso del derecho en el derecho argentino}

% TODO: contenido incompleto o ambiguo en las notas originales (la referencia al art. 279 / ex art. 953 figura solo en el cuadro Cornell del apunte)
\begin{longtable}{
    >{\raggedright\arraybackslash}p{0.27\textwidth}
    >{\raggedright\arraybackslash}p{0.63\textwidth}
}
\toprule
\textbf{Preguntas / palabras clave} & \textbf{Notas / respuestas} \\
\midrule
\endfirsthead
\toprule
\textbf{Preguntas / palabras clave} & \textbf{Notas / respuestas} \\
\midrule
\endhead
\bottomrule
\endfoot

\textbf{¿Cómo trató el tema el Código de Vélez?} &
El art. 1071 original decía que el ejercicio de un derecho propio no puede constituirse como ilícito. No recepcionó el abuso del derecho: para Vélez Sársfield, ejercer un derecho era siempre legítimo, por su concepción liberal e individualista del derecho subjetivo. \\[0.8em]

\textbf{¿Qué hizo la jurisprudencia antes de 1968?} &
Reconoció el instituto por vía pretoriana en distintos fallos, como casos de intereses excesivos, limitando la visión del derecho subjetivo como algo absoluto. \\[0.8em]

\textbf{¿Qué aportó la Constitución de 1949?} &
Su art. 35 establecía que los abusos de los derechos que perjudiquen a la comunidad o lleguen a cualquier forma de explotación del hombre por el hombre configuran delitos. Aunque fue derogada en 1956, marcó un punto de inflexión constitucional, con límites a las potestades individuales, sobre todo a la propiedad. \\[0.8em]

\textbf{¿Qué dijo la Corte en el fallo Reina (1956)?} &
Que la teoría del ejercicio abusivo del derecho es inherente a la noción del derecho mismo, porque todo derecho es relativo y tiene una finalidad. El instituto existe aunque no esté registrado en una norma. \\[0.8em]

\textbf{¿Cuáles son los 3 cambios de la Ley 17.711 (1968)?} &
1.° Se agrega la palabra ``regular'' al art. 1071. 2.° Se expresa que ``la ley no ampara el ejercicio abusivo de los derechos''. 3.° Se incorporan dos criterios objetivos: contrariedad con los fines que tuvo en mira la ley (Josserand), o exceso de los límites de la buena fe, la moral y las buenas costumbres (Borda). \\[0.8em]

\textbf{¿Por qué se usa ``exceder los límites'' y no ``contrario a''?} &
Porque si el derecho fuera directamente contrario, habría un problema de compatibilidad con el objeto del acto (art. 279 / ex art. 953 CC). ``Exceder los límites'' implica que el derecho en sí es válido, pero su ejercicio sobrepasa las fronteras de lo aceptable. \\[0.8em]

\textbf{¿Cuáles son los 3 cambios del Código Civil y Comercial (2015)?} &
1.° Se ubica en el Título Preliminar (art. 10), con alcance expansivo. 2.° Se amplía el criterio: de ``fines que tuvo la ley'' a ``fines del ordenamiento jurídico''. 3.° Se explicitan las atribuciones del juez: evitar los efectos del ejercicio abusivo, reponer al estado anterior e indemnizar. \\[0.8em]

\textbf{¿Qué otros artículos del CCC se relacionan con el ejercicio regular?} &
Arts. 9 (buena fe), 11 (abuso de posición dominante), 14 (derechos individuales y colectivos), 53, 235, 240 y 1.718 inc. a) (justificación del daño por ejercicio regular). \\

\midrule
\multicolumn{2}{p{\dimexpr0.90\textwidth+2\tabcolsep\relax}}{%
\textbf{Resumen:} el Código de Vélez Sársfield (1871) rechazaba el abuso del derecho al concebir el derecho subjetivo como potestad absoluta de la voluntad. La jurisprudencia lo incorporó por vía pretoriana; la Constitución de 1949 lo elevó al plano constitucional (aunque fue derogada en 1956); el fallo Reina lo consideró inherente a la noción del derecho; la Ley 17.711 (1968) lo positivizó en el art. 1071 con la palabra ``regular'' y dos criterios objetivos; y el Código Civil y Comercial (2015) lo amplió en el art. 10, ubicado en el Título Preliminar, extendiéndolo a los fines del ordenamiento jurídico y dotando al juez de atribuciones expresas para prevenir, reponer e indemnizar.
} \\
\end{longtable}
\partintro{Esta parte recorre la situación jurídica de la persona humana: las restricciones a su capacidad y el sistema de apoyos, los derechos personalísimos que protegen su dignidad, el fin de su existencia por la muerte y las figuras que regulan la incertidumbre sobre su vida (la ausencia y la presunción de fallecimiento).}
\part{El sujeto de la relación jurídica: la persona humana}
\chapter{Marco convencional y antecedentes}
\label{ch:bloque1}

\section{Presentación del tema}

Este capítulo y los tres siguientes reconstruyen el sistema de restricciones a la capacidad de las personas regulado en el Código Civil y Comercial argentino. Se aborda el contexto constitucional que ofrece la Convención sobre los Derechos de las Personas con Discapacidad (CDPD) y su artículo 12, los tres tipos de restricción que regula el código --capacidad restringida, incapacidad e inhabilitación por prodigalidad--, el procedimiento judicial que las rodea --legitimación, medidas cautelares, entrevista personal, dictamen interdisciplinario, sentencia, registro y revisión-- y el sistema de apoyos que reemplaza, como criterio general, a la antigua representación absoluta.

Quien ejerza la abogacía --en un estudio, en un juzgado, en el Ministerio Público o en cualquier área que involucre contratos, sucesiones, salud o familia-- tarde o temprano se encontrará con una persona cuya capacidad está restringida, o con la necesidad de evaluar si conviene promover esa restricción para proteger a alguien. Comprender este sistema permite redactar contratos que no terminen siendo nulos, asesorar correctamente a una familia que enfrenta el deterioro cognitivo de un mayor, intervenir con criterio en un proceso de determinación de la capacidad y, sobre todo, evitar dos errores opuestos: privar de autonomía a quien todavía puede decidir por sí mismo, o desproteger a quien necesita un apoyo real.

Estas reglas tampoco son ajenas a la vida cotidiana de las familias: un padre, una madre, un abuelo o un hermano pueden necesitar, en algún momento, ser acompañados con un apoyo para tomar decisiones. Entender este sistema ayuda a saber qué pedir, qué esperar de un juez y, fundamentalmente, cómo acompañar a esa persona sin sustituir su voluntad.

\begin{note}
El sistema puede organizarse en cuatro bloques que se relacionan entre sí. El primero explica por qué existe una convención internacional que exige no apartar a la persona de la toma de decisiones, sino acompañarla mediante apoyos. El segundo distingue qué tipo de restricción corresponde a cada situación: una restricción limitada para quien conserva gran parte de su autonomía (capacidad restringida), una restricción más intensa para quien no puede en absoluto interaccionar con su entorno (incapacidad), o una limitación específica para quien administra sus bienes de manera perjudicial para su familia (inhabilitación por prodigalidad). El tercero regula el procedimiento judicial que debe seguirse, con las garantías correspondientes, para decidir si corresponde imponer una restricción y de qué alcance. El cuarto explica qué sucede cuando una persona actúa en contra de lo dispuesto por la sentencia, cómo se recupera la plena capacidad y cómo funciona, en concreto, el sistema de apoyos que reemplaza a la antigua representación absoluta.
\end{note}

Al regular la capacidad jurídica se distingue entre capacidad de derecho y capacidad de ejercicio. Dentro de la capacidad de ejercicio, la regla siempre fue la capacidad, pero existen algunos casos de incapaces: las personas por nacer, las personas menores de edad --que van teniendo una capacidad progresiva-- y todo el universo de situaciones en que se presenta alguna discapacidad por razones de salud mental. Es en este último campo donde se ubica el tema desarrollado en esta sección: la Sección Tercera, Capítulo Segundo, Título Primero, Libro Primero del Código Civil y Comercial, que trata las restricciones a la capacidad.

\section{La Convención sobre los Derechos de las Personas con Discapacidad (CDPD)}

Para entender el tema es necesario tener presente el contexto de la Convención sobre los Derechos de las Personas con Discapacidad. Esta convención es del año 2006; la Argentina la ratifica en 2008 mediante la Ley 26.378.

\begin{important}
\textbf{Artículo 75, inciso 22 --- Constitución Nacional.} La jerarquía constitucional de ciertos tratados internacionales de derechos humanos se otorga mediante el procedimiento previsto en esta norma. En 2014, por este procedimiento, la CDPD obtuvo jerarquía constitucional, equiparándose a la primera parte de la Constitución Nacional (declaraciones, derechos y garantías).
\end{important}

En 2014 se le dio jerarquía constitucional: la Convención pasó a tener la misma jerarquía que la primera parte de la Constitución --las declaraciones, derechos y garantías--, incorporándose como uno de los tratados de derechos humanos con jerarquía constitucional.

\subsection{El artículo 12 de la CDPD: igual reconocimiento como persona ante la ley}

El artículo clave de la Convención para este tema es el \textbf{artículo 12}, titulado ``Igual reconocimiento como persona ante la ley''. La Convención parte de reafirmar que las personas con discapacidad son seres humanos que deben gozar de igualdad de condiciones y ser incorporadas e incluidas en la comunidad y en la sociedad, lo que supone superar barreras.

La discapacidad surge, según la Convención, de la interacción entre una deficiencia --que puede ser física, psíquica, sensorial o motriz-- y una barrera que impide la plena incorporación de la persona. Esta idea de barreras que deben eliminarse resulta especialmente clara en la dimensión arquitectónica: las rampas obligatorias en las construcciones o la adaptación de los colectivos con rampa son un modo de superar la barrera que representa, por ejemplo, una escalera.

Junto con la noción de barreras a eliminar, la Convención introduce dos nociones fundamentales: los \textbf{apoyos} y los \textbf{ajustes razonables}.

\begin{definition}{Apoyo (noción convencional)}
El apoyo consiste en brindar la ayuda necesaria para que la persona con discapacidad pueda integrarse por sí misma. Resulta fundamental en materia del ejercicio de la capacidad jurídica.
\end{definition}

\begin{definition}{Ajuste razonable}
Los ajustes razonables son las adaptaciones del procedimiento que permiten la plena integración y participación de la persona. Se encuentran, en el Código Civil y Comercial, en los artículos 35 y 36, referidos al proceso judicial vinculado con la capacidad. Por ejemplo, si se debe tomar una audiencia y la persona es sordomuda, corresponde disponer de un traductor de lenguaje de señas para que pueda entender lo que sucede y participar.
\end{definition}

\begin{important}
\textbf{Artículo 12, inciso 1 --- CDPD.} Las personas con discapacidad tienen derecho al reconocimiento de su personalidad jurídica.
\end{important}

\begin{important}
\textbf{Artículo 12, inciso 2 --- CDPD.} Las personas con discapacidad tienen capacidad jurídica en igualdad de condiciones con los demás en todos los aspectos de la vida.
\end{important}

Sobre el alcance de la expresión ``capacidad jurídica'' en este inciso existen dos posturas: quienes entienden que se refiere solo a la capacidad de derecho (ser titular de derechos) y quienes sostienen que comprende ambas dimensiones --capacidad de derecho y capacidad de ejercicio--. Agustina Palacios y Francisco Bariffi son referentes argentinos en el estudio de la Convención; Palacios sostiene, en un capítulo de libro sobre el tema, que el concepto de capacidad jurídica del artículo 12, inciso 2, comprende tanto la capacidad de derecho como la de ejercicio, es decir, que corresponde partir de un principio de capacidad. El propio Comité de los Derechos de las Personas con Discapacidad, creado por la Convención, sostiene la misma interpretación.

\begin{important}
\textbf{Artículo 12, inciso 3 --- CDPD.} Los Estados que han suscripto esta Convención deben garantizarle a las personas con discapacidad los apoyos necesarios para el ejercicio de su capacidad jurídica.
\end{important}

La palabra ``apoyo'' constituye un concepto clave del sistema del nuevo código. Se trata de un concepto amplio, que involucra un conjunto de medidas: por ejemplo, una persona con discapacidad visual puede utilizar apoyos informáticos o perros guía. En este contexto, la referencia es más específica: apoyos para el ejercicio de la capacidad jurídica, es decir, apoyos para la toma de decisiones.

\begin{example}{Apoyo para la venta de un inmueble}
Una persona con discapacidad mental que es dueña de un inmueble y quiere venderlo puede recibir del juez la designación de un apoyo que la asesore y preste su conformidad, a fin de evitar que alguien se aproveche de la situación y le compre el bien muy por debajo de su valor. La decisión de vender sigue siendo de la persona; el apoyo la asiste en la toma de esa decisión.
\end{example}

\begin{important}
\textbf{Artículo 12, inciso 4 --- CDPD (salvaguardias).} Los Estados deben asegurar que las medidas relativas al ejercicio de la capacidad jurídica respeten los derechos, la voluntad y las preferencias de la persona; que estén libres de conflicto de intereses e influencia indebida; que sean proporcionales y adaptadas a las circunstancias de la persona; que se apliquen en el plazo más corto posible; y que estén sujetas a exámenes periódicos por parte de una autoridad u órgano judicial competente, independiente e imparcial.
\end{important}

Las salvaguardias deben ser adecuadas y efectivas, y la medida que adopte un juez debe ajustarse a la persona concreta. En primer lugar, deben respetarse los derechos, la voluntad y las preferencias de la persona; a diferencia de la Convención de los Derechos del Niño, aquí no aparece la noción de ``interés superior'', sino la de respeto de la voluntad y las preferencias.

\begin{example}{Respeto de la voluntad y las preferencias}
Una persona dueña de un campo prefiere sembrar soja porque le gusta esa tarea. Si el apoyo designado considera más conveniente destinar el campo al ganado por rendir más, de todos modos debe respetar los derechos, la voluntad y las preferencias de la persona titular del campo.
\end{example}

Debe evitarse, además, todo conflicto de intereses o influencia indebida.

\begin{definition}{Conflicto de intereses e influencia indebida}
El conflicto de intereses es un concepto más objetivo: existe, por ejemplo, cuando quien es designado apoyo de una persona es, a la vez, quien le alquila un inmueble que administra o sobre el que debe decidir. La influencia indebida es más difícil de determinar, porque supone cierta manipulación de la voluntad de la persona, como puede ocurrir cuando quien ejerce una influencia decisiva convive con ella.
\end{definition}

\subsection{El Comité de la CDPD y el modelo de apoyo a la decisión}

Una observación general del Comité de Naciones Unidas sobre los Derechos de las Personas con Discapacidad, del año 2014, señala que los Estados deben reemplazar los regímenes basados en la sustitución de decisiones --lo que habitualmente se conoce como representación, en la cual otra persona sucede a la titular en su lugar-- por regímenes basados en el apoyo, de modo que la persona pueda decidir por sí misma. Esto es precisamente lo que instrumenta el nuevo código.

En definitiva, existe un delicado equilibrio entre autonomía y protección: la Convención tiende a enfatizar la autonomía, pero también prevé protecciones, de modo que corresponde articular ambos objetivos.

\section{El régimen del Código Civil derogado: incapaces e inhabilitados}

El código anterior, en su artículo 141, se refería a los incapaces. Eran declaradas incapaces de hecho las personas que, por demencia --también llamadas insanas--, padecieran una enfermedad mental que no les permitía administrar su persona ni sus bienes. Existían también los inhabilitados.

En su configuración posterior a la Ley 17.711, es decir, en la configuración inmediatamente anterior a la entrada en vigencia del nuevo código, el ordenamiento contemplaba dos figuras.

\subsection{Figura 1: los incapaces de hecho declarados judicialmente}

También llamados insanos o interdictos, eran personas declaradas incapaces de hecho por la concurrencia de dos factores: uno biológico (la presencia de una enfermedad mental) y uno jurídico (la ineptitud para administrar su persona o sus bienes). Se les designaba un curador y eran, básicamente, incapaces.

\subsection{Figura 2: los inhabilitados (artículo 152 bis)}

El inhabilitado, en cambio, partía de un principio de capacidad y se le nombraba, por sentencia judicial, un curador con funciones delimitadas por esa sentencia. Se preveían tres supuestos:

\begin{itemize}
\item \textbf{Embriaguez habitual o uso de estupefacientes}, cuando exponían al otorgamiento de actos perjudiciales para el patrimonio o la persona --por ejemplo, alguien con alcoholismo que puede aprovecharse y vender bienes a bajo precio o regalarlos.
\item \textbf{Disminución de facultades} que no llegara al supuesto de incapacidad del artículo 141, pero en la cual el juez estimara que podía producirse un daño a la persona o a su patrimonio.
\item \textbf{Prodigalidad}: quienes, por la prodigalidad en los actos de administración y disposición de sus bienes, exponían a su familia a la pérdida del patrimonio. Este supuesto solo procedía si la persona tenía familia y había dilapidado ya parte de sus bienes; es el caso típico de la persona con ludopatía.
\end{itemize}

\begin{note}
En el sistema del código anterior existían, en síntesis, dos grandes esquemas: el incapaz de hecho, totalmente incapaz y sujeto a curador, y el inhabilitado, capaz pero con restricciones puntuales, también sujeto a curador. En el inhabilitado la regla era la capacidad, con una excepción limitada a los tres supuestos mencionados.
\end{note}

\section{La Ley de Salud Mental 26.657 y el artículo 152 ter}

Dos años después de que Argentina aprobara la Convención, en 2010, se sancionó la Ley de Salud Mental (Ley 26.657), que incorporó al código anterior un artículo 152 ter con tres precisiones para el dictado de la sentencia de incapacidad o inhabilitación.

\begin{important}
\textbf{Artículo 152 ter --- Código Civil (texto incorporado por Ley 26.657).} Para declarar la incapacidad o la inhabilitación se requiere un examen de facultativos conformado por evaluaciones interdisciplinarias. La sentencia no podrá extenderse por más de tres años. La sentencia debe especificar las funciones y actos que se limitan.
\end{important}

Esta reforma constituyó un primer intento --parcial y no del todo coherente con el resto del sistema-- de recoger los principios de la Convención, ya que la figura de la incapacidad y su regla general seguían existiendo. Sin embargo, la exigencia de que la sentencia especificara qué funciones y actos se limitaban anticipaba una idea central del nuevo código: que la regla es la capacidad y que corresponde precisar qué es lo que no puede hacerse. Pese a ello, en la práctica muchos jueces continuaron dictando sentencias que declaraban a la persona ``incapaz para todos los actos de la vida civil''.

\section{Cuadro Cornell --- Bloque 1}

\begin{table}[H]
\centering
\begin{tabularx}{\textwidth}{
    >{\raggedright\arraybackslash}p{0.30\textwidth}
    >{\raggedright\arraybackslash}X
}
\cornellhead

\textbf{¿Por qué la CDPD cambió el paradigma de la capacidad jurídica?} &
Porque, a través de su artículo 12, reconoce la personalidad jurídica de las personas con discapacidad y su capacidad jurídica en igualdad de condiciones, reemplazando el modelo de sustitución de decisiones (representación) por un modelo de apoyo a la toma de decisiones, que promueve la autonomía y respeta la voluntad y las preferencias de la persona. \\

\textbf{¿Qué diferencia hay entre un apoyo y un ajuste razonable?} &
El apoyo es la ayuda que permite a la persona con discapacidad ejercer su capacidad jurídica y tomar decisiones por sí misma. El ajuste razonable es la adaptación del procedimiento judicial (por ejemplo, un intérprete de lengua de señas) que garantiza la accesibilidad y la participación efectiva de la persona en el proceso. \\

\textbf{¿Cómo regulaba la capacidad el Código Civil derogado?} &
Mediante dos figuras: los incapaces de hecho (declarados por un factor biológico y uno jurídico, con curador) y los inhabilitados del artículo 152 bis (embriaguez o toxicomanía, disminución de facultades, o prodigalidad), que partían de la capacidad con restricciones puntuales y también tenían curador. \\

\textbf{¿Qué cambios introdujo la Ley 26.657 antes del nuevo código?} &
Incorporó el artículo 152 ter, que exigió examen interdisciplinario para declarar incapacidad o inhabilitación, fijó un plazo máximo de tres años para la sentencia y obligó a especificar las funciones y actos limitados, anticipando el principio de que la regla es la capacidad. \\

\midrule
\multicolumn{2}{p{\textwidth}}{
\textbf{Resumen:} La CDPD, con jerarquía constitucional desde 2014, impone un modelo social de la discapacidad basado en apoyos y ajustes razonables, en reemplazo de la sustitución de la voluntad. El Código Civil derogado, incluso tras la reforma de la Ley 26.657, mantenía la incapacidad y la inhabilitación como ejes del sistema, con curadores en ambos casos; este es el punto de partida que el nuevo Código Civil y Comercial viene a transformar.
} \\
\bottomrule
\end{tabularx}
\end{table}

% =========================================================================
% CAPÍTULO 2
% =========================================================================
\chapter{Tipos de restricciones a la capacidad en el Código Civil y Comercial}
\label{ch:bloque2}

El nuevo Código Civil y Comercial regula el tema a partir de un concepto general de restricciones a la capacidad y de reglas generales aplicables a todos los supuestos.

\section{Reglas generales del artículo 31}

\begin{important}
\textbf{Artículo 31 --- Código Civil y Comercial (reglas generales).} La capacidad de ejercicio se presume, incluso cuando la persona se encuentra internada. Las limitaciones a la capacidad son de carácter excepcional y se imponen siempre en beneficio de la persona. La intervención estatal tiene carácter interdisciplinario, tanto en el tratamiento como en el proceso judicial. La persona tiene derecho a recibir información a través de medios y tecnologías adecuadas. La persona tiene derecho a participar en el proceso judicial con asistencia letrada, que debe ser proporcionada por el Estado si carece de ella. Deben priorizarse las alternativas terapéuticas menos restrictivas de los derechos y libertades.
\end{important}

La regla general es la capacidad: se presume incluso cuando la persona está internada. Por ello, la internación en un establecimiento psiquiátrico no implica necesariamente una restricción a la capacidad; puede haber internación sin restricción, restricción sin internación, o ambas situaciones combinadas.

Las limitaciones a la capacidad son excepcionales y se imponen en beneficio de la persona. En materia de salud mental se parte de que todas las personas son capaces, y la sentencia judicial cumple un rol central porque delimita con precisión qué actos y funciones se limitan.

La intervención estatal debe ser interdisciplinaria. Las disciplinas que habitualmente intervienen son la psiquiatría, la psicología y el trabajo social: la psiquiatría, por tratarse de una especialidad médica, permite diagnosticar y prescribir tratamiento; la psicología aporta la comprensión de los dinamismos psíquicos y el bienestar integral; el trabajo social aporta información sobre el entorno de la persona, dado que la intensidad de la restricción puede variar según cuente o no con una red de contención familiar y social.

\section{Personas con capacidad restringida (artículo 32, primer párrafo)}

\begin{important}
\textbf{Artículo 32, primer párrafo --- Código Civil y Comercial.} El juez puede restringir la capacidad para determinados actos de una persona mayor de trece años que padece una adicción o alteración mental permanente o prolongada de suficiente gravedad, siempre que estime que del ejercicio de su plena capacidad puede resultar un daño a su persona o a sus bienes.
\end{important}

La persona con capacidad restringida constituye el supuesto base del sistema. Se requiere, en primer lugar, sentencia judicial de restricción de la capacidad, en tanto la Convención establece que toda medida que afecte la capacidad de ejercicio debe ser dictada por una autoridad judicial independiente.

\subsection{Requisitos del artículo 32}

\begin{enumerate}
\item \textbf{Ser mayor de trece años.} Entre los 13 y los 18 años la regla es la incapacidad de ejercicio con capacidad progresiva; la sentencia dictada en esa etapa restringirá aquello vinculado con las atribuciones propias de un adolescente. Antes de los 13 años, el artículo 261 del código presume que la persona carece de discernimiento para los actos lícitos, por lo que no es necesaria una restricción judicial.
\item \textbf{Requisito intrínseco o biológico}: presencia de una adicción o de una alteración mental permanente o prolongada, de suficiente gravedad. El código no precisa el término ``adicción''; la doctrina tiende a interpretarlo a la luz del artículo 4 de la Ley de Salud Mental, que define la adicción como el consumo problemático de drogas, legales o ilegales. Una afición intrascendente --como cierta afición a una serie de televisión-- carece de la gravedad necesaria para justificar la limitación de la capacidad. La ``alteración mental permanente o prolongada'', también de suficiente gravedad, refiere a una situación en la que la persona está afectada en su funcionamiento o en sus capacidades mentales, a los fines de un proceso judicial civil, y no debe confundirse con el certificado de discapacidad, que es un concepto distinto.
\item \textbf{Requisito extrínseco}: que del ejercicio de la plena capacidad pueda resultar un daño a la persona o a sus bienes. El código no contempla el daño a terceros; si existiera riesgo de daño a terceros, correspondería otra regulación --la de la responsabilidad civil-- y, eventualmente, un supuesto de internación si el riesgo fuera cierto e inminente.
\end{enumerate}

\subsection{Régimen de las personas con capacidad restringida}

Las personas con capacidad restringida son, como regla, capaces; la sentencia limita únicamente los actos y funciones específicamente indicados, debe ser proporcional y ajustada a la situación concreta de la persona. Se les designa un \textbf{apoyo} --y no un curador--, cuya modalidad de actuación fija el juez: puede tratarse de un apoyo cuyo consentimiento deba concurrir junto con el de la persona, o de un apoyo de consejo. El artículo 101, inciso b, admite también el apoyo con funciones de representación, supuesto que la doctrina considera excepcional y que debe estar especialmente fundado, para no convertirse en una curatela con otro nombre.

\begin{example}{Apoyo con función de asistencia en la venta de un inmueble}
Si a una persona con una adicción se le designa un apoyo con funciones de asistencia para la venta de su inmueble, al momento de celebrarse la escritura concurren el escribano, la persona con capacidad restringida y el apoyo, y ambos firman en el ámbito de su respectiva intervención.
\end{example}

\section{Personas incapaces (artículo 32, último párrafo)}

\begin{important}
\textbf{Artículo 32, último párrafo --- Código Civil y Comercial.} Excepcionalmente, cuando la persona se encuentre absolutamente imposibilitada de interaccionar con su entorno y expresar su voluntad por cualquier modo, medio o formato adecuado, y el sistema de apoyos resulte ineficaz, el juez puede declarar la incapacidad y designar un curador.
\end{important}

La incapacidad constituye un supuesto excepcional, no la regla. Un análisis de las sentencias dictadas por la Cámara Nacional de Apelaciones en lo Civil durante los dos primeros años de vigencia del código mostró que los casos de capacidad restringida eran ampliamente mayoritarios frente a los de incapacidad, que resultaron poco frecuentes.

% TODO: contenido incompleto o ambiguo en las notas originales: no se identifican en el apunte los datos completos (autores, publicación) de la investigación referida sobre la Cámara Nacional de Apelaciones en lo Civil.

Se trata de casos muy severos: la persona no puede en absoluto entablar una relación con su entorno ni expresar su voluntad por ningún modo, medio o formato adecuado, y además el sistema de apoyos resulta ineficaz para asistirla. De no designarse un representante en estos supuestos extremos, la persona quedaría, como un ``muerto civil'', sin posibilidad de que se protejan sus derechos --por ejemplo, para participar en la asamblea de una sociedad anónima o para administrar bienes en caso de encontrarse en coma.

\subsection{Régimen de las personas incapaces}

La sentencia también debe establecer qué actos y funciones se limitan, aunque la limitación será mayor que en la capacidad restringida. Se designa un \textbf{curador} con funciones de representación, a diferencia del supuesto anterior, en el que se designaban apoyos. En el nuevo código, la figura del curador queda restringida a este supuesto excepcional del artículo 32, último párrafo.

\section{Inhabilitados por prodigalidad (artículo 48)}

\begin{important}
\textbf{Artículo 48 --- Código Civil y Comercial.} Pueden ser inhabilitados quienes por la prodigalidad en la gestión de sus bienes expongan a su cónyuge, conviviente o hijos menores de edad o con discapacidad a la pérdida del patrimonio. A estos inhabilitados se les designarán apoyos. Solo podrán otorgar los actos de disposición entre vivos previstos en la sentencia con la conformidad del apoyo. De ser necesario, el juez puede ampliar la nómina de actos que requieren la conformidad del apoyo.
\end{important}

La denominación de ``inhabilitados'' proviene del código anterior, pero en el nuevo código queda restringida exclusivamente al supuesto de prodigalidad; los otros dos supuestos que antes contemplaba el artículo 152 bis --embriaguez o toxicomanía, y disminución de facultades-- pasan a encuadrarse en la capacidad restringida.

\begin{definition}{Pródigo}
Pródigo es quien derrocha sus bienes --no debe confundirse con ``prodigio'', que designa a alguien extraordinario en sus capacidades--. La figura no protege solo a la persona pródiga, sino también a su familia frente al riesgo de pérdida del patrimonio.
\end{definition}

A diferencia del régimen anterior, que exigía la dilapidación efectiva de una parte del patrimonio, el artículo 48 no requiere que la dilapidación ya se haya producido: basta con que se ponga en riesgo la posibilidad de dilapidar el patrimonio.

\subsection{Régimen de los inhabilitados}

Las personas inhabilitadas son, como regla, capaces, pero tienen limitada la posibilidad de realizar actos de disposición de bienes --aquellos que impliquen enajenarlos-- y los demás actos que fije la sentencia. Se les designan apoyos.

\begin{example}{El pródigo y la protección del patrimonio familiar}
Ante una persona que dilapida bienes jugando en el casino, la familia puede iniciar la acción de inhabilitación; la sentencia puede disponer, por ejemplo, que para vender su automóvil o sus bienes registrables, o para hipotecar, necesite el consentimiento de su cónyuge. De ese modo se preserva el patrimonio familiar. Si la persona no tuviera familia, no podría ser declarada inhabilitada en estos términos.
\end{example}

\section{Cuadro Cornell --- Bloque 2}

\begin{table}[H]
\centering
\begin{tabularx}{\textwidth}{
    >{\raggedright\arraybackslash}p{0.30\textwidth}
    >{\raggedright\arraybackslash}X
}
\cornellhead

\textbf{¿Cuándo puede un juez restringir la capacidad de una persona?} &
Cuando la persona es mayor de trece años, padece una adicción o una alteración mental permanente o prolongada de suficiente gravedad, y del ejercicio de su plena capacidad puede resultar un daño a su persona o a sus bienes (art. 32, primer párrafo). \\

\textbf{¿Por qué la incapacidad es hoy un supuesto excepcional?} &
Porque solo procede cuando la persona está absolutamente imposibilitada de interaccionar con su entorno, no puede expresar su voluntad por ningún modo, medio o formato adecuado, y el sistema de apoyos resulta ineficaz; en la práctica, los casos de capacidad restringida resultan mucho más frecuentes que los de incapacidad. \\

\textbf{¿Qué protege realmente la inhabilitación por prodigalidad?} &
Protege al cónyuge, conviviente o hijos menores o con discapacidad frente al riesgo de pérdida del patrimonio familiar por la mala gestión de los bienes de la persona pródiga, limitando sus actos de disposición y exigiendo la conformidad de un apoyo. \\

\midrule
\multicolumn{2}{p{\textwidth}}{
\textbf{Resumen:} El Código Civil y Comercial distingue tres figuras según la intensidad de la limitación: la capacidad restringida (la regla, con apoyos), la incapacidad (excepcional, con curador) y la inhabilitación por prodigalidad (protección patrimonial familiar, con apoyos). En todos los casos rige el principio general de capacidad del artículo 31.
} \\
\bottomrule
\end{tabularx}
\end{table}

% =========================================================================
% CAPÍTULO 3
% =========================================================================
\chapter{Aspectos procesales}
\label{ch:bloque3}

El Código Civil y Comercial incorpora normas de tipo procesal. Aunque se trata de un código de fondo --y los códigos de procedimiento corresponden a las provincias y a la Nación, cada una para su propio fuero--, el código incluye normas procesales que apuntan a dar una regulación común a una materia que afecta directamente la capacidad de una persona.

\section{Juez competente y legitimados activos}

\subsection{Juez competente (artículo 36)}

El juez competente para intervenir en un proceso de determinación de la capacidad es el del domicilio o el del lugar de internación de la persona. La jurisprudencia tiende a aplicar un principio de inmediatez: si la persona tiene domicilio en un lugar pero está internada en otro, suele resultar competente el juez más cercano al lugar de la internación.

\subsection{Legitimados activos (artículo 33)}

\begin{important}
\textbf{Artículo 33 --- Código Civil y Comercial (legitimados para solicitar la declaración).} Están legitimados para solicitar la declaración de incapacidad y de capacidad restringida: a) el propio interesado; b) el cónyuge no separado de hecho y el conviviente mientras la convivencia no haya cesado; c) los parientes dentro del cuarto grado y, si fuera por afinidad, dentro del segundo grado; d) el Ministerio Público.
\end{important}

Una novedad del nuevo código es que el propio interesado puede iniciar el proceso; el código anterior no lo contemplaba. Esto cobra sentido, por ejemplo, cuando una persona a la que se le diagnosticó una enfermedad de deterioro cognitivo desea anticiparse y solicitar su propia restricción de la capacidad, proponiendo a alguien de su confianza como apoyo.

También puede solicitarlo el cónyuge, mientras no esté separado de hecho ni divorciado, y el conviviente. En torno al alcance del término ``conviviente'' existe cierta ambigüedad en el código, que la doctrina discute: si abarca a toda persona que convive o solo a quienes tienen una unión convivencial registrada; la interpretación parecería ser amplia.

Los parientes legitimados son los que se encuentran dentro del cuarto grado --quienes tendrían vocación hereditaria en caso de fallecimiento-- y, si el parentesco es por afinidad, hasta el segundo grado. El código anterior admitía como legitimado a cualquier vecino cuando la persona era considerada ``furiosa'' (peligrosa); esa posibilidad fue eliminada. Actualmente, quien no encuadre en los supuestos legitimados debe acudir al Ministerio Público (regulado en el artículo 103 del Código Civil y Comercial), que evaluará la situación y, en su caso, se presentará ante el juez para iniciar el proceso.

\section{Medidas cautelares y entrevista personal}

\subsection{Medidas cautelares (artículo 34)}

\begin{important}
\textbf{Artículo 34 --- Código Civil y Comercial (medidas cautelares).} Durante el proceso, el juez puede ordenar medidas de restricción al ejercicio de la capacidad, con el fin de garantizar los derechos personales y patrimoniales de la persona. El juez debe determinar con precisión para qué actos se disponen estas medidas y designar los apoyos necesarios, o excepcionalmente el curador para el caso de los incapaces.
\end{important}

Durante el trámite del proceso la persona es, en principio, capaz; sin embargo, el juez puede disponer restricciones cautelares, siempre de manera excepcional y justificada, respetando los principios generales del artículo 31 y de la Convención, para atender situaciones de urgencia --por ejemplo, la designación de un apoyo para decisiones referidas a un inmueble o a un fondo de inversión mientras la persona atraviesa un problema de salud mental.

\subsection{Entrevista personal (artículo 35)}

\begin{important}
\textbf{Artículo 35 --- Código Civil y Comercial (entrevista personal).} El juez debe garantizar la inmediatez con el interesado durante el proceso y entrevistarlo personalmente, antes de dictar resolución alguna. Debe asegurar la accesibilidad y los ajustes razonables del procedimiento de acuerdo a la situación de aquel. El Ministerio Público y, al menos, un letrado que preste asistencia al interesado durante el proceso deben estar presentes en las audiencias.
\end{important}

Si bien la entrevista personal ya era una práctica habitual en algunos fueros desde hace décadas, el nuevo código la establece como exigencia general: antes de adoptar una decisión tan significativa, el juez debe conocer personalmente a la persona involucrada, garantizando accesibilidad y ajustes razonables.

\section{La persona como parte y asistencia letrada}

\begin{important}
\textbf{Artículo 36 --- Código Civil y Comercial.} La persona en cuyo interés se lleva adelante el proceso es parte y puede aportar todas las pruebas que hacen a su defensa. Interpuesta la solicitud de declaración de incapacidad o de restricción a la capacidad ante el juez, éste debe nombrar un abogado que lo represente, si la persona afectada no tuviera letrado propio. En el caso de personas incapaces, también es necesaria la presencia del Ministerio Público.
\end{important}

El artículo 36 enfatiza que la persona interesada no puede quedar ausente de su propio proceso, situación que en ocasiones ocurría en el pasado. La persona reviste calidad de parte y tiene derecho a aportar pruebas; si carece de abogado propio, el Estado --normalmente a través del Ministerio Público-- debe proveerle uno, cuya función es específicamente la representación judicial, distinta de la del apoyo o de la del curador.

\section{Dictamen interdisciplinario y contenido de la sentencia}

\subsection{Artículo 37: el dictamen interdisciplinario}

\begin{important}
\textbf{Artículo 37 --- Código Civil y Comercial (dictamen e informe interdisciplinario).} Para la declaración de incapacidad o de capacidad restringida es necesario el dictamen de un equipo interdisciplinario. El dictamen debe contener: 1) diagnóstico y pronóstico; 2) época en que se manifestó la situación; 3) recursos personales, familiares y sociales existentes; 4) régimen aconsejable para la protección y asistencia del interesado, con la mayor autonomía posible. Estos puntos deben reflejarse en la sentencia.
\end{important}

La Convención no exige literalmente un análisis interdisciplinario, pero de su espíritu se desprende que el examen no debe limitarse a lo médico, sino integrar distintos saberes. La interdisciplina implica que las disciplinas involucradas --psiquiatría, psicología y trabajo social-- alcancen un conocimiento común, dentro de los límites de sus respectivas competencias profesionales.

% TODO: contenido incompleto o ambiguo en las notas originales: no queda resuelto en el apunte si el dictamen interdisciplinario debe ser conjunto (un único informe integrado) o secuencial (informes individuales de cada disciplina).

El equipo interdisciplinario debe expedirse sobre el diagnóstico y pronóstico --distinguiendo, por ejemplo, una crisis psiquiátrica aguda de una situación permanente o de un deterioro cognitivo severo--; la época en que se manifestó la situación, relevante para la eventual anulación de actos anteriores conforme al artículo 45; los recursos personales, familiares y sociales existentes, que permiten conocer la red de apoyo natural de la persona; y el régimen aconsejable para la protección, asistencia y promoción de la mayor autonomía posible.

\subsection{Artículo 38: el contenido de la sentencia}

\begin{important}
\textbf{Artículo 38 --- Código Civil y Comercial (alcance de la sentencia).} La sentencia debe determinar la extensión y alcance de la restricción y especificar las funciones y actos que se limitan, procurando que la afectación de la autonomía personal sea la menor posible. Asimismo, designa una o más personas de apoyo o curadores según el caso, y señala las condiciones de validez de los actos específicos sujetos a la restricción indicando si la persona debe actuar con la asistencia o representación del apoyo o curador.
\end{important}

El artículo 38 concreta la personalización que exige la Convención: a diferencia del artículo 152 ter del código anterior, que resultaba desajustado respecto del resto del ordenamiento, en el nuevo código todo el sistema parte de la capacidad, y la sentencia se limita a precisar qué actos y funciones se restringen y con qué modalidad de apoyo.

\section{Registro, revisión, internación y traslados}

\subsection{Registro (artículo 39)}

La sentencia debe registrarse en el Registro de Estado Civil, anotándose al margen del acta de nacimiento. A partir de la inscripción se producen efectos frente a terceros: los actos posteriores a la inscripción que sean contrarios a lo dispuesto por la sentencia son nulos. Cuando se levanta la restricción, corresponde cancelar la anotación mediante oficio al Registro Civil.

Este registro plantea una tensión entre la protección de la intimidad --por tratarse de un dato altamente sensible-- y la necesidad de dar publicidad suficiente, ya que quienes contratan deben poder conocer si su interlocutor tiene alguna restricción a la capacidad.

\subsection{Revisión (artículo 40)}

\begin{important}
\textbf{Artículo 40 --- Código Civil y Comercial (revisión).} La sentencia debe ser revisada por el juez en un plazo no superior a tres años, sobre la base de nuevos dictámenes interdisciplinarios y mediando la audiencia personal con el interesado. Puede ser revisada en cualquier momento a instancia del interesado. El Ministerio Público debe fiscalizar el cumplimiento efectivo de la revisión periódica establecida e instarla, si corresponde.
\end{important}

La revisión puede solicitarse en cualquier momento, pero no puede dejar de producirse por más de tres años. Corresponde un nuevo dictamen interdisciplinario, una nueva audiencia y una nueva sentencia; el Ministerio Público debe fiscalizar que el juez cumpla con esta obligación y, en su caso, instarla.

\subsection{Internación (artículo 41) y traslados (artículo 42)}

\begin{important}
\textbf{Artículo 41 --- Código Civil y Comercial (internación).} La internación sin consentimiento de una persona, tenga o no restringida su capacidad, procede solo si se cumplen los recaudos previstos en la legislación especial y se cumplen los requisitos de esta sección. En particular, la internación involuntaria solo es procedente cuando a criterio del equipo de salud mediare situación de riesgo cierto e inminente para la persona o para terceros.
\end{important}

No es posible internar a una persona contra su voluntad si no existe riesgo cierto e inminente de un daño de entidad para ella o para terceros. Este estándar, proveniente de la Ley de Salud Mental y ratificado por el nuevo código, busca evitar los abusos que en el pasado llevaron a internar a personas por carecer de vivienda o por resultar molestas; la internación involuntaria debe ser breve, acotada, orientada a la recuperación en un contexto de crisis y a la posterior reintegración, con supervisión periódica, debido proceso, control judicial y defensa --entre otros mecanismos, la Unidad de Letrados del artículo 22 de la Ley de Salud Mental, encargada de velar por los derechos de las personas internadas por razones de salud mental.

\section{Cuadro Cornell --- Bloque 3}

\begin{table}[H]
\centering
\begin{tabularx}{\textwidth}{
    >{\raggedright\arraybackslash}p{0.30\textwidth}
    >{\raggedright\arraybackslash}X
}
\cornellhead

\textbf{¿Quién puede iniciar un proceso de restricción de la capacidad?} &
El propio interesado, el cónyuge o conviviente no separado de hecho, los parientes dentro del cuarto grado (segundo grado por afinidad) y el Ministerio Público (art. 33). \\

\textbf{¿Por qué el juez debe entrevistar personalmente al interesado?} &
Porque el artículo 35 exige inmediatez y entrevista personal antes de dictar cualquier resolución, garantizando accesibilidad y ajustes razonables, de modo que el juez conozca directamente a la persona antes de restringir su capacidad. \\

\textbf{¿Qué garantiza que la persona no quede ausente de su propio proceso?} &
El artículo 36, que le reconoce calidad de parte, derecho a aportar pruebas y derecho a contar con un abogado --provisto por el Estado si no tuviera uno propio-- que la asista específicamente en el proceso judicial. \\

\textbf{¿Qué cuatro aspectos debe abordar toda sentencia?} &
Según el artículo 37, el dictamen y la sentencia deben contener: diagnóstico y pronóstico; época en que se manifestó la situación; recursos personales, familiares y sociales existentes; y régimen aconsejable para la protección y promoción de la mayor autonomía posible. \\

\textbf{¿Cuándo es legítimo internar a una persona contra su voluntad?} &
Solo ante riesgo cierto e inminente de daño para la persona o para terceros, fundado en evaluación interdisciplinaria, con control judicial inmediato, asistencia letrada y revisión periódica (art. 41). \\

\midrule
\multicolumn{2}{p{\textwidth}}{
\textbf{Resumen:} El proceso de determinación de la capacidad está atravesado por garantías procesales: legitimación amplia pero acotada, medidas cautelares excepcionales, entrevista personal obligatoria, calidad de parte del interesado con asistencia letrada, dictamen interdisciplinario con contenido mínimo obligatorio, registro con efectos frente a terceros, revisión periódica no mayor a tres años, e internación involuntaria limitada al riesgo cierto e inminente.
} \\
\bottomrule
\end{tabularx}
\end{table}

% =========================================================================
% CAPÍTULO 4
% =========================================================================
\chapter{Validez de los actos y sistema de apoyos}
\label{ch:bloque4}

\section{Nulidad de los actos posteriores y anteriores a la sentencia}

\subsection{Actos posteriores a la inscripción (artículo 44)}

\begin{important}
\textbf{Artículo 44 --- Código Civil y Comercial (actos posteriores a la sentencia).} Son nulos los actos de la persona incapaz o con capacidad restringida que contrarían lo dispuesto en la sentencia realizada con posterioridad a su inscripción en el Registro de Estado Civil y Capacidad de las Personas.
\end{important}

Si una persona sujeta a una sentencia que restringe su capacidad realiza un acto prohibido por esa sentencia, el acto es nulo. La nulidad es la sanción civil que deja sin efecto el acto y retrotrae la situación al estado anterior: si se entregó un inmueble, debe restituirse; si se entregó dinero, debe restituirse.

\begin{example}{Venta de un inmueble contraria a la sentencia}
Una persona con discapacidad mental que vive de la renta de varias cocheras tiene una sentencia que exige el apoyo de un hermano o hermana para vender inmuebles. Si, sin ese apoyo, vende un inmueble y el escribano no advierte la restricción, el acto es nulo: la venta no produce efectos y las partes deben restituirse lo entregado.
\end{example}

Se trata de una nulidad relativa. Para que proceda, el acto debe ser posterior a la inscripción de la sentencia y contrario a lo que ella dispone. Si la sentencia solo restringe, por ejemplo, la venta de inmuebles, el consentimiento que la misma persona preste para una intervención quirúrgica --materia no alcanzada por la sentencia-- es perfectamente válido.

\subsection{Actos anteriores a la sentencia (artículo 45)}

\begin{important}
\textbf{Artículo 45 --- Código Civil y Comercial (actos anteriores a la sentencia).} Los actos anteriores a la inscripción de la sentencia pueden ser declarados nulos, si perjudican a la persona incapaz o con capacidad restringida, y se cumple alguno de los siguientes requisitos: a) la enfermedad mental era ostensible al momento de la celebración del acto; b) quien contrató con ella era de mala fe; c) el acto fue a título gratuito.
\end{important}

Estos actos son, en principio, válidos, pero pueden anularse si concurre alguno de los requisitos del artículo 45. El dictamen interdisciplinario del artículo 37, al establecer la época en que se manifestó la situación, permite determinar si existían causales al momento de celebrarse el acto anterior.

\begin{example}{Contratante de mala fe}
Un profesional de la salud que atiende a una persona con capacidad restringida y, conociendo su situación, le compra un inmueble a bajo precio, actúa de mala fe. Si, en cambio, el acto fue a título gratuito, no existe perjuicio patrimonial porque la persona no recibió contraprestación alguna.
\end{example}

\section{Cese de las restricciones (artículo 47)}

\begin{important}
\textbf{Artículo 47 --- Código Civil y Comercial (cese de la incapacidad y de las restricciones a la capacidad).} El cese de la incapacidad o de la restricción a la capacidad debe decretarse por el juez que la declaró, previo examen de un equipo interdisciplinario integrado conforme a las pautas del artículo 37, que dictamine sobre el restablecimiento de la persona. Si el restablecimiento es total, el juez deja sin efecto la sentencia. Si el restablecimiento es parcial o el interesado solicita expresamente la revisión, el juez puede ampliar la nómina de actos que la persona puede realizar por sí o con la asistencia o representación de los apoyos designados.
\end{important}

Además de la revisión obligatoria del artículo 40, el interesado puede solicitar en cualquier momento el cese de la restricción, mediante el mismo procedimiento: un nuevo examen interdisciplinario conforme al artículo 37 (diagnóstico, pronóstico, época, recursos personales, familiares y sociales, y régimen aconsejable).

\section{El sistema de apoyos (artículo 43) y conclusiones}

\begin{important}
\textbf{Artículo 43 --- Código Civil y Comercial (concepto de apoyo).} Se entiende por apoyo cualquier medida de carácter judicial o extrajudicial que facilite a la persona la toma de decisiones para dirigir su persona, administrar sus bienes y celebrar actos jurídicos en general. Las medidas de apoyo tienen como función la de promover la autonomía y facilitar la comunicación, la comprensión y la manifestación de voluntad de la persona para el ejercicio de sus derechos. El interesado puede proponer al juez la designación de una o más personas de su confianza para que le presten apoyo. El juez debe evaluar los alcances de la designación y procurar la protección de la persona respecto de eventuales conflictos de intereses o influencia indebida. La resolución debe establecer la condición y la calidad de las medidas de apoyo y, de ser necesario, ser inscripta en el Registro de Estado Civil y Capacidad de las Personas.
\end{important}

El código toma de la Convención la noción de apoyo como cualquier medida judicial o extrajudicial que facilite la toma de decisiones. El apoyo judicial es el que se designa en la sentencia; el apoyo extrajudicial resulta más difícil de precisar, porque toda persona recibe apoyos informales de quienes la rodean, sin que ello implique control judicial. La cuestión se presenta cuando el apoyo adquiere relevancia jurídica.

\begin{example}{Apoyo extrajudicial sin relevancia jurídica}
Durante una situación de aislamiento, un vecino que hace las compras de una persona mayor que no puede salir --recibiendo el dinero, comprando y entregando la mercadería-- presta un apoyo de hecho que no requiere control judicial, en la medida en que no involucre actos como vender un inmueble, comprar un automóvil o hipotecar una propiedad.
\end{example}

\subsection{Salvaguardias: conflicto de intereses e influencia indebida}

Para designar apoyos, el juez debe valorar la eventual existencia de conflicto de intereses e influencia indebida, dos mecanismos de salvaguarda previstos por la Convención para evitar abusos. El conflicto de intereses es más objetivo --por ejemplo, cuando quien administra un inmueble como apoyo es, a la vez, quien lo alquila--; la influencia indebida es más difícil de probar, porque supone cierta manipulación de la voluntad, situación que puede presentarse, por ejemplo, cuando una persona mayor con discapacidad mental convive con uno de sus hijos y los demás hijos consideran que existe manipulación.

\subsection{Criterio de actuación de los apoyos}

El criterio de actuación de los apoyos es el respeto de la voluntad y las preferencias de la persona, principio que surge directamente del artículo 12 de la Convención. El apoyo puede ser propuesto por la propia persona interesada. La posibilidad de que una persona jurídica actúe como apoyo se encuentra en discusión doctrinaria y no consta que exista jurisprudencia que la haya admitido respecto de la institución donde se encuentra internada una persona.

% TODO: contenido incompleto o ambiguo en las notas originales: la posibilidad de que una persona jurídica sea designada apoyo queda planteada como debate abierto, sin resolución expresa en el material.

\subsection{Conclusiones generales}

El nuevo código se adapta a la Convención mediante conceptos y terminología propios de ella --apoyos, ajustes razonables, accesibilidad, autonomía, recursos personales y familiares--. Ratifica el criterio de la Ley de Salud Mental según el cual la regla es la capacidad y las limitaciones son excepcionales, por actos y funciones concretos. El sistema general es el de apoyos; se reconoce ampliamente el derecho de la persona a participar en el proceso y a ser escuchada; la sentencia debe revisarse periódicamente; y se consolida el criterio de intervención interdisciplinaria.

Subsiste, sin embargo, la cuestión de la incapacidad y de los supuestos de sustitución de la voluntad, en tanto se discute si esas formas de representación respetan plenamente la Convención. La doctrina mayoritaria entiende que no la vulneran, por tratarse de situaciones absolutamente excepcionales en las que la persona no puede expresarse de ningún modo y en las que, de no designarse representante, quedaría sin protección jurídica alguna. A ello se suma que la persona conserva el derecho a participar y puede dejar directivas anticipadas --instituto recibido de manera limitada en el artículo 139, que permite a una persona señalar anticipadamente quién podría ser su curador ante la eventual pérdida de conciencia.

\section{Cuadro comparativo: régimen anterior y régimen vigente}

\begin{longtable}{
    >{\raggedright\arraybackslash}p{0.22\textwidth}
    >{\raggedright\arraybackslash}p{0.36\textwidth}
    >{\raggedright\arraybackslash}p{0.36\textwidth}
}
\toprule
\textbf{Situación} & \textbf{Código Civil derogado} & \textbf{Código Civil y Comercial vigente} \\
\midrule
\endfirsthead
\toprule
\textbf{Situación} & \textbf{Código Civil derogado} & \textbf{Código Civil y Comercial vigente} \\
\midrule
\endhead

\textbf{Regla general} & Incapacidad como eje del sistema; el inhabilitado era la excepción acotada a tres casos. & La capacidad es la regla; toda restricción es excepcional y se impone en beneficio de la persona (art. 31). \\[0.4em]

\textbf{Figura de protección} & Curador con función de representación, para incapaces e inhabilitados. & Apoyos para personas con capacidad restringida e inhabilitados; curador solo para incapaces (art. 32, último párrafo). \\[0.4em]

\textbf{Intervención profesional} & Evaluación médica (luego interdisciplinaria, desde el art. 152 ter de 2010). & Dictamen interdisciplinario obligatorio: psiquiatría, psicología y trabajo social (art. 37). \\[0.4em]

\textbf{Duración de la sentencia} & Sin plazo, hasta el art. 152 ter (2010), que fijó un máximo de tres años. & Revisión obligatoria cada tres años como máximo, o antes a pedido del interesado (art. 40). \\[0.4em]

\textbf{Voz de la persona} & No estaba prevista la entrevista personal obligatoria ni la condición de parte. & Entrevista personal obligatoria (art. 35) y la persona es parte del proceso (art. 36). \\[0.4em]

\textbf{Sistema de apoyos} & No existía el sistema de apoyos; solo representación a través del curador. & Sistema de apoyos como regla; respeto de la voluntad y preferencias de la persona (art. 43). \\
\bottomrule
\end{longtable}

% =========================================================================
% CORNELL DE REPASO GENERAL (comprensivo, aportado por el propio apunte)
% =========================================================================

\section{Cuadro Cornell de repaso general}
\label{ch:cornell-general}

Este cuadro de repaso reúne las preguntas y conceptos clave desarrollados a lo largo de los cuatro bloques, pensado para el estudio activo previo a un examen.

\begin{longtable}{
    >{\raggedright\arraybackslash}p{0.30\textwidth}
    >{\raggedright\arraybackslash}p{0.62\textwidth}
}
\toprule
\textbf{Pregunta / concepto clave} & \textbf{Notas / respuesta / explicación} \\
\midrule
\endfirsthead
\toprule
\textbf{Pregunta / concepto clave} & \textbf{Notas / respuesta / explicación} \\
\midrule
\endhead

\textbf{¿Qué jerarquía tiene la CDPD en Argentina?} &
Fue ratificada en 2008 (Ley 26.378) y en 2014 obtuvo jerarquía constitucional mediante el procedimiento del artículo 75, inciso 22, de la Constitución Nacional, equiparándose a la primera parte de la Constitución. \\[0.4em]

\textbf{¿Qué establece el artículo 12 de la CDPD?} &
Reconoce la personalidad jurídica de las personas con discapacidad, su derecho a la capacidad jurídica en igualdad de condiciones, la obligación estatal de garantizar apoyos para el ejercicio de esa capacidad y la exigencia de salvaguardias que respeten la voluntad, evitando conflictos de intereses e influencia indebida. \\[0.4em]

\textbf{¿Cuáles son las reglas generales del artículo 31 del CCyC?} &
La capacidad se presume, incluso en caso de internación; las restricciones son excepcionales y en beneficio de la persona; la intervención es interdisciplinaria; hay derecho a información accesible y a asistencia letrada gratuita; y deben priorizarse las alternativas menos restrictivas. \\[0.4em]

\textbf{¿Cuáles son los tres tipos de restricción a la capacidad del CCyC?} &
Personas con capacidad restringida (art. 32, primer párrafo, la regla), personas incapaces (art. 32, último párrafo, supuesto excepcional) e inhabilitados por prodigalidad (art. 48). \\[0.4em]

\textbf{¿Qué requisitos exige el artículo 32 para restringir la capacidad?} &
Persona mayor de trece años; adicción o alteración mental permanente o prolongada de suficiente gravedad; y que el ejercicio de la plena capacidad pueda causarle un daño a su persona o a sus bienes. \\[0.4em]

\textbf{¿Cuándo corresponde declarar la incapacidad?} &
Solo cuando la persona esté absolutamente imposibilitada de interaccionar con su entorno, no pueda expresar su voluntad por ningún modo, medio o formato adecuado, y el sistema de apoyos resulte ineficaz. Se designa entonces un curador con funciones de representación. \\[0.4em]

\textbf{¿Qué protege la inhabilitación por prodigalidad?} &
Protege al cónyuge, conviviente o hijos menores o con discapacidad frente al riesgo de pérdida del patrimonio por la mala gestión de los bienes; limita los actos de disposición y exige la conformidad del apoyo para ciertos actos. \\[0.4em]

\textbf{¿Quiénes están legitimados para pedir la restricción de la capacidad?} &
El propio interesado, el cónyuge o conviviente no separado de hecho, los parientes dentro del cuarto grado (segundo grado por afinidad) y el Ministerio Público (art. 33). \\[0.4em]

\textbf{¿Qué exige la entrevista personal del artículo 35?} &
Que el juez garantice la inmediatez con el interesado y lo entreviste personalmente antes de dictar cualquier resolución, asegurando accesibilidad y ajustes razonables del procedimiento. \\[0.4em]

\textbf{¿Qué cuatro puntos debe contener el dictamen y reflejar la sentencia según el artículo 37?} &
1) Diagnóstico y pronóstico; 2) época en que se manifestó la situación; 3) recursos personales, familiares y sociales existentes; 4) régimen de protección y promoción de la mayor autonomía posible, sobre la base del dictamen interdisciplinario. \\[0.4em]

\textbf{¿Cuándo es válida una internación involuntaria?} &
Solo ante riesgo cierto e inminente de daño para la persona o para terceros, fundada en evaluación interdisciplinaria, con control judicial inmediato, asistencia letrada y revisión periódica (art. 41). \\[0.4em]

\textbf{¿Qué efectos produce un acto contrario a la sentencia?} &
Si es posterior a la inscripción de la sentencia, el acto es nulo (art. 44); si es anterior, puede anularse si la enfermedad era ostensible, hubo mala fe del contratante, o el acto fue a título gratuito (art. 45). \\[0.4em]

\textbf{¿Qué es un apoyo según el artículo 43?} &
Cualquier medida judicial o extrajudicial que facilite a la persona la toma de decisiones para dirigir su vida, administrar sus bienes y celebrar actos jurídicos, promoviendo su autonomía y respetando su voluntad y preferencias. \\[0.4em]

\textbf{¿Qué diferencia hay entre conflicto de intereses e influencia indebida?} &
El conflicto de intereses es objetivo y visible (el apoyo que administra un inmueble que él mismo alquila); la influencia indebida es más sutil y difícil de probar, porque supone una manipulación de la voluntad de la persona. \\
\midrule
\multicolumn{2}{p{0.94\textwidth}}{
\textbf{Síntesis final:} El Código Civil y Comercial argentino adapta el sistema de restricciones a la capacidad al modelo social de la discapacidad consagrado en la Convención sobre los Derechos de las Personas con Discapacidad (CDPD). Parte de un principio general de capacidad y reserva la restricción para casos excepcionales, debidamente fundados en un dictamen interdisciplinario y resueltos por sentencia judicial personalizada. Reconoce tres figuras --capacidad restringida, incapacidad e inhabilitación por prodigalidad-- que se diferencian por la intensidad de la limitación y por la naturaleza de la figura de protección: el apoyo, que asiste sin sustituir la voluntad, o el curador, reservado a los casos más extremos de imposibilidad absoluta de interacción. Todo el procedimiento está atravesado por garantías procesales --entrevista personal, calidad de parte, asistencia letrada, revisión periódica-- que buscan que la persona nunca quede ausente de las decisiones que la afectan. Finalmente, el sistema de apoyos del artículo 43, con sus salvaguardias contra el conflicto de intereses y la influencia indebida, condensa el espíritu de toda la reforma: acompañar para que la persona pueda decidir por sí misma, y proteger solo en la medida estrictamente necesaria.
} \\
\bottomrule
\end{longtable}
\chapter[Derechos personalísimos]{Derechos personalísimos: noción, caracteres y régimen}
\label{ch:personalisimos}

\section{Introducción}

Los derechos personalísimos son derechos subjetivos fundamentales que protegen la dignidad humana en sus manifestaciones físicas y espirituales. Estos derechos reconocen y garantizan a cada persona el respeto a su vida, integridad corporal, libertad, honor, intimidad e identidad.

En el ejercicio profesional de la abogacía, la comprensión de estos derechos resulta necesaria para:

\begin{itemize}
    \item Defender a clientes cuyos derechos personalísimos han sido vulnerados.
    \item Asesorar sobre el consentimiento informado en procedimientos médicos.
    \item Proteger la intimidad y la privacidad de datos en el entorno digital.
    \item Litigar casos de difamación, invasión de la privacidad e injurias.
    \item Orientar cuestiones relacionadas con la bioética y los experimentos humanos.
    \item Comprender la protección constitucional de la dignidad humana.
\end{itemize}

Estos derechos adquieren especial relevancia en el contexto de la tecnología digital avanzada, donde los datos personales, las imágenes y la información íntima pueden ser capturados, almacenados y difundidos sin consentimiento. El marco legal que se desarrolla a continuación permite proteger tanto los derechos propios como los ajenos frente a estos desafíos contemporáneos.

\section{Alcance y orden de exposición}

Los derechos personalísimos forman un sistema integrado en el que la dignidad humana constituye el principio rector. La exposición comienza con los conceptos generales y los antecedentes, continúa con las características y la clasificación, prosigue con la regulación en el Código Civil y Comercial y se detiene, finalmente, en el derecho a la libertad. Cada bloque se construye sobre el anterior, de forma análoga a las capas que envuelven un núcleo central de protección de la persona humana.

% TODO: verificar consistencia entre fuentes originales. El mapa de contenidos de la fuente anuncia el desarrollo posterior de los bloques de integridad física e integridad espiritual (por ejemplo, vida, integridad corporal, honor, imagen, intimidad e identidad), pero ese desarrollo no figura en los archivos recibidos; solo se desarrollan los fundamentos, las características, la clasificación, la regulación del Código y el derecho a la libertad.

\section{Noción de derechos personalísimos}

Los derechos personalísimos corresponden a la persona humana. Son derechos que pertenecen a todo ser humano por el solo hecho de ser tal.

\begin{definition}{Derechos personalísimos (Tobías)}
Los derechos personalísimos son derechos subjetivos que reconocen y garantizan a la persona humana el goce en su propia entidad, en su interioridad, en todas sus manifestaciones físicas y espirituales.
\end{definition}

Estos derechos no deben confundirse con los atributos de la personalidad, distinción que tampoco realiza el Código Civil y Comercial. Los atributos de la personalidad son cualidades esenciales que permiten que la persona despliegue su personalidad y entable relaciones jurídicas. Los derechos personalísimos, en cambio, son verdaderos derechos subjetivos que protegen dimensiones fundamentales de la persona.

Estos derechos protegen dos dimensiones fundamentales:

\begin{enumerate}
    \item La \textbf{dimensión de la integridad física}: la vida y la integridad corporal.
    \item La \textbf{dimensión espiritual}: la libertad, el honor, la imagen, la intimidad y la identidad de la persona.
\end{enumerate}

\section{Importancia y contexto actual}

A lo largo de la historia siempre ha existido una protección de los valores esenciales de la persona: la vida, la integridad física, la libertad, la imagen y el honor. Sin embargo, especialmente en la última mitad del siglo XX, se dieron circunstancias que motivaron a los juristas y al legislador a sancionar normas y a adoptar tratados internacionales orientados a proteger estas dimensiones de la persona humana.

Han surgido nuevas formas de amenaza que reclaman una respuesta jurídica. En el entorno de las biotecnologías, hoy es posible intervenir sobre el cuerpo humano de una manera mucho más poderosa: el conocimiento de la genética y de la biología ha dado lugar a enormes beneficios, pero también a nuevas formas de intervención que pueden no ser respetuosas de la integridad física.

\begin{note}
A lo largo de la historia se han realizado experimentos sobre personas humanas sin su consentimiento, aprovechando la condición de vulnerabilidad de determinados grupos, lo que ha dado lugar a fuertes reacciones éticas y jurídicas.
\end{note}

Las nuevas tecnologías inciden también en la transmisión de la vida humana, en los trasplantes de órganos y en las intervenciones biomédicas y biotecnológicas en general. Respecto de la investigación en seres humanos, la genética permite hoy intervenciones de gran alcance, como ocurrió en China en el año 2018, cuando se obtuvieron seres genéticamente modificados. Estas posibilidades tecnológicas son enormemente promisorias para el bienestar y el avance de las personas, pero presentan también aspectos problemáticos que requieren una reflexión cuidadosa.

\subsection{Datos personales y privacidad digital}

Otro ámbito relevante es el de la recolección, el uso, el tratamiento, la transmisión y el almacenamiento de los datos personales: datos vinculados con la identidad, con las características personales, con el consumo o con la salud. La cantidad de datos que hoy se almacena es enorme, fenómeno conocido como \textit{big data}: la tendencia a generar grandes colecciones de datos de las personas para elaborar perfiles, realizar seguimientos y dirigir ofertas publicitarias.

Los datos personales permiten tomar decisiones más racionales y conocer mejor a las personas, lo cual puede traer beneficios, pero también nuevas posibilidades de discriminación y de ataques a las personas en función de sus datos. Cuando se trata de datos sensibles existen nuevas posibilidades de manipulación; así, en determinadas situaciones políticas, las campañas se han dirigido a públicos específicos para persuadirlos e influir en su voto.

\begin{important}
El desarrollo de la informática ha transformado radicalmente la capacidad de procesamiento de datos personales —antes limitada a fichas escritas a mano— lo que exige nuevas formas de protección de la intimidad, la identidad y el honor de las personas. La posibilidad actual de captar la voz y la imagen de manera mucho más poderosa hace que estos elementos, una vez registrados, se conviertan también en datos personales sujetos a este mismo régimen de protección.
\end{important}

\section{Antecedentes en Argentina}

Entre los autores que en la Argentina han desarrollado el estudio de esta materia se encuentran Salvat, Cifuentes, Rivera, Navarro Floria y Tobías, entre otros. Cabe señalar que este tema desborda el contenido propio de la parte general del derecho civil y podría constituir por sí mismo una materia completa.

% TODO: contenido incompleto o ambiguo en las notas originales — el segundo autor mencionado aparece en el apunte como "el profesor Salto Cifuentes", lo que podría corresponder a dos autores distintos (Salvat y Cifuentes) mal unidos en la transcripción; se conserva la grafía más plausible sin poder confirmarlo con el material disponible.

\subsection{El Código de Vélez}

El Código de Vélez no contenía un tratamiento metódico y orgánico de los derechos personalísimos. No obstante, los autores reconocen la presencia del tema en algunos de sus artículos:

\begin{longtable}{p{0.30\textwidth} p{0.62\textwidth}}
\toprule
\textbf{Disposición} & \textbf{Contenido} \\
\midrule
\endfirsthead
\toprule
\textbf{Disposición} & \textbf{Contenido} \\
\midrule
\endhead
Artículos 428 y 2.126 & Hablan de derechos inherentes a la persona. \\
\addlinespace
Artículo 1.075 & Establece que toda afectación de un derecho que se confunda con la existencia de la persona puede ser materia de un delito. \\
\bottomrule
\end{longtable}

También en la nota al artículo 2.312, el Código de Vélez hacía referencia a ciertos derechos que tienen origen en la existencia del individuo mismo, como la libertad, el honor, el cuerpo y la patria potestad. Se trata de una mención presente en el código, aunque no sistemática ni ordenada.

\subsection{Evolución legislativa en Argentina}

\begin{longtable}{p{0.30\textwidth} p{0.62\textwidth}}
\toprule
\textbf{Norma} & \textbf{Contenido} \\
\midrule
\endfirsthead
\toprule
\textbf{Norma} & \textbf{Contenido} \\
\midrule
\endhead
Ley 11.723 (1933) — Ley de Propiedad Intelectual & Norma que continúa vigente. En sus artículos 31 a 35 regula la fotografía y el derecho a la imagen, incluyendo el consentimiento y los casos en que puede prescindirse de él. \\
\addlinespace
Ley 21.173 — incorporó el artículo 1.071 bis al código anterior & Protegía el derecho a la intimidad. Este artículo sería recogido posteriormente por el nuevo Código Civil y Comercial en su artículo 1.770. \\
% TODO: contenido incompleto o ambiguo en las notas originales — el apunte indica el año 1905 para esta incorporación, lo cual resulta inconsistente con la numeración de la ley citada; se conserva el dato tal como figura en el material original.
\addlinespace
Ley 24.193 (1992) — Ley de Trasplantes & Reemplazada en 2018 por la Ley 27.447. \\
\addlinespace
Ley 25.326 (2000) — Ley de Protección de Datos Personales & Aún vigente, con proyectos de reforma. \\
\addlinespace
Ley 26.529 (2009) — Ley de Derechos del Paciente & Reformada en 2012 por la Ley 26.742. \\
\bottomrule
\end{longtable}

Estas normas reflejan las dimensiones de los derechos personalísimos vinculadas con la integridad física, la integridad espiritual y la libertad.

\subsection{La Constitución Nacional de 1994}

La reforma constitucional de 1994 incorpora, en el artículo 75 inciso 22, diversos tratados de derechos humanos con jerarquía constitucional que consagran derechos correspondientes a la personalidad.

Se produce así una convergencia entre el derecho público y el derecho privado. Aunque los tratados de derechos humanos incorporados a la Constitución corresponden al derecho público, tienen resonancias directas en el derecho privado, en tanto corresponden a la persona humana, que se despliega fundamentalmente en el campo civil.

\begin{note}
La protección de estos derechos tiene también faceta penal: la protección del derecho al honor puede dar lugar al delito de calumnias e injurias, además de a la protección civil. Ambas protecciones suelen coexistir. El presente desarrollo se limita a las dimensiones civiles del asunto y a su incorporación al Código Civil y Comercial.
\end{note}

\section{Caracteres de los derechos personalísimos}

Según Tobías, los derechos personalísimos presentan los siguientes caracteres:

\begin{itemize}
    \item \textbf{Son derechos subjetivos privados:} corresponden a la esfera de las personas, dimensión donde operan el derecho privado y el derecho civil.
    \item \textbf{Son derechos innatos y originarios:} no son derechos concedidos por el legislador, sino reconocidos. Un ser humano, por el solo hecho de ser tal, tiene derecho a la vida, al honor, a la imagen y a la identidad.
    \item \textbf{Son vitalicios y necesarios:} acompañan a la persona durante toda su vida; una persona humana siempre debe tenerlos.
    \item \textbf{Son absolutos (o de exclusión):} tienen un sujeto pasivo universal. Frente a cualquier persona existe la obligación de respetarlos, sin que sea necesario haber celebrado un contrato ni haber tenido relación previa alguna.
    \item \textbf{Son inherentes a la persona:} no pueden separarse de la persona; forman parte de su ser mismo.
    \item \textbf{Son intransmisibles:} no es posible transmitir completamente, por ejemplo, el derecho a la integridad física, ni delegar en otro la toma de decisiones sobre estos derechos.
    \item \textbf{Son relativamente indisponibles:} en principio no se pueden disponer, salvo mediante disposiciones relativas, de carácter restrictivo y siempre con posibilidad de revocación.
    \item \textbf{Son irrenunciables:} no es posible renunciar a estos derechos. La manifestación de ceder o renunciar, por ejemplo, al derecho a la imagen, carece de validez.
\end{itemize}

\begin{important}
Estos caracteres protegen dimensiones fundamentales de la persona y, en última instancia, tienen como bien jurídico común la dignidad humana.
\end{important}

\section{Clasificación tripartita}

Se adopta aquí la clasificación tripartita habitualmente utilizada por el profesor Tobías, sin perjuicio de que otros autores propongan clasificaciones distintas.

\begin{table}[H]
\centering
\begin{tabular}{p{0.28\textwidth} p{0.28\textwidth} p{0.28\textwidth}}
\toprule
\textbf{Integridad física} & \textbf{Libertad} & \textbf{Integridad espiritual} \\
\midrule
Vida \newline Integridad corporal \newline Salud &
Concepto que comprende aspectos tanto corporales como espirituales &
Honor \newline Intimidad \newline Imagen \newline Identidad \\
\bottomrule
\end{tabular}
\end{table}

En torno a esta clasificación se organiza el desarrollo del capítulo, que aborda las disposiciones generales del Código Civil y Comercial y una primera aproximación al derecho a la libertad.

\section{Regulación en el Código Civil y Comercial}

Metodológicamente, el Código Civil y Comercial incorpora los derechos personalísimos en el Capítulo 3, Título Primero del Libro Primero. El Libro Primero —Parte General— se abre con el Título Primero, dedicado a la persona humana; dentro de este título, luego de regular el comienzo de la existencia de la persona y la capacidad de las personas, aparecen los derechos personalísimos en los artículos 51 a 62.

\begin{important}
La incorporación sistemática de los derechos personalísimos constituye una de las grandes novedades del Código Civil y Comercial, en la medida en que el Código de Vélez no traía un tratamiento sistemático de estos derechos. Esta incorporación otorga una mayor presencia e importancia a la dignidad humana en el tratamiento del Código Civil y Comercial.
\end{important}

\subsection{Disposiciones generales: artículo 51}

El capítulo se abre con el artículo 51, referido a la inviolabilidad de la persona humana:

\begin{formula}{Artículo 51 CCyC}
``La persona humana es inviolable y en cualquier circunstancia tiene derecho al reconocimiento y respeto de su dignidad.''
\end{formula}

El artículo no menciona explícitamente el derecho a la vida, como sí lo hacían proyectos anteriores, pero dicha protección puede reconocerse detrás de la noción de inviolabilidad de la persona: no es posible afectar su persona, ni su cuerpo, ni su integridad física, ni su integridad espiritual. La expresión ``reconocimiento'' señala que la dignidad no constituye una concesión del ordenamiento jurídico, sino algo preexistente que el ordenamiento se limita a reconocer.

\subsection{Medios de protección: prevención y reparación}

El artículo 52 establece que la enumeración de medios de protección no es taxativa, sin perjuicio de que exista una sistematización que distingue entre:

\begin{table}[H]
\centering
\begin{tabular}{p{0.30\textwidth} p{0.55\textwidth}}
\toprule
\textbf{Prevención} & \textbf{Reparación} \\
\midrule
Consiste en evitar que ocurra el hecho que menoscabaría la dignidad; se trata de impedir que se produzca el acto dañino. &
Si el hecho ocurriere, se trata de volver al estado anterior y, si ello no fuera posible, de indemnizar los daños producidos. \\
\bottomrule
\end{tabular}
\end{table}

El Código incorpora, además, nuevas formas de reparación, entre ellas la posibilidad de publicar la sentencia en la que se ha discutido la afectación de un derecho personalísimo. Esta forma de reparación, presente ya en artículos anteriores, se incorpora también en la parte de daños del Libro Tercero como posibilidad específica cuando se encuentra afectado el honor de la persona.

\subsection{Disponibilidad relativa: artículo 55}

El artículo 55 dispone que los derechos personalísimos no son disponibles, salvo mediante el consentimiento de la persona, consentimiento que se encuentra sujeto a una serie de límites.

\begin{formula}{Artículo 55 CCyC — Disponibilidad relativa}
El consentimiento no puede ser contrario a las buenas costumbres. El consentimiento no se presume; es de interpretación restrictiva y, ante la duda, se interpreta que no ha sido otorgado. Es libremente revocable en cualquier momento.
\end{formula}

\begin{important}
El consentimiento sobre los derechos personalísimos debe reunir las siguientes condiciones:
\begin{itemize}
    \item No puede ser contrario a la ley, la moral o las buenas costumbres.
    \item Debe ser expresado; no puede presumirse a partir de la mera falta de oposición. Puede manifestarse por escrito, en forma verbal o mediante un gesto comprensible, pero siempre debe existir una manifestación real.
    \item Es de interpretación restrictiva: ante la duda, se interpreta que no hubo consentimiento. Así, el consentimiento otorgado para el uso de la propia imagen en una campaña de bien común no se extiende, por interpretación restrictiva, a una campaña con fines comerciales.
    \item Es libremente revocable: la persona puede dar marcha atrás y revocar el consentimiento oportunamente otorgado.
\end{itemize}
\end{important}

\section{Derecho a la libertad}

El valor de la libertad constituye uno de los derechos personalísimos más importantes. Su fuente se ubica generalmente en el artículo 19 de la Constitución Nacional.

\begin{formula}{Artículo 19 de la Constitución Nacional}
``Las acciones privadas de los hombres que de ningún modo ofendan a la moral pública, ni perjudiquen a un tercero, están exentas de la autoridad de los magistrados. Ningún habitante de la Nación será obligado a hacer lo que la ley no manda, ni privado de lo que ella no prohíbe.''
\end{formula}

Este constituye un principio de libertad fundamental en el ordenamiento jurídico argentino. Desde la perspectiva civil, la protección de la libertad se vincula sobre todo con los supuestos en que una persona sufre violencia para celebrar un acto; en tal caso, la persona afectada tiene derecho a reclamar los daños y perjuicios correspondientes.

También existe protección de la libertad en el campo de los derechos constitucionales, a través del habeas corpus, de la acción de amparo y de la acción de inconstitucionalidad, instrumentos que permiten conocer la situación de una persona, su paradero y su situación respecto de su integridad y de su libertad.

\begin{note}
Al igual que en el resto de los derechos personalísimos, las formas de protección de la libertad se articulan a través del artículo 52, mediante la prevención y la reparación.
\end{note}

% ================================================================
% CUADRO CORNELL
% ================================================================
\section{Cuadro Cornell de repaso}

El siguiente cuadro permite repasar los conceptos centrales desarrollados en este capítulo mediante la técnica de recuperación activa propia del método Cornell.

\begin{longtable}{
    >{\raggedright\arraybackslash}p{0.30\textwidth}
    >{\raggedright\arraybackslash}p{0.60\textwidth}
}
\toprule
\textbf{Preguntas / palabras clave} & \textbf{Notas / respuestas} \\
\midrule
\endfirsthead
\toprule
\textbf{Preguntas / palabras clave} & \textbf{Notas / respuestas} \\
\midrule
\endhead

\textbf{¿Qué son los derechos personalísimos?} &
Son derechos subjetivos que reconocen y garantizan a la persona humana el goce en su propia entidad, en su interioridad, en todas sus manifestaciones físicas y espirituales. Protegen dimensiones fundamentales del ser humano inherentes a su naturaleza y dignidad. \\
\addlinespace

\textbf{¿Cuál es la diferencia entre derechos personalísimos y atributos de la personalidad?} &
Los atributos de la personalidad son cualidades esenciales que permiten que la persona despliegue su personalidad y entable relaciones jurídicas (como el nombre, el domicilio o la capacidad). Los derechos personalísimos son verdaderos derechos subjetivos que protegen dimensiones fundamentales de la persona (como la vida, el honor o la integridad física). \\
\addlinespace

\textbf{¿Por qué es importante estudiar los derechos personalísimos hoy?} &
Las circunstancias contemporáneas —biotecnologías, big data, privacidad digital— han generado nuevas amenazas a estos derechos que requieren una respuesta jurídica actualizada, orientada a proteger a las personas frente a formas de intromisión y vulneración que no existían en épocas anteriores. \\
\addlinespace

\textbf{¿Cuáles son los caracteres principales de los derechos personalísimos?} &
Son: subjetivos privados; innatos y originarios (reconocidos, no concedidos); vitalicios y necesarios; absolutos (sujeto pasivo universal); inherentes a la persona; intransmisibles; relativamente indisponibles; e irrenunciables. Todos estos caracteres buscan proteger la dignidad humana. \\
\addlinespace

\textbf{¿Cómo se clasifican los derechos personalísimos?} &
En forma tripartita: integridad física (vida, integridad corporal, salud); libertad (concepto que comprende aspectos corporales y espirituales); e integridad espiritual (honor, intimidad, imagen, identidad). \\
\addlinespace

\textbf{¿Dónde se regulan los derechos personalísimos en el Código Civil y Comercial?} &
En el Capítulo 3, Título Primero, Libro Primero (artículos 51 a 62). El artículo 51 abre el capítulo con el principio de inviolabilidad de la persona humana y su derecho al reconocimiento y respeto de su dignidad. \\
\addlinespace

\textbf{¿Cuáles son los medios de protección de los derechos personalísimos?} &
El artículo 52 establece: la prevención, que busca evitar que ocurra el acto dañino, y la reparación, que busca volver al estado anterior o indemnizar. También se prevé la publicación de la sentencia como forma específica de reparación en casos de daño al honor, la intimidad o la identidad. \\
\addlinespace

\textbf{¿Qué significa que los derechos personalísimos son relativamente indisponibles?} &
No pueden ser dispuestos en principio, salvo mediante consentimiento que: no sea contrario a las buenas costumbres; no se presuma (debe ser expresado); sea de interpretación restrictiva; y sea libremente revocable. \\
\addlinespace

\textbf{¿Cómo se protege la libertad en el derecho civil?} &
A través de acciones civiles cuando existe violencia o coacción para celebrar un acto (reclamo de daños y perjuicios). También existen protecciones constitucionales, como el habeas corpus, el amparo y la acción de inconstitucionalidad, más específicas respecto de la libertad física. \\
\addlinespace

\midrule
\multicolumn{2}{p{0.92\textwidth}}{
\textbf{Resumen:} Los derechos personalísimos constituyen un sistema integrado de protección de la dignidad humana en todas sus dimensiones. Fundados en el reconocimiento de que cada persona posee una dignidad inherente, protegen la vida, la integridad física y espiritual, la libertad, el honor, la intimidad y la identidad de todo ser humano. Aunque tienen raíces profundas en la historia del derecho, adquieren una importancia renovada en el contexto contemporáneo, en el que las biotecnologías, la digitalización y el almacenamiento masivo de datos han generado nuevas amenazas a estas esferas fundamentales de la persona. El Código Civil y Comercial supone un avance significativo al sistematizar estos derechos en su Parte General, reconociendo su carácter innato, vitalicio, necesario, absoluto, inherente, intransmisible e irrenunciable. Su protección se articula mediante dos vías principales: la prevención y la reparación.
} \\
\bottomrule
\end{longtable}
\chapter{La muerte: fin de la existencia de la persona humana}
\label{ch:muerte}

\section{Ubicación temática y relevancia práctica}
La regulación civil de la persona humana comienza con el primer momento de su existencia --que en el ordenamiento argentino, en consonancia con las normas constitucionales, se sitúa en el momento de la concepción-- y concluye con el fin de la existencia, es decir, con la muerte. El estudio de esta última etapa presenta dos grandes cuestiones: por un lado, la situación de la muerte propiamente dicha y su acreditación civil; por otro, el instituto de la presunción de fallecimiento, figura jurídica aplicable cuando una persona se encuentra ausente de su domicilio sin que se tengan noticias de ella durante el plazo fijado por la ley. En este capítulo se desarrolla exclusivamente el fin de la existencia por muerte comprobada; la presunción de fallecimiento se estudia en el capítulo~\ref{ch:ausencia}.

\begin{note}
Los distintos apartados de este capítulo se encuentran conectados entre sí. Ante el fallecimiento de una persona es necesario, en primer lugar, comprender qué es la muerte y cómo se comprueba. Cuando no se dispone de un cadáver, corresponde analizar cómo se acredita judicialmente el fallecimiento. Si dos o más personas fallecen simultáneamente, la conmoriencia determina las relaciones sucesorias entre ellas. Finalmente, la muerte produce consecuencias jurídicas concretas sobre el patrimonio y los derechos de la persona fallecida.
\end{note}

El fin de la existencia de la persona humana constituye uno de los pilares del derecho civil: el sistema de herencias, la donación de órganos, los registros civiles, los seguros de vida y las sucesiones dependen de determinar cuándo, cómo y con qué efectos jurídicos se produce la muerte. En el ejercicio profesional del derecho privado resulta indispensable gestionar partidas de defunción, iniciar sucesiones, asesorar a las familias en situaciones de presunción de fallecimiento, promover o impugnar declaraciones judiciales de muerte y conocer los límites de la transmisión hereditaria.

\section{La muerte: concepto, comprobación y normativa}

\subsection{Principio general}

El Código Civil y Comercial de la Nación (en adelante, CCyCN) establece el principio general que rige la finalización de la existencia de la persona humana:

\begin{formula}{Art. 93 CCyCN --- Principio general}
La existencia de la persona humana termina por su muerte.
\end{formula}

El Código Civil anterior contenía una disposición similar, aunque hacía referencia a la ``muerte natural''. El término ``natural'', en la redacción del Código de Vélez Sarsfield, no se oponía a ``provocada'' --como si una persona asesinada no muriera-- sino que aludía a la muerte biológica y, fundamentalmente, procuraba excluir cualquier forma de muerte civil.

\begin{note}
La muerte civil implicaba que el derecho consideraba muerta a una persona por el solo cumplimiento de determinada circunstancia --por ejemplo, una pena, o el ingreso a una comunidad religiosa, especialmente de votos solemnes, situación en la cual la persona ``moría al mundo'' para entrar en religión--. Esta institución nunca tuvo aplicación efectiva en el derecho argentino desde el Código Civil anterior, y el CCyCN simplifica la cuestión disponiendo, sin más, que la existencia de la persona termina con su muerte.
\end{note}

\subsection{¿Qué es la muerte?}

\begin{definition}{Muerte}
Cese irreversible de las funciones corporales; fin de la unidad de cuerpo y alma. El momento exacto en que la muerte ocurre permanece, en gran medida, en el plano del misterio: el derecho civil no define en qué instante preciso se produce, sino que regula el modo en que se \textbf{comprueba} su acaecimiento, comprobación que siempre resulta posterior al hecho biológico en sí.
\end{definition}

El momento exacto de la muerte de una persona resulta, por lo tanto, muy difícil de determinar. Aquello de lo que se ocupa el derecho civil es de su comprobación.

\subsection{Comprobación de la muerte}

\begin{formula}{Art. 94 CCyCN --- Comprobación de la muerte}
La comprobación de la muerte queda sujeta a los estándares médicos aceptados, aplicándose la legislación especial en el caso de ablación de órganos del cadáver.
\end{formula}

La comprobación de la muerte se realiza conforme a estándares médicos que van variando a lo largo del tiempo, a medida que avanza la ciencia, y que permiten detectar el cese irreversible de las funciones corporales --especialmente las cardiorrespiratorias y cerebrales, que son las que comandan el organismo humano.

Se habla de ``comprobación'' precisamente porque esta ocurre con posterioridad al hecho de la muerte. Lo que se obtiene es una certeza moral, no una certeza apodíctica: las certezas sobre la muerte se basan siempre en observaciones empíricas que, en principio, podrían tener alguna falla, pero que --con el acompañamiento de los estándares médicos-- se procuran hacer lo más rigurosas posible, de manera de no dejar duda alguna acerca del fallecimiento de la persona. De allí la confianza depositada en la intervención médica.

\subsection{Instrumentos para acreditar la muerte}

\begin{formula}{Ley 26.413 --- Registro del Estado Civil y Capacidad de las Personas}
Regula la forma en que deben ser expedidos el certificado médico de defunción y la partida de defunción. El certificado médico es expedido por un profesional de la salud siempre sobre la base de la presencia del cadáver. La partida de defunción es el instrumento público en el que consta el asiento en el registro civil, con la fecha, la hora y los demás elementos que rodean a la muerte.
\end{formula}

La normativa distingue con claridad dos instrumentos:

\begin{table}[H]
\centering
\caption{Certificado médico de defunción y partida de defunción}
\begin{tabularx}{\textwidth}{X X}
\toprule
\textbf{Certificado médico de defunción} & \textbf{Partida de defunción} \\
\midrule
Expedido por un profesional de la salud & Instrumento público del Registro Civil \\
Siempre requiere la presencia del cadáver & No requiere la presencia del cadáver (se basa en el certificado) \\
Documento de base para inscribir la muerte & Constancia del asiento en el Registro del Estado Civil \\
Contiene fecha, hora y elementos que rodean la muerte & Regulada también por la Ley 26.413 \\
\bottomrule
\end{tabularx}
\end{table}

\section{Donación de órganos y criterios de muerte encefálica}

\subsection{Consentimiento presunto}

Cuando una persona fallece sin haber expresado su voluntad respecto de la donación de órganos, la ley presume que es dadora.

\begin{formula}{Ley 27.447 (2018) --- Trasplante de Órganos y Tejidos}
Establece que cuando una persona fallece sin haber expresado su voluntad, se presume que es dadora de órganos. La negativa debe ser expresa, realizada ante los organismos señalados en la ley. El artículo 36 dispone que la muerte puede certificarse con la confirmación del cese irreversible de las funciones circulatorias o encefálicas. El artículo 37 remite a un protocolo fijado por el Ministerio de Salud con asesoramiento del INCUCAI (Resolución 716/2019), estableciendo que la certificación debe ser realizada por dos médicos, uno de los cuales debe ser neurólogo o neurocirujano y no puede integrar el equipo de ablación.
\end{formula}

\begin{note}
La Ley 27.447 presume que toda persona es donante de órganos, salvo que haya expresado directamente una negativa ante los organismos señalados en la norma. Esta solución ha generado controversias, tanto por suprimir el rol de la familia en la decisión como por presumir una voluntad que, para algunos, debería quedar estimulada en lugar de presumida.
\end{note}

\subsection{Antecedentes legislativos}

\begin{formula}{Ley 21.541 (1977) --- Primera ley de trasplantes}
Definió el fallecimiento como la disfunción total e irreversible de las funciones encefálicas, verificada por un equipo médico. La reglamentación establecía los signos mínimos que debían presentarse en su totalidad.
\end{formula}

\begin{formula}{Ley 24.193 (1993) --- Ley de trasplantes anterior a 2018}
Establecía que la muerte se constataba por: ausencia irreversible de respuesta cerebral con pérdida absoluta de conciencia; ausencia de respiración espontánea; ausencia de reflejos cefálicos; pupilas fijas no reactivas; e inactividad encefálica comprobada por medios técnicos o instrumentales (electroencefalograma plano y constante durante seis horas). No era necesario comprobar la inactividad encefálica en caso de paro cardiorrespiratorio total e irreversible.
\end{formula}

Esta definición de la Ley 24.193 estuvo vigente entre 1993 y 2018.

\subsection{Criterios actuales y Protocolo 2019}

A partir de 2018, la ley opta por no incorporar en su texto una descripción tan detallada de los signos, remitiéndola a un protocolo. Se dispone que el fallecimiento de una persona puede certificarse con la confirmación del cese irreversible de las funciones circulatorias o encefálicas --cualquiera de los dos supuestos, tanto el paro cardiorrespiratorio como el encefálico.

\begin{important}
La palabra clave del criterio vigente es la \textbf{irreversibilidad}. Si existe la posibilidad de reversión --como ocurre con un paro cardíaco que puede revertirse mediante reanimación-- no hay persona fallecida, sino una persona cuya función circulatoria puede ser terapéuticamente reanimada. El fallecimiento, en sentido propio, es un cese irreversible de las funciones.
\end{important}

Respecto de la certificación, el artículo 37 de la Ley 27.447 remite a un protocolo fijado por el Ministerio de Salud con asesoramiento del INCUCAI, concretado en la Resolución 716/2019. Dicho protocolo puede variar, en atención a la necesidad de actualización técnica de los medios biotecnológicos utilizados para el diagnóstico.

\subsection{Garantías y críticas al régimen vigente}

\begin{important}
Un sector de la doctrina ha criticado el sistema vigente por no establecer en el propio texto legal algunas garantías mínimas --como las que preveía la ley anterior--, delegando en una reglamentación un hecho tan decisivo como la certificación de la muerte. Esta cuestión ha merecido discusión doctrinaria; no obstante, la ley mantiene como criterio decisivo el del cese irreversible de las funciones.
\end{important}

Para la profundización de este debate cabe remitir a los trabajos del Dr.\ José Tobías y del Dr.\ Carlos Muñiz, quienes han escrito ampliamente sobre la determinación del momento de la muerte.

\section{Conmoriencia y prueba supletoria del fallecimiento}

\subsection{Conmoriencia}

\begin{formula}{Art. 95 CCyCN --- Conmoriencia}
Se presume que las personas que fallecen en un desastre común mueren al mismo tiempo. Salvo prueba en contrario, es una presunción iuris tantum.
\end{formula}

\begin{note}
A lo largo del tiempo han existido distintas tesis sobre la preeminencia --es decir, sobre quién moría primero, si los más jóvenes, los más ancianos, o de acuerdo con el sexo--. Los códigos civiles argentinos, ya desde la época de Vélez Sarsfield, establecen que todas las personas que se encuentran en esas circunstancias mueren al mismo tiempo, admitiéndose prueba en contrario.
\end{note}

\begin{example}{Accidente de tránsito}
Fallecen el padre y el hijo en un accidente de tránsito; sobrevive la madre. Se supone que no existen otros herederos.

Si el hijo fallece \textbf{después} del padre: hereda del padre y, al morir el propio hijo, lo hereda la madre.

Si el hijo fallece \textbf{antes} que el padre: el padre, al morir, no tendría herederos directos, y serían sus propios ascendientes quienes heredarían.

\textbf{Solución legal (Art. 95 CCyCN):} se presume que ambos fallecieron al mismo tiempo. No hay transmisión de derechos entre ellos, lo cual evita discusiones sobre quién murió primero y alteraciones en los órdenes sucesorios.
\end{example}

\subsection{Prueba supletoria del fallecimiento}

El segundo supuesto especial que corresponde analizar es el previsto en el artículo 98, segunda parte, del CCyCN, que no debe confundirse con la presunción de fallecimiento.

\begin{formula}{Art. 98 CCyCN --- Prueba supletoria del fallecimiento}
Cuando el cadáver no es hallado o no puede ser identificado, y la desaparición se produjo en circunstancias tales que la muerte debe tenerse por cierta, el juez puede tener por comprobada la muerte y disponer la correspondiente inscripción en el registro.
\end{formula}

Para inscribir una muerte se requiere, en principio, un certificado médico. Sin embargo, cuando el profesional de la salud no dispone de cadáver alguno, o no puede identificarlo, dicha vía resulta imposible. Ante la certeza de que la persona ha fallecido, se habilita entonces un trámite judicial.

El criterio decisivo para optar por este artículo, en lugar de la presunción de fallecimiento, es la certeza:

\begin{table}[H]
\centering
\caption{Prueba supletoria (Art. 98) frente a la presunción de fallecimiento}
\begin{tabularx}{\textwidth}{X X}
\toprule
\textbf{Prueba supletoria --- Art. 98} & \textbf{Presunción de fallecimiento} \\
\midrule
Hay \textbf{certeza} de que la persona falleció & Hay \textbf{duda o incertidumbre} sobre si falleció \\
Cadáver no hallado o no identificado & Persona ausente sin noticias durante el tiempo fijado por ley \\
Procedimiento judicial más breve & Procedimiento judicial mucho más largo \\
El juez ordena la inscripción de la muerte & Se declara presuntamente fallecida a la persona \\
\bottomrule
\end{tabularx}
\end{table}

\subsection{Ejemplos desarrollados}

\begin{example}{Avión que se desintegra en el aire (Art. 98)}
Un avión de Air France que se desintegró en el aire (vuelo AF 447), o el avión de LAPA en el Río Uruguay que se desintegró en el aire: no hubo sobrevivientes, y resulta imposible que los hubiera habido. Ante ese hecho, se aplica el artículo 98 --prueba supletoria-- porque existe certeza de que las personas fallecieron.
\end{example}

\begin{example}{Los Andes (presunción de fallecimiento)}
El caso del avión de la película ``Viven'', cuyos pasajeros --uruguayos-- caen en la cordillera de los Andes: no solo no se los encontraba, sino que tampoco se sabía si estaban muertos o no. No podía afirmarse que esas personas habían fallecido; correspondía, por lo tanto, aplicar la presunción de fallecimiento y esperar los plazos propios de ese instituto. De hecho, fueron encontrados a los tres meses.
\end{example}

\begin{example}{Alpinista en precipicio inaccesible (Art. 98)}
Un alpinista cae en un precipicio de acceso imposible, acompañado por compañeros de montañismo que certifican que la persona llegó a ese punto, tuvo un problema, se resbaló y cayó, resultando imposible que haya sobrevivido e imposible acceder al lugar. Con esos testimonios y las constancias del viaje puede realizarse la prueba supletoria: no corresponde permanecer en la incertidumbre, pues la persona falleció. El juez, teniendo por comprobada la muerte, dispone la inscripción en el registro.
\end{example}

\section{Efectos jurídicos de la muerte}

\subsection{Efectos extrapatrimoniales}

La muerte produce, en el campo extrapatrimonial, la extinción de los atributos de la persona: el nombre y la capacidad. Del domicilio subsiste el domicilio convencional --aquel que se elige para un contrato-- que se transmite a los sucesores conforme al artículo 1024 CCyCN. Se extingue el estado civil, pero subsisten los vínculos a los fines sucesorios: si existía cónyuge, y respecto del padre, la madre y los hijos, el parentesco se mantiene a esos efectos.

\subsection{Protección jurídica post mortem}

Se extinguen los derechos personalísimos, pero subsiste una protección jurídica sobre el honor de la persona fallecida, su intimidad y la inviolabilidad de sus papeles privados. Existe, asimismo, una protección jurídica del cadáver, orientada a que este sea tratado con dignidad.

\begin{note}
El cadáver posee un estatuto especial: no es la persona ni es el cuerpo, pero constituye una prolongación de la persona, y por ello existe un cuidado de esta y de su memoria. La protección de la memoria de los difuntos se vincula también con su intimidad y con la inviolabilidad de sus papeles privados.
\end{note}

\subsection{Efectos patrimoniales}

En el campo patrimonial, los derechos patrimoniales de la persona fallecida se transmiten a sus herederos conforme a las reglas dispuestas por el CCyCN, que pueden ser de dos tipos: testamentarias --cuando la propia persona ha determinado quiénes son sus herederos, existiendo porciones legítimas correspondientes a herederos forzosos sobre las cuales no puede disponerse, y una parte disponible según existan descendientes o solo ascendientes-- o, en su defecto, ab intestato.

\subsection{Derechos no transmisibles}

\begin{formula}{Art. 1024 CCyCN --- Efectos de los contratos: sucesores}
Las obligaciones que son inherentes a las personas no son transmisibles por causa de muerte.
\end{formula}

\begin{formula}{Art. 2280 CCyCN --- Situación de los herederos}
No se transmiten por sucesión los derechos que la ley dispone que no sean transmisibles.
\end{formula}

No se transmiten, por tanto, los derechos que la ley dispone como no transmisibles. Un ejemplo de ello son las obligaciones inherentes a las personas, que no son transmisibles por causa de muerte.

\begin{example}{El pintor famoso}
Se contrata a un pintor famoso para que realice el retrato de un familiar por cinco mil dólares. El pintor fallece. Su hija, también pintora, ofrece realizar el retrato en su lugar. La obligación no se transmite, porque era inherente a la persona --se habían considerado las cualidades específicas del pintor contratado-- y su transmisión resultaría incompatible con la naturaleza de la obligación pactada.
\end{example}

Los derechos reales de usufructo y de habitación constituyen otros ejemplos típicos de derechos que no se transmiten por causa de muerte: por su propia constitución, se extinguen con la muerte del beneficiario. La muerte produce, por regla general, la transmisión de los derechos patrimoniales, con excepción de aquellos que no se transmiten.

\section{Síntesis del capítulo}
El capítulo ha desarrollado las cuestiones relativas al fin de la existencia de la persona humana: la muerte, su comprobación, el caso de la conmoriencia, la prueba supletoria del fallecimiento y los efectos jurídicos de la muerte.

% ------------------------------------------------------------
% CUADRO CORNELL
% ------------------------------------------------------------
\section{Cuadro Cornell --- Repaso de la unidad}
\begin{longtable}{
    >{\raggedright\arraybackslash}p{0.30\textwidth}
    >{\raggedright\arraybackslash}p{0.60\textwidth}
}
\caption{Cuadro Cornell --- Fin de la existencia de la persona humana} \\
\toprule
\textbf{Preguntas / palabras clave} & \textbf{Notas / respuestas} \\
\midrule
\endfirsthead

\toprule
\textbf{Preguntas / palabras clave} & \textbf{Notas / respuestas} \\
\midrule
\endhead

\textbf{¿Cuándo termina la existencia de la persona? (Art. 93)} &
Con la muerte. El CCyCN simplifica la redacción anterior eliminando la referencia a la ``muerte natural'' y la figura de la muerte civil, ya inexistente en el derecho argentino desde el Código de Vélez. \\

\textbf{¿Qué es la muerte?} &
El cese irreversible de las funciones corporales. El momento exacto en que ocurre permanece en el plano del misterio; lo que el derecho regula es su comprobación posterior mediante estándares médicos. \\

\textbf{Art. 94 --- ¿Cómo se comprueba la muerte?} &
La comprobación queda sujeta a los estándares médicos aceptados, que evolucionan con la ciencia. Para la ablación de órganos rige la legislación especial (Ley 27.447 y Protocolo 2019). \\

\textbf{Certificado médico vs.\ partida de defunción} &
El certificado médico es expedido por un profesional de la salud, siempre ante la presencia del cadáver, y constituye el documento base. La partida de defunción es el instrumento público del Registro Civil donde consta la inscripción de la muerte (Ley 26.413). \\

\textbf{Ley 27.447 --- ¿Todos somos donantes de órganos?} &
Sí, por presunción legal. Quien no quiera serlo debe expresar una negativa ante los organismos designados. Se eliminó el rol de la familia que existía bajo la Ley 24.193 (Arts.\ 36 y 37). \\

\textbf{¿Qué exigía la Ley 24.193 (1993--2018) para la muerte encefálica?} &
Ausencia irreversible de respuesta cerebral; ausencia de respiración espontánea; ausencia de reflejos cefálicos; pupilas fijas no reactivas; encefalograma plano y constante durante seis horas. El paro cardiorrespiratorio total e irreversible eximía de la prueba encefálica. \\

\textbf{Ley 27.447 --- ¿cuál es el criterio actual de muerte?} &
Cese irreversible de las funciones circulatorias o encefálicas. Los signos específicos y los plazos se remiten al Protocolo del Ministerio de Salud (Res.\ 716/2019), elaborado con el INCUCAI. Debe certificarlo un equipo de dos médicos, uno de ellos neurólogo o neurocirujano ajeno al equipo de ablación. \\

\textbf{¿Qué es la conmoriencia? (Art. 95)} &
Presunción iuris tantum de que dos personas que fallecen en un desastre común mueren al mismo tiempo. Evita discusiones sobre quién murió primero y posibles alteraciones en los órdenes sucesorios; admite prueba en contrario. \\

\textbf{Prueba supletoria del fallecimiento (Art. 98, segunda parte)} &
Procede cuando el cadáver no es hallado o no puede identificarse, pero existe certeza de que la persona falleció. Un juez --no un médico-- ordena la inscripción de la muerte en el registro. Ejemplos: aviones desintegrados, alpinistas en precipicios inaccesibles. \\

\textbf{Diferencia entre el Art. 98 y la presunción de fallecimiento} &
Art.\ 98: hay certeza de la muerte, el procedimiento es más breve y el juez ordena la inscripción. Presunción de fallecimiento: hay duda o incertidumbre, el procedimiento es mucho más largo y se declara ``presuntamente fallecida'' a la persona. \\

\textbf{Efectos extrapatrimoniales de la muerte} &
Se extinguen el nombre, la capacidad, el estado civil y los derechos personalísimos. Subsisten el domicilio convencional (para contratos), los vínculos familiares a fines sucesorios y la protección jurídica del honor, la intimidad y la inviolabilidad del cadáver. \\

\textbf{Efectos patrimoniales de la muerte} &
Se transmiten a los herederos los derechos patrimoniales, conforme a las reglas testamentarias o ab intestato del CCyCN, respetando las porciones legítimas de los herederos forzosos. \\

\textbf{Derechos NO transmisibles por causa de muerte} &
Las obligaciones inherentes a la persona (intuitu personae) y los derechos reales de usufructo y de habitación (Arts.\ 1024 y 2280 CCyCN). Ejemplo: si se contrata a un pintor famoso y este fallece, su hija no está obligada a cumplir ni tiene derecho a exigir el pago. \\

\midrule
\multicolumn{2}{p{\dimexpr\textwidth-2\tabcolsep}}{
\textbf{Resumen:} El fin de la existencia de la persona humana se produce con la muerte, cuya única definición jurídica en el CCyCN es el cese irreversible de las funciones corporales (Art.\ 93). Como ese momento exacto es de difícil determinación, el derecho regula su comprobación: en el caso ordinario, mediante el certificado médico de defunción --que exige la presencia del cadáver-- y la posterior partida de defunción en el Registro Civil, ambos regulados por la Ley 26.413. Para la ablación de órganos, los criterios son más rigurosos y están regulados por la Ley 27.447 (2018) y el Protocolo del INCUCAI de 2019: basta el cese irreversible de las funciones circulatorias o encefálicas, certificado por dos médicos, uno de los cuales debe ser neurólogo o neurocirujano y ajeno al equipo de trasplante. Bajo esta ley, toda persona es donante presunta salvo negativa expresa. Cuando dos personas fallecen en un desastre común sin poder determinarse quién murió primero, la conmoriencia (Art.\ 95) presume que ambas murieron simultáneamente, evitando distorsiones en los órdenes sucesorios. Si el cadáver no es hallado pero hay certeza de la muerte, el juez puede tenerla por acreditada y ordenar su inscripción mediante la prueba supletoria del Art.\ 98, que se distingue de la presunción de fallecimiento por la ausencia de incertidumbre. Finalmente, la muerte produce efectos extrapatrimoniales --extinción del nombre, la capacidad y el estado civil, con subsistencia de los vínculos familiares a fines sucesorios y protección jurídica post mortem del honor y la intimidad-- y patrimoniales: transmisión del patrimonio a los herederos según las reglas del CCyCN, con excepción de los derechos inherentes a la persona (usufructo, habitación, obligaciones intuitu personae), que se extinguen con la muerte del titular.
} \\
\bottomrule
\end{longtable}
\chapter[Marco conceptual e histórico]{Ausencia y presunción de fallecimiento: marco conceptual e histórico}
\label{ch:ausencia}

El fin de la existencia de la persona admite, además de la muerte comprobada, una vertiente de incertidumbre sobre su vida o su muerte, regulada en el Código Civil y Comercial de la Nación (CCCN) mediante dos institutos diferenciados: la \textbf{ausencia} (arts.~79 a 84) y la \textbf{presunción de fallecimiento} (arts.~85 a 92). Los capítulos de esta sección tratan, sucesivamente, el marco conceptual e histórico, la ausencia, los tres casos de presunción de fallecimiento, el procedimiento y la sentencia, y las garantías y efectos de la reaparición.

\begin{note}
Tarde o temprano resulta habitual encontrarse con un cliente cuyo familiar desapareció: un accidente en ruta, una catástrofe natural, o una persona que simplemente dejó de dar señales de vida. Saber distinguir la ausencia de la presunción de fallecimiento no es un detalle académico: es la diferencia entre proteger el patrimonio familiar desde el primer mes o perderlo por falta de acción.

Estos institutos permiten: (1) asesorar a herederos sobre cuándo y cómo iniciar acciones judiciales; (2) gestionar la curatela de bienes de personas desaparecidas; (3) representar acreedores que necesitan cobrar deudas pendientes; (4) intervenir en procesos sucesorios abiertos tras una declaración de fallecimiento presunto; (5) proteger los derechos del ausente ante una eventual reaparición, gestionando la prenotación registral.

En síntesis, estas reglas existen para que el derecho privado pueda dar una respuesta ordenada, justa y temporalmente acotada a uno de los dramas más dolorosos de la vida real: no saber si un ser querido está vivo o muerto.
\end{note}

\textbf{¿Cómo se relacionan los temas?} Puede pensarse el proceso como una escalada progresiva: si no se tienen noticias de una persona y sus bienes quedan desprotegidos, corresponde solicitar primero la \textbf{declaración de ausencia} (arts.~79 a 84) para que un curador los cuide. Si transcurren tres años sin novedades, puede iniciarse la \textbf{presunción de fallecimiento} (arts.~85 a 92), que abre la sucesión pero con garantías por si la persona reaparece. Todo el proceso protege tanto a quienes esperan como a quien podría volver.

\section{Introducción: ausencia y presunción de fallecimiento}

El fin de la existencia de la persona está regulado en el Código Civil y
Comercial a través de dos figuras: la \textbf{ausencia}, por un lado, y la
\textbf{presunción de fallecimiento}, por el otro.

Ausencia y presunción de fallecimiento son dos figuras distintas. Es
habitual que se confundan bajo la expresión «ausencia con presunción de
fallecimiento», lo cual no resulta exacto: se trata de institutos
diferentes. Por supuesto, la presunción de fallecimiento supone que existe
una ausencia de la persona, pero lo que corresponde estudiar es,
específicamente, por un lado la \textbf{ausencia} y por el otro la
\textbf{presunción de fallecimiento}.

\section{El problema jurídico de fondo}

El problema jurídico de fondo puede formularse como la pregunta: ¿cuál es
la respuesta jurídica ante la ausencia prolongada de una persona de su
domicilio?

Esta cuestión ya encontraba explicación en el contexto de las guerras: la
persona que quedaba cautiva y de la que no se sabía nada. ¿Qué sucedía con
sus bienes? ¿Quién podía tomar medidas conservatorias? También pueden
pensarse situaciones catastróficas más recientes: la desaparición de un
avión, el tsunami que hubo en Tailandia a comienzos del año 2002, un
incendio, o cualquier catástrofe en la que una persona desaparece y no se
sabe si está viva o no.

Se produce una \textbf{incertidumbre}, y esa incertidumbre no puede durar
eternamente, porque para el derecho civil existe una cantidad de
relaciones y situaciones jurídicas que ameritan respuesta y toma de
decisiones.

No existe una institución estatal que actúe de oficio ante la mera falta
de contacto de una persona: este es un \textbf{sistema de derecho civil,
de derecho privado}, en el cual los particulares que tienen algún tipo de
interés en los bienes o en las relaciones jurídicas acuden a un proceso
judicial, ya sea para declarar la ausencia o para declarar la presunción
de fallecimiento.

Esto supone preguntarse: ¿cuándo se considera que una persona ha
fallecido? ¿Cuándo resulta razonable asumirlo así? ¿Quién puede pedirlo?
¿Qué sucede si la persona reaparece? ¿Qué resguardos jurídicos se toman en
el camino por si reaparece? Ese es el tema de fondo que estructura ambos
institutos.

\begin{important}
No debe confundirse la ausencia con la presunción de fallecimiento. Si
bien están ubicadas inmediatamente una al lado de la otra en el Código
Civil y Comercial, y tienen elementos comunes, son institutos distintos.
\end{important}

\section{Antecedentes históricos}

Los antecedentes de esta institución provienen del \textbf{Código de
Vélez Sársfield}, que regulaba la ausencia con presunción de fallecimiento
en los artículos 110 a 625. La doctrina señala con claridad cómo en Vélez
Sársfield se produce una \textbf{confusión de elementos}:

Por un lado, existen antecedentes del \textbf{derecho germánico}, donde
existía claramente el instituto de la declaración de muerte, que Vélez
recibe a través de Freitas. Por el otro lado, el \textbf{derecho
francés}, que no conocía esa declaración de muerte y simplemente regulaba
la ausencia, sin poner fin a la existencia de la persona. Vélez también
reconoce esta figura de la ausencia, de modo que el instituto queda
regulado con elementos mixtos.

Tal como fue regulado por Vélez, la institución finalmente recibió
críticas. En 1984 una ley ómnibus —una ley que reunía muchos temas en su
contenido—, la \textbf{Ley 14.394}, reguló por un lado la ausencia y por
otro lado la presunción de fallecimiento: dos instituciones distintas,
clarificando y resolviendo lo que Vélez había planteado. El régimen de
Vélez, entonces, quedó reemplazado por el sistema de la Ley 14.394.

\begin{important}
\textbf{CCCN — Estructura normativa actual.} El nuevo Código Civil y
Comercial deroga la Ley 14.394 en estas partes, pero sigue sus
lineamientos e incorpora tanto la ausencia como la presunción de
fallecimiento al articulado del Título I, Libro I:
\begin{itemize}
    \item \textbf{Ausencia} — Capítulo 6, Arts. 79-84 (no implica fin de la existencia de la persona).
    \item \textbf{Presunción de Fallecimiento} — Capítulo 7, Arts. 85-92 (implica inscripción en el Registro Civil y apertura de la sucesión).
\end{itemize}
Metodológicamente, el CCCN los incorpora con muy pocos cambios respecto de
la ley anterior.
\end{important}

% ======================================================================
% CAPÍTULO 2
% ======================================================================
\chapter{Ausencia (Arts. 79-84 CCCN)}

\section{Concepto de ausencia}

La ausencia es una institución relativamente sencilla. Si la persona
desapareció de su domicilio sin tenerse noticias de ella y sin dejar
apoderado, o si el apoderado no tiene atribuciones suficientes, puede
asignarse un \textbf{curador a los bienes}.

\begin{definition}{Ausencia}
La ausencia, en el proceso del Art. 79 CCCN, no busca declarar a la
persona presuntamente fallecida. Simplemente busca \textbf{garantizar la
adopción de medidas sobre los bienes} de una persona de la que no se
tienen noticias.
\end{definition}

\begin{example}{¿Qué NO es ausencia?}
Si alguien quedó varado en Qatar durante la pandemia, se tienen noticias
de esa persona: se sabe que está varada y que no puede regresar a su
domicilio. Eso no configura una ausencia. La ausencia es la situación de
la persona que desaparece, no se conecta a internet, no usa su tarjeta de
crédito, no pasa por migraciones, no tiene registro en su correo
electrónico, no tiene ningún tipo de actividad, y por lo tanto nadie de
las personas conocidas tiene noticias de ella.
\end{example}

\section{Legitimación y juez competente}

\begin{formula}{Art. 80 CCCN — Legitimación activa}
El Art. 80 establece quiénes pueden iniciar la acción: el Ministerio
Público y las personas que tengan interés legítimo en los bienes (por
ejemplo, los acreedores o los sucesores).
\end{formula}

Respecto del juez competente:

\begin{table}[H]
\centering
\begin{tabularx}{\textwidth}{>{\raggedright\arraybackslash}X >{\raggedright\arraybackslash}X}
\toprule
\textbf{Situación} & \textbf{Juez competente} \\
\midrule
Persona con domicilio en Argentina & Juez del domicilio \\
Sin domicilio en el país o lugar desconocido & Juez del lugar donde están los bienes que hay que cuidar \\
Bienes en distintas jurisdicciones & El primero que intervenga (principio de prevención) \\
\bottomrule
\end{tabularx}
\caption{Juez competente en el proceso de ausencia}
\end{table}

\section{Procedimiento y curatela}

Se cita al ausente por \textbf{cinco días por edictos}. Los edictos son
los avisos que se publican en el Boletín Oficial y en diarios de
circulación importante, mediante los cuales se convoca a la persona a que
se haga presente. Si no comparece, se le da intervención al
\textbf{defensor oficial}, que interviene en representación del ausente,
y pueden tomarse medidas precautorias para los bienes.

La prueba se orienta a demostrar que no hay noticias de esta persona: se
libran distintos oficios, se solicitan testigos, y se determina si existe
o no ausencia. El Ministerio Público procura defender al presunto
ausente, evaluando la verosimilitud de la situación alegada.

Determinada la ausencia, se nombra un curador. El procedimiento del
curador comprende actos de \textbf{conservación y administración de los
bienes}. No puede realizar actos extraordinarios: no puede vender ni
hipotecar un inmueble, porque en ese caso estaría comprometiendo el
patrimonio. Los actos extraordinarios que resulten de necesidad evidente e
impostergable requieren \textbf{autorización judicial}.

Los frutos —por ejemplo, la renta de un inmueble que se alquila— se
destinan al \textbf{sostenimiento de los familiares del ausente}.

\section{Conclusión de la curatela}

La curatela concluye:

\begin{enumerate}[label=(\alph*)]
    \item Porque reaparece el ausente, personalmente o por apoderado.
    \item Porque se constata su muerte: ya sea porque aparece el cadáver y se hace la inscripción del fallecimiento y la correspondiente partida, o bien porque se declara el fallecimiento presunto judicialmente, instituto que se analiza a continuación.
\end{enumerate}

% ======================================================================
% CAPÍTULO 3
% ======================================================================
\chapter{Presunción de fallecimiento (Arts. 85-92 CCCN)}

El sistema del Código es igual al sistema de la Ley 14.394 y se estructura
en torno a \textbf{tres casos}: un caso ordinario y dos casos
extraordinarios —el extraordinario genérico y el extraordinario
específico—.

\section{Caso ordinario (Art. 85 CCCN)}

\begin{formula}{Art. 85 CCCN — Caso ordinario}
La persona que se ausenta de su domicilio sin que se tengan noticias de
ella por el término de \textbf{tres (3) años}. Los elementos son: la
ausencia sin noticias durante tres años. No hay ningún elemento
particular adicional.
\end{formula}

El plazo es de \textbf{tres años}. Esto se vincula con situaciones como,
por ejemplo, una persona que desaparece —quizás víctima de un homicidio
no resuelto— y de la que no se sabe bien qué pasó. Al cabo de tres años
de incertidumbre, los familiares pueden solicitar su fallecimiento
presunto.

No es obligatorio: el cumplimiento del plazo de tres años no impone
iniciar el proceso de manera obligatoria. Esto se relaciona con la
certeza de las relaciones jurídicas. Dado que en la mayoría de los casos
se trata de un familiar y todo está rodeado de afectos y conmociones, hay
quienes siguen buscando y pueden dejar pasar cinco o seis años. Lo que
debe demostrarse en el proceso es que los supuestos legales se han
cumplido.

\begin{important}
\textbf{La última noticia.} El plazo de tres años —o el que corresponda
según el caso— se cuenta desde la \textbf{última noticia} que se tuvo de
la persona. Por ejemplo: si alguien cruzó por migraciones en una fecha
determinada por el paso fronterizo de Paso de los Libres, esa es la
última noticia que se tuvo. A partir de allí se cuentan los plazos.
\end{important}

\section{Caso extraordinario genérico (Art. 86, 1.\textsuperscript{er} párr.)}

El caso extraordinario se denomina así porque acontece algún hecho que
marca una \textbf{mayor probabilidad} de que haya habido un deceso de la
persona. El caso genérico comprende situaciones de incendio, terremoto,
acción de guerra u \textit{otro suceso semejante}, o bien la
participación de la persona en una actividad que implique el mismo
riesgo.

\begin{formula}{Art. 86, 1.\textsuperscript{er} párr. — Caso extraordinario genérico}
Si la persona hubiese estado en el lugar de un incendio, terremoto,
acción de guerra u otro suceso semejante susceptible de ocasionar la
muerte, o hubiese participado de una actividad que implique el mismo
riesgo, y no se tuviese noticias de ella durante \textbf{dos (2) años}
desde el día en que el suceso ocurrió o pudo haber ocurrido.
\end{formula}

\begin{example}{Casos de extraordinario genérico}
\begin{itemize}
    \item El alpinista que estaba escalando el Aconcagua o el Everest y dejó de dar señales.
    \item El tsunami de Tailandia.
    \item Un incendio, un terremoto, una acción bélica.
    \item Cualquier suceso que resulte muy probable que haya provocado la muerte.
\end{itemize}
\end{example}

En este caso, el plazo se adelanta a \textbf{dos años}, contados desde el
día en que el suceso ocurrió o pudo haber ocurrido. Por ejemplo, si se
trata de un alpinista y el último día que consta en un registro que
estaba subiendo fue el 10 de mayo de 2020, desde allí se cuentan los dos
años.

\section{Caso extraordinario específico (Art. 86, 2.\textsuperscript{do} párr.)}

El caso extraordinario específico es el supuesto más detallado por el
Código: el de un \textbf{buque naufragado o aeronave perdida}. Aquí el
acontecimiento resulta mucho más probable como causa de muerte, y el
plazo es de \textbf{seis meses} desde que el suceso ocurrió o pudo
ocurrir.

\begin{formula}{Art. 86, 2.\textsuperscript{do} párr. — Caso extraordinario específico}
Si el presunto fallecido estuvo en un buque naufragado o aeronave
siniestrada o perdida, y no se tuviese noticias de su suerte durante
\textbf{seis (6) meses} desde el día en que el suceso ocurrió o pudo
haber ocurrido.
\end{formula}

\section{Tabla comparativa de los tres casos}

\begin{table}[H]
\centering
\begin{tabularx}{\textwidth}{>{\raggedright\arraybackslash}X c c >{\raggedright\arraybackslash}X}
\toprule
\textbf{Tipo de caso} & \textbf{Artículo} & \textbf{Plazo} & \textbf{Circunstancia} \\
\midrule
Ordinario & Art. 85 & 3 años & Sin noticias \\
Extraordinario genérico & Art. 86, 1.\textsuperscript{er} párr. & 2 años & Catástrofe / actividad de riesgo \\
Extraordinario específico & Art. 86, 2.\textsuperscript{do} párr. & 6 meses & Buque naufragado / aeronave perdida \\
\bottomrule
\end{tabularx}
\caption{Comparación de los tres casos de presunción de fallecimiento}
\end{table}

\section{Casos reales ilustrativos}

\begin{example}{Caso ``Los Andes'' (1972) — Aeronave perdida}
En 1972, un avión militar uruguayo iba rumbo a Chile y cayó en la mitad
de la cordillera de los Andes. Durante 72 días se buscó a sus ocupantes
hasta que dos supervivientes lograron cruzar hacia el lado chileno y
dieron aviso. Esas personas estuvieron desaparecidas durante 72 días. Si
hubiera existido necesidad de declarar su fallecimiento, habría
correspondido el caso extraordinario específico (buque/aeronave), por lo
que hubiera sido necesario esperar seis meses.
% TODO: contenido incompleto o ambiguo en las notas originales — el apunte indica que "los encontraron antes de los 62 días", cifra que no coincide con los "72 días" mencionados previamente en el mismo párrafo; se conserva tal como figura en el material original.
Como los encontraron antes de los 62 días, ese plazo no llegó a cumplirse.
\end{example}

\begin{example}{Caso Malaysia Airlines MH370 (2014) — Aeronave perdida}
En 2014 desapareció un avión de Malaysia Airlines y nunca se supo con
certeza dónde terminó. El último contacto se toma como fecha para contar
los seis meses. Quien tuviera un familiar en ese vuelo debía esperar, como
mínimo, seis meses desde ese último contacto para poder iniciar el
proceso de presunción de fallecimiento bajo el caso extraordinario
específico.
\end{example}

% ======================================================================
% CAPÍTULO 4
% ======================================================================
\chapter{Procedimiento y sentencia}

\section{Legitimación y juez competente}

\begin{formula}{Art. 87 CCCN — Legitimación activa}
Cualquier persona que tenga un derecho subordinado a la muerte puede
pedir la declaración del fallecimiento presunto y debe justificar los
extremos legales: la ausencia, la falta de noticias y las diligencias
realizadas para averiguar la existencia del ausente.
\end{formula}

El \textbf{juez competente} es el juez del domicilio del ausente. Si bien
las normas procesales corresponden a cada provincia, en este caso el
Código Civil y Comercial las incorpora para dar una
\textbf{regulación común a todo el país}, evitando diferencias en el
tratamiento de una situación muy delicada, ya que se está declarando
fallecida a una persona para iniciar su sucesión.

\section{Diligencias probatorias}

Para justificar los extremos legales deben realizarse diligencias
tendientes a averiguar la existencia del ausente, entre ellas:

\begin{itemize}
    \item Oficios a distintas entidades.
    \item Pedido de informes a su tarjeta de crédito.
    \item Pedido de informes a su cuenta bancaria.
    \item Pedido a la empresa de telefonía celular.
    \item Pedido al club si es socio.
    \item Pedido a Migraciones para verificar entradas o salidas del país.
    \item Pedido de informe sobre actividad de su celular.
\end{itemize}

Se le nombra un \textbf{defensor al ausente} (o el defensor oficial). Se
lo cita por \textbf{edictos una vez por mes durante seis meses}: se
publican seis edictos, uno por mes. Se le nombra, asimismo, un curador de
los bienes si no hay quien los esté administrando.

\begin{important}
\textbf{Aclaración del Art. 79 (ausencia simple) frente a la presunción
de fallecimiento.} El Código aclara que el hecho de haber tramitado
previamente la simple ausencia (Art. 79) \textbf{no} es presupuesto
necesario ni suficiente para la presunción de fallecimiento. Es decir:
(a) puede iniciarse directamente la presunción de fallecimiento sin haber
tramitado la ausencia previa; y (b) haber tramitado la ausencia no exime
de tener que probar nuevamente, en el juicio de presunción de
fallecimiento, que se han realizado todas las diligencias para dar con el
ausente. Ambos procesos son independientes, aunque puedan darse en forma
sucesiva.
\end{important}

\section{Sentencia y día presuntivo de fallecimiento}

Pasados los seis meses de citación, recibida la prueba y oído el
defensor, el juez debe dictar sentencia. La sentencia contiene tres
elementos:

\begin{enumerate}
    \item Declara el fallecimiento presunto, comprobados que estén los extremos legales.
    \item Fija el día presuntivo de fallecimiento.
    \item Dispone la inscripción de la sentencia en el Registro Civil.
\end{enumerate}

\textbf{¿Cómo se fija el día presuntivo?}

\begin{table}[H]
\centering
\begin{tabularx}{\textwidth}{>{\raggedright\arraybackslash}X >{\raggedright\arraybackslash}X >{\raggedright\arraybackslash}X}
\toprule
\textbf{Caso} & \textbf{Día presuntivo} & \textbf{Ejemplo} \\
\midrule
Ordinario (Art. 85) & Último día del 1.\textsuperscript{er} año = mitad del plazo de 3 años & Última noticia: 31/12/2017 $\rightarrow$ día presuntivo: 30/06/2019 \\
Extraordinario genérico (Art. 86, 1.\textsuperscript{er} párr.) & El día en que ocurrió el suceso (o el término medio si se extendió en el tiempo) & Si el incendio duró varios días: el día del mediodía del período \\
Extraordinario específico (Art. 86, 2.\textsuperscript{do} párr.) & El último día en que se tuvo noticia del buque o aeronave & MH370: último contacto = 8 de marzo de 2014 $\rightarrow$ ese es el día presuntivo \\
\bottomrule
\end{tabularx}
\caption{Fijación del día presuntivo de fallecimiento según el caso}
\end{table}

El día presuntivo es de gran relevancia porque determina
\textbf{quiénes son los herederos} que entran en la sucesión del
presuntamente fallecido, declarado judicialmente. Si es posible, la
sentencia debe determinar en qué hora presuntiva se fijó el
fallecimiento; si no, se toma al final de ese día.

% ======================================================================
% CAPÍTULO 5
% ======================================================================
\chapter{Garantías y efectos de la reaparición}

\section{Las tres garantías adoptadas}

Con la declaración de fallecimiento presunto se abre la sucesión y los
herederos pueden tomar los bienes. Pero el Código adopta \textbf{tres
garantías} por si el ausente reapareciera:

\begin{table}[H]
\centering
\begin{tabularx}{\textwidth}{>{\raggedright\arraybackslash}p{0.28\textwidth} >{\raggedright\arraybackslash}X}
\toprule
\textbf{Garantía} & \textbf{Descripción} \\
\midrule
\textbf{Inventario} & Se hace un inventario detallado de todos los bienes del presuntamente fallecido (el acervo sucesorio / masa hereditaria). En base a ese inventario, si reaparece, los herederos deben rendir cuentas de todos esos bienes. \\
\textbf{Prenotación registral} & Al transferir los bienes a nombre de los herederos (por ejemplo, un departamento a nombre de dos hijos), se hace una prenotación: se deja constancia en el asiento registral de que el bien está sometido al régimen de declaración de fallecimiento presunto. \\
\textbf{Restricción para enajenar/gravar} & Para enajenar (vender) o gravar (hipotecar) los bienes heredados, los herederos necesitan autorización judicial. La prenotación avisa a terceros contratantes de que ese bien está sometido a este régimen. \\
\bottomrule
\end{tabularx}
\caption{Las tres garantías del sistema ante una eventual reaparición}
\end{table}

\section{Duración de la prenotación (Art. 91 CCCN)}

\begin{formula}{Art. 91 CCCN — Prenotación}
La prenotación caduca a los \textbf{cinco (5) años} desde el día
presuntivo del fallecimiento, o cuando la persona hubiera cumplido
\textbf{ochenta (80) años} de edad, lo que ocurra primero. Vencida la
prenotación, los herederos pueden enajenar libremente los bienes.
\end{formula}

\begin{important}
Son cinco años desde el \textbf{día presuntivo del fallecimiento}, no
desde el día en que la persona desapareció, y no desde la fecha de la
sentencia. Por eso el día presuntivo es tan relevante: es el que pone en
marcha el plazo de prenotación.
\end{important}

\section{Efectos de la reaparición (Art. 91 CCCN)}

Si el ausente reaparece \textbf{antes} de que venza la prenotación: se le
devuelven todos los bienes, porque el inventario asegura que se sabe
exactamente con qué se empezó, y la restricción de enajenación junto con
el requisito de autorización judicial protegieron el patrimonio.

Si el ausente reaparece \textbf{después} de vencida la prenotación
(cuando ya los bienes podían enajenarse libremente):

\begin{itemize}
    \item Se le entregan los bienes que aún existen.
    \item Se le entregan los bienes adquiridos con el valor de los que faltaron (subrogación real: por ejemplo, si se vendió un departamento y se compraron cuatro cocheras, se entregan las cuatro cocheras).
    \item Se le entrega el precio pendiente de cobro si se vendió un bien con saldo impago.
    \item Se le entregan los frutos no consumidos.
\end{itemize}

\begin{example}{Caso real: Oberá, Misiones (2011)}
Una persona se metió al río Uruguay y desapareció; la buscaron por tierra
y aire con helicópteros. Tiempo después, publicó en una red social una
foto desde Curitiba, en la playa. Se armó un escándalo: resultó que había
decidido hacerse ``desaparecer'' voluntariamente y fue descubierto con
una nueva familia. Si esa persona hubiera tenido bienes y se hubiera
iniciado la declaración, podría haber reaparecido y reclamado la
devolución de lo que le correspondía.
\end{example}

\begin{example}{Caso real: Inglaterra — reaparición tras cinco años}
Se conoció un caso en Inglaterra (noticias de diciembre de 2007) de una
persona que había desaparecido y reapareció cinco años después. Si
durante ese tiempo se hubiera declarado su fallecimiento presunto en
Argentina, al reaparecer tendría derecho a que se le restituyeran los
bienes —o su equivalente— conforme las reglas del Art. 91 CCCN.
\end{example}

% ======================================================================
% CUADRO CORNELL
% ======================================================================

\section{Cuadro de repaso Cornell}

El siguiente cuadro reúne, en formato Cornell, las preguntas y conceptos
centrales desarrollados a lo largo de los capítulos anteriores, con el
propósito de facilitar el repaso activo del material.

\begin{longtable}{
    >{\raggedright\arraybackslash}p{0.30\textwidth}
    >{\raggedright\arraybackslash}p{0.58\textwidth}
}
\toprule
\textbf{Preguntas / palabras clave} & \textbf{Notas / respuestas} \\
\midrule
\endfirsthead

\toprule
\textbf{Preguntas / palabras clave} & \textbf{Notas / respuestas} \\
\midrule
\endhead

\midrule
\multicolumn{2}{r}{\textit{Continúa en la página siguiente}} \\
\endfoot

\bottomrule
\endlastfoot

\textbf{¿Cuál es la diferencia entre ausencia y presunción de fallecimiento?} &
La ausencia (Arts. 79-84) es una medida conservatoria de bienes: nombra
un curador sin declarar fallecida a la persona. La presunción de
fallecimiento (Arts. 85-92) declara judicialmente el fallecimiento,
inscribe la sentencia en el Registro Civil y abre la sucesión. La
presunción supone ausencia, pero no es lo mismo. \\

\textbf{¿Quién puede pedir la declaración de ausencia?} &
El Ministerio Público y las personas que tengan interés legítimo en los
bienes (acreedores, sucesores). El juez competente es el del domicilio
del ausente; si no tiene domicilio en el país o es desconocido, el juez
del lugar donde están los bienes. \\

\textbf{¿Cuáles son los tres casos de presunción de fallecimiento?} &
1) Ordinario: ausencia sin noticias por 3 años (Art. 85). 2)
Extraordinario genérico: persona involucrada en catástrofe o actividad de
riesgo, plazo de 2 años (Art. 86, 1.\textsuperscript{er} párr.). 3)
Extraordinario específico: buque naufragado o aeronave perdida, plazo de
6 meses (Art. 86, 2.\textsuperscript{do} párr.). \\

\textbf{¿Cómo se fija el día presuntivo de fallecimiento?} &
En el ordinario: mitad del plazo (1 año y medio después de la última
noticia). En el extraordinario genérico: el día en que ocurrió el suceso
(o término medio si se extendió). En el extraordinario específico: el
último día en que se tuvo noticia del buque o aeronave. \\

\textbf{¿Cuáles son las tres garantías que protegen al ausente por si
reaparece?} &
1) Inventario detallado de todos los bienes. 2) Prenotación registral en
los asientos de los bienes. 3) Prohibición de enajenar o gravar sin
autorización judicial. La prenotación dura 5 años desde el día presuntivo
o hasta que la persona hubiera cumplido 80 años. \\

\textbf{¿Qué pasa si el ausente reaparece después de vencida la
prenotación?} &
Se le restituyen los bienes existentes, los adquiridos con el valor de
los que faltaron (subrogación real), el precio pendiente de cobro y los
frutos no consumidos. No puede recuperar bienes ya enajenados libremente
por los herederos una vez caduca la prenotación. \\

\textbf{¿Es necesario haber tramitado la ausencia antes de iniciar la
presunción de fallecimiento?} &
No. Según el Código, la declaración de ausencia (Art. 79) no es
presupuesto necesario ni suficiente para la presunción de fallecimiento.
Pueden tramitarse en forma sucesiva, pero en el juicio de presunción hay
que probar nuevamente las diligencias para dar con el ausente. \\

\textbf{¿Qué regula el CCCN sobre el procedimiento?} &
El CCCN incorpora normas procesales comunes para todo el país (citación
por edictos una vez por mes durante 6 meses, defensor oficial del
ausente, curador de bienes), para evitar diferencias entre jurisdicciones
en situaciones de gran sensibilidad jurídica y personal. \\

\midrule

\multicolumn{2}{p{\textwidth}}{
\textbf{Resumen:} El derecho civil argentino responde a la incertidumbre
generada por la desaparición de una persona mediante dos institutos
diferenciados que conviven en el CCCN (Arts. 79-92). La \textbf{ausencia}
(Arts. 79-84) es una herramienta de protección patrimonial preventiva:
permite nombrar un curador de bienes y adoptar medidas conservatorias sin
declarar fallecida a la persona. La \textbf{presunción de fallecimiento}
(Arts. 85-92) va más lejos: declara judicialmente el fallecimiento,
inscribe la sentencia en el Registro Civil y habilita la apertura de la
sucesión. Se estructura en tres casos según el grado de probabilidad de
muerte: el ordinario (3 años), el extraordinario genérico —catástrofes o
actividades de riesgo— (2 años) y el extraordinario específico —buque o
aeronave— (6 meses). La sentencia fija el día presuntivo, que determina
quiénes son los herederos y desde cuándo corre la prenotación registral
de cinco años. Durante ese período, los bienes no pueden enajenarse sin
autorización judicial, garantizando la restitución al ausente en caso de
reaparición. El sistema, con sus antecedentes en Vélez Sársfield y la Ley
14.394, busca equilibrar la certeza jurídica que exigen las relaciones
patrimoniales y sucesorias con la protección de los derechos del
ausente, a quien siempre se le reserva la posibilidad de recuperar su
patrimonio si regresa.
} \\

\bottomrule
\end{longtable}
\partintro{Junto a la persona humana, el ordenamiento reconoce como sujetos de derecho a las personas jurídicas. Esta parte estudia su concepto y naturaleza, el régimen de las personas jurídicas privadas, su responsabilidad civil, el fin de su existencia, las fundaciones y las asociaciones civiles y simples asociaciones.}
\part{El sujeto de la relación jurídica: las personas jurídicas}
\chapter{Introducción a las personas jurídicas}
\label{ch:pj}

Todo abogado, trabaje en derecho de familia, comercial, laboral, notarial o administrativo, se encuentra permanentemente con personas jurídicas. Saber identificar su tipo, leer su estatuto, comprender sus órganos de gobierno y conocer las reglas de responsabilidad es una habilidad transversal y cotidiana: desde redactar el estatuto de una ONG, hasta asesorar sobre la responsabilidad de los directivos de una sociedad anónima, gestionar la disolución de un consorcio o defender a una asociación civil en un proceso disciplinario. Este conocimiento también resulta útil en la vida cotidiana, dado que las personas se relacionan constantemente con personas jurídicas: el consorcio del edificio, el club al que se pertenece, la empresa en la que se trabaja, el Estado que regula.

\textbf{Ejemplo integrador que guía la exposición:} un grupo de vecinos forma un club de fútbol. Para ello necesitan constituir una persona jurídica. ¿Qué tipo eligen? ¿Qué nombre pueden usar? ¿Quién representa al club? ¿Pueden sancionar a un socio? ¿Qué pasa si el club se disuelve? Cada capítulo de esta parte responde a una de estas preguntas: el primero trata el concepto, la naturaleza y la clasificación; el segundo, el régimen de las personas jurídicas privadas; el tercero, su responsabilidad civil; el cuarto, el fin de su existencia y las fundaciones; y el quinto, las asociaciones civiles y simples asociaciones.

% ============================================================

\section{Presentación del tema}

El estudio del régimen de las personas jurídicas en el derecho civil y comercial se sitúa dentro de la Parte General del Código, que se ocupa de los elementos de la relación jurídica: el sujeto, el objeto y la causa. Habiendo estudiado a fondo la persona humana, corresponde ahora tratar la persona jurídica, es decir, los entes colectivos que entablan también relaciones jurídicas en el campo del derecho privado.

\section{El problema de fondo: la asociatividad humana}

El problema central que resuelve esta regulación es cómo dar tratamiento jurídico a la asociatividad, dado que el ser humano es un ser sociable por naturaleza que tiende a unirse con otros para fines útiles. Esa regulación jurídica es fundamental porque permite establecer con qué requisitos esos entes colectivos, morales o ideales pueden entablar relaciones jurídicas: qué tipos de personas jurídicas existen, cuándo hay personalidad, cuáles son sus atributos, cuándo empieza y cuándo finaliza su existencia, y qué consecuencias tiene.

Estas personas jurídicas apuntan muchas veces a facilitar la organización de la vida cotidiana: desde el Estado como organización política que trasciende a los ciudadanos y tiene una personalidad distinta a la de estos, hasta las sociedades, las mutuales, las cooperativas y las iglesias.

\begin{note}
Existe un tipo de persona jurídica de enorme trascendencia económica, las sociedades, que tiene su materia específica (el Derecho Societario); el Estado también tiene la suya (el Derecho Administrativo). En esta materia se estudian las grandes disposiciones generales del Código Civil y Comercial aplicables a todas las personas jurídicas.
\end{note}

\section{Regulación en el Código de Vélez (Ley 340, 1871)}

\begin{note}
\textbf{Código de Vélez --- artículos relevantes.}

\textbf{Art. 30:} ``Todos los entes susceptibles de adquirir derechos y contraer obligaciones son personas''.

\textbf{Art. 31:} Las personas se clasifican en personas físicas (o de existencia visible) y personas jurídicas (o de existencia ideal).

\textbf{Art. 32:} Las personas jurídicas son definidas por exclusión: son todos los entes susceptibles de adquirir derechos y contraer obligaciones que no son personas de existencia visible.

\textbf{Arts. 33 a 50:} Régimen general de las personas jurídicas.
\end{note}

\section{El nuevo Código Civil y Comercial: metodología}

El nuevo Código regula la materia de forma más organizada y sistemática. El Libro Primero, la Parte General, comienza con un Título Primero sobre la persona humana, y el Título Segundo trata de las personas jurídicas.

Dentro de los seis libros del Código, en el Libro Primero:
\begin{itemize}
    \item Persona humana: artículos 19 al 140.
    \item Personas jurídicas: a partir del artículo 141.
\end{itemize}

El Título Segundo sobre personas jurídicas se divide en:
\begin{itemize}
    \item Capítulo 1: Parte general (personalidad, composición y clasificación).
    \item Capítulo 2: Asociaciones civiles y simples asociaciones.
    \item Capítulo 3: Fundaciones.
\end{itemize}

\begin{note}
El nuevo Código adopta la denominación única de ``persona jurídica'', diferenciándose del Código anterior, que hablaba de ``persona de existencia ideal'' y ``persona jurídica''. Antes de la Ley 17.711 (1968) se discutía si el género era la persona de existencia ideal y la persona jurídica una especie; con dicha reforma se tendió a equipararlos. Otras denominaciones usadas en doctrina son ``personas morales'' o ``personas colectivas''.
\end{note}

\section{Definición del artículo 141 del CCyCN}

\begin{definition}{Persona jurídica (art. 141 CCyCN)}
Son personas jurídicas todos los entes a los cuales el ordenamiento jurídico les confiere aptitud para adquirir derechos y contraer obligaciones para el cumplimiento de su objeto y los fines de su creación.
\end{definition}

Esta definición presenta dos características destacables:

\begin{enumerate}
    \item \textbf{Se define por la positiva}: se incorpora como elemento al ordenamiento jurídico, que es quien confiere la aptitud. Esto marca la diferencia con las personas humanas, cuya personalidad es un dato de la naturaleza reconocido —no concedido— por el legislador.
    \item \textbf{Se incorpora el elemento finalista}: la aptitud está orientada al cumplimiento del objeto y los fines de la creación. Este es el fundamento del principio de especialidad.
\end{enumerate}

\section{Principio de especialidad}

A diferencia de la persona humana, que puede dedicarse a cualquier actividad lícita, las personas jurídicas tienen un objeto y un fin propios. El \textbf{principio de especialidad} significa que la persona jurídica, dentro de su posibilidad de adquirir derechos y contraer obligaciones, está acotada a sus fines, a su especialidad, a su objeto, y a aquello para lo cual fue constituida.

\begin{important}
\textbf{Artículo 156 CCyCN --- Objeto.} El objeto de la persona jurídica privada debe ser preciso y determinado.
\end{important}

Esto resulta especialmente relevante para la redacción de los estatutos: los constituyentes deben ser claros, explícitos y descriptivos sobre cómo piensan lograr sus fines y cuál es el objeto de la persona jurídica.

\section{Teorías sobre la naturaleza jurídica}

\subsection{Teoría de la ficción (Savigny)}

Savigny recoge elaboraciones que provenían del derecho canónico y desarrolla la idea de la ficción. Para esta teoría, la persona humana es la única dotada de voluntad; como para Savigny el derecho es una expresión de la voluntad, no existirían propiamente personas jurídicas, sino un artificio creado por el legislador, que puede ser ampliado o restringido, pero que carece de voluntad propia.

\begin{important}
Consecuencia fundamental de la teoría de la ficción: la persona jurídica no respondería por los daños causados —responderían solo las personas humanas—, y como se trataría de una representación con mandato, este mandato nunca podría ser conferido para el ilícito. Esto dio lugar a situaciones injustas y, con el tiempo, se modificó esta posición. El artículo 43 del Código de Vélez, en su redacción original, sostenía esta posición: las personas jurídicas no respondían por los daños causados por quienes las dirigían o administraban.
\end{important}

\subsection{Teorías negatorias de la personalidad}

Niegan que existan propiamente personas jurídicas. Sostienen que lo que existe son patrimonios de afectación o bienes colectivos, pero no algo que pueda llamarse persona jurídica.

\subsection{Teorías de la realidad}

Buscan encontrar un sustrato real detrás de lo que se llama persona jurídica:

\begin{itemize}
    \item \textbf{Teoría organicista}: la persona jurídica sería como un organismo vivo, con sus mecanismos de adopción de la voluntad (los órganos que la gobiernan). Existe una realidad que el legislador viene a tutelar, proteger y regular.
    \item \textbf{Teoría de la institución} (Hauriou, seguida en Argentina por Llambías): las personas jurídicas son instituciones. Una institución requiere tres elementos: (a) una idea, objeto o fines en torno a los cuales se nuclean personas; (b) un sustrato personal, es decir, las personas que se nuclean en torno a esa idea; y (c) un poder u organización que permite llevar adelante esos fines.
\end{itemize}

\subsection{Teorías normativistas (Kelsen)}

Las personas jurídicas son un centro de imputación de normas: una realidad a la que el derecho le asigna consecuencias jurídicas a través de normas dictadas por la autoridad competente.

\begin{note}
En el Código Civil y Comercial vigente se percibe una fuerte presencia de las teorías normativistas —el ordenamiento jurídico es quien confiere la aptitud—, pero también existe una dimensión institucional, en tanto se reconocen realidades organizativas como las simples asociaciones.
\end{note}

\section{Clasificación: personas jurídicas públicas y privadas}

El Código Civil y Comercial simplifica la clasificación entre públicas y privadas.

\subsection{Personas jurídicas públicas (arts. 146 y 147)}

\begin{definition}{Personas jurídicas públicas (art. 146 CCyCN)}
Son personas jurídicas públicas:
\begin{enumerate}[label=\alph*)]
    \item El Estado nacional, las provincias, la Ciudad Autónoma de Buenos Aires, los municipios, las entidades autárquicas y las demás organizaciones a las que el ordenamiento jurídico les atribuya ese carácter.
    \item Los estados extranjeros, las organizaciones a las que el derecho internacional público reconozca personalidad jurídica (como la ONU) y toda otra persona jurídica constituida en el extranjero cuyo carácter público resulte de su derecho aplicable.
    \item La Iglesia Católica.
\end{enumerate}
\end{definition}

La inclusión de la Iglesia Católica con personería jurídica pública se vincula con el artículo 2 de la Constitución Nacional, que establece que el gobierno federal sostiene el culto católico apostólico romano, en virtud de la preeminencia histórica de la Iglesia antes de la constitución del Estado argentino.

En cuanto a su reconocimiento, comienzo, capacidad, funcionamiento y organización, las personas jurídicas públicas se rigen por sus propias leyes y ordenamientos. Así, por ejemplo, si se litiga contra el Banco Central, habrá que revisar su Carta Orgánica y las leyes que lo regulan.

\subsection{Personas jurídicas privadas (art. 148)}

El artículo 148 enumera los tipos de personas jurídicas privadas.

\begin{table}[H]
\centering
\caption{Tipos de personas jurídicas privadas (art. 148 CCyCN)}
\begin{tabularx}{\textwidth}{
    >{\raggedright\arraybackslash}p{0.22\textwidth}
    >{\raggedright\arraybackslash}X
    >{\raggedright\arraybackslash}p{0.22\textwidth}
}
\toprule
\textbf{Tipo} & \textbf{Característica principal} & \textbf{Ley aplicable} \\
\midrule
Sociedades & Finalidad directa de lucro. Los socios participan en ganancias y pérdidas. & Ley 19.550 \\
Asociaciones civiles & No tienen fin de lucro. Unión de personas para un fin común (cultural, deportivo, social, etc.). & Arts. 168--186 CCyCN \\
Simples asociaciones & Similar a la asociación civil, pero sin autorización estatal para funcionar. & Arts. 187--192 CCyCN \\
Fundaciones & Patrimonio afectado por uno o varios fundadores a un fin de bien común. Sin miembros. & Arts. 193--224 CCyCN \\
Iglesias y comunidades religiosas (no católicas) & No tienen fin de lucro. Libertad de culto y conciencia. & Ley 21.745 \\
Mutuales & Ayuda mutua entre sus miembros. Socorros recíprocos. & Ley 20.321 \\
Cooperativas & Asociación de personas para fines comunes de producción, trabajo o consumo. & Ley 20.337 \\
Consorcios de propiedad horizontal & Conjunto de propietarios de unidades funcionales de un inmueble sometido al régimen de propiedad horizontal. & Art. 2044 CCyCN \\
\bottomrule
\end{tabularx}
\end{table}

\begin{note}
El artículo 149 establece que, si el Estado participa en una persona jurídica privada, ello no modifica el carácter privado de esta, aunque puede haber derechos y obligaciones diferenciadas en razón del interés público comprometido.
\end{note}

\section{Ley aplicable (art. 150)}

\begin{formula}{Ley aplicable a las personas jurídicas privadas (art. 150 CCyCN)}
Las personas jurídicas privadas se rigen:
\begin{enumerate}
    \item[1.°)] Por las normas imperativas de la ley especial o, en su defecto, de este Código;
    \item[2.°)] Por las normas del acto constitutivo con sus modificaciones y de los reglamentos, prevaleciendo las primeras en caso de divergencia;
    \item[3.°)] Por las normas supletorias de leyes especiales, o en su defecto, por las de este Título.
\end{enumerate}
\end{formula}

El estatuto tiene fuerza de norma jurídica. Por ello, siempre que se aborde una cuestión en la que interviene una persona jurídica, lo primero que corresponde hacer es leer su estatuto en su última versión, con todas las modificaciones, y también sus reglamentos.

\section{Personalidad diferenciada e inoponibilidad (art. 144)}

Un presupuesto fundamental es que la persona jurídica es distinta de sus miembros. Esto implica distintos patrimonios, distintas responsabilidades, y la posibilidad incluso de que el propio miembro demande a la persona jurídica o viceversa.

\begin{important}
\textbf{Regla general}: los miembros no responden por las obligaciones de la persona jurídica. Esta es una de las diferencias más importantes que debe considerar quien constituye una persona jurídica: responde con el capital que aporta, pero no con todo su patrimonio personal.
\end{important}

\begin{definition}{Inoponibilidad de la persona jurídica (art. 144 CCyCN)}
La actuación que esté destinada a la consecución de fines ajenos a la persona jurídica, que constituya un recurso para violar la ley, el orden público o la buena fe, o para frustrar derechos de cualquier persona, se imputa a quienes, a título de socios, asociados, miembros o controlantes directos o indirectos, la hicieron posible, quienes responderán solidaria e ilimitadamente por los perjuicios causados.
\end{definition}

Esta norma reconoce como antecedente el artículo 54 de la Ley de Sociedades (Ley 19.550), pero ahora se aplica a todo el sistema de personas jurídicas por estar ubicada en la Parte General.

La inoponibilidad significa que los efectos de la personería no pueden hacerse valer frente a quien ha sido dañado, quien puede accionar directamente contra los socios, asociados, miembros o controlantes que hicieron posible esa actuación dañina. Responden solidaria e ilimitadamente: todos responden por el todo y con todos sus patrimonios.

% ============================================================
\chapter{Régimen de las personas jurídicas privadas}
% ============================================================

El Código ofrece reglas generales sobre nombre, domicilio, capacidad, patrimonio, estatutos y comienzo de la existencia, que luego cada tipo de persona jurídica complementa con sus normas específicas.

\section{Nombre (art. 151)}

\begin{definition}{Nombre de la persona jurídica (art. 151 CCyCN)}
La persona jurídica debe tener un nombre que la identifique como tal, con el aditamento indicativo de la forma jurídica adoptada. La persona jurídica en liquidación debe aclarar esta circunstancia en la utilización de su nombre. El nombre debe satisfacer recaudos de veracidad, novedad y aptitud distintiva, tanto respecto de otros nombres como de marcas, nombres de fantasía u otras formas de referencia a bienes o servicios, se relacionen o no con el objeto de la persona jurídica.

No puede contener términos o expresiones contrarios a la ley, el orden público o las buenas costumbres, ni inducir a error sobre la clase u objeto de la persona jurídica.

Si se incluye el nombre de personas humanas, deben contar con la conformidad de estas, que se presume si son socios. Sus herederos pueden oponerse a la continuación del uso si acreditan perjuicios materiales o morales.
\end{definition}

Los requisitos pueden precisarse del siguiente modo:

\begin{itemize}
    \item \textbf{Indicación de la forma jurídica}: debe agregarse el tipo de persona jurídica (por ejemplo, ``El estudiante feliz S.A.''). Si está en liquidación, se agrega ``en liquidación''.
    \item \textbf{Veracidad}: concordancia entre el nombre y el objeto, para que no haya lugar a confusión. Esto se juega especialmente en ciertos ámbitos, como el de las universidades.
    \item \textbf{Novedad y aptitud distintiva}: el nombre debe distinguirse de otros nombres, marcas, nombres de fantasía u otras formas de referencia a bienes y servicios.
\end{itemize}

\begin{example}{Teatro Colón S.A.}
No se puede fundar una sociedad anónima con el nombre ``Teatro Colón S.A.'' porque se confundiría con el Teatro Colón preexistente. El nombre no tendría novedad ni aptitud distintiva.
\end{example}

\begin{example}{Fundación Lionel Messi / Fundación Borges}
Si se quiere usar el nombre de una persona humana, se necesita su conformidad. Si esa persona ya falleció, los herederos pueden oponerse si acreditan perjuicios materiales o morales. Así, si la Fundación Borges —supuestamente dedicada a las altas letras— organizara un concurso de letras de reggaetón, los herederos de Borges podrían oponerse alegando que ello afecta el buen nombre del escritor y resulta incompatible con los fines fundacionales.
\end{example}

\section{Domicilio y sede social (arts. 152 y 153)}

El domicilio es el asiento jurídico de la persona jurídica y permite ubicarla, determinar el juez competente y conocer su inscripción registral.

\begin{table}[H]
\centering
\caption{Domicilio y sede social}
\begin{tabularx}{\textwidth}{>{\raggedright\arraybackslash}X >{\raggedright\arraybackslash}X}
\toprule
\textbf{Domicilio} & \textbf{Sede social} \\
\midrule
Se fija en el estatuto. Puede expresarse simplemente con una ciudad (por ejemplo, ``Buenos Aires''). & Es la dirección física concreta donde tienen lugar la dirección y administración de la persona jurídica (por ejemplo, ``Av. Figueroa Alcorta 2263, CABA''). \\
Determina el juez competente y la jurisdicción registral. & Generalmente se inscribe ante el organismo de contralor (IGJ). Si está en el estatuto, cambiarla requiere modificar el estatuto. \\
\bottomrule
\end{tabularx}
\end{table}

\begin{important}
\textbf{Artículo 153 CCyCN --- Efectos de la sede social.} Se tienen por válidas y vinculantes para la persona jurídica todas las notificaciones efectuadas en la sede inscripta.
\end{important}

También existe el domicilio especial: cuando la persona jurídica posee establecimientos o sucursales, tiene su domicilio especial en el lugar de dichos establecimientos, pero solo para las obligaciones allí contraídas.

\begin{example}{Sucursal del Banco Santander en Mendoza}
Si un contratista de pintura firmó un contrato con la sucursal mendocina del Banco Santander y este no le paga, puede litigar en Mendoza. No necesita ir a Buenos Aires a litigar contra la casa matriz, porque la obligación se contrajo localmente.
\end{example}

\section{Capacidad y principio de especialidad (arts. 141, 156)}

La persona jurídica es capaz de derecho; la capacidad de ejercicio se realiza a través de sus órganos representativos. Pero esa capacidad tiene límites:

\begin{enumerate}
    \item \textbf{Principio de especialidad}: la capacidad está acotada al objeto y a los fines de la persona jurídica.
    \item \textbf{Limitaciones legales}: la ley prohíbe a las personas jurídicas adquirir ciertos derechos o los limita. Por ejemplo, el usufructo en favor de personas jurídicas tiene un plazo determinado, mientras que para personas humanas puede extenderse hasta su muerte.
    \item \textbf{Limitaciones por la naturaleza de las cosas}: las personas jurídicas no pueden adoptar niños, porque el derecho de familia se refiere a las personas humanas.
\end{enumerate}

\section{Patrimonio (arts. 143 y 154)}

La persona jurídica debe tener un patrimonio distinto del de sus miembros. Esto es una consecuencia directa de la personalidad diferenciada.

La persona jurídica tiene bienes a su nombre; si se trata de información registrable, se pueden inscribir preventivamente. Por ejemplo, si alguien dona un inmueble a una fundación y esta tarda en inscribirse, puede realizarse una anotación preventiva en el registro para proteger el bien mientras se completan los trámites.

\begin{note}
En cuanto a su duración, las personas jurídicas son en principio ilimitadas en el tiempo, excepto que el acto constitutivo disponga lo contrario. La prórroga está regulada en el artículo 165 y requiere decisión de los miembros o aprobación de la autoridad de contralor antes del vencimiento del plazo.
\end{note}

\section{Comienzo de la existencia (art. 142)}

\begin{definition}{Comienzo de la existencia (art. 142 CCyCN)}
La existencia de la persona jurídica privada comienza desde su constitución. No necesita autorización legal para funcionar, excepto disposición legal en contrario. En los casos en que se requiere autorización estatal, la persona jurídica no puede funcionar antes de obtenerla.
\end{definition}

En el Código anterior existía una discusión sobre si el comienzo de la existencia era el momento de la constitución o el de la autorización. El nuevo Código es claro: el comienzo es la constitución. Sin embargo, si la ley requiere autorización estatal para funcionar (como en el caso de las fundaciones), la persona jurídica no puede comenzar a operar hasta obtenerla.

\begin{example}{Fundación banco de alimentos}
Si hoy se constituye una fundación con objeto de banco de alimentos y se aporta un capital de \$100.000, la fundación existe desde ese momento. Pero no puede empezar a hacer compras, emitir facturas ni hacer donaciones hasta que se obtenga la autorización estatal. La constitución y la autorización son dos momentos distintos.
\end{example}

\section{Estatutos y gobierno (arts. 150, 158, 159, 160)}

El estatuto es una norma jurídica secundaria que emana de las partes que constituyen la persona jurídica y que la regula internamente. No debe confundirse con el acto constitutivo:

\begin{itemize}
    \item El \textbf{acto constitutivo} es el acto jurídico por el cual las partes deciden dar inicio a la persona jurídica. Generalmente, en ese acto se redacta el estatuto.
    \item El \textbf{estatuto} cobra independencia porque puede sufrir sucesivas modificaciones a lo largo del tiempo.
\end{itemize}

\begin{note}
Hay clubes que tienen más de cien años y cuyo estatuto ha cambiado muchísimas veces. Cuando se solicita el estatuto ante la Inspección de Justicia, se pide el actualizado con todos sus cambios.
\end{note}

\subsection{Contenido del estatuto}

El estatuto de una persona jurídica privada debe contener:

\begin{itemize}
    \item Nombre y domicilio.
    \item Objeto preciso y determinado (art. 156).
    \item Derechos y deberes de los miembros.
    \item Forma de gobierno: composición de los órganos, cómo se toman las decisiones, quórum y mayorías necesarias (incluidas las agravadas para decisiones trascendentes, como reforma del estatuto, expulsión de miembros o disolución).
    \item Atribuciones de cada órgano y sus miembros.
    \item Administración de bienes.
    \item Destino de los bienes en caso de disolución.
\end{itemize}

\subsection{Deber de lealtad y diligencia de los administradores (art. 159)}

\begin{important}
\textbf{Artículo 159 CCyCN --- Deber de lealtad y diligencia.} Los administradores de la persona jurídica deben obrar con lealtad y diligencia. No pueden perseguir ni favorecer intereses contrarios a los de la persona jurídica. Si en determinada operación tuvieran un interés contrario al de ella, deben hacerlo saber a los demás miembros del órgano de administración o, en su caso, al órgano de gobierno, y abstenerse de cualquier intervención relacionada con dicha operación.
\end{important}

Ser administrador de una persona jurídica implica una gran responsabilidad, porque se confían bienes económicos que no son propios.

% ============================================================
\chapter{Responsabilidad civil de las personas jurídicas}
% ============================================================

\section{Evolución histórica: del artículo 43 de Vélez al nuevo Código}

El Código de Vélez Sarsfield, adherente a la teoría de la ficción, establecía en el artículo 43 que las personas jurídicas no respondían por los daños que causaban quienes las dirigían o administraban: no se podía iniciar una acción de daños y perjuicios contra una persona jurídica.

La lógica era coherente con la teoría de la ficción: como la personería era solo ficción y el mandato nunca se confería para el ilícito, las actuaciones ilícitas se imputaban a las personas humanas, no a la persona jurídica. A lo largo del tiempo, esta solución fue generando la sensación de que constituía una injusticia.

% TODO: contenido incompleto o ambiguo en las notas originales — la referencia al profesor "Rivarola" y a la traducción del término utilizado por Freitas se transcribe tal como figura en el apunte, sin poder verificarse contra la fuente original.
Existieron distintas posturas doctrinarias para resolver el problema. Un profesor de apellido Rivarola acudió al diccionario que empleaba Vélez y encontró que, en el texto original de Freitas (en portugués), el término utilizado podía entenderse como ``cuando'' y no como ``aunque'', lo que cambiaría el sentido de la norma; sin embargo, esta interpretación no fue pacífica. El jurista Borda, en un fallo, llegó a declarar que el artículo 43 era de una injusticia intolerable.

\begin{important}
Todo cambió en 1968 con la Ley 17.711, que reformó el artículo 43 y estableció que la persona jurídica responde por los daños que causen quienes la dirigen o administran en ejercicio o con ocasión de sus funciones. El nuevo Código recoge esta solución en el artículo 1763.
\end{important}

\section{Tres supuestos de responsabilidad}

\subsection{Responsabilidad por actos de directivos y administradores (art. 1763)}

\begin{formula}{Responsabilidad de la persona jurídica (art. 1763 CCyCN)}
La persona jurídica responde por los daños que causen quienes las dirigen o administran en ejercicio o con ocasión de sus funciones.
\end{formula}

La expresión ``en ejercicio o con ocasión de sus funciones'' resulta determinante: si un presidente de la empresa está de vacaciones y comete un daño, la empresa no responde, porque el acto no ocurre en ejercicio ni con ocasión de sus funciones.

\subsection{Responsabilidad por cosas y actividades riesgosas (art. 1757)}

\begin{formula}{Hecho de las cosas y actividades riesgosas (art. 1757 CCyCN)}
Toda persona responde por el daño causado por el riesgo o vicio de las cosas, o de las actividades que sean riesgosas o peligrosas por su naturaleza, por los medios empleados o por las circunstancias de su realización.
\end{formula}

Esta norma se aplica a todas las personas, incluidas las jurídicas. Se trata de una responsabilidad objetiva: no depende de factores de atribución subjetivos (dolo o culpa); la responsabilidad se atribuye siempre.

\begin{example}{Sociedad anónima y accidente de tránsito}
Una sociedad anónima tiene un automóvil. Ese auto causa un daño. La persona jurídica responde, aunque no estuviera manejando: se le imputa la responsabilidad porque es dueña de la cosa riesgosa.
\end{example}

\subsection{Responsabilidad por los dependientes (art. 1753)}

\begin{formula}{Responsabilidad del principal por el hecho del dependiente (art. 1753 CCyCN)}
El principal responde objetivamente por los daños que causen los que están bajo su dependencia, o las personas de las cuales se sirve para el cumplimiento de sus obligaciones, cuando el hecho dañoso acaece en ejercicio o con ocasión de las funciones encomendadas.
\end{formula}

También esta norma engloba a las personas jurídicas. El término ``dependientes'' es genérico: incluye a cualquier persona de la que la persona jurídica se sirve para cumplir sus obligaciones, más allá de la relación de empleo formal.

\begin{example}{Personal de seguridad que golpea a alguien}
En una empresa, el personal de seguridad golpea a una persona de manera arbitraria y violenta. El daño causado se imputa tanto al empleado como a la persona jurídica. La empresa no puede escudarse en que ``fue el empleado''.
\end{example}

\section{Responsabilidad de los administradores (art. 160)}

\begin{formula}{Responsabilidad de los administradores (art. 160 CCyCN)}
Los administradores responden en forma ilimitada y solidaria frente a la persona jurídica, sus integrantes y terceros, por los daños causados por su culpa en el ejercicio o con ocasión de sus funciones, por acción u omisión.
\end{formula}

Esta disposición es la contracara del artículo 1763: si el presidente causó un daño, frente al tercero dañado responden tanto el presidente como la persona jurídica. Pero la persona jurídica, a su vez, puede iniciar acción contra el presidente por los daños que causó. Los administradores responden de forma ilimitada (con todos sus bienes) y solidaria (todos responden por el todo).

% ============================================================
\chapter{Fin de la existencia y fundaciones}
% ============================================================

\section{Causales de disolución (art. 163)}

\begin{formula}{Causales de disolución (art. 163 CCyCN)}
La persona jurídica se disuelve por:
\begin{enumerate}[label=\alph*)]
    \item La decisión de sus miembros adoptada por unanimidad o por la mayoría establecida por el estatuto o disposición especial;
    \item El cumplimiento de la condición resolutoria a la que el acto constitutivo subordinó su existencia;
    \item La consecución del objeto para el cual la persona jurídica se formó, o la imposibilidad sobreviniente de cumplirlo;
    \item El vencimiento del plazo;
    \item La declaración de quiebra;
    \item La fusión respecto de las personas jurídicas que se fusionan o la escisión respecto de la escindida;
    \item La reducción a uno del número de miembros, si la ley especial exige pluralidad de ellos y esta no es restablecida dentro de los tres meses;
    \item La denegatoria o revocación firmes de la autorización estatal para funcionar, cuando esta sea requerida;
    \item El agotamiento de los bienes destinados a sostenerla;
    \item Cualquier otra causa prevista en el estatuto o en disposiciones de este Título o de ley especial.
\end{enumerate}
\end{formula}

Entre las causales más relevantes cabe señalar:

\begin{itemize}
    \item \textbf{Inciso a) --- Decisión de los miembros}: por unanimidad si el estatuto no dice nada; si el estatuto establece una mayoría, se aplica esa mayoría.
    \item \textbf{Inciso b) --- Condición resolutoria}: cuando la existencia estaba subordinada a un hecho futuro e incierto que se cumple.
    \item \textbf{Inciso c) --- Consecución del objeto}: cuando la persona jurídica se creó para un fin específico y este se cumplió o se tornó imposible.
    \item \textbf{Inciso e) --- Quiebra}: cuando el pasivo supera al activo y no hay propuesta de acuerdo preventivo viable. La quiebra está regulada por ley especial y tiene por objeto liquidar el activo para pagar a los acreedores; excepcionalmente, puede haber avenimiento o conversión.
    \item \textbf{Inciso h) --- Revocación estatal}: el Estado puede revocar la autorización para funcionar, con las restricciones del artículo 164.
\end{itemize}

\begin{example}{Fundación condicionada a un programa estatal}
Una fundación se crea para ayudar a personas en situación de calle ``siempre y cuando no exista un programa especial del Gobierno de la Ciudad de Buenos Aires para esa problemática''. Si el gobierno implementa tal programa, se cumple la condición y la fundación se disuelve.
\end{example}

\begin{example}{Fundación para los cien años de un club}
Se constituye una fundación especial para organizar la celebración del centenario de un club. Pasada la celebración, el objeto se cumplió y la personería se disuelve.
\end{example}

\section{Revocación estatal --- garantías (art. 164)}

\begin{important}
\textbf{Artículo 164 CCyCN --- Revocación de la autorización estatal.} La revocación de la autorización estatal debe fundarse en la comisión de actos graves que violen la ley, el estatuto o el reglamento. La resolución que la disponga debe ser fundada y contener la indicación de los actos que la justifican. Se debe dar a los interesados la posibilidad de ser escuchados antes de que la resolución se dicte. Es apelable judicialmente. La apelación puede tener efectos suspensivos.
\end{important}

Las garantías que protegen a la persona jurídica frente a una revocación arbitraria pueden sintetizarse así:

\begin{enumerate}
    \item Solo procede por actos graves que violen la ley, el estatuto o el reglamento.
    \item La resolución debe ser fundada (no puede ser arbitraria).
    \item Se debe garantizar el derecho de defensa (descargo previo).
    \item Es apelable judicialmente, y la apelación puede tener efecto suspensivo (la persona jurídica sigue funcionando mientras se tramita).
\end{enumerate}

El fundamento constitucional de estas garantías es la libertad de asociación consagrada en el artículo 14 de la Constitución Nacional.

\section{Reconducción y liquidación}

\begin{itemize}
    \item \textbf{Reconducción}: si se decidió la disolución pero no terminó la liquidación, se puede revertir esa decisión siempre que la causa pueda ser removida. Por ejemplo, si se cumplió el plazo pero los miembros deciden ampliarlo y transformar la personería de tiempo limitado a indefinido.
    \item \textbf{Liquidación}: una vez dada la causal de disolución, se pagan las deudas y obligaciones pendientes con el activo, se abonan los gastos, y los bienes remanentes se destinan según lo establezca la ley o el estatuto. En asociaciones civiles y fundaciones generalmente se exige que los bienes vayan a otra institución de bien común, sin fin de lucro.
\end{itemize}

\section{Las fundaciones (arts. 193--224)}

\begin{definition}{Fundación (art. 193 CCyCN)}
Las fundaciones son personas jurídicas que se constituyen con una finalidad de bien común, sin propósito de lucro, mediante el aporte patrimonial de una o más personas, denominadas fundadores, destinando dicho aporte al cumplimiento de sus fines.
\end{definition}

Características distintivas de las fundaciones:

\begin{itemize}
    \item No tienen miembros: el fundador aporta un patrimonio, que queda afectado al cumplimiento de un fin de bien común.
    \item No tienen fin de lucro.
    \item Requieren autorización estatal para funcionar (art. 142).
    \item Pueden constituirse por acto entre vivos o por testamento.
\end{itemize}

\begin{example}{Fundación banco de alimentos (caso detallado)}
En un contexto de crisis, un fundador con capital decide constituir una fundación con \$100.000 para organizar un banco de alimentos. La fundación alquila un depósito, contrata empleados, establece convenios con instituciones que proveen alimentos y organiza campañas de donación. Todo eso es la fundación funcionando. No hay ``socios'': hay un fundador que dotó el patrimonio y un Consejo de Administración que lo gestiona.

También puede constituirse por testamento: un deportista que practique hockey puede establecer en su testamento que su parte disponible de la herencia vaya a una fundación para que personas sin recursos puedan practicar ese deporte.
\end{example}

\subsection{Patrimonio inicial suficiente}

El Código exige que la fundación cuente con un patrimonio inicial que posibilite razonablemente el cumplimiento de sus fines (recaudo no exigido para las asociaciones civiles). La Inspección de Personas Jurídicas no otorgará la autorización si los fondos son insuficientes.

El plan de acción a tres años (plan trienal) que se presenta ante el organismo de contralor puede incluir tanto los aportes iniciales como los ingresos futuros proyectados (eventos benéficos, campañas de donación, etc.).

\subsection{Órgano de gobierno: Consejo de Administración}

El único órgano de gobierno es el Consejo de Administración, dotado de todas las atribuciones para el gobierno y la administración. Debe tener al menos tres personas humanas como integrantes.

Puede delegar funciones en un Comité Ejecutivo (art. 205), que ejerce esas funciones entre las reuniones del Consejo. Los miembros del Comité Ejecutivo pueden ser remunerados, mientras que los del Consejo de Administración no.

\subsection{Control externo}

Las fundaciones no tienen órgano de fiscalización interna: el control es exclusivamente externo, a cargo de la autoridad de contralor (IGJ u organismo provincial equivalente). Entre las atribuciones del organismo de contralor (art. 222) se encuentran:

\begin{itemize}
    \item Aprobar estatutos y sus reformas.
    \item Fiscalizar el funcionamiento.
    \item Designar administradores interinos, suspender deliberaciones, remover administradores, convocar asambleas.
    \item Intervenir en la disolución y liquidación.
\end{itemize}

\subsection{Inmutabilidad del objeto (art. 216)}

No es posible reformar el objeto de la fundación, salvo que haya resultado de cumplimiento imposible. Esto se debe a que la fundación fue autorizada para ese objeto específico, y cambiarlo implicaría traicionar la voluntad inicial del fundador.

% ============================================================
\chapter{Asociaciones civiles y simples asociaciones}
% ============================================================

Las asociaciones civiles y las simples asociaciones son personas jurídicas privadas de enorme importancia social: clubes deportivos, ONG, asociaciones barriales, organizaciones de víctimas, entidades religiosas. La Constitución Nacional consagra la libertad de asociación como un derecho fundamental.

\section{Asociaciones civiles (arts. 168--186)}

\subsection{Concepto y objeto}

\begin{definition}{Objeto de las asociaciones civiles (art. 168 CCyCN)}
La asociación civil debe tener un objeto que no sea contrario al interés general o al bien común. El objeto no puede ser contrario a las disposiciones de este Código ni a la ley especial, ni referirse a cuestiones de interés particular de los asociados.
\end{definition}

La definición es por la negativa: cualquier finalidad que no afecte al interés general puede ser objeto de una asociación civil. La asociación civil es una corporación: unión de personas en torno a un fin común (cultural, artístico, deportivo, religioso, social).

\begin{important}
La asociación civil no tiene fin de lucro y no persigue el lucro como fin principal. Puede eventualmente generar ganancias si organiza un evento a beneficio, pero ese dinero se reinvierte en los fines de la asociación; no se distribuye entre los miembros.
\end{important}

\subsection{Diferencia fundamental con las sociedades}

En la sociedad, los socios aportan capital, participan de ganancias y pérdidas, y obtienen utilidades. En la asociación civil, los miembros aportan su cuota —fijada por el estatuto y las asambleas— que contribuye a los fines colectivos, sin distribuirse ganancias.

\begin{note}
La Resolución de la IGJ 7/2015 entiende por ``bien común'' a las finalidades que contribuyen al bien de la comunidad en general o a la mejora de las condiciones de vida, en contraposición al bien individual.
\end{note}

\subsection{Constitución y contenido del estatuto (art. 170)}

El instrumento constitutivo debe hacerse en escritura pública. El artículo 170 establece que debe contener:

\begin{itemize}
    \item Identificación de los constituyentes.
    \item Nombre, objeto y domicilio.
    \item Plazo de duración y causales de disolución.
    \item Composición del patrimonio inicial.
    \item Clases o categorías de socios (activos, cadetes, vitalicios, etc.).
    \item Régimen de admisión, renuncia y sanciones disciplinarias.
    \item Órganos de gobierno: composición, atribuciones, funcionamiento y duración de cargos.
    \item Cierre del ejercicio económico anual.
    \item Destino de los bienes en caso de liquidación: deben ir a otra entidad de bien común, pública o privada, sin fin de lucro, con domicilio en la República Argentina.
\end{itemize}

\section{Órganos de gobierno}

\subsection{Comisión Directiva (art. 171)}

Es el órgano de dirección y administración cotidiana de la asociación civil. Los cargos mínimos son Presidente, Secretario y Tesorero (los demás son Vocales). Cada estatuto puede agregar cargos adicionales: vicepresidentes primero y segundo, prosecretario, protesorero, etc.

Debe reunirse al menos una vez al mes. Sus atribuciones las fija el estatuto, y hay que distinguir entre las que tiene la Comisión como cuerpo colegiado y las que tiene el Presidente individualmente.

\begin{important}
No hay que confundir el hecho de que alguien sea presidente con que pueda resolver todo por sí mismo. El presidente es quien preside un órgano que es la Comisión Directiva.
\end{important}

\begin{example}{Venta de un inmueble por una asociación civil}
Si una asociación civil quiere vender un inmueble, hay que leer el estatuto para saber quién puede tomar esa decisión. A veces se exige que la asamblea autorice previamente la venta; otras veces basta con la Comisión Directiva. El abogado debe siempre revisar el estatuto antes de asesorar.
\end{example}

\subsection{Asamblea}

Es el órgano supremo: la reunión de todos los asociados. En asociaciones muy numerosas puede instrumentarse a través de representantes elegidos por los socios.

\begin{itemize}
    \item \textbf{Asamblea ordinaria}: generalmente una vez al año; aprueba la memoria y el balance presentado por la Comisión Directiva con el dictamen del órgano de fiscalización.
    \item \textbf{Asamblea extraordinaria}: se convoca cuando es necesario modificar el estatuto, elegir autoridades, aprobar reglamentos o tomar decisiones que requieran mayorías especiales.
\end{itemize}

\subsection{Órgano de fiscalización}

Puede recaer en personas no asociadas. Si la asociación cuenta con más de cien socios, debe tener una Comisión Revisora de Cuentas. Si la asociación exige una profesión determinada para ser socio, los integrantes del órgano de fiscalización no necesariamente deben contar con ese título, pero debe contratarse a profesionales independientes para el asesoramiento correspondiente.

Además del control interno, existe también la fiscalización estatal del artículo 174.

\section{Poder disciplinario}

La asociación civil puede sancionar a sus miembros. En esto se evidencia la distinción entre la persona jurídica y sus miembros.

\subsection{Principios que rigen el poder disciplinario}

\begin{enumerate}
    \item Debe estar regulado en el estatuto, con claridad sobre qué órgano interviene (asamblea, Comisión Directiva o Comisión de Disciplina).
    \item Se debe garantizar el derecho de defensa del afectado: debe darse vista al sumariado y permitirle hacer sus descargos.
    \item Gradualidad y proporcionalidad de la sanción: primero se amonesta, luego se suspende por períodos crecientes, y finalmente se expulsa. No se puede expulsar sin proceso.
    \item La decisión es apelable ante los tribunales judiciales si no se respetó el derecho de defensa, si fue desproporcionada o si no hubo debido proceso.
\end{enumerate}

\begin{example}{Destrucción de un vestuario}
Un socio muy enojado, luego de un partido, destruye el vestuario del club un domingo. El lunes, la Comisión Directiva puede suspenderlo transitoriamente como medida cautelar. Pero no puede expulsarlo sin proceso: debe abrirse un sumario, recopilarse pruebas y testimonios, notificarse al imputado, permitirle hacer su descargo y luego resolver. La sanción debe ser proporcional a la falta.
\end{example}

\subsection{Límites de la revisión judicial}

El juez puede revisar la forma del proceso disciplinario (si se respetó el derecho de defensa, si la sanción fue proporcional). Pero no puede revisar el fondo de la decisión —salvo grave violación de derechos humanos—, porque ello implicaría alterar el objeto y la autonomía de la persona jurídica.

\begin{example}{Club vegano y el vendedor de panchos}
Un club de personas veganas expulsa a un socio que insiste en vender panchos en las instalaciones. Si el juez le diera la razón al expulsado en cuanto al fondo, estaría alterando el objeto de la asociación y el derecho de todos sus miembros a mantener la identidad que los une. Por eso, el control judicial se centra en la forma, no en el fondo de la decisión disciplinaria.
\end{example}

\section{Simples asociaciones (arts. 187--192)}

Las simples asociaciones son personas jurídicas privadas que, a diferencia de las asociaciones civiles, no necesitan autorización estatal para funcionar. Existen muchísimas en la práctica.

\subsection{Diferencias con las asociaciones civiles}

\begin{table}[H]
\centering
\caption{Asociación civil frente a simple asociación}
\begin{tabularx}{\textwidth}{
    >{\raggedright\arraybackslash}p{0.24\textwidth}
    >{\raggedright\arraybackslash}X
    >{\raggedright\arraybackslash}X
}
\toprule
\textbf{Aspecto} & \textbf{Asociación civil} & \textbf{Simple asociación} \\
\midrule
Instrumento constitutivo & Escritura pública & Instrumento público o privado con firma certificada \\
Autorización estatal & Requerida & No requerida \\
Comienzo de existencia & Desde la constitución (previa autorización para funcionar) & Desde la constitución \\
Órgano de fiscalización & Obligatorio (Comisión Revisora de Cuentas si hay más de 100 socios) & Prescindible si hay menos de 20 socios; se certifica con contador externo \\
Responsabilidad en insolvencia & Separación plena entre persona jurídica y miembros & Administradores y miembros que participaron responden solidariamente con sus bienes propios \\
Régimen supletorio & No aplica & Se aplica el régimen de asociaciones civiles \\
\bottomrule
\end{tabularx}
\end{table}

\subsection{Responsabilidad en las simples asociaciones (art. 191)}

\begin{formula}{Responsabilidad en simples asociaciones (art. 191 CCyCN)}
Si la simple asociación no puede pagar a sus acreedores, los administradores y todos los miembros que hubieren participado de hecho en la administración responden solidariamente con sus propios bienes por las deudas de la simple asociación.

Los demás miembros que no participaron en la administración solo responden hasta el límite de las cuotas y contribuciones prometidas o debidas.
\end{formula}

En la simple asociación existe una menor diferencia entre la persona jurídica y sus miembros: quienes participaron de alguna forma en la administración y llevaron a un resultado de quebranto o insolvencia deben responder con sus propios bienes.

% ============================================================

\section{Cuadro Cornell de repaso general}
% ============================================================

El siguiente cuadro reúne, según el método Cornell, las preguntas centrales que atraviesan los cuatro bloques temáticos estudiados, con sus respuestas y una síntesis final para el repaso integral de la materia.

\begin{longtable}{
    >{\raggedright\arraybackslash}p{0.28\textwidth}
    >{\raggedright\arraybackslash}p{0.62\textwidth}
}
\cornellhead
\endfirsthead
\cornellhead
\endhead
\midrule
\multicolumn{2}{r}{\textit{continúa en la página siguiente}} \\
\midrule
\endfoot
\bottomrule
\endlastfoot

\textbf{¿Qué es una persona jurídica?} &
Son los entes a los cuales el ordenamiento jurídico les confiere aptitud para adquirir derechos y contraer obligaciones para el cumplimiento de su objeto y los fines de su creación (art. 141 CCyCN). A diferencia de la persona humana, cuya personalidad es reconocida por el legislador, la de las personas jurídicas es conferida por él. \\[0.4em]

\textbf{¿En qué consiste el principio de especialidad?} &
La capacidad de la persona jurídica está acotada a su objeto y sus fines específicos. No puede hacer cualquier cosa: solo aquello para lo que fue creada. De ahí la importancia de redactar el objeto con precisión y amplitud suficiente en el estatuto. \\[0.4em]

\textbf{¿Cuál es la diferencia entre la teoría de la ficción y las teorías de la realidad?} &
La teoría de la ficción (Savigny) entiende que la persona jurídica es un artificio del legislador, sin voluntad real. Las teorías de la realidad (organicista, de la institución) buscan un sustrato real. Las teorías normativistas las conciben como centros de imputación de normas. El CCyCN combina elementos normativistas e institucionales. \\[0.4em]

\textbf{¿Cómo clasifica el Código a las personas jurídicas?} &
Públicas: Estado nacional, provincias, municipios, entidades autárquicas, estados extranjeros, organismos internacionales, Iglesia Católica. Privadas: sociedades, asociaciones civiles, simples asociaciones, fundaciones, mutuales, cooperativas, consorcios de propiedad horizontal y otras previstas en leyes especiales. \\[0.4em]

\textbf{¿Qué norma prevalece según el artículo 150?} &
Orden de prevalencia: (1) normas imperativas de la ley especial; (2) normas imperativas del CCyCN; (3) estatuto y reglamentos; (4) normas supletorias. El estatuto tiene fuerza de norma jurídica en el ámbito de la persona jurídica. \\[0.4em]

\textbf{¿Cuándo se aplica la inoponibilidad del artículo 144?} &
Cuando la persona jurídica es usada para violar la ley, el orden público, la buena fe o para frustrar derechos de terceros, la personería se vuelve inoponible y se puede accionar directamente contra los socios, miembros o controlantes que hicieron posible esa actuación. Responden solidaria e ilimitadamente. \\[0.4em]

\textbf{¿Qué requisitos debe cumplir el nombre de una persona jurídica?} &
Debe indicar la forma jurídica adoptada; debe ser veraz, novedoso y con aptitud distintiva (diferenciarse de nombres, marcas y fantasías); no puede inducir a error sobre el objeto; no puede contener términos contrarios a la ley, el orden público o las buenas costumbres. Si incluye el nombre de una persona humana, se necesita su conformidad. \\[0.4em]

\textbf{¿Cuál es la diferencia entre domicilio y sede social?} &
Domicilio: ciudad indicada en el estatuto; determina jurisdicción y juez competente. Sede social: dirección física inscripta ante el organismo de contralor; las notificaciones allí enviadas se tienen por válidas. Domicilio especial: en el lugar de las sucursales, para las obligaciones allí contraídas. \\[0.4em]

\textbf{¿Cómo responde civilmente la persona jurídica?} &
Tres supuestos: (1) daños causados por directivos/administradores en ejercicio o con ocasión de sus funciones (art. 1763); (2) daños por cosas o actividades riesgosas, responsabilidad objetiva (art. 1757); (3) daños causados por dependientes en ejercicio o con ocasión de funciones (art. 1753). Los administradores también responden ilimitada y solidariamente ante la propia persona jurídica (art. 160). \\[0.4em]

\textbf{¿Cuáles son las causales de disolución del artículo 163?} &
Decisión de los miembros; condición resolutoria; consecución o imposibilidad del objeto; vencimiento del plazo; quiebra; fusión o escisión; reducción a uno del número de miembros; revocación o denegatoria de autorización estatal; agotamiento de bienes; otras causas estatutarias o legales. \\[0.4em]

\textbf{¿Qué garantías protegen frente a la revocación estatal del artículo 164?} &
Solo procede por actos graves que violen la ley, el estatuto o el reglamento. Requiere resolución fundada, garantía del derecho de defensa, y es apelable judicialmente con posible efecto suspensivo. El fundamento es la libertad de asociación constitucional. \\[0.4em]

\textbf{¿Qué caracteriza a las fundaciones?} &
Persona jurídica sin miembros, con fin de bien común y sin lucro, creada por el aporte patrimonial de uno o varios fundadores. Requiere autorización estatal para funcionar. Necesita patrimonio inicial suficiente (plan trienal). Órgano de gobierno: Consejo de Administración (mínimo tres personas). Solo fiscalización externa. El objeto es inmutable, salvo imposibilidad sobreviniente. \\[0.4em]

\textbf{¿Qué caracteriza a las asociaciones civiles?} &
Corporación de personas con fin común no lucrativo. Instrumento constitutivo: escritura pública. Tres órganos: Comisión Directiva, Asamblea (suprema) y órgano de fiscalización (Comisión Revisora de Cuentas si hay más de 100 socios). Poder disciplinario: proceso con derecho de defensa, gradualidad y proporcionalidad. Apelación judicial solo sobre la forma, salvo grave violación de derechos humanos. \\[0.4em]

\textbf{¿En qué se diferencian las simples asociaciones?} &
No requieren autorización estatal. Pueden constituirse por instrumento privado con firma certificada. Comienzan su existencia desde la constitución. Si hay menos de 20 socios, pueden prescindir del órgano de fiscalización. Principal diferencia: en caso de insolvencia, los administradores y miembros que participaron en la gestión responden solidariamente con sus propios bienes. \\[0.4em]

\midrule
\multicolumn{2}{p{0.94\textwidth}}{
\textbf{Resumen:}
Las personas jurídicas son entes colectivos a los que el ordenamiento jurídico --no la naturaleza-- otorga aptitud para ser titulares de derechos y obligaciones, con el fin de cumplir el objeto para el que fueron creadas (principio de especialidad). El Código Civil y Comercial las clasifica en públicas --el Estado y sus organismos, los estados extranjeros, los organismos internacionales y la Iglesia Católica-- y privadas --sociedades, asociaciones civiles, simples asociaciones, fundaciones, mutuales, cooperativas y consorcios--. Su existencia comienza desde la constitución, aunque algunas requieren autorización estatal para funcionar. Son entidades con personalidad diferenciada de sus miembros (patrimonios y responsabilidades distintas), salvo los supuestos de inoponibilidad del artículo 144, donde la personería cede frente al uso abusivo o fraudulento. Responden civilmente por los daños que causen sus directivos en ejercicio de sus funciones (art. 1763), por las cosas y actividades riesgosas que posean (art. 1757) y por los hechos dañosos de sus dependientes (art. 1753); además, los propios administradores responden ilimitada y solidariamente frente a la persona jurídica (art. 160). Su disolución puede deberse a la voluntad de los miembros, al cumplimiento o imposibilidad del objeto, al vencimiento del plazo, a la quiebra o a la revocación estatal --esta última con estrictas garantías de fundamentación y defensa--. Entre los tipos especiales, las fundaciones se distinguen por la ausencia de miembros, la afectación de un patrimonio a un fin altruista y el control externo exclusivo de la autoridad de contralor; las asociaciones civiles, por su carácter corporativo, sus tres órganos de gobierno y el poder disciplinario sobre sus miembros con plenas garantías procesales; y las simples asociaciones, por no requerir autorización estatal y por la responsabilidad solidaria que recae sobre los administradores y miembros activos ante la insolvencia.
} \\

\end{longtable}
\partintro{Esta parte conecta dos preguntas fundamentales del derecho privado: ¿quién responde? y ¿con qué responde? El patrimonio es el puente entre ambas. Se estudia el patrimonio, la clasificación de los derechos, los bienes y las cosas que lo componen y, finalmente, la garantía patrimonial de los acreedores, los bienes que quedan excluidos de ella y las acciones destinadas a conservarla.}
\part{El objeto de la relación jurídica: patrimonio, derechos y bienes}
\chapter{El patrimonio}

\section{Planteo general}

Esta parte conecta dos preguntas fundamentales del Derecho Privado: \textbf{¿quién responde?} y \textbf{¿con qué responde?} El patrimonio es el puente entre ambas: es el atributo que le da contenido económico a la persona jurídica y, al mismo tiempo, la garantía sobre la que los acreedores ejecutarán sus créditos. Entender qué son los bienes y las cosas ---y cómo se clasifican--- es entender de qué está hecho ese patrimonio.

\begin{note}
\textbf{Relevancia práctica.} Estas nociones constituyen la base de la práctica profesional jurídica cotidiana:
\begin{itemize}
    \item Al redactar una escritura de compraventa, es necesario distinguir si el bien es mueble o inmueble para determinar la forma requerida (escritura pública o simple contrato).
    \item Al asesorar a un cliente endeudado, permite determinar qué bienes puede conservar por ser inembargables.
    \item En un litigio, permite identificar qué acciones utilizar para recuperar el patrimonio de un deudor que intenta eludir a sus acreedores.
\end{itemize}
El patrimonio como garantía es el concepto central del crédito: no se presta dinero ni se celebra un contrato sin entender sobre qué responde el deudor.
\end{note}

\section{El patrimonio como atributo de la personalidad}

En este capítulo confluyen dos aspectos distintos y complementarios. El primero es el estudio del \textbf{patrimonio como atributo de la personalidad}. Existen distintas posturas: algunas señalan que el patrimonio es un atributo de la personalidad y otras no lo entienden así. Se adopta aquí ---junto con el profesor Lozano Roy--- la posición que entiende el patrimonio como atributo de la personalidad.

La materia aborda el patrimonio desde dos planos:

\begin{tabla}
\begin{tabularx}{\textwidth}{B{3.4cm}YY}
\toprule
Perspectiva & \textbf{Qué estudia} & \textbf{Dónde se ubica} \\
\midrule
Atributo de la personalidad & El patrimonio como cualidad inherente al sujeto de derecho & Dentro del estudio del sujeto de la relación jurídica \\
Objeto de la relación jurídica & Los bienes y cosas que conforman el patrimonio & Dentro del estudio del objeto de la relación jurídica \\
\bottomrule
\end{tabularx}
\end{tabla}

La materia recorre tres grandes ejes: el \textbf{sujeto} (persona humana y jurídica), el \textbf{objeto} (bienes y cosas) y la \textbf{causa} de la relación jurídica. El patrimonio aparece en los dos primeros ejes: como atributo del sujeto y como continente de los objetos.

\begin{norma}{Artículo 15 CCCN --- Titularidad de derechos individuales}
Las personas son titulares de los derechos individuales sobre los bienes que integran su patrimonio conforme a lo que se establece en este Código.
\end{norma}

\section{El patrimonio como garantía común de los acreedores}

El artículo 242 constituye el núcleo funcional del patrimonio:

\begin{norma}{Artículo 242 CCCN --- Garantía común}
Todos los bienes del deudor están afectados al cumplimiento de sus obligaciones y constituyen la garantía común de sus acreedores, con excepción de aquellos que este Código o leyes especiales declaran inembargables e inejecutables.
\end{norma}

\begin{example}{La deuda de un millón de pesos}
Supóngase que una persona debe un millón de pesos. El acreedor la intima a pagar y el deudor no paga. El acreedor no puede limitarse a solicitar cortésmente el pago: en algún momento deberá ejecutar compulsivamente al deudor.

Si el deudor no entrega voluntariamente el dinero, no realiza una transferencia bancaria ni cumple de ninguna manera la obligación, el acreedor se cobra \textbf{ejecutando su patrimonio}. Si el deudor tiene un auto a su nombre y un departamento, el acreedor traba embargo sobre esos bienes y, si el deudor no paga, los ejecuta judicialmente; con el producido de los bienes se cobra.
\end{example}

\section{El patrimonio como universalidad de derecho}

\begin{definition}{Patrimonio (universalidad de derecho)}
Según la teoría clásica ---Llambías y otros autores---, el patrimonio es una \textbf{universalidad de derecho}: un conjunto de bienes y deudas que conforman una unidad porque la ley se la otorga, con independencia de su contenido concreto en cada momento.
\end{definition}

Se destacan dos notas esenciales:

\begin{tabla}
\begin{tabularx}{\textwidth}{B{3.0cm}YY}
\toprule
Nota & \textbf{Descripción} & \textbf{Consecuencia práctica} \\
\midrule
Unicidad & Cada persona tiene un solo patrimonio & No puede fragmentarlo a voluntad, salvo casos legales \\
Mutabilidad del contenido & Los bienes entran y salen, pero el patrimonio subsiste & Los acreedores se cobran sobre los bienes actuales, no sobre los que existían al contraer la deuda \\
\bottomrule
\end{tabularx}
\end{tabla}

Los acreedores pueden cobrarse con los bienes actuales del deudor. El hecho de que el deudor haya vendido un departamento no implica que los acreedores hayan perdido su garantía: la garantía es el \textbf{patrimonio como tal}, con independencia de su contenido. Esto indica también que, si el contenido del patrimonio de la persona va cambiando con el tiempo, sigue siendo su patrimonio.

\subsection{Los patrimonios especiales}

El Código Civil y Comercial reconoce la existencia de patrimonios especiales, disipando la ambigüedad que existía en el código anterior. Son universalidades autónomas que conviven con el patrimonio general:

\begin{tabla}
\begin{tabularx}{\textwidth}{B{4.6cm}Y}
\toprule
Tipo de patrimonio especial & \textbf{Descripción} \\
\midrule
Fondo de comercio & Conjunto de bienes afectados a una explotación comercial \\
Masa de la quiebra & Bienes del fallido afectados al proceso concursal \\
Masa hereditaria & Bienes del causante antes de la partición \\
Herencia del presunto fallecido & Patrimonio sujeto a prenotación durante el período legal \\
\bottomrule
\end{tabularx}
\end{tabla}

\section{Cuadro Cornell: El patrimonio}

\cornellinicio
\begin{longtable}{|Q|Z|}
\hline
\rowcolor{gray!15}\textbf{Preguntas / palabras clave} & \textbf{Notas / respuestas} \\ \hline
\endfirsthead
\hline
\rowcolor{gray!15}\textbf{Preguntas / palabras clave} & \textbf{Notas / respuestas} \\ \hline
\endhead

\cornellpregunta{¿Qué es el patrimonio y por qué se dice que es una universalidad de derecho?}{%
Es una universalidad de derecho: el conjunto de bienes (activo) y deudas (pasivo) de una persona, al que el ordenamiento jurídico le otorga unidad. Es una unidad porque la ley se la otorga, con independencia de su contenido concreto en cada momento (teoría clásica: Llambías y otros autores). En la materia se lo entiende, junto con el profesor Lozano Roy, como atributo de la personalidad.}

\cornellpregunta{¿Cuáles son las notas esenciales del patrimonio?}{%
\textbf{Unicidad:} cada persona tiene un solo patrimonio, que no puede fragmentar a voluntad salvo casos legales.

\textbf{Mutabilidad del contenido:} los bienes entran y salen, pero el patrimonio subsiste; por eso los acreedores se cobran sobre los bienes actuales y no sobre los que existían al contraer la deuda.}

\cornellpregunta{¿Puede existir más de un patrimonio para una misma persona?}{%
Cada persona tiene un solo patrimonio (unicidad), salvo casos legales. El CCCN reconoce patrimonios especiales, universalidades autónomas que conviven con el patrimonio general: fondo de comercio, masa de la quiebra, masa hereditaria y herencia del presunto fallecido.}

\cornellpregunta{¿Por qué se dice que el patrimonio es la garantía común de los acreedores?}{%
Porque todos los bienes del deudor responden por sus obligaciones (Arts. 242 y 743 CCCN). Si el deudor no paga voluntariamente, el acreedor puede ejecutar compulsivamente los bienes que integran su patrimonio, con excepción de los declarados inembargables.}

\cornellpregunta{¿Desde qué perspectivas se estudia el patrimonio?}{%
Como atributo de la personalidad (dentro del estudio del sujeto de la relación jurídica) y como conjunto de los bienes y cosas que lo conforman (dentro del estudio del objeto de la relación jurídica). Los tres ejes de la materia son el sujeto, el objeto y la causa; el patrimonio aparece en los dos primeros.}

\cornellresumen{El patrimonio es una universalidad de derecho ---conjunto de bienes y deudas--- y un atributo de la personalidad. Es único y de contenido cambiante, aunque el CCCN admite patrimonios especiales. Es la garantía común de los acreedores (Art. 242 CCCN): todos los bienes del deudor responden por sus obligaciones, salvo los declarados inembargables e inejecutables, y el acreedor puede cobrarse ejecutándolos cuando el deudor no paga voluntariamente.}

\end{longtable}
\cornellfin

% =====================================================================
% CAPÍTULO 2 — CLASIFICACIÓN DE LOS DERECHOS
% =====================================================================
\chapter{Clasificación de los derechos}

\section{Derechos patrimoniales y extrapatrimoniales}

La gran clasificación de los derechos distingue los que componen el patrimonio de los que quedan fuera de él:

\begin{tabla}
\begin{tabularx}{\textwidth}{B{4.0cm}YY}
\toprule
Categoría & \textbf{Subcategorías} & \textbf{Caracteres} \\
\midrule
Extrapatrimoniales & Derechos personalísimos; derechos de familia & Sin contenido económico directo; ligados a la persona o al vínculo familiar \\
Patrimoniales --- Reales & Dominio, condominio, usufructo, uso, habitación, servidumbre, hipoteca, prenda & Relación directa entre la persona y la cosa; oponibles \textit{erga omnes} \\
Patrimoniales --- Personales & Obligaciones y contratos & Relación entre acreedor y deudor respecto de una prestación \\
Patrimoniales --- Intelectuales & Obras científicas, literarias, artísticas, invenciones, software & Recaen sobre creaciones inmateriales \\
\bottomrule
\end{tabularx}
\end{tabla}

\section{Los derechos reales}

\begin{definition}{Derechos reales}
Los derechos reales son aquellos que conceden un \textbf{señorío inmediato sobre una cosa}: existe una relación directa de la persona con la cosa.
\end{definition}

El derecho real por antonomasia es el \textbf{dominio}, que es el derecho más perfecto y confiere el uso y goce de una cosa. Técnicamente, en el Código Civil y Comercial ---cuya tradición proviene de los romanos---, el \textit{dominus} es el señor de una \textit{res}; por eso estos derechos son «reales»: el término viene del latín \textit{res}, que significa «cosa».

\begin{tabla}
\begin{tabularx}{\textwidth}{B{4.4cm}P{4.2cm}Y}
\toprule
Tipo & \textbf{Ejemplos} & \textbf{Nota} \\
\midrule
Sobre cosa propia --- plenos & Dominio & El más amplio: uso, goce y disposición \\
Sobre cosa propia --- limitados & Condominio & Dos o más titulares sobre la misma cosa \\
Sobre cosa ajena --- de disfrute & Usufructo, uso, habitación, servidumbre & El titular goza de la cosa de otro \\
Sobre cosa ajena --- de garantía & Hipoteca (inmuebles), prenda (muebles) & Garantizan el cumplimiento de una obligación \\
\bottomrule
\end{tabularx}
\end{tabla}

\section{Los derechos personales}

\begin{important}
No deben confundirse los derechos \textbf{personales} con los derechos \textbf{personalísimos}.
\end{important}

Los derechos se llaman «personales» porque confieren prerrogativas para exigir de \textbf{otro} una determinada prestación: la relación no es directa entre una cosa y su dueño, sino entre un acreedor y un deudor en torno a una prestación. Estos derechos se plasman fundamentalmente en obligaciones y contratos.
% TODO: contenido incompleto o ambiguo en las notas originales (la frase original sobre el sujeto que «se plasma fundamentalmente en obligaciones y contratos» resulta ininteligible; se conservó únicamente la referencia a obligaciones y contratos, que coincide con el cuadro de clasificación)

\section{Los derechos intelectuales}

\begin{definition}{Derechos intelectuales}
Los derechos intelectuales confieren un \textbf{derecho de explotación, de uso y de goce sobre una obra intelectual}: una obra científica, la creación de una invención, la obra literaria o artística, el software.
\end{definition}

Se trata de distintos tipos de creaciones que tienen un campo de inmaterialidad, porque no son una casa, un auto o una computadora, sino obras intelectuales.

\section{Cuadro Cornell: Clasificación de los derechos}

\cornellinicio
\begin{longtable}{|Q|Z|}
\hline
\rowcolor{gray!15}\textbf{Preguntas / palabras clave} & \textbf{Notas / respuestas} \\ \hline
\endfirsthead
\hline
\rowcolor{gray!15}\textbf{Preguntas / palabras clave} & \textbf{Notas / respuestas} \\ \hline
\endhead

\cornellpregunta{¿Cómo se clasifican los derechos según su relación con el patrimonio?}{%
\textbf{Extrapatrimoniales} (personalísimos y de familia): sin contenido económico directo, ligados a la persona o al vínculo familiar; quedan fuera del patrimonio.

\textbf{Patrimoniales:} reales, personales e intelectuales.}

\cornellpregunta{¿Qué son los derechos reales y cuál es el derecho real por antonomasia?}{%
Conceden un señorío inmediato sobre una cosa: hay una relación directa de la persona con la cosa, y son oponibles \textit{erga omnes}. El derecho real por antonomasia es el dominio, el derecho más perfecto, que confiere el uso y goce de una cosa. Se los llama «reales» por el latín \textit{res} (cosa).}

\cornellpregunta{¿Qué diferencia hay entre ser dueño de una cosa y tener un derecho sobre ella?}{%
El dominio es un derecho real sobre cosa propia, pleno y el más amplio (uso, goce y disposición). Los derechos reales sobre cosa ajena son de disfrute (usufructo, uso, habitación, servidumbre: el titular goza de la cosa de otro) o de garantía (hipoteca sobre inmuebles, prenda sobre muebles). El condominio es un derecho sobre cosa propia limitado: dos o más titulares sobre la misma cosa.}

\cornellpregunta{¿Por qué se llaman «personales» los derechos personales si no son personalísimos?}{%
Porque confieren prerrogativas para exigir de otro una determinada prestación: la relación no es entre una cosa y su dueño, sino entre un acreedor y un deudor en torno a una prestación. Se plasman en obligaciones y contratos. No deben confundirse con los derechos personalísimos, que son extrapatrimoniales.}

\cornellpregunta{¿Qué son los derechos intelectuales?}{%
Derechos de explotación, de uso y de goce sobre una obra intelectual (obra científica, invención, obra literaria o artística, software). Recaen sobre creaciones con un campo de inmaterialidad.}

\cornellresumen{Los derechos se clasifican en extrapatrimoniales (personalísimos y de familia, sin contenido económico directo) y patrimoniales (reales, personales e intelectuales). Los derechos reales otorgan un señorío inmediato sobre una cosa y son oponibles \textit{erga omnes}; el dominio es el derecho real por antonomasia. Los derechos personales confieren la facultad de exigir de otro una prestación y se plasman en obligaciones y contratos; no deben confundirse con los personalísimos. Los derechos intelectuales recaen sobre creaciones inmateriales.}

\end{longtable}
\cornellfin

% =====================================================================
% CAPÍTULO 3 — BIENES Y COSAS
% =====================================================================
\chapter{Bienes y cosas}

\section{Distinción entre bienes y cosas}

\begin{norma}{Artículo 16 CCCN --- Bienes y cosas}
Los derechos referidos en el primer párrafo del artículo 15 pueden recaer sobre bienes susceptibles de valor económico. Los bienes materiales se llaman cosas. Las disposiciones referentes a las cosas son aplicables a la energía y a las fuerzas naturales susceptibles de ser puestas al servicio del hombre.
\end{norma}

\begin{tabla}
\begin{tabularx}{\textwidth}{B{3.4cm}YY}
\toprule
Concepto & \textbf{Definición} & \textbf{Ejemplos} \\
\midrule
Bien (sentido amplio) & Todo objeto susceptible de valor económico, material o inmaterial & Un crédito, un derecho de autor, una casa, dinero \\
Cosa & Objeto material susceptible de valor económico & Casa, auto, computadora, ganado \\
Bien inmaterial & Bien sin corporalidad física & Un crédito, una marca, un derecho de usufructo \\
\bottomrule
\end{tabularx}
\end{tabla}

\section{Muebles e inmuebles}

La distinción entre cosas inmuebles y cosas muebles es la clasificación más importante de las cosas en el Código.

\subsection{Inmuebles}

\begin{tabla}
\begin{tabularx}{\textwidth}{B{2.8cm}YY}
\toprule
Tipo de inmueble & \textbf{Descripción} & \textbf{Ejemplos} \\
\midrule
Por naturaleza & Fijados al suelo de manera natural o que se encuentran bajo él sin hecho del hombre & El suelo, los minerales, los árboles arraigados \\
Por accesión & Cosas muebles inmovilizadas por su adhesión física al suelo con carácter perdurable & Paredes, pisos, ventanas, puertas, techos de un edificio \\
\bottomrule
\end{tabularx}
\end{tabla}

\begin{example}{Límite de la accesión}
Los pisos, las paredes, las ventanas, las puertas y los techos de una casa están compuestos por cosas que eran muebles pero que fueron incorporadas y fijadas al piso con carácter perdurable, por lo que se convierten en inmuebles. La ventana es separable, pero es parte del inmueble.

En cambio, una biblioteca colocada en un departamento no es parte del inmueble aunque esté puesta con carácter permanente, porque está afectada a la actividad de su titular: si se vende el departamento, se lo vende sin la biblioteca, salvo que se trate de un placar construido directamente e incorporado con las paredes.
\end{example}

\subsection{Cosas muebles}

\begin{tabla}
\begin{tabularx}{\textwidth}{B{3.8cm}YY}
\toprule
Tipo de cosa mueble & \textbf{Descripción} & \textbf{Ejemplos} \\
\midrule
Por desplazamiento propio (semovientes) & Se mueven por sí mismas & Ganado \\
Por fuerza externa & Se desplazan mediante fuerza humana o mecánica & Muebles del hogar, electrodomésticos \\
Locomóviles & Vehículos con propulsión propia & Automóviles, motos, camiones, aeronaves \\
Muebles registrables & Muebles que por su importancia la ley manda registrar & Automotores, buques, aeronaves \\
Muebles no registrables & Sin obligación de registro & La mayoría de los objetos cotidianos \\
\bottomrule
\end{tabularx}
\end{tabla}

\subsection{Importancia práctica de la distinción entre muebles e inmuebles}

\begin{tabla}
\begin{tabularx}{\textwidth}{B{3.4cm}YY}
\toprule
Aspecto jurídico & \textbf{Cosas inmuebles} & \textbf{Cosas muebles} \\
\midrule
Forma de transmisión & Escritura pública obligatoria & Mayor libertad de formas \\
Ley aplicable (DIPr) & Ley del territorio donde está el bien & Ley del domicilio del propietario \\
Prescripción adquisitiva & 20 años (10 con justo título y buena fe) & Plazo más corto; sin registro: posesión vale título \\
Garantías reales & Hipoteca & Prenda \\
\bottomrule
\end{tabularx}
\end{tabla}

\section{Cosas divisibles e indivisibles}

\begin{definition}{Cosas divisibles}
Son divisibles las cosas que pueden dividirse en porciones reales sin ser destruidas, y cada una de las cuales forma un todo homogéneo y análogo tanto a las otras partes como a la cosa original.
\end{definition}

\begin{example}{División de una cosa}
Si deben hacerse, a partir de diez quintales de soja, dos entregas de cinco quintales, la cosa es perfectamente divisible: cada parte es un todo homogéneo. En cambio, si se divide un auto, no quedan dos autos: queda un auto destruido.
\end{example}

La divisibilidad tiene un límite legal: la división no puede tornar antieconómico el uso y aprovechamiento de la cosa.

\begin{example}{División de un lote}
Dividir un lote de diez metros por diez en pequeños lotes de un metro cuadrado sería absolutamente inconveniente y antieconómico.
\end{example}

Por eso la reglamentación en materia de inmuebles fija unidades mínimas ---en algunos lugares, de 866 metros cuadrados---. Las autoridades locales (municipios o provincias) tienen las potestades para reglamentar las divisiones en materia de inmuebles.

\section{Cosas principales y accesorias}

\begin{definition}{Cosas principales y accesorias}
Son \textbf{principales} las cosas que pueden existir por sí mismas; \textbf{accesorias}, aquellas cuya existencia y naturaleza son determinadas por otra de la cual dependen o están adheridas.
\end{definition}

El régimen jurídico de la cosa accesoria es el de la principal.

\begin{example}{El cuadro y su marco}
Sea un cuadro compuesto por la tela, la pintura y un marco: el marco es accesorio con relación a la tela y a la pintura, y el régimen jurídico aplicable es el de la principal. Pero si se trata de una lámina que se compra en cualquier kiosco y se la coloca en un marco enormemente valioso de oro, lo principal pasa a ser el marco, y la lámina colocada en él resulta absolutamente accesoria.
\end{example}

\begin{norma}{Artículo 230 CCCN --- Criterio residual}
Si no es posible determinar la relación de principalidad y accesoriedad entre las cosas, se considera principal la de mayor valor. Si los valores son iguales, no hay cosa principal ni accesoria.
\end{norma}

\section{Cosas consumibles y no consumibles; fungibles y no fungibles}

\begin{tabla}
\begin{tabularx}{\textwidth}{B{3.0cm}YYY}
\toprule
Clasificación & \textbf{Definición} & \textbf{Ejemplos} & \textbf{Importancia} \\
\midrule
Consumibles & Su existencia termina con el primer uso & Alimentos, combustible & Relevante en obligaciones de restitución \\
No consumibles & No dejan de existir por el uso habitual & Muebles, ropa (se deterioran, pero subsisten) & Base de contratos de uso y préstamo \\
Fungibles & Un individuo equivale a otro de la misma especie; pueden sustituirse & Toneladas de trigo de igual calidad, billetes & Admite sustitución en el cumplimiento \\
No fungibles & No pueden sustituirse por otro individuo & Cuadro de artista reconocido, botella de vino numerada & Incumplimiento si se entrega cosa distinta \\
\bottomrule
\end{tabularx}
\end{tabla}

Las obras artísticas no son fungibles; las cosas producidas industrialmente, por lo general, sí lo son.

\section{Frutos y productos}

\begin{tabla}
\begin{tabularx}{\textwidth}{B{3.2cm}YYY}
\toprule
Categoría & \textbf{Definición} & \textbf{Renovación} & \textbf{Ejemplos} \\
\midrule
Frutos naturales & Producidos espontáneamente por la naturaleza & Se renuevan periódicamente & Manzanas de un árbol \\
Frutos industriales & Producidos con intervención humana (agricultura, industria) & Se renuevan con la actividad & Cosechas, producción fabril \\
Frutos civiles & Rentas producidas por el bien a través de actos jurídicos & Se renuevan periódicamente & Alquileres, intereses del dinero prestado \\
Productos & Objetos no renovables que al separarse alteran o disminuyen la sustancia & No se renuevan & Minerales extraídos de una cantera \\
\bottomrule
\end{tabularx}
\end{tabla}

\begin{example}{Producto y fruto}
Si se extrae cobre de una cantera, se trata de un \textbf{producto}: la cantera no volverá a producir ese mineral sino en dos millones de años, de modo que se extraen productos y se disminuye el valor de la cosa. En cambio, las manzanas de un árbol se renuevan: son un \textbf{fruto}.
\end{example}

Esta distinción adquiere importancia cuando debe restituirse algo con sus frutos y sus productos.

\begin{example}{Restitución de una casa por un poseedor de mala fe}
Quien recibe una casa y es de mala fe debe restituirla con las rentas: si la casa producía un alquiler mensual de veinte mil pesos, debe devolver la casa con las rentas que hubiese producido durante esos años.
\end{example}

Los artículos 1935 y 1936 CCCN (y el artículo 234, sobre bienes fuera del comercio) regulan los deberes de restitución de frutos y productos según la buena o mala fe del poseedor.
% TODO: contenido incompleto o ambiguo en las notas originales (no se explica por qué el Art. 234 se menciona junto a los Arts. 1935 y 1936; se conservó la referencia tal como figura)

\section{Bienes según las personas}

\begin{tabla}
\begin{tabularx}{\textwidth}{B{2.9cm}P{1.7cm}YY}
\toprule
Categoría & \textbf{Artículo CCCN} & \textbf{Caracteres} & \textbf{Ejemplos} \\
\midrule
Dominio público del Estado & Art. 235 & Inalienables, inembargables, imprescriptibles; uso de todos & Ríos, lagos, playas, plazas, calles \\
Dominio privado del Estado & Art. 236 & El Estado es titular como particular; pueden enajenarse & Inmueble fiscal no afectado a uso público \\
Bienes de los particulares & Art. 238 & Todos los bienes no enumerados en los artículos anteriores & Cualquier bien en manos de un privado \\
Bienes fuera del comercio & Art. 234 & Prohibida su enajenación por ley o por acto jurídico & Bienes de uso cultual, bienes embargados \\
Bienes de incidencia colectiva & Art. 240 & Interés colectivo en su protección; no apropiables individualmente & Ecosistema, fauna, biodiversidad, agua, valores culturales \\
\bottomrule
\end{tabularx}
\end{tabla}

\subsection{Bienes de incidencia colectiva}

\begin{definition}{Bienes de incidencia colectiva}
Son bienes que no son estrictamente individuales, sino que corresponden de forma homogénea a un conjunto de personas o a la comunidad como tal. Por ejemplo: el ecosistema, la flora, la fauna, la biodiversidad, el agua y los valores culturales.
\end{definition}

\section{Cuadro Cornell: Bienes y cosas}

\cornellinicio
\begin{longtable}{|Q|Z|}
\hline
\rowcolor{gray!15}\textbf{Preguntas / palabras clave} & \textbf{Notas / respuestas} \\ \hline
\endfirsthead
\hline
\rowcolor{gray!15}\textbf{Preguntas / palabras clave} & \textbf{Notas / respuestas} \\ \hline
\endhead

\cornellpregunta{¿Qué diferencia hay entre un bien y una cosa?}{%
Todo bien es un objeto susceptible de valor económico (material o inmaterial). Una cosa es un bien material (objeto corporal). Todos los objetos inmateriales con valor económico son bienes, pero no cosas (Art. 16 CCCN).}

\cornellpregunta{¿Cuál es la clasificación más importante de las cosas?}{%
Muebles e inmuebles. Los inmuebles pueden ser por naturaleza (el suelo y lo adherido orgánicamente) o por accesión (cosas muebles incorporadas con carácter perdurable). Los muebles pueden desplazarse por sí mismos (semovientes) o por fuerza externa; algunos son registrables (automotores).}

\cornellpregunta{¿Por qué importa distinguir muebles de inmuebles?}{%
Define la ley aplicable (en DIPr), el plazo de prescripción adquisitiva (20/10 años para inmuebles \textit{vs.} plazos menores para muebles), la forma de transmisión (escritura pública para inmuebles) y el tipo de garantía real (hipoteca \textit{vs.} prenda).}

\cornellpregunta{¿Cuándo una cosa es divisible?}{%
Cuando puede dividirse en porciones reales sin ser destruida y cada porción forma un todo homogéneo y análogo a las otras partes y a la cosa original (diez quintales de soja en dos entregas de cinco; un auto no es divisible: dividido queda destruido). Límite: la división no puede tornar antieconómico el uso y aprovechamiento; en inmuebles, las autoridades locales fijan unidades mínimas.}

\cornellpregunta{¿Cuándo una cosa es principal o accesoria?}{%
Es principal la que puede existir por sí sola; accesoria, la que depende de otra. El régimen jurídico sigue a lo principal. Si no se puede determinar cuál es principal, prevalece la de mayor valor (Art. 230 CCCN); si los valores son iguales, no hay cosa principal ni accesoria.}

\cornellpregunta{¿Qué diferencia hay entre cosas consumibles y no consumibles, y entre fungibles y no fungibles?}{%
\textbf{Consumibles:} su existencia termina con el primer uso (alimentos, combustible). \textbf{No consumibles:} no dejan de existir por el uso habitual.

\textbf{Fungibles:} un individuo equivale a otro de la misma especie y pueden sustituirse, lo que admite sustitución en el cumplimiento (trigo de igual calidad, billetes). \textbf{No fungibles:} no pueden sustituirse; entregar una cosa distinta es incumplimiento (cuadro de artista reconocido).}

\cornellpregunta{¿Qué diferencia a un fruto de un producto?}{%
El fruto se produce de modo renovable sin alterar ni disminuir la sustancia de la cosa (manzanas, intereses, alquileres). El producto es no renovable y su separación altera o disminuye la sustancia (minerales extraídos de una cantera).}

\cornellpregunta{¿Por qué importa distinguir entre frutos y productos?}{%
Interesa cuando debe restituirse algo con sus frutos y productos: los Arts. 1935 y 1936 CCCN regulan los deberes de restitución según la buena o mala fe del poseedor. Quien recibe una casa de mala fe debe restituirla con las rentas que hubiese producido.}

\cornellpregunta{¿Cómo se clasifican los bienes según las personas?}{%
\begin{itemize}[nosep,leftmargin=1.2em]
\item \textbf{Dominio público del Estado} (Art. 235): inalienables, inembargables, imprescriptibles; uso de todos.
\item \textbf{Dominio privado del Estado} (Art. 236): el Estado es titular como particular; pueden enajenarse.
\item \textbf{Bienes de los particulares} (Art. 238): los no enumerados en los artículos anteriores.
\item \textbf{Bienes fuera del comercio} (Art. 234): prohibida su enajenación por ley o por acto jurídico.
\item \textbf{Bienes de incidencia colectiva} (Art. 240): interés colectivo en su protección; no apropiables individualmente.
\end{itemize}}

\cornellpregunta{¿Qué son los bienes de incidencia colectiva?}{%
Bienes no estrictamente individuales, que corresponden de forma homogénea a un conjunto de personas o a la comunidad como tal (ecosistema, flora, fauna, biodiversidad, agua, valores culturales).}

\cornellresumen{Todo bien es un objeto susceptible de valor económico; las cosas son los bienes materiales (Art. 16 CCCN). Las cosas se clasifican en muebles e inmuebles (la distinción más importante), divisibles e indivisibles, principales y accesorias, consumibles y no consumibles, fungibles y no fungibles, y se distingue entre frutos y productos. Según las personas, los bienes pueden ser del dominio público o privado del Estado, de los particulares, fuera del comercio o de incidencia colectiva. Cada clasificación tiene consecuencias jurídicas precisas.}

\end{longtable}
\cornellfin

% =====================================================================
% CAPÍTULO 4 — GARANTÍA PATRIMONIAL, ACREEDORES Y ACCIONES
% =====================================================================
\chapter{Garantía patrimonial, acreedores y acciones}
\label{ch:garantia}

\section{El patrimonio como garantía: repaso y profundización}

El artículo 743 precisa cuáles son los bienes que constituyen la garantía:

\begin{norma}{Artículo 743 CCCN --- Bienes que constituyen la garantía}
Son bienes del deudor sujetos a ejecución todos los que integran su patrimonio a la fecha de la sentencia o en el que se hubiese determinado, y los que adquiera hasta la extinción de la obligación. El acreedor puede exigir la venta judicial de los bienes del deudor, pero sólo en la medida necesaria para satisfacer su crédito. Todos los acreedores pueden ejecutar estos bienes en posición igualitaria, excepto que exista una causa legal de preferencia.
\end{norma}

\section{Tipos de acreedores y preferencias de cobro}

Cuando hay concurrencia de acreedores ---especialmente en situaciones de insolvencia--- el orden de cobro no es necesariamente igualitario.

\begin{tabla}
\begin{tabularx}{\textwidth}{B{2.7cm}YYY}
\toprule
Tipo de acreedor & \textbf{Fuente de preferencia} & \textbf{Alcance} & \textbf{Ejemplo} \\
\midrule
Con privilegio especial & La ley & Sobre determinados bienes & Acreedor laboral sobre maquinarias del empleador en quiebra \\
Con privilegio general & La ley & Sobre la generalidad del patrimonio & Créditos del fisco en algunos supuestos \\
Con garantía real & Autonomía de la voluntad (contrato) & Sobre el bien gravado & Hipotecario que ejecuta el inmueble hipotecado \\
Quirografario (común) & Ninguna preferencia & Sobre el remanente, en igualdad & Acreedor sin ninguna garantía \\
\bottomrule
\end{tabularx}
\end{tabla}

\begin{example}{Concurrencia de acreedores}
Supóngase que el deudor debe 10.000 dólares a un acreedor, 20.000 a otro y 30.000 a un tercero. Tiene un único bien: un departamento valuado en 50.000 dólares.

\textbf{Escenario A (todos quirografarios).} El pasivo total es de 60.000 y el activo de 50.000: hay insolvencia. Cada acreedor cobrará proporcionalmente sobre los 50.000.

\textbf{Escenario B (uno tiene hipoteca por 50.000).} El acreedor hipotecario ejecuta primero el departamento, cobra sus 50.000 y los otros dos no cobran nada, porque el bien con garantía real tiene preferencia.
\end{example}
% TODO: contenido incompleto o ambiguo en las notas originales (en el escenario B la hipoteca es «por 50.000», pero las deudas individuales planteadas son de 10.000, 20.000 y 30.000; se conservó el planteo tal como figura)

\section{Bienes inembargables (Art. 744 CCCN)}

Razones humanitarias justifican la exclusión de ciertos bienes de la garantía común.

\begin{tabla}
\begin{tabularx}{\textwidth}{P{8.4cm}Y}
\toprule
\textbf{Bien inembargable} & \textbf{Fundamento} \\
\midrule
Ropa, muebles e inmuebles de uso indispensable del deudor, cónyuge, conviviente e hijos & Necesidades básicas del hogar \\
Instrumentos necesarios para ejercer la profesión, arte u oficio & Sin ellos, la persona no tendría manera de sobrevivir \\
Sepulcros afectados a su destino (salvo que la deuda sea por su construcción o reparación) & Respeto a la memoria de los difuntos \\
Bienes afectados a cualquier religión reconocida por el Estado (templos, objetos de culto) & Libertad religiosa \\
Derechos de usufructo, uso y habitación (en los términos del Código) & Carácter personalísimo de estos derechos \\
Indemnizaciones por daño moral o lesiones a la integridad psicofísica & Compensación de daños personalísimos \\
Indemnizaciones por alimentos al cónyuge, conviviente e hijos en caso de homicidio & Protección de la familia del fallecido \\
Otros declarados inembargables por leyes especiales (por ejemplo, el salario mínimo vital y móvil) & Sustento mínimo para la vida \\
\bottomrule
\end{tabularx}
\end{tabla}

\begin{important}
El telón de fondo de todas estas exclusiones es proteger la dignidad humana en sus dimensiones más fundamentales.
\end{important}

\section{La protección de la vivienda (Arts. 244 y ss. CCCN)}

El Código Civil y Comercial reemplazó la denominación «bien de familia» (Ley 14.394) por el concepto de «protección de la vivienda».

\begin{norma}{Artículo 244 CCCN --- Afectación}
Puede afectarse al régimen previsto en este Capítulo, un inmueble destinado a vivienda, por su totalidad o hasta una parte de su valor. Esta protección no excluye la concurrencia de otros acreedores.
\end{norma}

\begin{tabla}
\begin{tabularx}{\textwidth}{B{3.4cm}Y}
\toprule
Aspecto & \textbf{Régimen} \\
\midrule
Requisito & Inscripción en el Registro de la Propiedad Inmueble (trámite sencillo) \\
Efecto & El inmueble queda inembargable frente a las deudas posteriores a la inscripción \\
Deudas anteriores & Los acreedores con deudas anteriores no se ven afectados \\
Límite cuantitativo & Puede afectarse la totalidad o solo parte del valor \\
Límite de unidades & No puede afectarse más de un inmueble \\
\bottomrule
\end{tabularx}
\end{tabla}

\section{Acciones de los acreedores}

\subsection{Medidas preventivas}

Antes de llegar a la ejecución, los acreedores disponen de medidas preventivas para conservar la garantía:

\begin{tabla}
\begin{tabularx}{\textwidth}{B{3.4cm}YY}
\toprule
Medida & \textbf{Descripción} & \textbf{Ejemplo} \\
\midrule
Embargo preventivo & Afecta un bien determinado para que no salga del patrimonio del deudor & Trabar embargo sobre la casa del deudor que está por venderla antes de obtener sentencia \\
Inhibición general de bienes & Impide al deudor enajenar cualquier bien cuando no se conocen bienes concretos & Anotar en el registro una inhibición general cuando el deudor no tiene bienes identificables \\
\bottomrule
\end{tabularx}
\end{tabla}

Los deudores pueden plantearse la estrategia de sacar bienes de su patrimonio y transformarlos en dinero en efectivo, de tal manera que los acreedores, cuando vengan a cobrar, no encuentren bienes a su nombre. Estas medidas sirven para contrarrestar esa estrategia.

\subsection{Acción directa (Arts. 736 y 737 CCCN)}
% TODO: contenido incompleto o ambiguo en las notas originales (el índice del apunte menciona «Art. 737 y 739 CCCN» para las acciones directa y subrogatoria; el desarrollo cita «Art. 736 y ss.» para la acción directa y «Art. 739» para la subrogatoria)

\begin{definition}{Acción directa}
Es la acción que se le da al acreedor para \textbf{cobrar lo que un tercero debe a su deudor}, hasta el importe del propio crédito.
\end{definition}

Sus caracteres son los siguientes:

\begin{tabla}
\begin{tabularx}{\textwidth}{B{3.6cm}Y}
\toprule
Carácter & \textbf{Descripción} \\
\midrule
Propia & El acreedor la ejerce por derecho propio \\
Exclusiva & En beneficio exclusivo del acreedor que la ejerce, no de todos \\
Excepcional & Solo procede en los casos expresamente previstos por la ley \\
Restrictiva & De interpretación restrictiva \\
Presupone homogeneidad & Los créditos deben ser de la misma naturaleza (por ejemplo, dinerarios) \\
\bottomrule
\end{tabularx}
\end{tabla}

\begin{example}{Acción directa con créditos homogéneos}
Supóngase que se prestó dinero a una persona y que esta, a su vez, se lo prestó a una tercera. Los créditos son homogéneos, porque el propio dinero del acreedor pasó de uno a otro. El acreedor va directamente contra el deudor de su deudor y le ejecuta la deuda, pero solo hasta el importe de su crédito: si prestó doscientos mil pesos a su deudor y este le prestó a otro trescientos mil, puede reclamarle al tercero doscientos mil; los cien mil restantes se los reclamará su deudor.
\end{example}

Esta figura aparece también denominada en las fuentes originales «acción de delegación directa».
% TODO: verificar referencia presente en las fuentes originales (denominación «acción de delegación directa»).

\begin{formula}{Art.~737 CCCN --- Requisitos de la acción directa}
\begin{itemize}
    \item Existencia de un crédito exigible del acreedor contra su propio deudor.
    \item Existencia de una deuda correlativa exigible del tercero demandado a favor del deudor.
    \item Homogeneidad de ambos créditos entre sí.
    \item Ninguno de los dos créditos debe estar embargado ni ser indisponible.
    \item El deudor debe ser notificado de la demanda.
\end{itemize}
\end{formula}

\subsection{Acción subrogatoria (Arts. 739 y 740 CCCN)}

\begin{definition}{Acción subrogatoria}
El acreedor ocupa el lugar de su deudor y ejerce un derecho que le es propio a este, cuando el deudor está inactivo y esa omisión perjudica el cobro del crédito.
\end{definition}

\begin{example}{La cochera heredada}
Supóngase que se prestaron doscientos mil pesos a una persona. Esta tiene un familiar que fallece y le deja en herencia una cochera que vale cuatrocientos mil pesos. El deudor, considerando que el acreedor se quedaría con la cochera apenas la tuviera, decide no iniciar la sucesión y permanece inactivo.

El acreedor, que advierte que su deudor está inactivo y tiene la posibilidad de exigir judicialmente un crédito, y que esa inactividad lo perjudica, se subroga: ocupa su lugar e inicia la sucesión. La cochera ingresa al patrimonio del deudor y el acreedor traba embargo sobre ese bien para cobrarse. Pero aquí no tiene ningún privilegio especial: cualquier otro acreedor embargante tendrá también preferencia.
\end{example}

\begin{formula}{Arts.~739 y 740 CCCN --- Acción subrogatoria}
\textbf{Art.~739.} El acreedor de un crédito cierto, exigible o no, puede ejercer judicialmente los derechos patrimoniales de su deudor si este es remiso en hacerlo y esa omisión afecta el cobro de su acreencia. El acreedor no goza de preferencia alguna sobre los bienes obtenidos por ese medio.

\smallskip
\textbf{Art.~740.} Los derechos y acciones del deudor excluidos de la subrogación son los inherentes a su persona, los no patrimoniales, los patrimoniales no susceptibles de ejecución forzada y los excluidos por disposición legal.
\end{formula}

El crédito o la acción recuperados ingresan al patrimonio del deudor; por eso la acción subrogatoria beneficia a todos los acreedores, mientras que en la acción directa resulta beneficiado únicamente el acreedor que actúa.

\subsubsection{Comparación entre la acción directa y la acción subrogatoria}

\begin{tabla}
\begin{tabularx}{\textwidth}{B{3.4cm}YY}
\toprule
Criterio & \textbf{Acción directa} & \textbf{Acción subrogatoria} \\
\midrule
Nombre en el que se actúa & Nombre propio del acreedor & Nombre del deudor (subrogación) \\
Beneficiario & Solo el acreedor que la ejerce & Todos los acreedores del deudor \\
Privilegio & Sí: cobra directamente hasta su crédito & No: el bien ingresa al patrimonio; todos pueden ejecutarlo \\
Presupuesto & Créditos homogéneos y caso legal expreso & Inactividad del deudor que perjudica la acreencia \\
Límite & Hasta el importe del propio crédito & Todo lo recuperado ingresa al patrimonio del deudor \\
Normas & Arts. 736 y 737 CCCN & Arts. 739 y 740 CCCN \\
\bottomrule
\end{tabularx}
\end{tabla}

\begin{important}
La acción subrogatoria, a diferencia de la acción directa, no otorga ningún privilegio y beneficia a todos los acreedores.
\end{important}

\subsubsection{Derechos excluidos de la acción subrogatoria}

No se pueden subrogar derechos vinculados con la personalidad ni de ejercicio absolutamente discrecional por el titular (Arts. 741--742 CCCN). Como ejemplo se menciona el reconocimiento de un hijo o las cuestiones de puro ejercicio de la voluntad familiar. Según otra de las fuentes, los derechos excluidos son los inherentes a la persona del deudor, los no patrimoniales, los patrimoniales no susceptibles de ejecución forzada y los excluidos por disposición legal (art. 740 CCCN).
% TODO: verificar consistencia entre fuentes originales: una fuente cita los arts. 741 y 742 y otra el art. 740 para los derechos excluidos de la subrogación.


\subsection{Acción de simulación y acción revocatoria (inoponibilidad)}
% TODO: contenido incompleto o ambiguo en las notas originales (el índice del apunte asocia esta sección con «nulidad relativa», concepto que no se desarrolla en el cuerpo de las notas)

\begin{tabla}
\begin{tabularx}{\textwidth}{B{2.6cm}P{3.2cm}YY}
\toprule
Acción & \textbf{Denominación en el CCCN} & \textbf{Presupuesto} & \textbf{Efecto} \\
\midrule
De simulación & Nulidad por simulación & El deudor aparenta sacar bienes pero los conserva (por ejemplo, venta ficticia al hermano/a sin contraprestación) & Nulidad del acto; la cosa vuelve al patrimonio del deudor \\
Revocatoria / pauliana & Inoponibilidad (CCCN) & El deudor vende realmente su bien a un cómplice para defraudar a sus acreedores & El acto es válido entre partes pero inoponible al acreedor; puede ejecutar el bien \\
\bottomrule
\end{tabularx}
\end{tabla}

Ambas acciones se distinguen mediante casos paralelos.

\begin{example}{Simulación}
El deudor llama a su hermana y le dice «simulemos que te vendo mi casa». La hermana no tiene dinero ni bienes, pero aparece como compradora. Cuando el acreedor va a ejecutar al deudor, la casa no está a su nombre. El acreedor alega que el acto fue simulado; el juez decreta la nulidad, las cosas vuelven al estado anterior y el acreedor puede cobrarse.
\end{example}

\begin{example}{Inoponibilidad}
El deudor vende efectivamente su casa ---el dinero ingresa a su patrimonio---, pero el acreedor demuestra que el comprador era cómplice del deudor. Al revocar la venta, esta no se anula: es inoponible al acreedor defraudado, que ejecuta el bien como si la venta no hubiera existido. El adquirente cómplice pierde el bien, pero tiene acción para cobrarse contra el deudor originario.
\end{example}

\section{Cuadro Cornell: Garantía, acreedores y acciones}

\cornellinicio
\begin{longtable}{|Q|Z|}
\hline
\rowcolor{gray!15}\textbf{Preguntas / palabras clave} & \textbf{Notas / respuestas} \\ \hline
\endfirsthead
\hline
\rowcolor{gray!15}\textbf{Preguntas / palabras clave} & \textbf{Notas / respuestas} \\ \hline
\endhead

\cornellpregunta{¿Qué bienes del deudor están sujetos a ejecución?}{%
Según el Art. 743 CCCN, todos los que integran su patrimonio a la fecha de la sentencia (o en el que se hubiese determinado) y los que adquiera hasta la extinción de la obligación. El acreedor puede exigir la venta judicial solo en la medida necesaria para satisfacer su crédito, y todos los acreedores ejecutan en posición igualitaria salvo causa legal de preferencia.}

\cornellpregunta{¿Qué tipos de acreedores existen y cómo afecta eso al cobro?}{%
Con privilegio especial (sobre determinados bienes, por ejemplo el laboral), con privilegio general (sobre todo el patrimonio), con garantía real (hipoteca o prenda, elegida voluntariamente) y quirografarios (sin preferencia, cobran en igualdad sobre el remanente). En caso de insolvencia, cobran en ese orden.
% TODO: contenido incompleto o ambiguo en las notas originales (el orden de cobro «en ese orden» figura solo en el cuadro Cornell del apunte; el desarrollo no precisa la prelación entre las distintas preferencias)
}

\cornellpregunta{¿Qué bienes son inembargables? (Art. 744 CCCN)}{%
Ropa y muebles indispensables del hogar, instrumentos del oficio o profesión, sepulcros (salvo deuda de construcción), bienes de culto, ciertos derechos de usufructo, uso y habitación, indemnizaciones por lesiones o daño moral, alimentos por homicidio, y los declarados así por leyes especiales (salario mínimo). El fundamento común es proteger la dignidad humana en sus dimensiones más fundamentales.}

\cornellpregunta{¿Cómo se protege la vivienda frente a las deudas?}{%
Inscribiendo el inmueble en el Registro de la Propiedad Inmueble (Arts. 244 y ss. CCCN). Desde la inscripción, las deudas posteriores no pueden ejecutarlo. Solo se protege un inmueble, en su totalidad o hasta una parte de su valor. Las deudas anteriores a la inscripción no se ven afectadas.}

\cornellpregunta{¿Qué medidas preventivas tienen los acreedores antes de la ejecución?}{%
\textbf{Embargo preventivo:} afecta un bien determinado para que no salga del patrimonio del deudor. \textbf{Inhibición general de bienes:} impide al deudor enajenar cualquier bien cuando no se conocen bienes concretos. Sirven para contrarrestar la estrategia del deudor de convertir sus bienes en dinero en efectivo para que los acreedores no encuentren bienes a su nombre.}

\cornellpregunta{¿Qué es la acción directa?}{%
Permite al acreedor cobrar directamente lo que un tercero le debe a su deudor (Art. 736 CCCN). La ejerce en nombre propio, solo lo beneficia a él y opera hasta el límite de su crédito. Requiere homogeneidad de créditos y un caso legal expreso.}

\cornellpregunta{¿Qué es la acción subrogatoria?}{%
El acreedor ocupa el lugar de su deudor inactivo y ejerce los derechos patrimoniales de este en su nombre (Art. 739 CCCN). El bien recuperado ingresa al patrimonio del deudor y beneficia a todos los acreedores, no solo al subrogante.}

\cornellpregunta{¿Qué diferencias hay entre la acción directa y la acción subrogatoria?}{%
\begin{itemize}[nosep,leftmargin=1.2em]
\item \textbf{Directa:} se actúa en nombre propio; beneficia solo al acreedor que la ejerce; con privilegio (cobra directamente hasta su crédito); requiere créditos homogéneos y caso legal expreso.
\item \textbf{Subrogatoria:} se actúa en nombre del deudor; beneficia a todos los acreedores; sin privilegio (el bien ingresa al patrimonio y todos pueden ejecutarlo); requiere inactividad del deudor que perjudique la acreencia.
\end{itemize}}

\cornellpregunta{¿Qué derechos no pueden ejercerse mediante la acción subrogatoria?}{%
Los derechos vinculados con la personalidad y los de ejercicio absolutamente discrecional por el titular (Arts. 741--742 CCCN); por ejemplo, el reconocimiento de un hijo o cuestiones de puro ejercicio de la voluntad familiar.}

\cornellpregunta{¿Qué diferencia hay entre la acción de simulación y la de inoponibilidad?}{%
En la simulación el acto es ficticio (el deudor aparenta vender pero conserva el bien); el efecto es la nulidad. En la inoponibilidad (acción revocatoria o pauliana) el acto es real pero fraudulento: el acreedor lo hace inoponible a sí mismo y puede ejecutar el bien aunque el tercero lo haya adquirido realmente.}

\cornellresumen{Los bienes del deudor responden por sus obligaciones y los acreedores los ejecutan en posición igualitaria salvo causa legal de preferencia (Art. 743 CCCN). Los acreedores pueden tener privilegio especial, privilegio general, garantía real o ser quirografarios. La garantía tiene límites: los bienes inembargables (Art. 744) y la protección de la vivienda (Art. 244). Para conservar y recuperar la garantía, los acreedores cuentan con medidas preventivas (embargo preventivo, inhibición general de bienes) y con las acciones directa, subrogatoria, de simulación y de inoponibilidad.}

\end{longtable}
\cornellfin

% =====================================================================
% CAPÍTULO 5 — SÍNTESIS GENERAL DE LA UNIDAD
% =====================================================================

\section{Síntesis de la parte}
\subsection{Recorrido de la parte}

Esta parte recorrió el patrimonio, el patrimonio como garantía, los distintos tipos de acreedores y los distintos tipos de acciones de los acreedores. Desde el patrimonio como garantía existe la posibilidad de ejercer acciones; asimismo, hay bienes que quedan excluidos de ella, y que el Código enuncia en los artículos 744 y 244.

\subsection{Síntesis final}

Esta parte construyó la arquitectura conceptual sobre la que descansa buena parte del Derecho Privado argentino. El patrimonio es el puente entre la persona y el mercado: como atributo de la personalidad, le da contenido económico al sujeto; como garantía común (Art. 242 CCCN), le da sustento a todo el sistema crediticio.

Para entender de qué está hecho ese patrimonio, el Código clasifica los bienes y cosas según múltiples criterios ---muebles e inmuebles, fungibles y no fungibles, divisibles e indivisibles, principales y accesorias, frutos y productos---, cada uno con consecuencias jurídicas precisas que atraviesan el Derecho Real, Obligacional, Contractual, de Familia y de Sucesiones.

El sistema de protección articula dos tensiones: la del acreedor, que necesita ejecutar el patrimonio para cobrar, y la del deudor, que no puede ser reducido a la indigencia. Esta tensión se resuelve mediante los bienes inembargables y la protección de la vivienda.

Finalmente, las acciones directa, subrogatoria, de simulación y de inoponibilidad son los instrumentos procesales que permiten al acreedor conservar y recuperar la garantía frente a estrategias de vaciamiento patrimonial.
\partintro{Si el sujeto es quien entabla la relación y el objeto es la materia sobre la que recae, la causa explica cómo surgen las relaciones jurídicas. Esta parte estudia los hechos jurídicos y el acto voluntario, la imputabilidad, el acto jurídico y sus elementos (objeto, causa y forma), los instrumentos, la clasificación y las modalidades de los actos y, por último, los efectos de los actos jurídicos respecto de las partes, los representantes, los sucesores y los terceros.}
\part{La causa de la relación jurídica: hechos y actos jurídicos}
\chapter{Fundamentos de los Hechos y Actos Jurídicos}

\begin{note}
Los primeros capítulos de esta parte abordan el tercer elemento de la relación jurídica: la causa, es decir, el origen del cual emergen las relaciones jurídicas. El recorrido conceptual parte de la clasificación general de los hechos jurídicos, avanza hacia la distinción entre hechos voluntarios e involuntarios (discernimiento, intención y libertad) y culmina con el estudio de la imputabilidad, es decir, de las consecuencias que se atribuyen a quien realiza un acto voluntario, así como del tratamiento excepcional que reciben los actos involuntarios a través del principio de equidad.
\end{note}

\section{Introducción: la causa como tercer elemento de la relación jurídica}

La teoría de los hechos y actos jurídicos constituye el tercer elemento de la relación jurídica, luego de haberse estudiado el sujeto y el objeto. Mientras el sujeto es quien entabla la relación y el objeto es la materia sobre la que ésta recae, la causa explica cómo surgen las relaciones jurídicas.

\begin{table}[H]
\centering
\caption{Elementos de la relación jurídica}
\begin{tabularx}{\textwidth}{l l X}
\toprule
\textbf{Elemento} & \textbf{Definición} & \textbf{Contenido} \\
\midrule
Sujeto & Quien entabla la relación & Persona humana o jurídica \\
Objeto & Materia sobre la que recae la relación & Bienes y cosas \\
Causa & De dónde surge la relación jurídica & \textbf{Hechos y actos jurídicos} \\
\bottomrule
\end{tabularx}
\end{table}

Esta teoría proyecta efectos sobre el conjunto del derecho civil y comercial, en la medida en que toda disciplina jurídica se ocupa, en definitiva, de hechos que dan origen a relaciones y situaciones jurídicas.

\section{Ubicación metodológica en el Código Civil y Comercial}

En el Código de Vélez, los hechos y actos jurídicos se regulaban en el Libro Segundo, Sección Segunda, dedicada a los derechos personales. El artículo 896 marcaba el inicio de la regulación de los hechos jurídicos, distanciada de los derechos reales y de otras áreas del código.

\begin{table}[H]
\centering
\caption{Cambio metodológico: Código de Vélez frente al CCCN}
\begin{tabularx}{\textwidth}{X X}
\toprule
\textbf{Código de Vélez (anterior)} & \textbf{CCCN actual} \\
\midrule
Libro Segundo -- Derechos personales & Libro Primero -- Parte General (metodológicamente más correcto) \\
Art. 896: iniciaba los hechos jurídicos & Desarrolla la teoría en nueve capítulos integrados a la parte general \\
Distanciado de los derechos reales y otras áreas & Integra la Parte General y fue bien recibido por la doctrina \\
\bottomrule
\end{tabularx}
\end{table}

El Código Civil y Comercial de la Nación (CCCN) desarrolla la teoría en nueve capítulos: disposiciones generales, los vicios (error, dolo, violencia), los actos jurídicos propiamente dichos, los vicios de los actos, las modalidades, la representación y la ineficacia. Este cambio metodológico constituye un acierto del nuevo Código, en la medida en que integra la parte general del derecho civil. Concretamente, la regulación de los hechos y actos jurídicos se ubica en el Libro Primero, Título Cuarto del Código.

\section{Definición de hecho jurídico en el CCCN}

\begin{definition}{Art. 257 CCCN -- Hecho jurídico}
El hecho jurídico es el acontecimiento que, conforme al ordenamiento jurídico, produce el nacimiento, modificación y extinción de relaciones o situaciones jurídicas.
\end{definition}

\begin{table}[H]
\centering
\caption{Comparación entre el art. 257 CCCN y el art. 896 del Código de Vélez}
\begin{tabularx}{\textwidth}{l X X}
\toprule
\textbf{Aspecto} & \textbf{CCCN (Art. 257)} & \textbf{Código de Vélez (Art. 896)} \\
\midrule
Término clave & Produce (certeza) & Susceptible de producir (posibilidad) \\
Referencia al orden jurídico & Sí: ``conforme al ordenamiento jurídico'' & No hacía esta referencia expresa \\
Efectos abarcados & Relaciones o situaciones jurídicas & Derechos u obligaciones solamente \\
\bottomrule
\end{tabularx}
\end{table}

El CCCN incorpora la expresión ``conforme al ordenamiento jurídico'', señalando que es el ordenamiento jurídico el que genera estas consecuencias. Esta afirmación no debe entenderse exclusivamente desde una visión positivista, como si únicamente el ordenamiento jurídico las generase, pues existe al respecto un debate doctrinario abierto.

El reemplazo del término ``susceptible de producir'' por ``produce'' resulta más preciso, en tanto al derecho le interesan los acontecimientos que efectivamente producen consecuencias jurídicas.

\begin{example}{El hecho jurídico frente al acontecimiento cotidiano}
Tomar un vaso de agua no constituye un hecho jurídico. Arrojarlo por la ventana y dañar a un transeúnte sí lo es, en tanto genera una consecuencia jurídica relevante para el ordenamiento.
\end{example}

\section{Clasificación general de los hechos jurídicos}
\label{sec:clasifhechos}

El cuadro general de clasificación de los hechos jurídicos resulta central para el desarrollo de la materia.

\begin{table}[H]
\centering
\caption{Hechos jurídicos: clasificación general}
\begin{tabularx}{\textwidth}{X X}
\toprule
\textbf{Hechos de la naturaleza} & \textbf{Hechos humanos} \\
\midrule
Granizo, terremoto, aluvión. Relevan por sus efectos, no por su origen. &
\textbf{Involuntarios}: sin discernimiento, intención o libertad. \\
\cmidrule(l){2-2}
 & \textbf{Voluntarios} (con discernimiento, intención y libertad): \\
 & \quad -- \textbf{Lícitos}: simples actos lícitos / actos jurídicos \\
 & \quad -- \textbf{Ilícitos}: delitos (con dolo) / cuasidelitos (con culpa) \\
\bottomrule
\end{tabularx}
\end{table}

\section{Hechos voluntarios e involuntarios: importancia de la distinción}

\begin{example}{Caso del ataque epiléptico en la estación de subte}
Una persona sufre un ataque de epilepsia en una estación de subte y se desvanece, golpeando a otra persona, que a su vez golpea a una tercera que casi cae a las vías del tren.

Conclusión jurídica: el hecho se origina de manera involuntaria. No hubo voluntad de empujar; la persona sufrió un ataque de epilepsia y se desvaneció. Desde el punto de vista jurídico, este hecho es tratado de manera diferente al supuesto en que existe voluntad, pues quienes tienen dominio de sus actos se hacen cargo de ellos, existiendo mayor responsabilidad cuando media voluntad.
\end{example}

\section{Actos ilícitos: dolo y culpa}

Los actos ilícitos son los prohibidos por la ley que generan daños. A diferencia del derecho penal, donde los delitos se encuentran tipificados, en el derecho civil los ilícitos comprenden todos los actos que causan un daño, presumiéndose que dicho daño es siempre antijurídico.

\begin{table}[H]
\centering
\caption{Delito civil y cuasidelito civil}
\begin{tabularx}{\textwidth}{l X}
\toprule
\textbf{Figura} & \textbf{Definición} \\
\midrule
Delito civil & Acto ilícito cometido con dolo (intención de dañar o manifiesta indiferencia por los intereses ajenos) \\
Cuasidelito civil & Acto ilícito cometido con culpa (negligencia, imprudencia o impericia, sin intención de dañar) \\
\bottomrule
\end{tabularx}
\end{table}

\begin{note}
Los términos ``delito civil'' y ``cuasidelito civil'' no aparecen literalmente en el texto del CCCN. Se trata de una clasificación histórica elaborada por la doctrina.

Cabe señalar, asimismo, una crítica metodológica: los actos ilícitos se encuentran regulados en el Libro Tercero del CCCN (Derechos Personales -- Responsabilidad Civil), cuando, al menos en sus elementos centrales, deberían haberse tratado en el Libro Primero, junto a la teoría general de los hechos jurídicos.
\end{note}

\begin{definition}{Art. 1724 CCCN -- Factores subjetivos de atribución (dolo y culpa)}
El dolo se configura por la producción de un daño de manera intencional o con manifiesta indiferencia por los intereses ajenos.

La culpa consiste en la omisión de la diligencia debida según la naturaleza de la obligación y las circunstancias de las personas, el tiempo y el lugar. Comprende la imprudencia, la negligencia y la impericia en el arte o profesión.
\end{definition}

El CCCN amplía la noción clásica romanista del dolo (intención de dañar), incorporando también la ``manifiesta indiferencia''.

\begin{example}{Caso del conductor que ignora a la anciana}
Un conductor circula a 120 km/h y advierte que una anciana cruza correctamente con el semáforo en verde para peatones. Pese a verla, no realiza ningún intento de esquivarla ni de frenar, y la mata.

Análisis: no puede determinarse con certeza que el conductor haya querido matarla con premeditación, pero sí mostró una ``manifiesta indiferencia'' ante los intereses ajenos. Por lo tanto, se configura el dolo por indiferencia, que el CCCN incorpora expresamente.
\end{example}

\begin{table}[H]
\centering
\caption{Formas de culpa}
\begin{tabularx}{\textwidth}{l X}
\toprule
\textbf{Forma de culpa} & \textbf{Definición} \\
\midrule
Negligencia & Torpeza, falta de cuidado o diligencia debida. Ej.: un médico que no se lava las manos antes de una intervención. \\
Imprudencia & Realizar algo osado que no correspondía, asumiendo un riesgo innecesario sin intención de dañar. \\
Impericia & Falta de conocimientos técnicos propios de la profesión u oficio. Ej.: un médico que interviene sin los conocimientos necesarios. \\
\bottomrule
\end{tabularx}
\end{table}

\section{Actos lícitos: simples actos lícitos y actos jurídicos}

\begin{definition}{Art. 258 CCCN -- Simple acto lícito}
La acción voluntaria no prohibida por la ley, de la que resulta alguna adquisición, modificación o extinción de relaciones o situaciones jurídicas.
\end{definition}

\begin{definition}{Art. 259 CCCN -- Acto jurídico}
El acto jurídico es el acto voluntario lícito que tiene por fin inmediato la adquisición, modificación o extinción de relaciones o situaciones jurídicas.
\end{definition}

\begin{table}[H]
\centering
\caption{Simple acto lícito frente a acto jurídico}
\begin{tabularx}{\textwidth}{l X X}
\toprule
\textbf{Criterio} & \textbf{Simple acto lícito (Art. 258)} & \textbf{Acto jurídico (Art. 259)} \\
\midrule
Voluntad & Sí, voluntario & Sí, voluntario \\
Licitud & Sí, no prohibido por la ley & Sí, lícito \\
Fin inmediato & No busca producir consecuencias jurídicas directamente & Sí: fin inmediato de producir consecuencias jurídicas \\
Ejemplo & Pasear por la calle (puede derivar en consecuencias jurídicas imprevistas) & Contrato de mutuo (préstamo de dinero con plazo e interés pactados) \\
\bottomrule
\end{tabularx}
\end{table}

\begin{example}{El contrato de mutuo como acto jurídico}
En un contrato de mutuo, una persona decide prestarle dinero a otra, fijando el plazo y, eventualmente, una tasa de interés. Todo ello se establece en un contrato firmado por personas humanas que resuelven voluntariamente realizar un acto lícito --prestarse dinero-- con el fin inmediato de generar una relación jurídica.
\end{example}

% =====================================================================
% CAPÍTULO 2
% =====================================================================
\chapter{Teoría del Acto Voluntario}

\section{Las condiciones del acto voluntario}

\begin{definition}{Art. 260 CCCN -- Acto voluntario}
El acto voluntario es el ejecutado con discernimiento, intención y libertad, que se manifiesta por un hecho exterior.
\end{definition}

\begin{table}[H]
\centering
\caption{Condiciones del acto voluntario}
\begin{tabularx}{\textwidth}{l l X}
\toprule
\textbf{Condición} & \textbf{Tipo} & \textbf{Descripción} \\
\midrule
Discernimiento & Interna & Aptitud para diferenciar lo bueno de lo malo, lo verdadero de lo falso, lo justo de lo injusto y sus consecuencias \\
Intención & Interna & Propósito de llevar a cabo un acto determinado, de manera consciente y dirigida a ese fin \\
Libertad & Interna & Posibilidad de actuar de acuerdo a la propia conveniencia, deseo y convicción, sin coacción \\
Exteriorización & Externa & Manifestación visible del acto: expresa (oral, escrita, signos inequívocos) o tácita \\
\bottomrule
\end{tabularx}
\end{table}

\section{El discernimiento}

El discernimiento es la aptitud para diferenciar lo bueno de lo malo, lo verdadero de lo falso, lo justo de lo injusto y sus consecuencias. Se trata de una comprensión genérica, una capacidad o aptitud para comprender lo que está sucediendo.

\begin{table}[H]
\centering
\caption{Causas obstativas del discernimiento (Art. 261 CCCN)}
\begin{tabularx}{\textwidth}{l X}
\toprule
\textbf{Causa} & \textbf{Detalle} \\
\midrule
Privación de razón (permanente o transitoria) & Ej.: golpe en la cabeza que genera un estado de inconsciencia durante el cual se realizó un acto perjudicial. No requiere sentencia judicial de incapacidad. \\
Menor de diez (10) años -- para actos ilícitos & El código presume que antes de esa edad no hay discernimiento para comprender lo injusto y sus consecuencias. \\
Menor de trece (13) años -- para actos lícitos & Para los actos lícitos se exige mayor madurez cognitiva; la edad mínima es superior. \\
\bottomrule
\end{tabularx}
\end{table}

\begin{important}
El discernimiento no debe confundirse con la capacidad. La capacidad es un concepto jurídico integral que se refiere a la aptitud para entablar relaciones jurídicas, mientras que el discernimiento es un componente del acto voluntario. Puede haber discernimiento sin capacidad: es el caso de las personas menores de edad.
\end{important}

\begin{example}{Discernimiento sin capacidad jurídica}
Un niño de diez años puede tener discernimiento suficiente para un acto lícito, pero no necesariamente capacidad jurídica para celebrar ese acto.
\end{example}

\begin{table}[H]
\centering
\caption{Discernimiento frente a capacidad jurídica}
\begin{tabularx}{\textwidth}{X X}
\toprule
\textbf{Discernimiento (Art. 261 CCCN)} & \textbf{Capacidad jurídica} \\
\midrule
Elemento del acto voluntario. Se evalúa en cada caso concreto. No depende de sentencia judicial. & Concepto jurídico integral. Puede requerir sentencia judicial. Se adquiere con la mayoría de edad (18 años). \\
\bottomrule
\end{tabularx}
\end{table}

\section{La intención y sus vicios: error y dolo}

La intención es el propósito de llevar a cabo un acto determinado con consciencia. Mientras el discernimiento es la característica general que puede estar presente o no en cada acto, la intención es el propósito de realizar ese acto concreto, preciso y específico.

\begin{table}[H]
\centering
\caption{Vicios de la intención}
\begin{tabularx}{\textwidth}{l X X}
\toprule
\textbf{Vicio} & \textbf{Definición} & \textbf{Ejemplo} \\
\midrule
Error & Ignorancia o falso conocimiento de la realidad que vicia la voluntad & Creer que se vende una casa cuando la otra parte entiende que es una donación. Error sobre la naturaleza del acto. \\
Dolo (vicio de la voluntad) & Maniobra engañosa por la que se induce a la otra parte a celebrar el acto & Alquilar una propiedad mostrando fotos falsas; al llegar, la casa está en mal estado y lejos de la playa. \\
\bottomrule
\end{tabularx}
\end{table}

La palabra ``dolo'' posee al menos dos significados en el derecho civil: (1) como intención de dañar, elemento subjetivo de la responsabilidad civil, y (2) como maniobra engañosa que vicia la voluntad.

El error puede recaer sobre los objetos, las personas, las causas o las cualidades de la cosa, dando lugar a distintas circunstancias que se estudian en el capítulo~\ref{ch:error}.

Respecto del dolo como vicio, la cuestión central consiste en determinar cuándo el engaño tiene entidad suficiente para afectar la intención, es decir, cuándo puede afirmarse que el acto no se habría realizado de no mediar el engaño.

\section{La libertad y el vicio de violencia}

La libertad es la posibilidad de llevar a cabo el acto de acuerdo con la propia conveniencia, deseo y convicción.

\begin{table}[H]
\centering
\caption{Formas de violencia como causa obstativa de la libertad}
\begin{tabularx}{\textwidth}{l X}
\toprule
\textbf{Forma de violencia} & \textbf{Definición} \\
\midrule
Fuerza física irresistible & Uso de la fuerza física de manera irresistible sobre el sujeto para que realice el acto. Ej.: torturar a alguien para que firme. \\
Intimidación (violencia moral) & Amenaza de sufrir un mal grave e inminente que no puede contrarrestarse, y que lleva a la persona a realizar el acto sin libertad real. \\
\bottomrule
\end{tabularx}
\end{table}

Es posible comprender, querer y realizar un acto con esa intención, y sin embargo no ser suficientemente libre por encontrarse la persona bajo amenaza. Un acto realizado bajo amenaza no es, por tanto, un acto voluntario.

\section{La exteriorización de la voluntad}

La voluntad debe exteriorizarse para ser considerada tal. Esto puede ocurrir de manera expresa o tácita.

\begin{table}[H]
\centering
\caption{Modalidades de exteriorización de la voluntad}
\begin{tabularx}{\textwidth}{l X X}
\toprule
\textbf{Modalidad} & \textbf{Descripción} & \textbf{Ejemplo} \\
\midrule
Expresa oral & La voluntad se comunica verbalmente & Pactar verbalmente un préstamo \\
Expresa escrita & Instrumentada en documento público o privado & Contrato firmado; escritura pública \\
Signos inequívocos & Gestos que no admiten otra interpretación & Hacer el gesto de pedir un café en un bar \\
Ejecución de hecho material & Realización de un acto que revela la voluntad & Girar dinero como pago de una deuda \\
Tácita & Se infiere de conductas que permiten conocer con certidumbre la voluntad & Serie de actos que evidencian aceptación de una oferta \\
\bottomrule
\end{tabularx}
\end{table}

\begin{definition}{Art. 264 CCCN -- Límite a la voluntad tácita}
Si la ley o la convención exigen que la voluntad sea expresa, no puede suplirse con una expresión tácita.
\end{definition}

Un ejemplo de esta regla es el consentimiento exigido para los actos de disposición de derechos, respecto de los cuales la ley exige que sea expreso, no pudiendo interpretarse un consentimiento tácito en ese contexto.

\section{El silencio como (no) manifestación de voluntad}

\begin{definition}{Art. 263 CCCN -- El silencio como manifestación de voluntad}
El silencio opuesto a actos o a una interrogación no es considerado como una manifestación de voluntad conforme al acto o la interrogación, excepto en los casos en que haya un deber de expedirse que puede resultar de la ley o de la voluntad de las partes.
\end{definition}

En el derecho civil y comercial, el silencio no constituye, en principio, una manifestación de voluntad: quien calla no otorga.

\begin{example}{El silencio frente a una oferta no solicitada}
Si una persona recibe un mensaje de texto que indica ``te has ganado un auto; a partir de mañana se descuentan \$20.000 de tu tarjeta de crédito durante 60 meses, salvo que contestes que no lo querés en 24 horas'', y guarda silencio, ese silencio no puede considerarse aceptación. De lo contrario, sería necesario responder permanentemente a las publicidades para aclarar que no se las desea.
\end{example}

\begin{table}[H]
\centering
\caption{Excepciones: cuándo el silencio sí es manifestación de voluntad}
\begin{tabularx}{\textwidth}{l X}
\toprule
\textbf{Excepción} & \textbf{Descripción} \\
\midrule
Deber legal de expedirse & La ley obliga a pronunciarse. Ej.: intimación para reconocer o desconocer la firma de un instrumento. Si se guarda silencio, se tiene por reconocida. \\
Deber convencional de expedirse & Las partes acordaron que el silencio implica determinada manifestación. Ej.: contrato que establece que se renueva automáticamente si la parte no comunica lo contrario con un mes de anticipación. \\
Relaciones entre declaraciones precedentes & Cuando una serie de declaraciones anteriores llevaron razonablemente a la otra parte a esperar una determinada respuesta, y el silencio posterior confirma esa dirección. \\
\bottomrule
\end{tabularx}
\end{table}

% =====================================================================
% CAPÍTULO 3
% =====================================================================
\chapter{Imputabilidad de los Actos Voluntarios}

\section{El deber de reparar: fundamento y presupuestos}

Quienes realizan un acto voluntario deben hacerse cargo de sus consecuencias, en virtud de la relación de causalidad existente entre la acción y el resultado. La cuestión central consiste en determinar hasta dónde se extiende esa responsabilidad.

\begin{definition}{Art. 1716 CCCN -- Deber de reparar}
La violación del deber de no dañar a otro, o el incumplimiento de una obligación, da lugar a la reparación del daño causado, conforme con las disposiciones de este Código.
\end{definition}

\begin{definition}{Art. 1717 CCCN -- Antijuridicidad}
Cualquier acción u omisión que causa un daño a otro es antijurídica si no está justificada.
\end{definition}

\begin{table}[H]
\centering
\caption{Presupuestos de la responsabilidad civil}
\begin{tabularx}{\textwidth}{c l X}
\toprule
\textbf{N.\textdegree{}} & \textbf{Presupuesto} & \textbf{Descripción} \\
\midrule
1 & Daño & Debe existir un daño efectivo a una persona o a su patrimonio. \\
2 & Antijuridicidad & El daño debe ser antijurídico (se presume siempre, salvo causa de justificación -- Art. 1718). \\
3 & Factor de atribución & Subjetivo: dolo o culpa (Art. 1724). Objetivo: la ley establece que se responde independientemente de dolo o culpa (ej.: art. 1757 -- responsabilidad por cosas y actividades riesgosas). \\
4 & Relación de causalidad & El acto voluntario debe ser la causa adecuada del daño producido (nexo causal adecuado). \\
\bottomrule
\end{tabularx}
\end{table}

\begin{example}{Concausalidad: auto que choca a motociclista sin casco}
Un auto choca a una moto cuyo conductor no lleva casco. El accidente le provoca una lesión cerebral severa, agravada precisamente por la falta de casco.

Análisis: quien chocó a la moto tiene una parte fundamental en la causación del daño, pues sin el choque el motociclista no se habría caído. Pero la negligencia del conductor que no llevaba casco también incide en el resultado dañoso, ya que con casco se habrían evitado o reducido las consecuencias. Esto no exime al conductor del auto de responder, pero obliga a analizar cómo incidió cada causa (concausalidad).
\end{example}

\section{Tipos de consecuencias: inmediatas, mediatas y casuales}

\begin{table}[H]
\centering
\caption{Clasificación de las consecuencias del acto}
\begin{tabularx}{\textwidth}{l X X}
\toprule
\textbf{Tipo} & \textbf{Definición} & \textbf{Ejemplo (choque de auto)} \\
\midrule
Inmediatas & Las que acostumbran suceder según el curso normal y ordinario de las cosas (Art. 1726 CCCN). & Auto choca a persona $\rightarrow$ fractura de cadera. Resultado directo del impacto. \\
Mediatas & Resultan de la conexión del hecho con un acontecimiento distinto (Art. 1726 CCCN). & Internación $\rightarrow$ pérdida de ingresos como carpintero durante 20 días. \\
Casuales & No pueden ser previstas, o resultan de circunstancias extraordinarias (Art. 1727 CCCN). & Pérdida de ingresos $\rightarrow$ depresión $\rightarrow$ divorcio $\rightarrow$ hijos con adicciones. Imprevisibles. \\
\bottomrule
\end{tabularx}
\end{table}

\begin{formula}{Arts. 1726 y 1727 CCCN -- ¿Por cuáles consecuencias se responde?}
Se reparan las consecuencias inmediatas y las mediatas previsibles. No se responde por las consecuencias casuales (mediatas no previsibles).
\end{formula}

\section{El criterio de previsibilidad y la graduación de la responsabilidad}

\begin{definition}{Art. 1725 CCCN -- Valoración de la conducta}
Cuanto mayor sea el deber de obrar con prudencia y pleno conocimiento de las cosas, mayor es la diligencia exigible al agente y la valoración de la previsibilidad de las consecuencias.
\end{definition}

Cuanto mayor es lo que una persona conoce, y cuanto mayor es su deber de conocer y obrar con prudencia, mayor es la diligencia que se le exige y, en consecuencia, mayor es también la valoración de la previsibilidad de las consecuencias por las que deberá responder.

\begin{table}[H]
\centering
\caption{Graduación de la diligencia exigible según la situación}
\begin{tabularx}{\textwidth}{X X}
\toprule
\textbf{Situación} & \textbf{Exigencia de diligencia} \\
\midrule
Profesional de la salud tratando un paciente & Alta: se le exige la diligencia propia de su especialidad y el estándar de cuidado de la comunidad médica. \\
Trabajador en planta con material nuclear o contaminante & Muy alta: el riesgo elevado implica mayor previsibilidad exigida. \\
Trabajador en tareas de bajo riesgo & Estándar ordinario de diligencia. \\
\bottomrule
\end{tabularx}
\end{table}

La condición especial o la facultad intelectual particular del agente no se toma en cuenta a los fines de esta graduación, a menos que haya sido específicamente contemplada en un contrato.

\section{El caso fortuito y la fuerza mayor como eximentes}

\begin{definition}{Art. 1730 CCCN -- Caso fortuito y fuerza mayor}
Se considera caso fortuito o fuerza mayor al hecho que no ha podido ser previsto, o que, habiendo sido previsto, no ha podido ser evitado. El caso fortuito o fuerza mayor exime de responsabilidad, excepto disposición en contrario.
\end{definition}

\begin{important}
Cuando algo no puede ser previsto, el supuesto excede el ámbito de lo voluntario y se enmarca dentro del caso fortuito o la fuerza mayor. Se trata de un principio clásico del derecho civil: cuando algo no es voluntario, no se responde.
\end{important}

% =====================================================================
% CAPÍTULO 4
% =====================================================================
\chapter{Responsabilidad por Actos Involuntarios}

\section{Regla general: no responsabilidad por actos involuntarios}

Los actos involuntarios, a diferencia de los voluntarios, no generan responsabilidad en principio. Ello obedece a que, si el agente no obró con discernimiento, intención y libertad, no puede imputársele el deber de responder por las consecuencias del hecho.

\begin{important}
En principio, los actos involuntarios no generan responsabilidad, salvo las excepciones expresamente previstas por el ordenamiento.
\end{important}

El antiguo artículo 907 del Código de Vélez establecía que en ningún caso se respondía por los actos involuntarios.

\section{La evolución: el principio de equidad}

A partir, especialmente, de la jurisprudencia francesa, comenzó a abrirse camino la idea de que la regla de no responsabilidad por actos involuntarios podía resultar injusta en determinados casos.

\begin{example}{El caso que motivó la incorporación de la equidad}
Una persona privada de la razón por una enfermedad mental, con un gran patrimonio, causa un daño enorme a una persona muy pobre con familia numerosa.

Dilema: desde el punto de vista teórico, el hecho es involuntario y no puede imputarse. Pero cabe preguntarse si es justo que una persona muy rica, cuya conducta involuntaria arruina a una familia pobre, no repare absolutamente nada. Aquí aparece el concepto de equidad.
\end{example}

\begin{table}[H]
\centering
\caption{Dos concepciones de equidad}
\begin{tabularx}{\textwidth}{l X}
\toprule
\textbf{Concepción} & \textbf{Contenido} \\
\midrule
Equidad romana / igualitaria & Buscar que las personas tengan lo indispensable para la vida; igualdad material o distributiva. \\
Epieikeia aristotélica (equidad correctiva) & Rectificación de lo justo legal en el caso concreto. Cuando la aplicación estricta de la ley conduce a un resultado manifiestamente injusto, el juez puede apartarse de ella y disponer algo distinto para ese caso. Esta es la concepción que opera en el artículo 1750 CCCN. \\
\bottomrule
\end{tabularx}
\end{table}

\section{El artículo 1750 del CCCN: régimen vigente}

\begin{definition}{Art. 1750 CCCN -- Daños causados por actos involuntarios}
El autor de un daño causado por un acto involuntario responde por razones de equidad. Se aplica lo dispuesto en el artículo 1742. El acto realizado por quien sufre fuerza irresistible no genera responsabilidad para su autor, sin perjuicio de la que corresponde a título personal a quien ejerce esa fuerza.
\end{definition}

\begin{table}[H]
\centering
\caption{Evolución normativa: del art. 907 (mod. 1968) al art. 1750 CCCN}
\begin{tabularx}{\textwidth}{l X X}
\toprule
\textbf{Aspecto} & \textbf{Anterior (Art. 907 mod. 1968)} & \textbf{CCCN (Art. 1750)} \\
\midrule
Regla general & Nunca se respondía & El autor responde por razones de equidad \\
Facultad judicial & ``Los jueces podrán disponer'' & Texto más imperativo; el autor ``responde'' \\
Quantum & Patrimonio del autor y situación de la víctima & Se remite al art. 1742 (atenuación) \\
Debate doctrinario & Incorporado en 1968, debatido desde entonces & Jornadas Nacionales de Derecho Civil (Santa Fe): ¿es posible una indemnización de cero? \\
\bottomrule
\end{tabularx}
\end{table}

\begin{note}
Constituye una cuestión doctrinaria debatida si el juez podría considerar que la equidad implica no establecer ninguna indemnización, frente a la posición de otros autores que sostienen que siempre debe establecerse alguna. La mayoría de la doctrina se inclina por esta última posición, aunque se trata de un debate abierto en el que se ha señalado que podría el juez, por razones de equidad, no fijar indemnización alguna.

Se menciona al Dr. Juan Carlos Palmero, jurista cordobés, como el autor que más ha profundizado este tema en Argentina, con una publicación específica sobre el daño involuntario.
\end{note}

\section{Criterios de la equidad para la atenuación (Art. 1742 CCCN)}

\begin{definition}{Art. 1742 CCCN -- Atenuación de la responsabilidad}
El juez, al fijar la indemnización, puede atenuarla si es equitativo en función del patrimonio del deudor, la situación personal de la víctima y las circunstancias del hecho. Esta facultad no es aplicable en caso de dolo del responsable.
\end{definition}

\begin{table}[H]
\centering
\caption{Factores para la determinación de la indemnización equitativa}
\begin{tabularx}{\textwidth}{l X}
\toprule
\textbf{Factor} & \textbf{Descripción} \\
\midrule
Patrimonio del causante del daño & ¿Cuánta capacidad económica tiene quien causó el daño involuntario? \\
Situación personal de la víctima & ¿Cuáles son las circunstancias de la persona que sufrió el daño? \\
Circunstancias del hecho & El contexto particular del caso concreto que rodea al evento dañoso. \\
\bottomrule
\end{tabularx}
\end{table}

El juez debe considerar estos factores, pudiendo apartarse de la regla general por razones de equidad y disponer el monto de la indemnización, llegando incluso a la posibilidad de que no exista indemnización, según una posición doctrinaria minoritaria; la mayoría de los autores sostiene, en cambio, que siempre corresponde otorgar algún tipo de indemnización.

Con el desarrollo de la imputabilidad de los actos voluntarios y del régimen de los actos involuntarios queda cerrada la presentación general de la teoría de los hechos y actos jurídicos.

% =====================================================================
% CORNELL
% =====================================================================

\section{Cuadro Cornell de Repaso General}
\textit{Hechos y Actos Jurídicos · Actos Voluntarios · Imputabilidad}

\begin{longtable}{
    >{\raggedright\arraybackslash}p{0.30\textwidth}
    >{\raggedright\arraybackslash}p{0.60\textwidth}
}
\toprule
\textbf{Preguntas / palabras clave} &
\textbf{Notas / respuestas} \\
\midrule
\endfirsthead
\multicolumn{2}{l}{\textit{Cuadro Cornell de Repaso General (continuación)}} \\
\toprule
\textbf{Preguntas / palabras clave} &
\textbf{Notas / respuestas} \\
\midrule
\endhead
\midrule
\multicolumn{2}{r}{\textit{Continúa en la página siguiente}} \\
\endfoot
\bottomrule
\endlastfoot

\textbf{¿Qué es un hecho jurídico según el CCCN?} &
El hecho jurídico es el acontecimiento que, conforme al ordenamiento jurídico, produce el nacimiento, modificación y extinción de relaciones o situaciones jurídicas (Art. 257 CCCN). El CCCN mejoró la redacción anterior: usa ``produce'' en lugar de ``susceptible de producir'', y habla de ``relaciones o situaciones jurídicas'' en lugar de solo ``derechos u obligaciones''. \\

\textbf{¿Cómo se clasifican los hechos jurídicos?} &
Hechos de la naturaleza (sin intervención humana: granizo, terremoto, aluvión) y hechos humanos (con intervención de la voluntad). Los humanos se subclasifican en voluntarios (con discernimiento, intención y libertad) e involuntarios. \\

\textbf{¿Cuándo un acto humano es voluntario?} &
Cuando reúne tres condiciones internas --discernimiento (aptitud para distinguir lo bueno de lo malo), intención (propósito dirigido al acto concreto) y libertad (ausencia de coacción)-- y se exterioriza (Art. 260 CCCN). \\

\textbf{¿Qué son las causas obstativas y cuáles corresponden a cada condición?} &
Son los vicios o impedimentos que eliminan cada condición. Discernimiento: privación de razón, minoría de edad (menor de 10 años para actos ilícitos; menor de 13 para actos lícitos). Intención: error y dolo. Libertad: violencia (fuerza o intimidación). \\

\textbf{¿En qué se diferencia el acto jurídico del simple acto lícito?} &
Ambos son voluntarios y lícitos. La diferencia radica en el fin inmediato: el acto jurídico tiene por fin inmediato producir consecuencias jurídicas (Art. 259 CCCN); el simple acto lícito no busca ese resultado directamente, aunque puede producirlo (Art. 258 CCCN). \\

\textbf{¿Cuáles son los dos significados del ``dolo'' en el Derecho Civil?} &
(1) Como factor de atribución: intención de dañar o manifiesta indiferencia por los intereses ajenos (Art. 1724 CCCN). (2) Como vicio de la voluntad: maniobra engañosa que induce a la otra parte a celebrar un acto que de otro modo no habría realizado. \\

\textbf{¿Qué diferencia al dolo de la culpa como factores de atribución?} &
Dolo: el agente produce el daño de manera intencional o con manifiesta indiferencia. Culpa: omisión de la diligencia debida, sin intención de dañar. Se subdivide en negligencia (torpeza), imprudencia (temeridad) e impericia (falta de conocimiento técnico). \\

\textbf{¿Cuándo el silencio es manifestación de voluntad?} &
En principio no lo es (Art. 263 CCCN). Excepcionalmente sí, cuando: (a) la ley obliga a expedirse; (b) las partes pactaron que el silencio implica determinada consecuencia; o (c) existe una serie de declaraciones precedentes que permiten inferir el sentido del silencio. \\

\textbf{¿Cuáles son los cuatro presupuestos de la responsabilidad civil?} &
1) Daño (efectivo). 2) Antijuridicidad (se presume, salvo causa de justificación -- Art. 1717). 3) Factor de atribución (subjetivo: dolo/culpa; objetivo: la ley). 4) Relación de causalidad adecuada entre el acto y el daño. \\

\textbf{¿Por cuáles consecuencias responde el agente?} &
Por las consecuencias inmediatas (curso natural y ordinario -- Art. 1726) y las mediatas previsibles. No se responde por las casuales (mediatas imprevisibles -- Art. 1727). El caso fortuito --imprevisto o inevitable-- exime de responsabilidad (Art. 1730). \\

\textbf{¿Qué pasa con los actos involuntarios y el deber de reparar?} &
En principio no generan responsabilidad. Pero el Art. 1750 CCCN establece que el autor de un daño involuntario responde por razones de equidad. El juez pondera el patrimonio del causante, la situación de la víctima y las circunstancias del hecho (remisión al Art. 1742). Quien actúa bajo fuerza irresistible no responde; sí responde quien ejerció esa fuerza. \\

\textbf{¿Qué es la equidad en este contexto?} &
En su sentido aristotélico (epieikeia): la rectificación de lo justo legal cuando la aplicación estricta de la norma conduce a un resultado manifiestamente injusto en el caso concreto. Permite al juez apartarse de la regla general de no responsabilidad por actos involuntarios para establecer (o no) una indemnización justa. \\

\midrule

\multicolumn{2}{p{0.94\textwidth}}{
\vspace{0.2em}
\textbf{Resumen:}
La teoría de los hechos y actos jurídicos constituye el tercer elemento de la relación jurídica, es decir, la fuente de donde emergen derechos y obligaciones. El punto de partida es la clasificación general: hechos de la naturaleza y hechos humanos; dentro de los humanos, voluntarios e involuntarios según concurran discernimiento, intención y libertad. Los actos voluntarios se dividen en lícitos --simples actos lícitos y actos jurídicos (Art. 259 CCCN), cuyo elemento diferenciador es el fin inmediato de producir consecuencias jurídicas-- e ilícitos --delitos (con dolo) y cuasidelitos (con culpa)--. La exteriorización de la voluntad puede ser expresa o tácita; el silencio, como regla general, no es manifestación de voluntad salvo excepciones legales o convencionales. La segunda parte abordó la imputabilidad: el deber de reparar (Arts. 1716-1717 CCCN) nace cuando confluyen cuatro presupuestos --daño, antijuridicidad, factor de atribución y nexo causal adecuado--. Las consecuencias inmediatas y mediatas previsibles son reparables; las casuales no. El caso fortuito exime de responsabilidad (Art. 1730). Finalmente, los actos involuntarios, si bien no generan responsabilidad en principio, pueden dar lugar a una indemnización por equidad en virtud del Art. 1750 CCCN, aplicando los criterios de atenuación del Art. 1742: patrimonio del causante, situación personal de la víctima y circunstancias del hecho. Esta tensión entre la regla de la no responsabilidad y la equidad como válvula correctora constituye uno de los debates doctrinarios más actuales del derecho civil argentino.
\vspace{0.2em}
} \\

\end{longtable}
\chapter[El acto jurídico]{El acto jurídico: concepto y elementos}
\label{ch:actojuridico}

El acto jurídico es, dentro de la clasificación general de los hechos jurídicos (sección~\ref{sec:clasifhechos}), un hecho humano voluntario lícito que tiene por fin inmediato producir la adquisición, modificación o extinción de derechos, relaciones o situaciones jurídicas. Es la categoría central de esta parte.

\section{Acto jurídico y contrato: alcance de la teoría general}

El Código Civil y Comercial argentino contiene un sistema general de los actos jurídicos, a diferencia del Código francés, que no trata los actos jurídicos en general sino que se limita al contrato. \textbf{El acto jurídico es un concepto más amplio que el contrato}; el contrato es una especie dentro del género acto jurídico.

\begin{definition}{Contrato}
El contrato es el acto jurídico en el cual hay acuerdo de voluntades para el surgimiento de obligaciones en el campo patrimonial.
\end{definition}

\begin{example}{Contratos como especie de acto jurídico}
\begin{itemize}
    \item \textbf{Contrato de compraventa}: acuerdo de voluntades para transferir la propiedad de una cosa contra el pago de un precio.
    \item \textbf{Contrato de locación}: acuerdo para transferir la tenencia de un bien —mueble o inmueble— contra la entrega de una renta mensual (canon).
\end{itemize}
\end{example}

No todos los actos jurídicos son contratos. Existen actos jurídicos que no responden a esta estructura de acuerdo patrimonial de voluntades:

\begin{itemize}
    \item \textbf{Testamento}: acto que tiene por finalidad disponer de los bienes para el supuesto de la propia muerte, asignándolos según la voluntad del testador. Produce efectos luego de la muerte, pero no constituye un contrato.
    \item \textbf{Matrimonio}: acto jurídico de carácter especial que crea relaciones jurídicas de naturaleza personal.
\end{itemize}

\section{Definición y características del acto jurídico}

\begin{definition}{Acto jurídico (Artículo 259, Código Civil y Comercial de la Nación)}
Acto jurídico es el acto voluntario lícito que tiene por fin inmediato la adquisición, modificación o extinción de derechos, relaciones o situaciones jurídicas.
\end{definition}

De esta definición se desprenden tres elementos fundamentales: que el acto sea \textbf{voluntario}, que sea \textbf{lícito}, y que tenga por \textbf{fin inmediato} producir consecuencias jurídicas.

\begin{important}
La idea de ``fin inmediato'' es crucial: significa que el acto jurídico busca de manera directa la producción de consecuencias jurídicas. Esto es lo que lo diferencia del simple acto lícito, que puede producir consecuencias jurídicas, pero sin que ello constituya su finalidad.
\end{important}

\section{Acto jurídico frente al simple acto lícito}

La distinción entre el acto jurídico y el simple acto lícito reside en la intención del sujeto respecto del efecto jurídico producido.

\begin{example}{El pescador}
Un pescador que sale a pescar, al apropiarse de lo pescado, realiza un simple acto lícito. Se produce un efecto jurídico —el pescador pasa a ser dueño de lo que pescó—, pero su intención original al pescar no es adquirir propiedad, sino pescar. El efecto jurídico es una consecuencia accesoria o secundaria del acto, no su fin inmediato.

Por el contrario, cuando se celebra un contrato de compraventa, el fin inmediato del sujeto es adquirir la propiedad de una cosa. El efecto jurídico es precisamente el objetivo buscado, la intención directa. Por ello, se trata de un acto jurídico.
\end{example}

\section{Los elementos del acto jurídico}

El acto jurídico se compone de los siguientes elementos esenciales:

\begin{enumerate}
    \item \textbf{Sujeto}: la persona que concurre al acto para su celebración.
    \item \textbf{Objeto}: la materia sobre la que recae el acto, es decir, los bienes o hechos que constituyen su contenido.
    \item \textbf{Forma}: la manera en que se exterioriza el acto.
    \item \textbf{Causa}: la finalidad perseguida por las partes al realizar el acto.
\end{enumerate}

% TODO: contenido incompleto o ambiguo en las notas originales — el tratamiento detallado del elemento "sujeto" y del elemento "forma" no fue desarrollado en esta clase; el apunte indica expresamente que la forma se abordará en una clase posterior, y el sujeto no vuelve a ser tratado en el material disponible.

El sujeto de los actos jurídicos se estudió en las Partes~II y III. Los capítulos siguientes desarrollan el \textbf{objeto} y la \textbf{causa} del acto jurídico; la \textbf{forma} se estudia en el capítulo~\ref{ch:forma}, y los elementos accidentales o modalidades (condición, plazo y cargo), que pueden estar o no presentes en un acto determinado, en el capítulo~\ref{ch:modalidades}.
\chapter{El objeto del acto jurídico}

\section{Fundamento normativo}

El artículo 279 del Código Civil y Comercial es central en materia de objeto de los actos jurídicos. Este artículo reconoce como antecedente el artículo 953 del Código Civil de Vélez, disposición que tuvo enorme influencia a lo largo de toda la vigencia del código anterior, en tanto marcaba el contenido moral del acto.

\begin{definition}{Objeto del acto jurídico (Artículo 279, Código Civil y Comercial de la Nación)}
El objeto del acto jurídico debe ser un hecho o un bien. Los hechos no pueden ser imposibles (ni jurídica ni prácticamente), prohibidos por la ley, o contrarios a la moral y a las buenas costumbres. Los bienes que son objeto del acto no pueden ser prohibidos, ni ser de carácter extrapatrimonial a menos que exista una autorización legal.
\end{definition}

\section{Contenido moral y límite de la autonomía de la voluntad}

Un acto jurídico debe ser conforme a la ley, a la moral y al orden público. No existe una libertad absoluta en la autonomía de la voluntad: esta se encuentra limitada y moldeada por razones de distinto tipo —la ley, el orden público, la moral y la dignidad de la persona.

\begin{note}
A lo largo de la historia, la disposición del artículo 953 del Código de Vélez originó numerosas limitaciones y aplicaciones jurisprudenciales, entre ellas:
\begin{itemize}
    \item Antes de la sanción de la Ley 17.711, dicha disposición dio lugar a la aplicación del vicio de \textbf{lesión}, que permitía anular un acto cuando existía una notable desproporción entre las prestaciones.
    \item Limitación en la aplicación de \textbf{cláusulas penales}, entendidas como la cláusula que se impone a quien incumple un contrato, obligándolo a pagar, por ejemplo, intereses punitorios.
    \item Limitación de los intereses cuando resultan \textbf{usurarios}: cobrar una tasa de interés excesiva o desproporcionada constituye usura, considerada inmoral. Aunque ello depende de la situación y el contexto (por ejemplo, en escenarios inflacionarios), cobrar tasas del 200 al 300 por ciento de interés es claramente inmoral.
\end{itemize}
\end{note}

Existe una tensión ideológica sobre este punto: una mirada más liberal sostiene que cada parte es libre de pactar las tasas de interés que desee. Sin embargo, en Argentina —y en la mayoría de las legislaciones modernas— aun cuando no exista un límite normativo expreso a los intereses, el objeto del acto jurídico debe ser conforme a la moral y a las buenas costumbres, principio que conserva plena vigencia en la actualidad.

\section{Requisitos del objeto}

El artículo 279 divide el objeto del acto jurídico en dos grandes categorías: los hechos y los bienes.

\subsection{Los hechos como objeto del acto}

Los hechos, en tanto acontecimientos, no pueden ser:

\begin{enumerate}
    \item \textbf{Imposibles} (ni jurídica ni prácticamente): no puede pactarse lo imposible. Cruzar el océano Atlántico nadando, por ejemplo, no puede ser objeto de un acto jurídico por tratarse de un hecho prácticamente imposible.
    \item \textbf{Prohibidos por la ley}: si la ley prohíbe determinado hecho, este no puede ser objeto de un acto jurídico.
    \item \textbf{Contrarios a la moral y a las buenas costumbres}: constituye una exigencia fundamental que el acto jurídico se adecue a los principios morales básicos.
\end{enumerate}

\subsection{Los bienes como objeto del acto}

El artículo 279 hace referencia genérica a prohibiciones que pueden surgir de la ley. Así, por ejemplo, ciertos bienes no pueden ser objeto de determinados contratos: los bienes de dominio público no pueden ser objeto de un acto jurídico de disposición.

\section{Objeto ilícito: el caso de la venta de influencias}

\begin{example}{Venta de influencias y acto de corrupción}
Una persona se compromete a realizar tratativas para obtener un privilegio en sede administrativa. En el fondo, se trata de una venta de influencias, un acto de corrupción.

Este acto tiene un objeto ilícito porque vender influencias es contrario a la ética pública y constituye un acto de corrupción que afecta el orden público.

Si posteriormente esa persona pretende hacer valer la obligación —por ejemplo, cobrar lo que le fue prometido—, la consecuencia es la \textbf{nulidad del acto}. La nulidad surge precisamente porque el objeto no responde a los requisitos del artículo 279. Se trata de una nulidad de tipo \textbf{absoluta}, regulada en el artículo 386 del Código Civil y Comercial.
\end{example}
\chapter{La causa del acto jurídico}
\label{ch:causa}

\section{Causa fuente y causa fin}

El concepto de causa admite varias acepciones en el derecho civil. En particular, se distingue entre la causa fuente y la causa fin.

\begin{definition}{Causa fuente}
Es la fuente de las obligaciones, es decir, el origen del cual surgen. Las obligaciones surgen de los contratos, en general de los actos jurídicos, y también de los actos ilícitos. No puede haber una obligación sin una causa, sin un elemento que le dé origen.
\end{definition}

\begin{definition}{Causa fin}
Es el fin perseguido por las partes al momento de celebrar un acto: qué se propusieron las partes al celebrarlo y cuál es la finalidad que persiguen.
\end{definition}

El desarrollo que sigue se concentra en el problema de la \textbf{causa fin}, y no en la causa fuente.

\section{Evolución histórica de la teoría de la causa}

El concepto de causa ha sido objeto de debate a lo largo de la historia del derecho, dando lugar a distintas posturas doctrinarias. Elementos vinculados a la valoración de los actos aparecen desde la antigüedad, especialmente con los canonistas y los glosadores.

\subsection{Teoría causalista}

Esta teoría, atribuida a un autor francés, sostiene que la causa fin de un contrato es la prestación de la otra parte. Se presenta como algo evidente que, al celebrar un contrato sinalagmático —perfectamente bilateral, en el que una parte tiene una prestación que la obliga y la otra, a su vez, otra prestación que también la obliga—, cada parte persigue como causa la prestación de la contraria.

\begin{example}{Compraventa bajo la teoría causalista}
En una compraventa, la causa por la cual una parte celebra el contrato es adquirir la cosa, y para la otra parte, la causa es recibir el precio.
\end{example}

\subsection{Crítica anticausalista}

Un autor belga —mencionado posteriormente por otros autores, entre ellos Planiol— critica la tesis causalista. El argumento principal sostiene que, si la causa fin es la prestación de la otra parte, entonces la causa se identifica con el objeto mismo del contrato, sin ofrecer ninguna novedad explicativa.

La crítica anticausalista señala que la idea de causa fin no agrega nada que no esté ya detrás del objeto: si el objeto es la transferencia de la propiedad (en una compraventa) y la causa es adquirir esa propiedad, causa y objeto se superponen. La postura anticausalista afirma, en consecuencia, que la causa no constituye un elemento independiente.

\subsection{Teoría moderna de la causa}

Enrique Capitant, hacia comienzos del siglo XX, retoma la idea de la causa y la profundiza. Para Capitant, la causa relevante es la finalidad perseguida por las partes, una finalidad que estas incorporan y que buscan en el momento de celebrar el acto.

\section{La causa en el Código de Vélez}

El Código Civil de Vélez no menciona la causa entre los elementos del acto de manera explícita en la sección específica de actos jurídicos. Sin embargo, el tema de la causa aparecía en algunos artículos ubicados en la parte de obligaciones (artículos 499 a 502), en el Libro Segundo, que se abría con la sección de obligaciones.

\begin{definition}{Artículo 499, Código Civil de Vélez}
No hay obligación sin causa. Se entiende por tal el hecho o acto lícito o ilícito en la relación de familia, o entre extraños, que es la causa generadora del derecho que faculta al acreedor a exigir del deudor la prestación debida.
\end{definition}

Este artículo se refiere claramente a la \textbf{causa fuente} (el hecho o acto que genera la obligación). Sin embargo, en los artículos 500 y 502, y especialmente en las notas que Vélez Sarsfield colocaba al pie de la ley —donde se critica la confusión entre causa fuente y causa fin—, aparece la idea de causa fin vinculada a la finalidad perseguida.

\begin{definition}{Artículo 502, Código Civil de Vélez (principio de la causa ilícita)}
De ningún efecto es el acto constituido con una causa ilícita, si esa causa es contraria a las leyes o al orden público.
\end{definition}

Aunque la causa no estaba expresamente mencionada como elemento en el Código de Vélez, existía la presunción de que, si las partes no expresaban la causa, esta debía existir igualmente, y la obligación se mantenía válida aun cuando la causa expresada fuera falsa.

\begin{important}
Una obligación con causa ilícita carece de todo efecto. Si la causa es ilícita, si contraviene las leyes o el orden público, el acto es nulo.
\end{important}

% TODO: contenido incompleto o ambiguo en las notas originales — el apunte menciona una observación del docente sobre el "artículo 429 del Código de Vélez", indicando que, cuando dicho artículo se refiere a una causa ilícita, en realidad remite a la causa fuente y no a la causa fin (ejemplo dado: quien causa un daño a otro debe repararlo, lo cual es causa fuente de la obligación de reparar). El apunte no precisa con exactitud el contenido ni la ubicación de dicho artículo dentro del sistema de Vélez, por lo que se conserva la referencia tal como fue consignada en clase.

\section{La causa en el Código Civil y Comercial}

El nuevo Código Civil y Comercial incorpora expresamente la causa como elemento del acto jurídico. El artículo 281 es la disposición central de esta regulación.

\begin{definition}{Artículo 281, Código Civil y Comercial}
La causa es el fin inmediato perseguido por las partes que ha sido determinante de la voluntad. Cuando las partes no la expresan, se presume que existe. Se presume también su licitud, a menos que resulte ilícita del acto mismo. También integran la causa los motivos exteriorizados cuando sean lícitos y hayan sido incorporados al acto en forma expresa o tácita, y sean esenciales para ambas partes.
\end{definition}

En rigor, lo que resulta relevante de esta definición es que la causa es el fin perseguido por las partes que ha sido determinante de la voluntad.

\subsection{Dimensiones objetiva y subjetiva de la causa}

El artículo 281 incorpora dos dimensiones de la causa:

\begin{enumerate}
    \item \textbf{Dimensión objetiva}: el fin inmediato autorizado por el ordenamiento jurídico. Es la finalidad típica que el derecho reconoce como válida para ese tipo de acto. En una compraventa, la causa objetiva es transferir la propiedad de una cosa a cambio de un precio.
    \item \textbf{Dimensión subjetiva}: los motivos o intenciones de las partes. Estos integran la causa cuando son exteriorizados (expresados o manifestados), lícitos, y esenciales para ambas partes; es decir, cuando ambas partes conocen y aceptan esos motivos como parte de la razón por la cual celebran el acto.
\end{enumerate}

\begin{example}{El caso de la profesora de francés}
Una profesora de francés contrata a un escritor para que redacte un ensayo que ella presentará, traducido, en un concurso destinado a obtener una beca para viajar a Francia. La profesora presentará el ensayo como propio.

En este caso se distingue con claridad entre objeto y causa:
\begin{itemize}
    \item El \textbf{objeto} del acto es contratar a alguien para que realice un ensayo o un trabajo. Esto es, en sí mismo, perfectamente válido y lícito.
    \item La \textbf{causa} (finalidad perseguida) es presentar el ensayo como propio para engañar a un concurso y obtener una beca. Esto es ilícito.
\end{itemize}

Si una parte persigue una causa ilícita pero la otra parte lo ignora, esta última no resulta afectada: el motivo ilícito no integra el acto. Pero si se interpreta que ambas partes conocen la causa y esta es esencial para ambas, entonces dicha causa integra el acto y lo vicia.
\end{example}

\subsection{Síntesis de las características de la causa}

\begin{enumerate}
    \item El Código Civil y Comercial incorpora la causa como elemento del acto jurídico en el artículo 281.
    \item La causa comprende dos dimensiones: una \textbf{objetiva} (el fin inmediato autorizado por el ordenamiento, determinante de la voluntad) y una \textbf{subjetiva} (los motivos o intenciones de las partes).
    \item Solo integran la causa los motivos o intenciones cuando: (a) son exteriorizados, (b) son lícitos, y (c) son esenciales para ambas partes.
    \item Quedan excluidos de la causa los puros motivos o intenciones personales que no estén exteriorizados ni sean esenciales para ambas partes.
\end{enumerate}

\section{Régimen de la causa}

\begin{definition}{Artículo 282, Código Civil y Comercial}
Cuando la causa no está expresada, se presume que existe. Se presume también su licitud, a menos que resulte ilícita del acto mismo.
\end{definition}

El artículo 282 retoma los principios ya contenidos en los artículos 500, 501 y 502 del Código de Vélez:

\begin{itemize}
    \item Si la causa no está expresada, se presume que existe (presunción de causa).
    \item El acto vale aunque la causa expresada sea falsa. Si lo que se aparenta es falso pero existe una causa verdadera, se estará ante una simulación, aunque el acto siga siendo válido.
    \item Una causa ilícita vicia el acto: si la causa es ilícita —contraria a las leyes y al orden público—, el acto es nulo o de ningún efecto.
\end{itemize}

\section{Actos abstractos}

\begin{definition}{Artículo 1014, Código Civil y Comercial (en materia de contratos)}
La inexistencia, falsedad o ilicitud de la causa no son discutibles en el acto abstracto.
\end{definition}

% TODO: contenido incompleto o ambiguo en las notas originales — el apunte atribuye esta regla al "artículo 283" en el desarrollo oral de la clase, mientras que el recuadro normativo transcripto en el propio material la identifica como "artículo 1014". Se conserva la referencia tal como figura en las notas originales, sin resolver la discrepancia mediante conocimiento externo.

Un \textbf{acto abstracto} es aquel que se independiza de su causa y recibe un tratamiento autónomo. El caso típico es el pagaré.

\begin{example}{El pagaré}
Una persona firma un pagaré porque contrae una deuda con un hospital. Al ingresar al hospital, el pagaré adquiere autonomía propia: puede circular y transferirse a terceros. Si posteriormente esa persona ya pagó al hospital y alguien pretende ejecutar el pagaré, no puede discutirse la causa (el hecho de que ya se pagó). La causa por la cual se emitió el pagaré no se discute en los títulos abstractos.
\end{example}

\section{Repaso mediante el método Cornell}

El siguiente cuadro reproduce, en formato de tabla Cornell, las preguntas y respuestas centrales para el repaso activo de los contenidos desarrollados en estos capítulos.

\begin{center}
\begin{longtable}{
    >{\raggedright\arraybackslash}p{0.28\textwidth}
    >{\raggedright\arraybackslash}p{0.62\textwidth}
}
\toprule
\textbf{Preguntas / palabras clave} & \textbf{Notas / respuestas} \\
\midrule
\endfirsthead
\toprule
\textbf{Preguntas / palabras clave} & \textbf{Notas / respuestas} \\
\midrule
\endhead
\midrule
\multicolumn{2}{r}{\textit{Continúa en la página siguiente}} \\
\midrule
\endfoot
\bottomrule
\endlastfoot

\textbf{¿Qué es un acto jurídico?} &
Es el acto voluntario lícito con fin inmediato de producir la adquisición, modificación o extinción de derechos, relaciones o situaciones jurídicas (art.\ 259, Código Civil y Comercial). Lo fundamental es que la intención directa de generar consecuencias jurídicas distingue al acto jurídico del simple acto lícito. \\

\addlinespace

\textbf{¿Cuál es la diferencia entre acto jurídico y simple acto lícito?} &
En el acto jurídico, la producción de consecuencias jurídicas es el fin inmediato e intencional. En el simple acto lícito, las consecuencias jurídicas se producen, pero no son el fin buscado (por ejemplo, el pescador que se apropia de lo que pesca). Lo determinante es la intención o finalidad del sujeto respecto del efecto jurídico. \\

\addlinespace

\textbf{¿Cuáles son los elementos del acto jurídico?} &
1.\ Sujeto: la persona que realiza el acto. 2.\ Objeto: la materia o los bienes sobre los que recae el acto. 3.\ Forma: la manera de exteriorizar el acto. 4.\ Causa: la finalidad perseguida por las partes. \\

\addlinespace

\textbf{¿Qué requisitos debe cumplir el objeto?} &
Debe ser posible (física y jurídicamente); no debe estar prohibido por la ley; debe ser lícito, es decir, conforme a la moral, a las buenas costumbres y al orden público; y no debe ser de carácter extrapatrimonial, salvo autorización legal. \\

\addlinespace

\textbf{¿Qué es la causa del acto?} &
Es el fin inmediato perseguido por las partes que ha sido determinante de la voluntad (art.\ 281, Código Civil y Comercial). Incorpora una dimensión objetiva (el fin autorizado por el derecho) y una dimensión subjetiva (los motivos lícitos, exteriorizados y esenciales para ambas partes). Se distingue de la causa fuente (origen de la obligación), que resulta menos relevante a estos efectos. \\

\addlinespace

\textbf{¿Qué diferencia hay entre objeto ilícito y causa ilícita?} &
En el objeto ilícito, la materia misma sobre la que recae el acto es ilícita (por ejemplo, vender un bien prohibido o contratar un acto ilegal). En la causa ilícita, el objeto es lícito, pero la finalidad o el motivo que persiguen las partes es ilícito (por ejemplo, contratar la redacción de un ensayo cuyo fin es el fraude académico). Ambas vician el acto, aunque por razones distintas. \\

\addlinespace

\textbf{¿Qué es un acto abstracto?} &
Es el acto que se independiza de su causa y recibe un tratamiento autónomo. El ejemplo típico es el pagaré: una vez emitido, puede circular sin vincularse a la causa original, y dicha causa no puede discutirse al momento de su ejecución. La inexistencia, falsedad o ilicitud de la causa no son discutibles en el acto abstracto. \\

\midrule
\multicolumn{2}{p{0.90\textwidth}}{
\textbf{Resumen:} Los actos jurídicos constituyen el fenómeno jurídico fundamental mediante el cual las personas, voluntariamente, crean, modifican o extinguen derechos y obligaciones. Su estudio requiere comprender que un acto jurídico no es simplemente cualquier conducta lícita: su característica esencial es que su fin inmediato es la producción intencional de consecuencias jurídicas. Para que un acto sea válido debe reunir cuatro elementos: un sujeto (persona con capacidad), un objeto (materia lícita, posible y conforme a la moral), una forma (manera de exteriorización) y una causa (finalidad perseguida). El objeto y la causa constituyen los límites más importantes a la autonomía de la voluntad: no puede haber objeto imposible, prohibido o contrario a la moral, ni causa ilícita. La distinción entre objeto ilícito y causa ilícita es crucial, porque permite comprender que un acto puede tener un objeto perfectamente válido pero una finalidad viciada. El Código Civil y Comercial introduce innovaciones relevantes: incorpora la causa expresamente como elemento (art.\ 281), reconoce sus dimensiones objetiva y subjetiva, y establece presunciones de existencia y licitud de la causa. Salvo en los actos abstractos (como el pagaré), la causa es siempre discutible y determinante: una causa ilícita torna nulo el acto, sin importar cuán lícito sea su objeto.
} \\

\end{longtable}
\end{center}
\chapter{La forma del acto jurídico}
\label{ch:forma}
% ============================================================

\section{Concepto de forma}

\begin{definition}{Forma del acto jurídico}
La forma es la exteriorización de la voluntad de los sujetos respecto del objeto, en orden a conseguir el fin jurídico y su protección al celebrar el acto. Ningún acto puede carecer de forma; en este sentido, la forma es un elemento esencial de los actos jurídicos.
\end{definition}

La forma del acto jurídico puede ser muy variada. El ordenamiento contempla desde formas simples —como un gesto que expresa la voluntad de adquirir algo— hasta actos de gran solemnidad, como el matrimonio y el testamento, para los cuales la ley establece formas determinadas. Entre estos extremos, existen los contratos verbales, los contratos electrónicos celebrados mediante formularios por internet y los instrumentos privados y públicos.

La forma se vincula con dos cuestiones centrales del derecho civil:

\begin{itemize}
    \item La \textbf{validez} del acto: en ciertos casos, la forma es requisito de existencia del acto y su incumplimiento genera la nulidad.
    \item La \textbf{prueba} del acto: en otros casos, la forma es exigida por la ley con el único fin de facilitar la demostración de que el acto fue celebrado.
\end{itemize}

\section{Regulación en el Código Civil y Comercial}

El Código Civil y Comercial de la Nación (CCyCN) regula la forma del acto jurídico dentro del Capítulo Quinto dedicado a los actos jurídicos, en las siguientes secciones del Libro Primero, Título Cuarto:

\begin{itemize}
    \item \textbf{Sección Tercera:} Forma y prueba del acto jurídico.
    \item \textbf{Sección Cuarta:} Instrumento público.
    \item \textbf{Sección Quinta:} Escritura pública.
    \item \textbf{Sección Sexta:} Instrumentos privados y particulares.
\end{itemize}

\section{El formalismo en el derecho: ventajas y desventajas}

La regulación jurídica de las formas no es neutral. Exige ponderarse en función de sus ventajas y desventajas para el tráfico jurídico.

\subsection{Ventajas del formalismo}

\begin{table}[H]
\centering
\begin{tabularx}{\textwidth}{>{\bfseries\raggedright\arraybackslash}p{0.35\textwidth}
                              >{\raggedright\arraybackslash}X}
\toprule
Ventaja & Explicación \\
\midrule
Facilita la prueba del acto & El acto queda documentado y puede acreditarse con mayor certeza. \\[4pt]
Publicita el acto & El acto es conocido por los demás; ingresa al conocimiento público. \\[4pt]
Colabora en la percepción de impuestos & Como ocurre con los contratos de alquiler que deben informarse a la AFIP. \\[4pt]
Permite distinguir el acto concluido de los preparatorios & La forma cristaliza el acto jurídico. Los borradores de un testamento no son un testamento; éste queda concluido con la forma mandada por la ley. \\[4pt]
Protege a terceros & Los derechos adquiridos respecto de terceros resultan tutelados cuando se cumple la forma. El tercero conoce su situación jurídica; esto es especialmente relevante en materia de inmuebles. \\
\bottomrule
\end{tabularx}
\end{table}

\begin{example}{Testamento — cristalización del acto}
Una persona puede estar pensando en hacer un testamento, consultar a distintas personas y elaborar borradores. En un momento determinado, la forma cristaliza ese acto jurídico. Los borradores no sirven: no son un testamento. El testamento queda concluido con la forma mandada por la ley: por escritura pública o por instrumento ológrafo.
\end{example}

\subsection{Desventajas del formalismo}

Las formas más estrictas presentan también inconvenientes:

\begin{itemize}
    \item Las formas son más costosas y engorrosas.
    \item Existe mayor riesgo de que el acto resulte inválido por algún defecto formal.
    \item Le dan \textbf{menor celeridad} a las transacciones.
    \item Generan incomodidades indudables en el tráfico jurídico cotidiano.
\end{itemize}

Por estas razones, el ordenamiento jurídico impone formas más exigentes únicamente a los actos más importantes o trascendentes, donde la mayor seguridad justifica la mayor incomodidad.

% ============================================================
\chapter{Clasificación de los Actos Jurídicos según la Forma}
% ============================================================

\section{Forma esencial y forma instrumental}

Es necesario distinguir dos sentidos del concepto de forma:

\begin{itemize}
    \item \textbf{Forma esencial:} es la manifestación de la voluntad que realizan las partes para celebrar el acto jurídico.
    \item \textbf{Forma instrumental:} es cuando la ley impone que algunos actos deben celebrarse mediante documentos escritos, ya sean públicos o privados.
\end{itemize}

El Código Civil y Comercial se concentra en la forma instrumental. Los instrumentos podrán ser públicos o privados.

\section{Primera clasificación: actos formales y no formales}

\begin{definition}{Actos formales}
Son aquellos respecto a los cuales la ley impone el cumplimiento de ciertas formas o solemnidades que deben cumplirse para la validez del acto. Ejemplo: el testamento, el matrimonio.
\end{definition}

\begin{definition}{Actos no formales}
Son aquellos que no tienen ninguna solemnidad prescripta por la ley; las partes son libres de realizarlos con libertad de formas.
\end{definition}

\section{Segunda clasificación: actos solemnes y no solemnes}

Los actos formales, a su vez, se clasifican en solemnes y no solemnes:

\begin{definition}{Actos formales solemnes}
Son aquellos en los que el rigor formal impuesto por la ley se vincula con la validez misma del acto. Si no se cumple con la forma, el acto es inválido.
\end{definition}

\begin{definition}{Actos formales no solemnes (ad probationem)}
Son aquellos en los que la forma es prescripta por la ley únicamente para facilitar la prueba del acto, sin afectar su validez.
\end{definition}

\section{Solemnidad absoluta y solemnidad relativa}

Dentro de los actos formales solemnes, el Código establece una distinción fundamental:

\begin{table}[H]
\centering
\begin{tabularx}{\textwidth}{>{\raggedright\arraybackslash}X
                              >{\raggedright\arraybackslash}X}
\toprule
\textbf{Solemnidad absoluta} & \textbf{Solemnidad relativa} \\
\midrule
La solemnidad es parte esencial del acto. Si no se cumple, el acto es ineficaz por nulidad de forma y no produce ningún efecto. & El acto no es válido (no genera los efectos propios del acto), pero genera la obligación de otorgar la forma prescripta. \\[6pt]
\textit{Ejemplo:} El testamento ológrafo que no está escrito y firmado de puño y letra por el testador. & \textit{Ejemplo:} El boleto de compraventa de un inmueble. No produce el efecto de la transferencia de la propiedad, pero sí genera la obligación de escriturar. \\
\bottomrule
\end{tabularx}
\end{table}

\subsection{Conversión del acto jurídico}

\begin{formula}{Art. 285 CCyCN — Conversión del acto jurídico}
\textit{El acto que no se otorga en la forma exigida por la ley no queda concluido como tal, pero vale como acto en el que las partes se obligan a cumplir con la expresada formalidad, salvo que ella se exija bajo sanción de nulidad.}
\end{formula}

La \textbf{conversión del acto jurídico} opera en los supuestos de solemnidad relativa: el acto no produce los efectos propios del negocio jurídico que se procuraba celebrar, pero se convierte en otro acto —la obligación de otorgar el instrumento proyectado—. Cuando la forma es exigida bajo sanción de nulidad (solemnidad absoluta), el incumplimiento genera la nulidad y no produce ningún otro efecto.

\begin{example}{Boleto de compraventa de inmueble}
Una compraventa de inmueble celebrada mediante boleto de compraventa (instrumento privado) no produce el efecto traslativo de dominio, ya que la ley exige escritura pública para ello. Sin embargo, en virtud de la conversión del artículo 285, genera la obligación de las partes de escriturar, es decir, de otorgar el instrumento proyectado.
\end{example}

% ============================================================

\section{Principio de Libertad de Formas}
% ============================================================

\begin{formula}{Art. 284 CCyCN — Libertad de formas}
\textit{Si la ley no designa una forma determinada para la exteriorización de la voluntad, las partes pueden utilizar la que estimen conveniente. Las partes pueden convenir una forma más exigente que la impuesta por la ley.}
\end{formula}

El principio de libertad de formas es la regla general del sistema. Si la ley no designa una forma determinada, las partes son libres de elegir la que consideren conveniente. Esta libertad tiene dos dimensiones:

\begin{itemize}
    \item \textbf{Libertad de elección:} ante el silencio de la ley sobre la forma, las partes determinan libremente el modo de exteriorizar su voluntad.
    \item \textbf{Autonomía para elevar la forma:} las partes pueden convenir una forma más exigente que la impuesta por la ley. Por ejemplo, celebrar un contrato de locación por escritura pública, aunque la ley no lo requiera.
\end{itemize}

\begin{example}{Locación por escritura pública}
Las partes de un contrato de alquiler pueden convenir celebrarlo por escritura pública, aunque la ley no exija esa forma para la locación. Ello implica costos adicionales (honorarios del escribano), pero resulta perfectamente válido.
\end{example}

Por supuesto, cuando la ley manda una forma determinada, las partes deben partir de la forma legalmente impuesta. El artículo 284 es el punto de partida para aquellos actos respecto de los cuales la ley no designa forma especial.

% ============================================================
\chapter{Los Instrumentos — Clasificación General}
% ============================================================

\section{Concepto de instrumento}

El concepto de instrumento refiere a la forma escrita. Es la expresión de la voluntad a través de un documento que da mayor seguridad y certeza acerca de los derechos y obligaciones asumidos por las partes.

\begin{formula}{Art. 286 CCyCN — Expresión escrita}
\textit{La expresión escrita puede tener lugar por instrumentos públicos, o por instrumentos particulares firmados o no firmados. Puede hacerse constar en cualquier soporte, siempre que su contenido sea representado con texto inteligible, aunque su lectura exija medios técnicos.}
\end{formula}

\begin{important}
La referencia a medios técnicos contempla los soportes informáticos y electrónicos. El Código extiende la categoría de instrumento a todo soporte inteligible, independientemente del medio requerido para su lectura.
\end{important}

\section{Clasificación de los instrumentos}

\begin{table}[H]
\centering
\begin{tabularx}{\textwidth}{>{\raggedright\arraybackslash}p{0.30\textwidth}
                              >{\raggedright\arraybackslash}p{0.30\textwidth}
                              >{\raggedright\arraybackslash}X}
\toprule
\textbf{Instrumentos públicos} & \textbf{Instrumentos privados} & \textbf{Particulares no firmados} \\
\midrule
Interviene un oficial público. Tienen validez y fuerza probatoria por sí mismos con los alcances que les concede el Código. & Surgen de las partes. Son instrumentos particulares firmados. La \textbf{firma} es un requisito esencial. & Todo escrito no firmado: impresos, registros visuales o auditivos de cosas o hechos, registros de la palabra y de la información, cualquiera sea el medio empleado. \\
\bottomrule
\end{tabularx}
\end{table}

\begin{formula}{Art. 287 CCyCN — Instrumentos particulares}
\textit{Los instrumentos particulares pueden estar firmados o no. Si lo están, se llaman instrumentos privados. Si no lo están, se los denomina instrumentos particulares no firmados; esta categoría comprende todo escrito no firmado, entre otros, los impresos, los registros visuales o auditivos de cosas o hechos y, cualquiera que sea el medio empleado, los registros de la palabra y de la información.}
\end{formula}

El Código Civil y Comercial amplió el tratamiento jurídico de esta materia respecto del Código anterior, incorporando la categoría de instrumentos particulares no firmados para dar cobertura a la realidad del tráfico contemporáneo, caracterizado por la informatización: formularios web, constancias electrónicas, mensajes de texto, registros digitales.

\begin{example}{Constancia de compra por internet}
Una constancia generada por una plataforma de comercio electrónico al completar un formulario de compra en línea —que registra los datos del contrato, el precio y el método de pago— es un instrumento particular no firmado. No tiene firma en sentido estricto, por lo que no es instrumento privado con su fuerza probatoria plena, pero puede tener valor como prueba en caso de conflicto.
\end{example}

\section{La firma como disposición general}

\begin{formula}{Art. 288 CCyCN — Firma}
\textit{La firma prueba la autoría de la declaración de voluntad expresada en el texto al cual corresponde. La firma consiste en el nombre del firmante o en un signo; en ambos casos se realiza con el trazo habitual y de tal manera que identifique a la persona.}
\end{formula}

En el Código Civil y Comercial, la regulación de la firma fue trasladada de la sección de instrumentos privados —donde se encontraba en el Código de Vélez— a las disposiciones generales sobre la forma, lo que significa que rige tanto para los instrumentos públicos (incluyendo la escritura pública) como para los privados. La firma:

\begin{itemize}
    \item Prueba la \textbf{autoría} de la declaración de voluntad.
    \item Expresa la \textbf{conformidad} con el contenido del documento.
    \item Consiste en el nombre o en un signo habitual e identificatorio de la persona.
\end{itemize}

\subsection{Firma digital y firma electrónica}

El artículo 288 contempla también la cuestión de los instrumentos generados por medios electrónicos y remite a la normativa específica sobre firma digital.

\begin{important}
La \textbf{firma digital} no debe confundirse con el agregado de datos personales al pie de un correo electrónico (nombre, cargo, teléfono). Este último no constituye firma digital en sentido jurídico.
\end{important}

\begin{definition}{Firma digital (Ley 25.506 — año 2001)}
Se entiende por firma digital al resultado de aplicar a un documento digital un procedimiento matemático (algoritmo) que requiere información de exclusivo conocimiento del firmante y que se encuentra bajo su absoluto control. La firma digital debe haber sido: (a) creada durante el período de vigencia del certificado digital válido del firmante; (b) verificada por referencia a los datos de verificación de firma; y (c) generada por un certificado emitido por un certificador licenciado.
\end{definition}

\begin{definition}{Firma electrónica}
Es un conjunto de datos electrónicos, ligados o asociados de manera lógica a otros datos electrónicos, que el signatario utiliza con fines de identificación, pero que no alcanza el estándar de la firma digital. No cuenta con certificación indubitable de autoría e integridad.
\end{definition}

La diferencia entre ambas figuras es fundamental en materia probatoria:

\begin{table}[H]
\centering
\begin{tabularx}{\textwidth}{>{\raggedright\arraybackslash}X
                              >{\raggedright\arraybackslash}X}
\toprule
\textbf{Firma digital} & \textbf{Firma electrónica} \\
\midrule
Certificación digital indubitable de \textbf{autoría} e \textbf{integridad} mediante tercero certificador licenciado. El documento firmado digitalmente no puede ser modificado con posterioridad sin que ello resulte detectable. & Identificación electrónica sin certificación indubitable. Puede ser desconocida; quien la invoca debe acreditar su validez mediante prueba. \\[6pt]
\textbf{Mayor valor probatorio.} Garantizadas la autoría y la integridad. & \textbf{Menor valor probatorio.} La autoría y la integridad no están garantizadas por estas certificaciones. \\
\bottomrule
\end{tabularx}
\end{table}

% ============================================================
\chapter{Instrumentos Públicos}
% ============================================================

\section{Concepto y caracteres}

\begin{definition}{Instrumento público}
Son los documentos a los cuales la ley les asigna un valor de autenticidad por sí mismos y que, como tales, no requieren de un reconocimiento previo de la firma para hacer plena fe de su contenido. La autenticidad del instrumento público proviene de la ley, no de la voluntad de las partes.
\end{definition}

En la mayoría de los instrumentos públicos interviene un oficial público —funcionario público o escribano— que certifica la validez del acto. Sin embargo, lo que otorga la fuerza probatoria propia al instrumento público es la ley: es ella la que, ante la intervención de determinado funcionario revestido de esa credibilidad social e institucional, concede al instrumento el carácter de auténtico.

Los caracteres esenciales del instrumento público son:

\begin{itemize}
    \item \textbf{Autenticidad:} hace plena fe de su contenido sin necesidad de reconocimiento previo.
    \item \textbf{Fuerza probatoria propia:} su valor probatorio proviene de la ley.
    \item \textbf{Oponibilidad erga omnes:} tiene valor frente a todos, no sólo entre las partes.
\end{itemize}

\section{Enumeración legal}

\begin{formula}{Art. 289 CCyCN — Enumeración de instrumentos públicos}
\textit{Son instrumentos públicos: a) Las escrituras públicas y sus copias o testimonios; b) Los instrumentos que extienden los escribanos o los funcionarios públicos con los requisitos que establecen las leyes; c) Los títulos emitidos por el Estado nacional, provincial o la Ciudad Autónoma de Buenos Aires, conforme a las leyes que autorizan su emisión.}
\end{formula}

Cada categoría comprende:

\begin{enumerate}
    \item \textbf{Escrituras públicas y sus copias o testimonios:} instrumento matriz del protocolo notarial y sus reproducciones autenticadas.

    \item \textbf{Instrumentos de escribanos o funcionarios públicos:} comprende, entre otros, las actas notariales, las sentencias judiciales homologadas, las actas celebradas en sede judicial, y los convenios celebrados ante funcionario competente (como un convenio de alimentos labrado en sede judicial).

    \item \textbf{Títulos emitidos por el Estado:} títulos de deuda pública nacional, provincial o de la Ciudad Autónoma de Buenos Aires, emitidos conforme a las leyes que autorizan su emisión.
\end{enumerate}

\section{Requisitos de validez}

\begin{formula}{Art. 290 CCyCN — Requisitos de validez del instrumento público}
\textit{Son requisitos de validez del instrumento público: a) La actuación del oficial público en los límites de sus atribuciones y de su competencia territorial, excepto que el lugar sea generalmente tenido como comprendido en ella; b) Las firmas del oficial público, de las partes, y en su caso, de sus representantes; si alguno de ellos no firma por sí o no puede hacerlo, otro de los intervinientes lo hará a su ruego, haciéndose constar esta circunstancia.}
\end{formula}

Los dos requisitos esenciales son:

\begin{itemize}
    \item \textbf{Competencia del oficial público:} comprende la competencia \emph{funcional} (atribuciones para celebrar ese tipo de acto) y la competencia \emph{territorial} (actuación dentro de la jurisdicción asignada). La actuación fuera de estos límites invalida el instrumento.
    \item \textbf{Firmas:} deben firmar el oficial público, las partes y, en su caso, los representantes. La falta de firma de cualquiera de los intervinientes invalida el instrumento para todos.
\end{itemize}

\begin{important}
Es presupuesto de validez que el oficial público se encuentre efectivamente en funciones. Los actos otorgados antes de la notificación de la suspensión o cesación de funciones son válidos. Si la persona ejercía efectivamente un cargo existente bajo apariencia de legitimidad, la falta de investidura no afecta al acto ni al instrumento dentro de los límites de la buena fe.
\end{important}

\section{Defectos de forma}

Los defectos de forma en el instrumento público —enmiendas, agregados, borraduras, entrelineados, alteraciones— deben ser salvados al final del instrumento antes de las firmas. Si no se salvan, el instrumento carece de validez en esa parte.

\begin{example}{Enmienda en número de DNI}
Si en una escritura pública se consignó erróneamente el número del documento nacional de identidad de uno de los comparecientes y se corrigió mediante una enmienda manuscrita sobre el texto impreso, esa corrección debe ser salvada al final del instrumento antes de las firmas. De lo contrario, la enmienda no salvada afecta la validez del instrumento.
\end{example}

\begin{note}
Si el instrumento contiene las firmas de las partes pero adolece de un defecto que le impide valer como instrumento público (por ejemplo, por incompetencia del escribano), puede igualmente valer como \textbf{instrumento privado}, en la medida en que el acto de que se trate no exija forma pública bajo sanción de nulidad.
\end{note}

\section{Eficacia probatoria del instrumento público}

Este es uno de los aspectos de mayor relevancia práctica. El instrumento público hace plena fe, pero el Código distingue dos situaciones según el tipo de contenido:

\begin{formula}{Art. 296 CCyCN — Eficacia probatoria del instrumento público}
\textit{El instrumento público hace plena fe: a) En cuanto a que se ha realizado el acto, la fecha, el lugar y los hechos que el oficial público enuncia como cumplidos por él o ante él hasta que sea declarado falso en juicio civil o criminal; b) En cuanto al contenido de las declaraciones sobre convenciones, disposiciones, pagos, reconocimientos y enunciaciones de hechos directamente relacionados con el objeto principal del acto instrumentado, hasta que se produzca prueba en contrario.}
\end{formula}

\subsection{Hechos percibidos directamente por el oficial público}

Los \textbf{hechos que el oficial público enuncia como cumplidos por él o ante él} —la fecha, el lugar, la celebración del acto, los hechos ocurridos en su presencia— hacen plena fe y sólo pueden ser cuestionados mediante un \textbf{juicio de redargución de falsedad} (acción civil o criminal tendiente a demostrar la falsedad del instrumento). Quien pretenda impugnarlos debe acusar al oficial público de haber faltado a la verdad, lo que supone una acusación de suma gravedad dada la credibilidad que la sociedad deposita en la función notarial.

\begin{example}{Acta notarial — constancia de hechos ante el escribano}
Si un escribano labra un acta en la que consigna: ``En el día de la fecha me constituí en el domicilio de la calle X y constaté que se realizaba la actividad Y, en presencia de las personas Z'', quien pretenda desacreditar ese hecho debe iniciar un juicio de redargución de falsedad, dado que estaría afirmando que el escribano mintió en su función oficial.
\end{example}

\subsection{Declaraciones de las partes}

En cuanto al \textbf{contenido de las declaraciones de las partes} —convenciones, disposiciones, pagos, reconocimientos y enunciaciones de hechos directamente relacionados con el objeto del acto—, el instrumento hace plena fe pero \textbf{admite prueba en contrario}. En este caso, no es necesario iniciar un juicio de redargución de falsedad: basta con demostrar por otros medios que el contenido de esas declaraciones no es veraz.

\begin{example}{Reconocimiento de deuda simulado}
Si en una escritura pública una de las partes reconoce adeudar una suma de dinero, ese reconocimiento es el contenido de su declaración de voluntad. Si posteriormente se pretende demostrar que ese reconocimiento fue simulado, no es necesaria la redargución de falsedad —el escribano no mintió: certificó que el compareciente hizo esa declaración—, sino que puede acreditarse la simulación por prueba en contrario.
\end{example}

% ============================================================
\chapter{La escritura pública y las actas notariales}
% ============================================================

\section{Concepto y protocolo notarial}

\begin{formula}{Art. 299 CCyCN — Escritura pública. Definición}
\textit{La escritura pública es el instrumento matriz extendido en el protocolo de un escribano público o de otro funcionario autorizado para ejercer las mismas funciones, que contiene uno o más actos jurídicos.}
\end{formula}

La escritura pública es el tipo de instrumento público más importante en la regulación civil y comercial. Sus características esenciales son:

\begin{itemize}
    \item \textbf{Instrumento matriz:} la escritura original es aquella que queda integrada al protocolo del escribano. Lo que las partes conservan en sus domicilios es una \textbf{copia} o \textbf{testimonio}, que también tiene carácter de instrumento público y hace plena fe tanto como la escritura matriz.
    \item \textbf{Contenido:} la escritura contiene \textbf{actos jurídicos}. Esto la distingue del acta notarial, cuyo objeto es la comprobación de \textbf{hechos}.
    \item \textbf{Función notarial:} los escribanos son regulados en cada provincia y en la Ciudad Autónoma de Buenos Aires. El Código Civil y Comercial establece reglas comunes para este instrumento con validez en todo el país.
\end{itemize}

\begin{important}
Si existe discordancia entre el texto de la copia o testimonio y el de la escritura matriz, prevalece el contenido de la escritura matriz.
\end{important}

\subsection{El protocolo notarial}

El protocolo está conformado por folios numerados correlativamente. No pueden quedar hojas en blanco —lo que implicaría la posibilidad de antedatar documentos. Los asientos deben realizarse en orden cronológico. La reglamentación del libro del protocolo es de competencia local (provincial o de la Ciudad Autónoma de Buenos Aires).

\section{Requisitos de la escritura pública}

La escritura pública debe ser recibida personalmente por el escribano, quien tiene el \textbf{deber de calificación} de todos los presupuestos y elementos del acto: verifica la identidad de los comparecientes, su capacidad, la legitimidad del acto y la regularidad de los títulos.

\subsection{Estructura: encabezado, cuerpo y pie}

El contenido de la escritura pública (artículo 305 CCyCN) se organiza en tres grandes partes:

\begin{table}[H]
\centering
\begin{tabularx}{\textwidth}{>{\bfseries\raggedright\arraybackslash}p{0.20\textwidth}
                              >{\raggedright\arraybackslash}X}
\toprule
Parte & Contenido \\
\midrule
Encabezado & Lugar y fecha de otorgamiento (con posibilidad de incluir la hora). Datos de los comparecientes: nombre, apellido, documento nacional de identidad, domicilio real, estado civil. En el caso de segundas nupcias, se hace constar. Nombre del cónyuge, relevante por cuestiones vinculadas a la vivienda y al régimen patrimonial. Si comparece una persona jurídica: denominación completa, domicilio y datos del representante y de su mandato. \\[4pt]
Cuerpo & Naturaleza del acto e individualización de los bienes que constituyen su objeto. En las transmisiones de derechos reales sobre inmuebles, incluye el \textbf{``corresponde''}: la relación de antecedentes de titularidad del bien transmitido. \\[4pt]
Pie / Cierre & Constancia de la lectura del instrumento por el escribano ante los comparecientes. Salvado de enmiendas, borrones o entrelineados realizados antes de las firmas. Firma de los otorgantes, del escribano y de los testigos, si los hubiera. Si algún compareciente no puede firmar, otro firmará a su ruego, dejándose constancia del impedimento y de la impresión digital del otorgante. \\
\bottomrule
\end{tabularx}
\end{table}

\subsection{El ``corresponde'' en las transmisiones de inmuebles}

El \textbf{``corresponde''} es una cláusula obligatoria en las escrituras que instrumentan transmisiones de derechos reales sobre inmuebles. Contiene la relación de todos los antecedentes en virtud de los cuales el transmitente llegó a ser titular del derecho que transmite.

\begin{example}{Cadena de antecedentes}
Si el transmitente adquirió el inmueble por herencia, el corresponde consigna: ``Le corresponde por declaratoria de herederos de fecha X y partición de fecha Y. A su vez, el causante lo había adquirido por escritura de fecha Z, que obra en el folio N del Registro de Propiedad Inmueble.'' Esta cláusula permite encadenar las distintas escrituras y verificar la titularidad del bien a lo largo del tiempo.
\end{example}

La omisión del ``corresponde'' no genera la invalidez de la escritura, pero puede acarrear sanciones para el escribano en el ejercicio de su función profesional.

\section{Copias y testimonios}

Las copias o testimonios de la escritura matriz deben entregarse a las partes. Pueden reproducirse por cualquier medio que asegure su permanencia. Si alguna de las partes requiere una nueva copia, debe entregársele, salvo que exista una obligación de dar dinero pendiente de cumplimiento.

% ============================================================

\section{Actas Notariales}
% ============================================================

\begin{formula}{Art. 310 CCyCN — Actas notariales}
\textit{Se denominan actas los documentos notariales que tienen por objeto la comprobación de hechos.}
\end{formula}

\begin{definition}{Acta notarial}
Es el documento expedido por un escribano a requerimiento de parte, cuyo objeto es la comprobación de \textbf{hechos}. A diferencia de la escritura pública —que contiene actos jurídicos (Art. 299)—, el acta sólo constata hechos, sin crear, modificar ni extinguir relaciones jurídicas.
\end{definition}

El acta notarial implica la concurrencia del escribano a un lugar determinado para constatar un hecho y dejar constancia de él. Entre sus aplicaciones habituales se encuentran:

\begin{itemize}
    \item Constatación del estado de un bien al momento de su restitución (locaciones).
    \item Intimaciones de pago al deudor.
    \item Comprobación de que un acreedor se negó a recibir un pago.
    \item Constancia de hechos relevantes para el ejercicio de derechos.
\end{itemize}

\begin{example}{Acta de constatación al recuperar un local comercial}
Al recuperar la posesión de un local comercial después de una locación, el locador puede requerir al escribano que concurra al acto de entrega de llaves y labre un acta notarial que constate el estado en que se encuentra el inmueble. El escribano describe en el acta todo lo que constata, y ese documento tendrá valor probatorio en caso de reclamarse daños o el incumplimiento de las obligaciones de restitución.
\end{example}

Los requisitos del acta notarial están regulados en el artículo 311 CCyCN. El acta tiene valor probatorio respecto de:

\begin{itemize}
    \item Los \textbf{hechos} que el notario tuvo a la vista: su existencia, características y estado.
    \item Las \textbf{personas} que concurrieron, si quisieron identificarse.
    \item Las \textbf{declaraciones} efectuadas por los presentes.
\end{itemize}

\begin{note}
La Ciudad Autónoma de Buenos Aires regula el ejercicio de la función notarial mediante la Ley 404 (año 2000). El escribano puede ser titular o adscripto. El ingreso al notariado requiere la aprobación de un procedimiento de selección competitivo y una investidura formal.
\end{note}

% ============================================================
\chapter{Instrumentos particulares y fecha cierta}
% ============================================================

\section{Clasificación y jerarquía}

El Código Civil y Comercial sistematiza los instrumentos particulares en dos categorías:

\begin{itemize}
    \item \textbf{Instrumentos particulares firmados} (instrumentos privados): requieren firma de las partes como requisito esencial.
    \item \textbf{Instrumentos particulares no firmados}: todo escrito no firmado, registro visual o auditivo, o cualquier otro soporte inteligible sin firma.
\end{itemize}

La categoría más importante es la de los instrumentos privados (particulares firmados). Los instrumentos particulares no firmados constituyen una categoría incorporada por el nuevo Código para dar cobertura jurídica a la realidad del tráfico contemporáneo de bienes y servicios.

\section{Instrumentos privados: el requisito de la firma}

\begin{important}
El \textbf{requisito esencial} de los instrumentos privados es la \textbf{firma} de las partes (Art. 288). Sin firma, el documento es un instrumento particular no firmado, con menor fuerza probatoria.
\end{important}

El Código Civil y Comercial eliminó el requisito del \textbf{doble ejemplar} que imponía el Código anterior (Código de Vélez). En virtud de ello, ya no es jurídicamente exigible que los instrumentos se celebren en tantos ejemplares como partes intervinientes. No obstante, la buena práctica profesional recomienda mantener este recaudo, especialmente cuando una de las partes asume obligaciones relevantes.

\section{Reconocimiento de la firma}

El instrumento privado tiene una firma que prueba la declaración de voluntad, pero para que se produzca la plenitud de la fuerza probatoria entre las partes, la firma debe ser \textbf{reconocida}.

\begin{formula}{Art. 314 CCyCN — Reconocimiento de la firma}
\textit{Todo aquel contra quien se presente un instrumento cuya firma se le atribuye debe manifestarse al respecto. El silencio se tendrá por reconocimiento de la firma.}
\end{formula}

\begin{important}
El silencio frente a la intimación de reconocimiento de firma equivale al reconocimiento. Esto constituye un supuesto de silencio como manifestación de la voluntad (Art. 263 CCyCN): ante la obligación legal de expedirse, el silencio importa aceptación tácita.
\end{important}

El reconocimiento de la firma puede ser:

\begin{itemize}
    \item \textbf{Expreso:} la persona declara que la firma le pertenece.
    \item \textbf{Tácito:} guarda silencio ante la intimación, o realiza gestiones sin desconocer la autoría de la firma.
    \item \textbf{Forzoso:} se acredita mediante prueba, especialmente a través de la \textbf{pericia caligráfica} (cotejo de firmas realizado por un perito calígrafo designado judicialmente).
\end{itemize}

\begin{example}{Negación de contrato de locación}
Si el locatario niega haber suscripto el contrato de alquiler, el locador puede exhibirle el instrumento en juicio e intimarlo a que reconozca la firma. Si la niega, podrá acreditarse la autoría mediante pericia caligráfica, que establecerá con alta probabilidad la autenticidad de la firma atribuida al demandado.
\end{example}

El reconocimiento de la firma importa el reconocimiento de todo el cuerpo del instrumento: las cláusulas, disposiciones, convenciones y reconocimientos que contiene. Lo que el instrumento público tiene por sí, el instrumento privado lo adquiere con el reconocimiento de firma.

\section{Firma a ruego e impresión digital}

\begin{formula}{Art. 313 CCyCN — Firma a ruego e impresión digital}
\textit{Si alguno de los firmantes no sabe o no puede firmar, puede dejarse constancia de la impresión digital o bien de la firma de otra persona a ruego, en presencia de dos testigos que también lo suscriban.}
\end{formula}

Cuando una parte no sabe o no puede firmar, el Código prevé dos alternativas:

\begin{enumerate}
    \item \textbf{Firma a ruego:} otra persona firma en nombre del que no puede hacerlo, en presencia de dos testigos que también suscriben el instrumento. Esta modalidad tiene pleno valor probatorio, similar a un mandato.

    \item \textbf{Impresión digital:} la estampación de la huella dactilar del signatario. Sin embargo, conforme al artículo 314 CCyCN, el documento signado con impresión digital vale únicamente como \textbf{principio de prueba por escrito} y puede ser impugnado. Su menor valor probatorio obedece a que la impresión digital puede obtenerse sin el pleno consentimiento de la persona.
\end{enumerate}

\section{Documento firmado en blanco}

\begin{formula}{Art. 315 CCyCN — Documento firmado en blanco}
\textit{El firmante de un documento en blanco puede impugnar su contenido mediante la prueba de que no responde a sus instrucciones, pero no puede valerse de testigos si no existe principio de prueba por escrito. El desconocimiento del documento no puede invocarse para afectar al tercero que de buena fe celebró un acto jurídico basándose en él.}
\end{formula}

El documento firmado en blanco implica una suerte de mandato implícito: quien firma en blanco autoriza al receptor del documento a completarlo dentro de las instrucciones acordadas. Si el documento fue completado en forma contraria a esas instrucciones, el firmante puede impugnar su contenido, pero con importantes limitaciones probatorias:

\begin{itemize}
    \item No puede valerse exclusivamente de testigos si no existe principio de prueba por escrito.
    \item Si el instrumento fue utilizado de buena fe por un tercero que celebró un acto jurídico basándose en él, ese tercero queda protegido: el firmante no puede invocar el desconocimiento del documento en su perjuicio.
\end{itemize}

\begin{important}
La firma de documentos en blanco es una práctica de alto riesgo, ya que quien tiene el documento puede completarlo con obligaciones enormemente perjudiciales para el patrimonio del firmante.
\end{important}

\section{Enmiendas en instrumentos privados}

Las enmiendas, raspaduras o entrelineados que afecten partes esenciales del instrumento privado deben ser salvadas al final, con la firma de las partes. Diferencia con la escritura pública: en los instrumentos privados, las enmiendas no salvadas no generan la nulidad del instrumento, sino que el juez determinará en qué medida el defecto excluye o reduce su fuerza probatoria.

% ============================================================

\section{Fecha Cierta y Oponibilidad a Terceros}
% ============================================================

\subsection{El instrumento privado frente a terceros}

El instrumento privado tiene fuerza probatoria entre las partes desde el reconocimiento de la firma. Frente a terceros, sin embargo, rige una regla diferente: el instrumento privado \textbf{no es oponible} a terceros mientras no adquiera \textbf{fecha cierta}.

\begin{formula}{Art. 317 CCyCN — Fecha cierta}
\textit{La eficacia probatoria de los instrumentos privados reconocidos se extiende a los terceros desde su fecha cierta. Adquieren fecha cierta el día en que acontece un hecho del que resulta como consecuencia ineludible que el documento ya estaba firmado o no pudo ser firmado después.}
\end{formula}

\begin{definition}{Fecha cierta}
No es la fecha consignada en el instrumento, sino el momento a partir del cual existe certeza objetiva de que ese instrumento ya estaba firmado o no pudo ser firmado con posterioridad. Sirve para hacer oponible el instrumento privado a terceros.
\end{definition}

\subsection{Medios para adquirir fecha cierta}

El Código de Vélez (artículo 1035) enumeraba taxativamente los casos de adquisición de fecha cierta:

\begin{itemize}
    \item Exhibición del instrumento en juicio o ante una repartición pública donde quedara archivado.
    \item Reconocimiento ante un escribano y testigos.
    \item Transcripción en cualquier registro público.
    \item Fallecimiento de la parte que lo firmó.
\end{itemize}

El Código Civil y Comercial \textbf{abandona la enumeración taxativa} y establece que la fecha cierta puede acreditarse por cualquier medio de prueba, que el juez apreciará de manera rigurosa.

\begin{example}{Fecha cierta por exhibición en juicio}
Cuando un instrumento privado se incorpora a un expediente judicial, la mesa de entradas de tribunales le estampa un sello que acredita la fecha de presentación. Desde ese momento existe certeza de que el instrumento ya existía: puede oponerse a terceros como instrumento firmado con anterioridad a esa fecha.
\end{example}

\begin{example}{Fecha cierta por fallecimiento del firmante}
Si uno de los firmantes de un instrumento privado fallece, desde ese momento se tiene certeza de que el instrumento fue firmado antes de la muerte: con posterioridad al fallecimiento, esa persona ya no pudo haberlo suscripto. El instrumento adquiere fecha cierta desde el día del fallecimiento.
\end{example}

% ============================================================
\chapter{Instrumentos Particulares no Firmados y Correspondencia}
% ============================================================

\section{Concepto y amplitud}

Los instrumentos particulares no firmados comprenden todos los escritos no firmados, impresos, registros visuales o auditivos de cosas o hechos, cualquiera sea el medio empleado, y los registros de la palabra y de la información. Esta amplísima categoría abarca tanto los soportes en papel como los medios electrónicos: mensajes de texto, correos electrónicos, publicaciones en redes sociales, tickets electrónicos, registros de plataformas web.

\begin{note}
A diferencia de los instrumentos privados, los instrumentos particulares no firmados no cuentan con firma. La firma prueba la autoría de la declaración de voluntad; su ausencia implica que el instrumento tiene menor fuerza probatoria y que la valoración de su alcance queda librada a la apreciación judicial.
\end{note}

\section{Fuerza probatoria — criterios del juez}

\begin{formula}{Art. 319 CCyCN — Valor probatorio del instrumento particular}
\textit{El valor probatorio de los instrumentos particulares debe ser apreciado por el juez ponderando, entre otras pautas, la congruencia entre lo sucedido y narrado, la precisión y claridad técnica del texto, los usos y prácticas del tráfico, las relaciones precedentes y la confiabilidad de los soportes utilizados y de los procedimientos técnicos que se apliquen.}
\end{formula}

Las pautas que el juez debe ponderar son:

\begin{enumerate}
    \item \textbf{Congruencia entre lo sucedido y lo narrado:} correspondencia entre lo que dice el instrumento y los hechos ocurridos y acreditados por otros medios.

    \item \textbf{Precisión y claridad técnica del texto:} un formulario web completo y preciso tendrá mayor valor que un mensaje informal con redacción ambigua.

    \item \textbf{Usos y prácticas del tráfico:} la forma habitual en que se concluyen los actos jurídicos en ese sector de actividad o negocio.

    \item \textbf{Relaciones precedentes entre las partes:} el historial de intercambios previos que contextualice el instrumento en cuestión.

    \item \textbf{Confiabilidad de los soportes y procedimientos técnicos empleados:} la seguridad e integridad de los sistemas informáticos utilizados para generar o transmitir el instrumento.
\end{enumerate}

\begin{example}{Ticket electrónico como prueba}
En un conflicto sobre el pago del precio en una compraventa de bienes, el comprador presenta el ticket generado por la plataforma de comercio electrónico y el resumen de tarjeta de crédito que registra el débito correspondiente. La congruencia entre ambos elementos es un indicio de peso que el juez puede ponderar para tener por acreditado el pago.
\end{example}

\section{La correspondencia}

\begin{formula}{Art. 318 CCyCN — Correspondencia}
\textit{La correspondencia, cualquiera sea el medio empleado para crearla o transmitirla, puede presentarse como prueba por el destinatario, pero la que es confidencial no puede ser utilizada sin consentimiento del remitente.}
\end{formula}

La correspondencia —cartas, correos electrónicos, mensajes— pertenece, una vez recibida, al destinatario. Éste puede utilizarla como prueba en juicio contra el remitente, con la excepción de la \textbf{correspondencia confidencial}: si la comunicación tiene carácter confidencial, no puede ser utilizada sin el consentimiento del remitente.

La confidencialidad puede surgir:
\begin{itemize}
    \item De la manifestación expresa de las partes (por ejemplo, que el remitente la haya designado así).
    \item De la naturaleza misma de la comunicación (una carta dirigida a un amigo de confianza sobre asuntos personales).
\end{itemize}

Esta protección tiene fundamento constitucional: la \textbf{inviolabilidad de la correspondencia y los papeles privados} reconocida por la Constitución Nacional como garantía vinculada a la intimidad y privacidad de las personas.

\begin{example}{Uso de correspondencia confidencial}
Si alguien envía una carta a un amigo en la que reconoce adeudar sumas de dinero a terceros, y ese amigo entrega la carta a los acreedores para que la usen en un juicio, el uso de esa correspondencia puede ser cuestionado por su carácter confidencial. Sin el consentimiento del remitente, los destinatarios originales no pueden utilizarla.
\end{example}

% ============================================================
% CUADRO CORNELL
% ============================================================
\section{Cuadro Cornell — Repaso Integral}

\begin{table}[H]
\centering
\begin{tabularx}{\textwidth}{
    >{\raggedright\arraybackslash}p{0.30\textwidth}
    >{\raggedright\arraybackslash}X
}
\toprule
\textbf{Pregunta / Concepto clave} &
\textbf{Notas / Respuesta} \\
\midrule

\textbf{¿Qué es la forma del acto jurídico? ¿Es un elemento esencial?} &
Es la exteriorización de la voluntad de los sujetos respecto del objeto, en orden a conseguir el fin jurídico. Ningún acto puede carecer de forma: en ese sentido, la forma es un elemento esencial. \\[6pt]

\textbf{¿Cuáles son las dos dimensiones de la forma?} &
La forma tiene incidencia sobre la \textbf{validez} del acto (en ciertos casos, su incumplimiento produce nulidad) y sobre la \textbf{prueba} del acto (en otros casos, la forma es exigida sólo para facilitar su demostración). \\[6pt]

\textbf{¿Cuáles son las ventajas y desventajas del formalismo?} &
\textit{Ventajas:} facilita la prueba, publicita el acto, colabora en la percepción impositiva, permite distinguir el acto concluido de los preparatorios y protege a los terceros. \textit{Desventajas:} mayor costo, mayor riesgo de invalidez por defecto formal, menor celeridad en las transacciones, incomodidades prácticas. \\[6pt]

\textbf{¿Cuál es el principio general que rige la forma? (Art. 284)} &
El \textbf{principio de libertad de formas}: si la ley no designa una forma determinada, las partes pueden utilizar la que estimen conveniente. Las propias partes pueden pactar formas más exigentes que las legales. \\[6pt]

\textbf{¿Cómo se clasifican los actos según la forma?} &
Actos \textbf{formales} (con forma legal impuesta) y \textbf{no formales} (libertad de formas). Los formales se subdividen en \textbf{solemnes} (la forma es requisito de validez) y \textbf{no solemnes} o \emph{ad probationem} (la forma es exigida sólo para la prueba). \\[6pt]

\textbf{¿Qué diferencia hay entre solemnidad absoluta y solemnidad relativa?} &
En la \textbf{solemnidad absoluta}, el incumplimiento produce la nulidad del acto, sin ningún otro efecto. En la \textbf{solemnidad relativa}, el acto no produce los efectos propios del negocio, pero genera la obligación de otorgar la forma prescripta (conversión del acto jurídico, Art. 285). \\[6pt]

\textbf{¿Qué es la conversión del acto jurídico? (Art. 285)} &
Es el efecto propio de la solemnidad relativa: el acto no concluido en la forma requerida por la ley no queda concluido como tal, pero vale como acto por el cual las partes se obligan a cumplir con la formalidad, siempre que ésta no sea exigida bajo sanción de nulidad. Ejemplo: el boleto de compraventa genera la obligación de escriturar. \\[6pt]

\textbf{¿Cuáles son las categorías de instrumentos en el CCyCN? (Arts. 286--287)} &
(1) Instrumentos \textbf{públicos}: con autenticidad propia por ley. (2) Instrumentos \textbf{privados}: particulares firmados; requisito esencial: firma (Art. 288). (3) Instrumentos \textbf{particulares no firmados}: todo escrito, registro visual, auditivo o electrónico sin firma. \\[6pt]

\textbf{¿Qué prueba la firma? ¿Qué diferencia hay entre firma digital y firma electrónica? (Art. 288, Ley 25.506)} &
La firma prueba la \textbf{autoría} de la declaración de voluntad. La \textbf{firma digital} garantiza autoría e integridad mediante algoritmo matemático y certificado emitido por certificador licenciado. La \textbf{firma electrónica} es sólo un conjunto de datos identificatorios sin esa certificación; puede ser desconocida y su validez debe acreditarse. \\[6pt]

\textbf{¿Qué son los instrumentos públicos y cuáles son sus caracteres? (Art. 289)} &
Son los documentos a los que la ley les atribuye autenticidad propia, sin necesidad de reconocimiento previo. Hacen plena fe de su contenido. Enumeración: escrituras públicas y copias, instrumentos de escribanos o funcionarios públicos, y títulos emitidos por el Estado. \\[6pt]

\textbf{¿Cuáles son los requisitos de validez del instrumento público? (Art. 290)} &
(a) Actuación del oficial público dentro de sus \textbf{atribuciones} (competencia funcional) y \textbf{competencia territorial}. (b) \textbf{Firmas} del oficial, las partes y sus representantes. La falta de cualquier firma invalida el instrumento para todos. \\[6pt]

\textbf{¿Cómo funciona la eficacia probatoria del instrumento público? (Art. 296)} &
Distingue dos regímenes: (a) Los \textbf{hechos percibidos por el oficial} (fecha, lugar, actos cumplidos ante él) sólo pueden cuestionarse mediante \textbf{juicio de redargución de falsedad}. (b) El \textbf{contenido de las declaraciones de las partes} hace plena fe pero admite \textbf{prueba en contrario} sin necesidad de redargución. \\[6pt]

\textbf{¿Qué es la escritura pública y en qué se diferencia del acta notarial? (Arts. 299 y 310)} &
La \textbf{escritura pública} es el instrumento matriz extendido en el protocolo notarial que contiene \textbf{actos jurídicos}. El \textbf{acta notarial} es el documento notarial que tiene por objeto la comprobación de \textbf{hechos}. La distinción es esencial: la escritura crea, modifica o extingue derechos; el acta sólo constata hechos. \\[6pt]

\textbf{¿Qué estructura tiene la escritura pública? (Art. 305)} &
Tres partes: (1) \textbf{Encabezado}: lugar, fecha y datos de los comparecientes. (2) \textbf{Cuerpo}: naturaleza del acto, individualización de bienes y, en transmisiones de inmuebles, el ``corresponde''. (3) \textbf{Pie}: constancia de lectura, salvado de enmiendas, firmas. \\[6pt]

\textbf{¿Qué es el ``corresponde'' y para qué sirve?} &
Es la cláusula que, en las escrituras de transmisión de derechos reales sobre inmuebles, relaciona todos los antecedentes de titularidad del bien. Permite encadenar las sucesivas escrituras y verificar la historia dominial. Su omisión no invalida la escritura pero puede generar sanciones para el escribano. \\[6pt]

\textbf{¿Cuál es el requisito esencial del instrumento privado y cómo se reconoce la firma? (Arts. 288 y 314)} &
El requisito esencial es la \textbf{firma}. Para que el instrumento tenga plena fuerza probatoria entre las partes, la firma debe ser \textbf{reconocida}: expresa, tácita (el silencio frente a la intimación equivale a reconocimiento, Art. 314) o forzosa (pericia caligráfica). El reconocimiento de firma implica el reconocimiento de todo el contenido del instrumento. \\[6pt]

\textbf{¿Qué es la firma a ruego y cuál es la diferencia con la impresión digital? (Art. 313)} &
La \textbf{firma a ruego} es la firma de un tercero en nombre de quien no sabe o no puede firmar, en presencia de dos testigos. Equivale a la firma y otorga plena fuerza probatoria. La \textbf{impresión digital} sólo vale como \emph{principio de prueba por escrito} (Art. 314) y puede ser impugnada, por su menor seguridad. \\[6pt]

\textbf{¿Qué es la fecha cierta y para qué sirve? (Art. 317)} &
Es el momento a partir del cual existe certeza objetiva de que el instrumento ya estaba firmado o no pudo firmarse después. Desde la fecha cierta, el instrumento privado es \textbf{oponible a terceros}. El CCyCN abandona la enumeración taxativa del Código anterior: la fecha cierta puede acreditarse por cualquier medio, que el juez apreciará rigurosamente. \\[6pt]

\textbf{¿Cómo valora el juez un instrumento particular no firmado? (Art. 319)} &
Debe ponderar: congruencia entre lo narrado y lo ocurrido; precisión y claridad técnica del texto; usos y prácticas del tráfico; relaciones precedentes entre las partes; y confiabilidad de los soportes y procedimientos técnicos empleados. \\[6pt]

\textbf{¿Puede utilizarse la correspondencia como prueba en juicio? (Art. 318)} &
Sí: el destinatario puede presentarla como prueba contra el remitente. Sin embargo, si la correspondencia es \textbf{confidencial} —por manifestación expresa o por la naturaleza de la comunicación—, no puede utilizarse sin consentimiento del remitente. Esta protección tiene fundamento constitucional (inviolabilidad de la correspondencia). \\

\midrule

\multicolumn{2}{p{\textwidth}}{%
\textbf{Resumen:} La forma del acto jurídico es el modo en que la voluntad de las partes se exterioriza y produce efectos jurídicos. El CCyCN parte del principio de libertad de formas (Art. 284), pero impone formas determinadas para los actos más trascendentes, con distintas consecuencias según se trate de solemnidad absoluta (nulidad del acto) o relativa (conversión: obligación de otorgar la forma prescripta, Art. 285). En materia de instrumentos, el Código distingue tres grandes categorías: los \textbf{instrumentos públicos}, que gozan de autenticidad propia y hacen plena fe, con dos regímenes de impugnación según se trate de hechos percibidos por el oficial (redargución de falsedad) o de declaraciones de las partes (prueba en contrario); la \textbf{escritura pública}, instrumento matriz del protocolo notarial que contiene actos jurídicos y requiere una estructura formal precisa; y los \textbf{instrumentos particulares}, que incluyen los privados (requisito esencial: firma; fuerza probatoria plena tras el reconocimiento) y los particulares no firmados (fuerza probatoria sujeta a la apreciación judicial según las pautas del Art. 319). La \textbf{firma digital} (Ley 25.506) garantiza autoría e integridad en el mundo digital, con mayor valor probatorio que la firma electrónica. La \textbf{fecha cierta} (Art. 317) determina desde cuándo el instrumento privado puede oponerse a terceros, y puede acreditarse por cualquier medio apreciado rigurosamente por el juez. Estas reglas configuran un sistema orientado a dar seguridad jurídica al tráfico de bienes y servicios, equilibrando la certeza del formalismo con la agilidad que demanda la vida contemporánea.
}\\

\bottomrule
\end{tabularx}
\end{table}

% ============================================================

% ============================================================
\chapter{Clasificación de los Actos Jurídicos}
% =====================================================

En el régimen del Código de Vélez se enunciaba expresamente una serie de clasificaciones de los actos jurídicos --unilaterales y bilaterales, entre otras--. El Código Civil y Comercial de la Nación (CCyC) vigente, en cambio, optó metodológicamente por no incluir una clasificación sistematizada de los actos jurídicos: no existen artículos dedicados a definir de manera unificada las distintas clases de actos. Sin embargo, estas clasificaciones subsisten y se manifiestan a lo largo de todo el articulado, en la medida en que el Código atribuye consecuencias jurídicas distintas según el acto sea bilateral, entre vivos o de última voluntad, oneroso o gratuito, entre otras categorías.

\begin{note}
Estas clasificaciones no son puramente teóricas, sino que tienen consecuencias prácticas concretas y se encuentran presentes en distintos lugares del Código. A lo largo de las distintas asignaturas del plan de estudios de abogacía es habitual encontrar aplicaciones de estas categorías.
\end{note}

Su relevancia práctica es amplia: saber si un acto es formal o no formal determina si se requiere escritura pública; saber si es oneroso o gratuito modifica las reglas de buena fe aplicables (art. 392, CCyC); y saber si el acto está sujeto a condición, plazo o cargo indica en qué momento nace, se suspende o se extingue la obligación asumida. Estas categorías atraviesan los contratos, las sucesiones, el derecho de familia, las obligaciones y los derechos reales, y constituyen el marco conceptual necesario para interpretar el Código en su conjunto.

\begin{example}{Aplicación integradora de las clasificaciones}
En la redacción de una compraventa de un inmueble, dicho acto resulta simultáneamente bilateral, oneroso, formal (requiere escritura pública), patrimonial y de disposición. Si además las partes acuerdan que la entrega quede sujeta a que el banco apruebe el crédito del comprador, se configura una condición suspensiva; si se pacta que el contrato se resuelve en caso de que el banco no se expida dentro de los sesenta días, se configura un plazo resolutorio; y si el vendedor se compromete a pintar el departamento antes de la entrega, se configura un cargo. De este modo, las clasificaciones y las modalidades estudiadas en este capítulo convergen en un único documento cotidiano de la práctica profesional.
\end{example}

\section{Actos Unilaterales y Bilaterales}

El criterio de esta clasificación es el número de partes que concurren a la celebración del acto.

\begin{definition}{Actos unilaterales y bilaterales}
El acto es \textbf{unilateral} cuando una sola parte concurre a su celebración. Es \textbf{bilateral} cuando dos o más partes concurren a su celebración.
\end{definition}

\begin{table}[H]
\centering
\begin{tabularx}{\textwidth}{l X l}
\toprule
\textbf{Tipo de acto} & \textbf{Definición} & \textbf{Ejemplo típico} \\
\midrule
Unilateral & Una sola parte concurre a la celebración del acto. & El testamento \\
Bilateral & Dos o más partes concurren a la celebración del acto. & Un contrato; el matrimonio \\
\bottomrule
\end{tabularx}
\caption{Actos unilaterales y bilaterales}
\end{table}

\section{Actos entre Vivos y de Última Voluntad}

\begin{definition}{Actos entre vivos y de última voluntad}
Los actos \textbf{entre vivos} son otorgados por personas vivas para producir efectos en vida de los otorgantes. Los actos \textbf{de última voluntad} son aquellos que producen sus efectos luego de la muerte de quien los otorga.
\end{definition}

\begin{table}[H]
\centering
\begin{tabularx}{\textwidth}{l X X}
\toprule
\textbf{Tipo de acto} & \textbf{Definición} & \textbf{Cuándo produce efectos} \\
\midrule
Entre vivos & Son otorgados por personas vivas para producir efectos en vida. & Desde el momento de su celebración \\
De última voluntad & Producen sus efectos luego de la muerte del otorgante. & Desde la muerte del otorgante \\
\bottomrule
\end{tabularx}
\caption{Actos entre vivos y de última voluntad}
\end{table}

\section{Actos a Título Oneroso y a Título Gratuito}

Esta clasificación resulta especialmente relevante a lo largo de todo el Código, dado que el CCyC tiende a proteger de manera diferenciada a quienes adquieren a título oneroso y de buena fe frente a quienes adquieren a título gratuito.

\begin{definition}{Actos onerosos y gratuitos}
El acto es \textbf{oneroso} cuando existen prestaciones recíprocas: cada parte entrega algo a cambio de lo que recibe. El acto es \textbf{gratuito} (o a título de liberalidad) cuando una de las partes recibe una prestación sin dar nada a cambio, lo cual enriquece a esa parte.
\end{definition}

\begin{table}[H]
\centering
\begin{tabularx}{\textwidth}{l X X}
\toprule
\textbf{Tipo de acto} & \textbf{Definición} & \textbf{Ejemplo} \\
\midrule
Oneroso & Hay prestaciones recíprocas: cada parte da algo a cambio. & Compraventa \\
Gratuito (liberalidad) & Una de las partes recibe sin dar nada a cambio; enriquece a esa parte. & Donación; renuncia a un derecho \\
\bottomrule
\end{tabularx}
\caption{Actos a título oneroso y a título gratuito}
\end{table}

El caso típico del título gratuito es la donación, pero también puede configurarse a título gratuito la renuncia a un derecho: quien tiene derecho a cobrar y renuncia a ese cobro está beneficiando a la otra parte sin recibir nada a cambio.

\begin{formula}{Art. 392, CCyC --- Efectos respecto de terceros en cosas registrables}
El Código diferencia la protección al subadquirente según haya adquirido a título oneroso o gratuito. El subadquirente de buena fe y a título oneroso goza de mayor protección.
\end{formula}

\section{Actos Principales y Accesorios}

\begin{definition}{Actos principales y accesorios}
Los actos \textbf{principales} son aquellos que no dependen, para su existencia, de otro acto. Los actos \textbf{accesorios} son aquellos que siguen la suerte del acto principal, del cual dependen para existir.
\end{definition}

\begin{table}[H]
\centering
\begin{tabularx}{\textwidth}{l X X}
\toprule
\textbf{Tipo de acto} & \textbf{Definición} & \textbf{Ejemplo} \\
\midrule
Principal & No depende de otro acto para existir. & Contrato de alquiler \\
Accesorio & Sigue la suerte del principal; depende de él para existir. & La garantía del contrato de alquiler \\
\bottomrule
\end{tabularx}
\caption{Actos principales y accesorios}
\end{table}

A modo de ejemplo, una garantía que acompaña a un contrato de alquiler constituye un acto accesorio respecto del principal y sigue su suerte, sin perjuicio de los matices particulares que se estudian en cada contrato en concreto.

\section{Actos Formales y No Formales}

Esta es una clasificación de particular importancia, que se desarrolla en profundidad en el capítulo~\ref{ch:forma}.

\begin{definition}{Actos formales y no formales}
Los actos \textbf{formales} son aquellos en los que la ley prescribe cierta solemnidad. Dentro de estos, la solemnidad puede hacer a la validez del acto (\textbf{formal solemne} o absolutamente formal) o hacer únicamente a su prueba (\textbf{formal no solemne} o relativamente formal). Los actos \textbf{no formales} son aquellos respecto de los cuales la ley no exige una forma determinada, quedando las partes en libertad de elegir la que prefieran.
\end{definition}

\begin{table}[H]
\centering
\begin{tabularx}{\textwidth}{l X X}
\toprule
\textbf{Tipo de acto} & \textbf{Definición} & \textbf{Consecuencia del incumplimiento de la forma} \\
\midrule
Formal solemne (absolutamente formal) & La ley prescribe una solemnidad cuyo incumplimiento invalida el acto. & Nulidad del acto \\
Formal no solemne (relativamente formal) & La ley prescribe una forma, pero su incumplimiento afecta solo la prueba, no la validez. & Puede probarse por otro medio \\
No formal & La ley no manda ninguna forma determinada; las partes pueden elegir la que prefieran. & No hay consecuencia: el acto es válido \\
\bottomrule
\end{tabularx}
\caption{Actos formales y no formales}
\end{table}

\begin{example}{Acto formal solemne}
El matrimonio es un acto formal solemne: se encuentra rodeado de una serie de solemnidades marcadas por el Código Civil y no puede celebrarse sino de acuerdo con los rituales que este establece. Cualquier otra forma determina que el matrimonio no sea válido.
\end{example}

\begin{example}{Acto no formal}
El alquiler de un inmueble es, en principio, un acto no formal: la ley no exige que se celebre por escritura pública. No obstante, las partes pueden optar por una forma más solemne que la exigida legalmente --por ejemplo, contratando a un escribano-- por razones de seguridad jurídica, sin que ello sea un requisito legal.
\end{example}

\section{Actos Patrimoniales y Extrapatrimoniales}

\begin{definition}{Actos patrimoniales y extrapatrimoniales}
Los actos \textbf{patrimoniales} tienen contenido económico y están vinculados con el patrimonio. Los actos \textbf{extrapatrimoniales} no tienen contenido económico directo y se vinculan con derechos personalísimos o con el derecho de familia.
\end{definition}

\begin{table}[H]
\centering
\begin{tabularx}{\textwidth}{l X X}
\toprule
\textbf{Tipo de acto} & \textbf{Definición} & \textbf{Ejemplo} \\
\midrule
Patrimonial & Tiene contenido económico; está vinculado con el patrimonio. & Un contrato de compraventa \\
Extrapatrimonial & No tiene contenido económico directo; vinculado con derechos personalísimos o de familia. & El reconocimiento de un hijo \\
\bottomrule
\end{tabularx}
\caption{Actos patrimoniales y extrapatrimoniales}
\end{table}

El reconocimiento de un hijo es un acto jurídico unilateral mediante el cual una persona expresa su voluntad de reconocerlo. Se trata de un acto típico del derecho de familia, extrapatrimonial en sí mismo, aun cuando pueda tener consecuencias patrimoniales indirectas, como el derecho a alimentos. Los actos vinculados con derechos personalísimos también son extrapatrimoniales.

\section{Actos de Disposición, de Administración y de Conservación}

\begin{definition}{Actos de disposición, administración y conservación}
Los actos de \textbf{disposición} producen una variación sustancial en la composición del patrimonio. Los actos de \textbf{administración} no producen un cambio sustancial en el patrimonio. Los actos de \textbf{conservación} están dedicados al mantenimiento de los bienes.
\end{definition}

\begin{table}[H]
\centering
\begin{tabularx}{\textwidth}{l X X}
\toprule
\textbf{Tipo de acto} & \textbf{Definición} & \textbf{Ejemplo} \\
\midrule
Disposición & Se produce una variación sustancial en la composición del patrimonio. & Venta de un inmueble: sale el bien, entra dinero \\
Administración & No produce un cambio sustancial en el patrimonio. & Cobrar un alquiler, recibir una renta \\
Conservación & Dedicados al mantenimiento de los bienes. & Reparaciones, actos de preservación \\
\bottomrule
\end{tabularx}
\caption{Actos de disposición, administración y conservación}
\end{table}

\begin{note}
Esta clasificación resulta especialmente relevante en el derecho de familia y de las personas, en materia de tutela, curatela y régimen de poderes. El representante legal de una persona incapaz, por ejemplo, necesita autorización judicial para los actos de disposición, pero no para los de administración ordinaria.
\end{note}

\section{Actos Puros y Simples, y Actos Modales}

\begin{definition}{Actos puros y simples, y actos modales}
Los actos \textbf{puros y simples} no están sujetos a ninguna modalidad y producen sus efectos de inmediato. Los actos \textbf{modales} son aquellos en los cuales las partes eligen sujetarlo a alguno de los tres modos posibles: condición, plazo o cargo.
\end{definition}

\begin{table}[H]
\centering
\begin{tabularx}{\textwidth}{l X}
\toprule
\textbf{Tipo de acto} & \textbf{Definición} \\
\midrule
Puro y simple & No está sujeto a ninguna modalidad. Produce efectos de inmediato. \\
Modal & Las partes eligen sujetarlo a alguno de los tres modos: condición, plazo o cargo. \\
\bottomrule
\end{tabularx}
\caption{Actos puros y simples, y actos modales}
\end{table}

\section{Actos Causales y Abstractos}

La noción de acto abstracto se desarrolla también al estudiar el régimen de la causa (capítulo~\ref{ch:causa}).


\begin{definition}{Actos causales y abstractos}
El acto \textbf{causal} permanece vinculado a su causa, la cual influye en su validez y en sus efectos. El acto \textbf{abstracto} se independiza de su causa.
\end{definition}

\begin{table}[H]
\centering
\begin{tabularx}{\textwidth}{l X}
\toprule
\textbf{Tipo de acto} & \textbf{Definición} \\
\midrule
Causal & El acto permanece vinculado a su causa; la causa influye en su validez y efectos. \\
Abstracto & El acto se independiza de su causa. \\
\bottomrule
\end{tabularx}
\caption{Actos causales y abstractos}
\end{table}

% TODO: contenido incompleto o ambiguo en las notas originales -- el apunte remite a un desarrollo previo de los actos abstractos ("se estudiaron previamente en la clase") que no forma parte de este documento.

\section{Actos Positivos y Negativos}

\begin{definition}{Actos positivos y negativos}
El acto \textbf{positivo} involucra una acción, es decir, un hacer. El acto \textbf{negativo} involucra una omisión, es decir, un no hacer.
\end{definition}

\begin{table}[H]
\centering
\begin{tabularx}{\textwidth}{l X X}
\toprule
\textbf{Tipo de acto} & \textbf{Definición} & \textbf{Ejemplo} \\
\midrule
Positivo & Involucra una acción (un hacer). & Entregar una cosa, pagar una suma \\
Negativo & Involucra una omisión (un no hacer). & Comprometerse a no competir en un rubro comercial \\
\bottomrule
\end{tabularx}
\caption{Actos positivos y negativos}
\end{table}

\begin{example}{Acto negativo: cláusula de no competencia}
Es habitual que quienes se dedican a una determinada actividad --por ejemplo, la producción de helados-- al vender la marca y la empresa se comprometan a no generar una nueva empresa ni emprendimientos vinculados con ese rubro durante los próximos cinco años. Este acto jurídico se encuadra dentro de los actos negativos, en tanto las partes se obligan a un ``no hacer'', es decir, a una omisión.
\end{example}

% =====================================================
\chapter{Modalidades del acto jurídico}
\label{ch:modalidades}
% =====================================================

Las modalidades del acto jurídico son la condición, el plazo y el cargo. En el régimen del Código de Vélez, las modalidades estaban reguladas dentro de la sección de obligaciones del Libro Segundo. En el Código Civil y Comercial vigente, en cambio, se las incorporó al Libro Primero, Título Cuarto, dentro del capítulo de hechos y actos jurídicos, en un capítulo específico dedicado a las modalidades del acto jurídico.

\section{La Condición}

\begin{formula}{Art. 343, CCyC --- Alcances}
Se denomina condición a la cláusula de los actos jurídicos por la cual las partes subordinan su plena eficacia o resolución a un hecho futuro e incierto.
\end{formula}

\begin{definition}{Condición}
La condición es una cláusula por la cual se subordina la adquisición o la resolución de un acto a un hecho futuro e incierto. Lo esencial de la noción es la incertidumbre: el hecho condicionante puede o no llegar a suceder.
\end{definition}

\begin{example}{Condición suspensiva (1)}
``Te voy a regalar un auto si Argentina sale campeón del mundo en el próximo mundial.'' Se trata de un hecho futuro e incierto: puede ocurrir o no. El efecto --la donación del auto-- queda suspendido hasta que ocurra, o no, el hecho condicionante.
\end{example}

\begin{example}{Condición suspensiva (2)}
``Te voy a regalar mi biblioteca jurídica si te recibís de abogado/a.'' El hecho es incierto: puede ocurrir o no. La obligación existe desde la constitución del acto, pero no es exigible --no tiene plena eficacia-- hasta el momento en que ocurre el hecho condicionante.
\end{example}

\begin{example}{Condición resolutoria}
``Te dono mi auto, pero siempre y cuando mi hermano no vuelva de España. Si mi hermano vuelve, tenés que devolverme el auto.'' Aquí el acto ya produce efectos desde el principio; el hecho condicionante produce la resolución, es decir, la devolución de los efectos hacia atrás.
\end{example}

\begin{table}[H]
\centering
\begin{tabularx}{\textwidth}{l X X}
\toprule
\textbf{Tipo de condición} & \textbf{Punto de partida} & \textbf{Efecto del hecho condicionante} \\
\midrule
Suspensiva & El acto existe pero NO produce efectos todavía. & Si ocurre el hecho, el acto produce sus efectos; si no ocurre, el acto no produce efectos. \\
Resolutoria & El acto YA produce efectos desde su celebración. & Si ocurre el hecho, los efectos cesan (se resuelven), con retroactividad. \\
\bottomrule
\end{tabularx}
\caption{Diferencia entre condición suspensiva y resolutoria}
\end{table}

\subsection{La Condición Impropia o Suposición}

El nuevo Código incorporó también la llamada ``condición impropia'' o ``suposición'': aquella en la que las partes condicionan el acto a un hecho presente o pasado que, al momento del acto, es ignorado por ellas.

\begin{formula}{Art. 343, CCyC (segundo párrafo)}
Se equipara a la condición la cláusula por la cual las partes sujetan la eficacia del acto a la existencia de cierto estado de cosas al tiempo de celebrarlo.
\end{formula}

\begin{example}{Condición impropia o suposición}
Una persona ignora la situación patrimonial de un familiar que vive lejos y, sin saberlo, celebra un acto sujeto a que ese familiar no haya caído en quiebra: por ejemplo, ``te presto la plata en la medida en que tal familiar mío no haya terminado en quiebra''. Aunque el hecho ya ocurrió o no ocurrió al momento de celebrar el acto, la incertidumbre subjetiva de las partes --que lo ignoran-- hace que este supuesto también valga como condición.
\end{example}

\subsection{Clasificación de las Condiciones}

\begin{table}[H]
\centering
\begin{tabularx}{\textwidth}{l l X}
\toprule
\textbf{Criterio} & \textbf{Tipo} & \textbf{Descripción} \\
\midrule
Según el hecho & Positiva & Implica una acción (que algo ocurra). \\
Según el hecho & Negativa & Implica una omisión (que algo no ocurra). \\
Según el efecto & Suspensiva & Subordina la eficacia del acto a un hecho futuro e incierto. \\
Según el efecto & Resolutoria & Subordina la extinción del acto a un hecho futuro e incierto. \\
Según la fuente & Causal & El hecho no depende de las partes; depende del azar o de un tercero. \\
Según la fuente & Potestativa & Depende de la voluntad del deudor o del acreedor. \\
Según la fuente & Mixta & Depende en parte del deudor y en parte de un tercero o del azar. \\
\bottomrule
\end{tabularx}
\caption{Clasificación de las condiciones}
\end{table}

La condición \textbf{causal} no depende de las partes, sino del azar o de un tercero --por ejemplo, que la selección nacional resulte campeona--. La condición \textbf{potestativa} depende del acreedor o del deudor --por ejemplo, ``te voy a regalar mi biblioteca si te recibís de abogado/a'' es potestativa del acreedor--. La condición \textbf{mixta} depende en parte del deudor y en parte de un hecho ajeno a él.

\begin{important}
La condición puramente potestativa del deudor está prohibida: ``te comprometo a entregarte algo si yo decido hacerlo'' no puede depender completamente de la voluntad del obligado.
\end{important}

\subsection{Condiciones Prohibidas}

\begin{formula}{Art. 344, CCyC --- Condiciones prohibidas}
Es nula la condición que: (1) sea un hecho imposible, (2) sea contraria a las buenas costumbres, (3) esté prohibida por el ordenamiento jurídico, (4) dependa exclusivamente de la voluntad del obligado (puramente potestativa). La condición de no hacer algo imposible no perjudica la validez de la obligación.
\end{formula}

\begin{formula}{Art. 345, CCyC --- Condiciones que gravan libertades fundamentales}
Las condiciones cuyo cumplimiento depende de que el beneficiario cambie de domicilio, religión u otras libertades fundamentales se tienen por no escritas.
\end{formula}

\begin{example}{Condición que grava una libertad fundamental}
``Te dono la casa si cambiás de religión'' constituye una condición prohibida, en tanto grava la libertad fundamental de profesar la religión que la persona elija.
\end{example}

\begin{table}[H]
\centering
\begin{tabularx}{\textwidth}{X X}
\toprule
\textbf{Tipo de condición prohibida} & \textbf{Consecuencia jurídica} \\
\midrule
Hecho imposible / contraria a buenas costumbres / puramente potestativa (art. 344, párr. 1) & Nulidad absoluta del acto \\
Que grava libertades fundamentales (art. 345, párr. 3) & Se tiene por no escrita (la obligación principal podría subsistir) \\
Imposible en testamentos (excepción) & La condición se tiene por no puesta; la disposición vale \\
\bottomrule
\end{tabularx}
\caption{Condiciones prohibidas y sus consecuencias}
\end{table}

Existe una diferencia entre el primer y el tercer párrafo del artículo 344: las condiciones ilícitas del primer párrafo producen la nulidad absoluta del acto, mientras que las condiciones comprendidas en el tercer párrafo se tienen por no escritas, discutiéndose en la doctrina si la obligación principal subsiste o no. La excepción se presenta en los testamentos: las condiciones prohibidas incorporadas en un testamento no afectan la validez de las disposiciones testamentarias de las que forman parte.

\section{El Plazo}

A diferencia de la condición, el plazo es la modalidad que posterga o limita temporalmente los efectos del acto. En el plazo hay un hecho futuro pero \textbf{cierto}: se sabe que va a ocurrir, aunque no siempre se sepa exactamente cuándo.

\begin{example}{Plazo suspensivo y plazo resolutorio}
Plazo suspensivo: ``Te voy a donar estas casas dentro de un año.'' No hay incertidumbre sobre si llegará el día, aunque pueda haberla sobre el momento exacto en el caso del plazo incierto. Plazo resolutorio: ``Te presto el uso de la casa por un año.'' Los efectos comienzan de inmediato, pero cesan cuando vence el plazo.
\end{example}

\begin{table}[H]
\centering
\begin{tabularx}{\textwidth}{l X X}
\toprule
\textbf{Tipo de plazo} & \textbf{Definición} & \textbf{Ejemplo} \\
\midrule
Suspensivo & Posterga los efectos del acto hasta el vencimiento del plazo. & ``Te dono la casa dentro de un año.'' \\
Resolutorio & Los efectos del acto cesan al vencimiento del plazo. & ``Te presto la casa por un año.'' \\
Cierto & Se conoce la fecha exacta del vencimiento. & ``En 30 días'', ``el 31 de diciembre'' \\
Incierto & No se sabe exactamente cuándo vencerá, pero es seguro que ocurrirá. & ``Cuando muera tal persona'' \\
Determinado & Las partes, la ley o el juez fijaron la fecha o el momento. & La ley o el contrato establece una fecha \\
Indeterminado & No se precisó fecha, pero existe un plazo posible. & Cláusula ``a mejor fortuna'' (art. 889, CCyC) \\
Esencial & El cumplimiento fuera de término no satisface el interés del acreedor. & Catering para una fiesta: si llega dos semanas tarde, ya no sirve \\
Accidental & El cumplimiento tardío puede igualmente satisfacer el interés del acreedor. & Entrega de materiales de construcción con demora \\
\bottomrule
\end{tabularx}
\caption{Clasificación del plazo}
\end{table}

\begin{example}{Plazo esencial}
Una persona contrata un servicio de catering para una fiesta y, por un error, el servicio se entrega dos semanas después del evento. El plazo era esencial --la fiesta misma-- por lo que la prestación tardía no satisface el interés del acreedor y el incumplimiento es total.
\end{example}

El plazo puede surgir de la convención de las partes, de la ley, puede ser judicial, o bien indeterminado, cuando se establece la fecha de inicio pero no la de finalización.

\begin{formula}{Art. 889, CCyC --- Cláusula ``a mejor fortuna''}
El Código admite la cláusula ``a mejor fortuna'': ``yo te presto plata y me devolvés cuando estés bien económicamente.'' Se trata de un plazo indeterminado. El juez, a solicitud del acreedor, puede fijar un plazo y hacer cumplir la obligación.
\end{formula}

\section{El Cargo}

\begin{definition}{Cargo}
El cargo es una obligación accesoria impuesta al adquirente de un derecho, que no impide los efectos del acto pero que, sin embargo, puede exigirse.
\end{definition}

\begin{example}{Cargo simple}
Una persona dona un campo a otra con el cargo de que mantenga en la casa a los caseros. Si el donatario no cumple, los caseros podrán exigir el cumplimiento del cargo o reclamar la indemnización de los daños y perjuicios ocasionados por el incumplimiento; sin embargo, en principio, el acto de donación no se cae.
\end{example}

\begin{example}{Cargo con finalidad benéfica}
Se dona a una fundación un terreno con el cargo de que el producto de la explotación económica de ese terreno se destine a obras para los barrios vulnerables. La fundación que recibe la donación debe administrar el bien conforme a ese cargo.
\end{example}

\begin{table}[H]
\centering
\begin{tabularx}{\textwidth}{X X X}
\toprule
\textbf{Característica} & \textbf{Condición} & \textbf{Cargo} \\
\midrule
Suspende efectos del acto & SÍ (si es suspensiva) & NO \\
Coercible (exigible) & Cuando se cumple el hecho: sí & Sí, siempre \\
Si falla en cumplirlo, revoca el acto & Sí (si es resolutoria) & Solo si fue impuesto como cargo condicional \\
Quién puede exigirlo & Las partes & Los beneficiarios y quien constituyó el derecho \\
\bottomrule
\end{tabularx}
\caption{Diferencias entre condición y cargo}
\end{table}

El cargo es coercitivo pero no suspensivo, a diferencia de la condición suspensiva, que sí suspende los efectos. El cargo siempre se impone al adquirente como una obligación que debe cumplir y, en tanto obligación, debe reunir los requisitos propios de las obligaciones: ser lícito, posible, determinado y de contenido patrimonial. Surge de la autonomía de la voluntad, es excepcional y accidental --puede o no estar presente--, constituye una obligación accesoria del derecho transmitido, debe ser expreso y es restrictivo del derecho adquirido, en la medida en que recibir algo con el cargo de cumplir una obligación restringe el derecho de quien lo recibe.

\subsection{Cargo Simple y Cargo Condicional}

\begin{table}[H]
\centering
\begin{tabularx}{\textwidth}{l X X}
\toprule
\textbf{Tipo de cargo} & \textbf{Definición} & \textbf{Consecuencia del incumplimiento} \\
\midrule
Cargo simple & Obligación accesoria que no condiciona la adquisición del derecho. & Se exige el cumplimiento o se reclaman daños. El acto no se revoca. \\
Cargo condicional & El cumplimiento del cargo fue puesto como condición del acto. & Si no se cumple, se puede revertir la adquisición del derecho. Opera como condición resolutoria. \\
\bottomrule
\end{tabularx}
\caption{Cargo simple y cargo condicional}
\end{table}

El cargo simple no tiene efecto resolutorio sobre el derecho principal. El cargo condicional, en cambio, opera como una condición resolutoria (o suspensiva), de modo que su incumplimiento afecta a la propia adquisición del derecho: si el cargo fue puesto como condicional, la adquisición puede revertirse en caso de incumplimiento.

\subsection{Cargos Prohibidos}

\begin{formula}{Art. 357, CCyC --- Cargos prohibidos}
Son de aplicación a los cargos las normas establecidas para las condiciones prohibidas (art. 344, CCyC). Sin embargo, cuando se impone un cargo prohibido, no se produce la nulidad del acto, sino que el cargo se tiene por no escrito.
\end{formula}

Los cargos prohibidos son, en su contenido, los mismos supuestos previstos para las condiciones prohibidas: el artículo 357 remite al artículo 344. La diferencia radica en la consecuencia: cuando se impone un cargo prohibido, este se tiene por no escrito, sin que ello acarree la nulidad del acto principal.

% =====================================================

\section{Cierre y Síntesis}
% =====================================================

Al finalizar el desarrollo de las clasificaciones y las modalidades del acto jurídico --condición, plazo y cargo-- corresponde una aclaración fundamental sobre la naturaleza de estas últimas.

\begin{important}
Las modalidades --condición, plazo y cargo-- \textbf{no} son una clasificación de los actos jurídicos. Son modos, cláusulas, elementos accidentales de un acto jurídico: pueden o no estar presentes. Si las partes las pactan, rigen las reglas de su regulación civil; si no las incorporan, el acto es puro y simple y produce efectos de inmediato.
\end{important}

Los actos jurídicos son manifestaciones de voluntad orientadas a producir efectos jurídicos. El Código Civil y Comercial no los clasifica en un artículo único, sino que sus categorías atraviesan todo el texto legal. Estas clasificaciones --unilateral/bilateral, entre vivos/última voluntad, oneroso/gratuito, principal/accesorio, formal/no formal, patrimonial/extrapatrimonial, disposición/administración/conservación, puro/modal, causal/abstracto, positivo/negativo-- no son puramente teóricas: determinan las reglas aplicables en cada situación jurídica concreta. Sobre esa base, el acto puede quedar sujeto a modalidades accidentales: la condición (hecho futuro e incierto que suspende o resuelve los efectos), el plazo (hecho futuro cierto que posterga o limita los efectos en el tiempo) y el cargo (obligación accesoria que no suspende el acto, pero que es coercitiva). Cada modalidad presenta variantes propias, reglas de validez específicas y artículos particulares dentro del CCyC (arts. 343 a 357).

% =====================================================

\section{Cuadro Cornell de Repaso}
% =====================================================

El siguiente cuadro reúne, según el método Cornell, las preguntas centrales que permiten repasar activamente el contenido desarrollado en los capítulos anteriores, junto con sus respuestas y una síntesis final integradora.

\begin{longtable}{
    >{\raggedright\arraybackslash}p{0.30\textwidth}
    >{\raggedright\arraybackslash}p{0.62\textwidth}
}
\toprule
\textbf{Preguntas / palabras clave} & \textbf{Notas / respuestas} \\
\midrule
\endfirsthead
\toprule
\textbf{Preguntas / palabras clave} & \textbf{Notas / respuestas} \\
\midrule
\endhead

\textbf{¿Qué son las clasificaciones de los actos jurídicos y para qué sirven?} &
Son categorías doctrinarias --y en gran parte también legales-- que agrupan los actos según criterios como el número de partes, el momento de producción de efectos, la existencia de prestaciones recíprocas, la dependencia respecto de otro acto, la forma, el contenido patrimonial o el nivel de variación del patrimonio. No son puramente teóricas: determinan las reglas aplicables al acto en cada situación concreta del Código y de la práctica profesional. \\

\textbf{Diferencia entre acto unilateral y bilateral. Ejemplos.} &
Unilateral: una sola parte (el testamento). Bilateral: dos o más partes (el contrato, el matrimonio). Esta diferencia resulta clave porque el régimen de perfeccionamiento, revocación y efectos es distinto en cada caso. \\

\textbf{¿Por qué importa si un acto es oneroso o gratuito?} &
El CCyC protege de manera distinta a los subadquirentes según hayan adquirido a título oneroso o gratuito. Conforme el art. 392, CCyC, quien adquiere de buena fe y a título oneroso goza de mayor protección frente a la evicción y otras acciones de terceros que quien adquirió gratuitamente. \\

\textbf{Acto formal solemne vs.\ formal no solemne vs.\ no formal.} &
Formal solemne: la forma es requerida para la validez; si no se cumple, el acto es nulo (testamento, matrimonio). Formal no solemne: la forma prescrita afecta solo la prueba; el acto igual es válido. No formal: la ley no exige ninguna forma; las partes eligen cómo exteriorizar su voluntad. \\

\textbf{¿Qué es la condición? Definición legal.} &
Cláusula por la cual se subordina la adquisición o la resolución de un acto a un hecho futuro e incierto (art. 343, CCyC). Lo esencial es la incertidumbre: el hecho puede o no ocurrir. Si el hecho es cierto, se trata de un plazo, no de una condición. \\

\textbf{Condición suspensiva vs.\ resolutoria.} &
Suspensiva: los efectos del acto no se producen hasta que ocurra el hecho condicionante; el acto existe, pero no es exigible. Resolutoria: el acto ya produce efectos; si ocurre el hecho, los efectos cesan y se resuelven retroactivamente. Ejemplo de suspensiva: ``te dono la biblioteca si te recibís''. Ejemplo de resolutoria: ``te dono el auto, pero si vuelve mi hermano, me lo devolvés''. \\

\textbf{¿Qué condiciones están prohibidas? Consecuencias.} &
Son nulas las condiciones que: (1) son imposibles, (2) son contrarias a las buenas costumbres, (3) son puramente potestativas del deudor (art. 344, CCyC). Se tienen por no escritas las que gravan libertades fundamentales, como la religión o el domicilio (art. 345). En los testamentos, las condiciones prohibidas no afectan la validez de la disposición testamentaria. \\

\textbf{Diferencia entre condición y plazo.} &
En ambos casos el hecho es futuro. La diferencia clave es la certeza: en la condición el hecho es incierto (puede no ocurrir); en el plazo el hecho es cierto (se sabe que ocurrirá, aunque no siempre se sepa exactamente cuándo, como en el plazo incierto: ``cuando muera tal persona''). \\

\textbf{Tipos de plazo según su certeza y determinación.} &
Cierto: se conoce la fecha exacta. Incierto: no se sabe cuándo, pero es seguro que ocurrirá. Determinado: fijado por las partes, la ley o el juez. Indeterminado: sin fecha precisada (por ejemplo, la cláusula ``a mejor fortuna'', art. 889, CCyC). Esencial: el cumplimiento tardío no satisface el interés del acreedor (catering para una fiesta). Accidental: el cumplimiento tardío igualmente puede satisfacer el interés del acreedor. \\

\textbf{¿Qué es el cargo? ¿Suspende los efectos del acto?} &
El cargo es una obligación accesoria impuesta al adquirente de un derecho que NO suspende los efectos del acto: el acto produce efectos de inmediato. Sin embargo, el cargo es coercitivo: puede exigirse su cumplimiento. Si fue impuesto como condición (cargo condicional), su incumplimiento puede revertir la adquisición del derecho. Los cargos prohibidos se tienen por no escritos (art. 357, que remite al art. 344, CCyC), sin que ello invalide el acto principal. \\

\textbf{¿Son las modalidades una clasificación de los actos jurídicos?} &
No. Las modalidades (condición, plazo, cargo) son elementos accidentales del acto jurídico: pueden o no estar presentes. Si las partes las incorporan, el acto queda sujeto a esas reglas; si no las incorporan, el acto es puro y simple y produce efectos de inmediato, sin condicionamiento alguno. \\

\midrule
\multicolumn{2}{p{\dimexpr\textwidth-2\tabcolsep}}{
\textbf{Resumen:} Los actos jurídicos admiten múltiples clasificaciones --según el número de partes, el momento de sus efectos, la existencia de contraprestación, su dependencia de otro acto, su forma, su contenido patrimonial y su incidencia en el patrimonio-- que determinan las reglas aplicables en cada caso concreto del Código Civil y Comercial. Además de estas clasificaciones, el acto jurídico puede quedar sujeto a modalidades accidentales: la condición (hecho futuro e incierto), el plazo (hecho futuro cierto) y el cargo (obligación accesoria y coercitiva que no suspende el acto). Cada modalidad tiene reglas propias de validez, variantes específicas y condiciones o cargos expresamente prohibidos por la ley (arts. 343 a 357, CCyC). Comprender estas categorías es la base del razonamiento jurídico: permiten leer el Código, redactar documentos, asesorar a clientes y resolver casos con precisión técnica en cualquier rama del derecho privado.
} \\
\bottomrule
\end{longtable}
\chapter[Principios generales]{Principios generales de los efectos del acto jurídico}

\section{Introducción: los efectos del acto jurídico}

El estudio de los efectos del acto jurídico se organiza como un sistema concéntrico. El núcleo es el principio \textbf{\textit{nemo plus iuris}} (nadie transmite más de lo que tiene), que luego se proyecta hacia las partes, sus representantes, los sucesores y, finalmente, los terceros. Puede pensarse como las ondas que produce una piedra en el agua: el acto jurídico es la piedra, y cada onda alcanza a un grupo diferente de personas con reglas propias.

\begin{definition}{Acto jurídico}
Acto \textbf{voluntario lícito} que tiene por fin inmediato producir el \textbf{nacimiento, la modificación, la transferencia o la extinción} de relaciones e instituciones jurídicas.
\end{definition}

Por definición, entonces, el acto jurídico busca producir efectos jurídicos. La cuestión central consiste en determinar \textbf{a quiénes abarcan esos efectos}. Los interrogantes que estructuran el tema son los siguientes:

\begin{itemize}
    \item ¿Los efectos abarcan a las partes?
    \item ¿Cómo están abarcados los sucesores?
    \item ¿Cómo están abarcados los acreedores?
    \item ¿Quiénes son representantes?
    \item ¿Qué reglas rigen estos efectos de los actos jurídicos?
\end{itemize}

\subsection{Importancia práctica}

Estos conceptos constituyen el fundamento de cada operación jurídica que realiza un abogado o una abogada. Se proyectan, por ejemplo, en las siguientes situaciones:

\begin{itemize}
    \item \textbf{Redacción de un contrato de compraventa:} es necesario saber si el cliente puede transmitir válidamente el derecho que pretende vender (\textit{nemo plus iuris}).
    \item \textbf{Otorgamiento o recepción de un poder notarial:} debe determinarse si el representante actuó dentro de sus facultades.
    \item \textbf{Asesoramiento a herederos:} se distingue qué obligaciones del causante los alcanzan y cuáles se extinguen.
    \item \textbf{Cobro de un crédito impago:} se elige entre la acción directa —para cobrar directamente— y la acción subrogatoria —para ingresar bienes al patrimonio del deudor—.
\end{itemize}

No se trata de temas abstractos: son herramientas cotidianas de la abogacía en materia contractual, sucesoria y ejecutiva.

\section{Nemo plus iuris (Art.~399 CCCN)}

\subsection{La transmisibilidad de los derechos como regla general}

Los derechos son \textbf{transmisibles}, salvo estipulación válida de las partes o prohibición legal. Sobre esta base se establece la regla central que organiza todo el tema: \textbf{nadie puede transmitir un derecho mejor que el que tiene}.

Esta regla ya se encontraba en el Código anterior, en el art.~3270: nadie puede transmitir a otro un derecho mejor o más extenso que el que tiene, salvo las excepciones legalmente dispuestas. Proviene del Derecho Romano y en algunos libros se la designa como \textit{nemo plus iuris}, por sus primeras palabras. Ya estaba presente en el Digesto.

\begin{formula}{Art.~399 CCCN --- Regla general (\textit{nemo plus iuris})}
Nadie puede transmitir a otro un derecho mejor o más extenso que el que tiene, sin perjuicio de las excepciones legalmente dispuestas.
\end{formula}

% TODO: contenido incompleto o ambiguo en las notas originales
% (las notas indican "en 398 estamos en el Título Quinto del Libro Primero" al ubicar el art. 399; se conserva la ubicación por Título y Libro, sin precisar el número de artículo)
\begin{note}
El \textit{nemo plus iuris} se encuentra en el art.~399, hacia el final del Libro Primero del Código Civil y Comercial de la Nación, dentro del Título Quinto, denominado «Transmisión de los derechos». Esto se relaciona con el estudio de los efectos del acto: quien transmite derechos no puede transmitirlos de un modo mejor que el que tiene.
\end{note}

\begin{example}{el inquilino que no puede vender}
Quien es inquilino de un inmueble no puede venderlo, porque no puede transmitir un derecho mejor que el que tiene. Su derecho sobre el inmueble es simplemente el uso y goce en los términos de un contrato de locación, y no el derecho en los términos del dominio del dueño de la cosa. Esto constituye claramente un \textbf{límite a la transmisibilidad de los derechos}.
\end{example}

\subsection{Caso práctico: Juan, Verónica y Daniela}

El principio se ilustra con un caso que involucra a tres personas.

\begin{example}{transmisión de propiedad viciada por intimidación}
Juan ejerce violencia (intimidación) sobre Verónica y logra que ella le transfiera su propiedad bajo amenazas. Juan aparece como propietario, pero su título está viciado por intimidación. Luego, Juan le vende el inmueble a Daniela.

\smallskip
\textbf{Pregunta:} ¿Daniela recibió la propiedad con el vicio o sin él?

\smallskip
Conforme al \textit{nemo plus iuris}, Juan transmite el derecho con el mismo alcance que lo tiene, es decir, con el vicio. Por lo tanto, Daniela, como sucesora, recibiría un derecho afectado por ese vicio de intimidación.
\end{example}

\section[Excepciones al nemo plus iuris]{Excepciones al \textit{nemo plus iuris}}

El mismo caso plantea una \textbf{excepción crucial} al principio general: es necesario determinar si Daniela es o no de buena fe. Pueden presentarse, así, excepciones al \textit{nemo plus iuris}.

\subsection{Adquirente de buena fe y a título oneroso (Art.~392 CCCN)}

Una de esas excepciones se encuentra en el art.~392. Si Daniela es de \textbf{buena fe} y a \textbf{título oneroso}, adquiere la propiedad, y quien sale perdiendo es Verónica —quien transmitió bajo intimidación—, porque Daniela se convierte en adquirente de buena fe a título oneroso. \textbf{Quien es de buena fe queda protegido.}

\begin{formula}{Art.~392 CCCN --- Efectos respecto de terceros en casos de nulidad}
El subadquirente de buena fe y a título oneroso queda a salvo de cualquier acción de nulidad que afecte el antecedente de su título, así como de los actos de los que fueran parte terceros de buena fe, con anterioridad a la inscripción de la sentencia que no fueren supuestos de restricciones a la capacidad, simulación o fraude.
\end{formula}

\subsection{Otros supuestos de protección de terceros}

Los supuestos en que los terceros quedan protegidos frente al principio general son los siguientes:

% TODO: contenido incompleto o ambiguo en las notas originales
% (la enumeración original menciona dos veces los "cambios en los poderes que no son oponibles" y los formula de manera imprecisa;
%  además, el texto del art. 392 transcripto excluye supuestos de simulación o fraude, mientras que la enumeración los incluye entre los supuestos de protección de terceros)
\begin{itemize}
    \item los \textbf{cambios en los poderes que no son oponibles}, según dependa de la posibilidad o no de hacerlos conocer a los terceros (cuestión vinculada con la protección de los terceros y la confirmación);
    \item el \textbf{poseedor de buena fe de cosas muebles}, que queda protegido;
    \item los supuestos de \textbf{simulación} (Art.~337);
    \item los supuestos de \textbf{fraude} (Art.~340).
\end{itemize}

\section{Cuadro Cornell: principios generales}

\begin{table}[H]
\centering
\small
\renewcommand{\arraystretch}{1.3}
\begin{tabularx}{\textwidth}{
    >{\raggedright\arraybackslash}p{0.28\textwidth}
    >{\raggedright\arraybackslash}X
}
\toprule
\textbf{Preguntas / palabras clave} &
\textbf{Notas / respuestas} \\
\midrule

\textbf{¿Qué busca producir el acto jurídico y cuál es la cuestión central sobre sus efectos?} &
El acto jurídico es un acto voluntario lícito cuyo fin inmediato es producir el nacimiento, la modificación, la transferencia o la extinción de relaciones e instituciones jurídicas; por definición, busca producir efectos jurídicos. La cuestión central es a quiénes abarcan esos efectos: las partes, los sucesores, los acreedores y los representantes, y qué reglas los rigen. \\

\midrule

\textbf{¿Qué es el \textit{nemo plus iuris}?} &
Regla del art.~399 CCCN: nadie puede transmitir a otro un derecho mejor o más extenso que el que tiene, sin perjuicio de las excepciones legalmente dispuestas. Proviene del Derecho Romano (estaba ya en el Digesto) y se encontraba en el Código anterior (art.~3270). Se ubica en el Título Quinto del Libro Primero, «Transmisión de los derechos». Es un principio cardinal del derecho privado. Los derechos son transmisibles, salvo estipulación válida de las partes o prohibición legal. \\

\midrule

\textbf{¿Cómo opera el \textit{nemo plus iuris} en el caso de Juan, Verónica y Daniela?} &
Juan obtuvo la propiedad de Verónica mediante intimidación y luego se la vendió a Daniela. Juan transmite el derecho con el mismo alcance que lo tiene, es decir, con el vicio; por lo tanto, Daniela, como sucesora, recibe un derecho afectado por ese vicio, salvo que sea adquirente de buena fe y a título oneroso (art.~392 CCCN). \\

\midrule

\textbf{¿Cuándo queda protegido el adquirente frente al \textit{nemo plus iuris}?} &
Según el art.~392 CCCN, el subadquirente de buena fe y a título oneroso queda a salvo de cualquier acción de nulidad que afecte el antecedente de su título; en el caso, Daniela adquiere y quien pierde es Verónica. También quedan protegidos el poseedor de buena fe de cosas muebles, los terceros en los supuestos de simulación (art.~337) y de fraude (art.~340), y los cambios en los poderes que no son oponibles a terceros. \\

\midrule

\textbf{¿Por qué un inquilino no puede vender el inmueble que alquila?} &
Porque no puede transmitir un derecho mejor que el que tiene: su derecho es solo al uso y goce en los términos de un contrato de locación, y no en los términos del dominio del dueño de la cosa. Es un límite a la transmisibilidad de los derechos. \\

\midrule

\multicolumn{2}{p{\dimexpr\textwidth-2\tabcolsep\relax}}{
\textbf{Resumen:}
El acto jurídico busca producir efectos, y el primer principio para determinar su alcance es el \textit{nemo plus iuris} (art.~399 CCCN): nadie transmite un derecho mejor o más extenso que el que tiene, de modo que el adquirente recibe el derecho con sus vicios. La excepción estudiada es la del adquirente de buena fe y a título oneroso (art.~392 CCCN), a la que se suman otros supuestos de protección de terceros.
} \\

\bottomrule
\end{tabularx}
\end{table}

% ============================================================
% CAPÍTULO 2
% ============================================================
\chapter[Partes y efecto relativo]{Las partes y el efecto relativo de los actos jurídicos}

\section{Efecto relativo de los actos jurídicos (Art.~1021 CCCN)}

El segundo gran marco teórico es el efecto relativo. Los actos jurídicos producen efectos entre las partes que los celebran; es decir, tienen un \textbf{efecto relativo}: sus efectos afectan a las partes que celebran el acto.

\begin{formula}{Art.~1021 CCCN --- Regla general (efecto relativo)}
El contrato solo tiene efectos entre las partes contratantes; no lo tiene con respecto a terceros, excepto en los casos previstos por la ley.
\end{formula}

Nadie puede verse afectado ni beneficiado por un acto del que no fue parte. Los actos jurídicos tienen valor entre las partes; como regla general, así lo establece el art.~1021.

\section{Definición de parte (Art.~1023 CCCN)}

La noción de «parte» resulta decisiva para determinar quién queda alcanzado por los efectos del acto. Esto conduce a considerar cuál es la definición de «parte», que se encuentra en el art.~1023.

\begin{definition}{Parte}
\textbf{Parte} es quien concurre a un acto ejerciendo prerrogativas jurídicas propias.

\smallskip
\textbf{Art.~1023 CCCN --- Parte del contrato.} Se considera parte del contrato a quien:
\begin{enumerate}[label=\alph*)]
    \item lo otorga a nombre propio, aunque lo haga en interés ajeno;
    \item es representado por un otorgante que actúa en su nombre o en su interés;
    \item manifiesta la voluntad contractual, aunque sea transmitida por un corredor o por un agente sin representación.
\end{enumerate}
\end{definition}

\subsection{Distinción entre parte y representante}

Quien actúa en nombre propio, o por medio de un representante que actúa en su nombre, es la parte, aunque en el acto conste que concurre el señor Fulanito de Tal en representación de Menganito de Tal. En ese caso, la parte es el \textbf{representado} (Menganito de Tal). El \textbf{representante} no ejerce una prerrogativa jurídica propia, sino la de otro.

\section{Cuadro Cornell: partes y efecto relativo}

\begin{table}[H]
\centering
\small
\renewcommand{\arraystretch}{1.3}
\begin{tabularx}{\textwidth}{
    >{\raggedright\arraybackslash}p{0.28\textwidth}
    >{\raggedright\arraybackslash}X
}
\toprule
\textbf{Preguntas / palabras clave} &
\textbf{Notas / respuestas} \\
\midrule

\textbf{¿Cuál es el efecto relativo de los actos jurídicos?} &
Los contratos solo producen efectos entre las partes (art.~1021 CCCN); no los producen respecto de terceros, excepto en los casos previstos por la ley. Un tercero ni puede beneficiarse ni perjudicarse por un acto en el que no intervino, salvo excepciones legales (estipulación a favor de tercero, acciones directa y subrogatoria). \\

\midrule

\textbf{¿Quién es «parte» en un acto jurídico?} &
Quien concurre al acto ejerciendo prerrogativas jurídicas propias (art.~1023 CCCN). Se considera parte a quien: a) otorga el contrato a nombre propio, aunque lo haga en interés ajeno; b) es representado por un otorgante que actúa en su nombre o en su interés; c) manifiesta la voluntad contractual, aunque sea transmitida por un corredor o por un agente sin representación. \\

\midrule

\textbf{¿Es parte el representante?} &
No. El representante no ejerce una prerrogativa jurídica propia, sino la de otro: actúa en nombre del representado, quien sí es la parte, aunque en el acto conste que concurre otra persona en su representación. El mandante de un poder, por ejemplo, es la parte aunque no esté presente. \\

\midrule

\multicolumn{2}{p{\dimexpr\textwidth-2\tabcolsep\relax}}{
\textbf{Resumen:}
Los actos jurídicos tienen efecto relativo: producen efectos entre las partes y, como regla, no alcanzan a terceros (art.~1021 CCCN). Es parte quien concurre ejerciendo prerrogativas jurídicas propias (art.~1023 CCCN); por eso, cuando hay representación, la parte es el representado y no el representante.
} \\

\bottomrule
\end{tabularx}
\end{table}

% ============================================================
% CAPÍTULO 3
% ============================================================
\chapter{Representación}

El Código regula, en el Título Cuarto del Libro Primero, todo lo relativo a la representación. Esta materia es muy importante para la vida de los abogados, ya que generalmente quienes ejercen la profesión cumplen funciones de representación, dadas por poderes o por distintos mecanismos.

\section{Tipos de representación (Art.~358 CCCN)}

El art.~358 del CCCN clasifica la representación en tres tipos:

\begin{center}
\renewcommand{\arraystretch}{1.3}
\begin{tabularx}{\textwidth}{
    >{\raggedright\arraybackslash\bfseries}p{2.4cm}
    >{\raggedright\arraybackslash}p{4.4cm}
    >{\raggedright\arraybackslash}X
}
\toprule
\textbf{Tipo} & \textbf{Origen} & \textbf{Ejemplo} \\
\midrule
Voluntaria & Acto jurídico (poder notarial) & Juana le otorga poder a Daniela para celebrar un contrato de mutuo en su nombre. \\
\addlinespace
Legal & Regla de derecho (la ley) & Padres que representan a sus hijos menores (Art.~101 CCCN); tutores; curadores. \\
\addlinespace
Orgánica & Estatuto de persona jurídica & El presidente de una sociedad que firma contratos porque el estatuto le atribuye esa función. \\
\bottomrule
\end{tabularx}
\end{center}

\section{Representación voluntaria (Arts.~362--381 CCCN)}

Los arts.~362 a 381 regulan la forma, la capacidad, los vicios, la actuación y la extinción de la representación voluntaria. En cuanto a la capacidad:

\begin{itemize}
    \item \textbf{Representante:} basta con el \textbf{discernimiento}. Puede serlo desde que tiene discernimiento para los actos lícitos, a los trece años.
    \item \textbf{Representado:} se requiere capacidad para otorgar el poder; debe ser una persona capaz.
\end{itemize}

\section{Facultades expresas (Art.~375 CCCN)}

El art.~375 es de máxima relevancia práctica: regula los casos en que se requieren \textbf{facultades expresas}. Quien, como abogado, vaya a realizar alguno de estos actos de mayor relevancia necesita lo que se denomina un \textbf{poder especial} de su representado.

\begin{formula}{Art.~375 CCCN --- Facultades expresas}
Las facultades para ciertos actos deben ser expresas. Entre ellos:
\begin{itemize}
    \item reconocer hijos;
    \item aceptar herencias;
    \item contraer matrimonio;
    \item hacer donaciones;
    \item avalar;
    \item constituir hipotecas;
    \item otorgar testamentos;
    \item transar;
    \item comprometer en árbitros;
    \item y otros actos de especial trascendencia.
\end{itemize}
\end{formula}

\section{Cuadro Cornell: representación}

\begin{table}[H]
\centering
\small
\renewcommand{\arraystretch}{1.3}
\begin{tabularx}{\textwidth}{
    >{\raggedright\arraybackslash}p{0.28\textwidth}
    >{\raggedright\arraybackslash}X
}
\toprule
\textbf{Preguntas / palabras clave} &
\textbf{Notas / respuestas} \\
\midrule

\textbf{¿Qué tipos de representación existen y cuál es el origen de cada una?} &
Arts.~358 y ss. CCCN. 1) \textbf{Voluntaria}: surge de un acto jurídico (poder notarial). 2) \textbf{Legal}: surge de la ley (padres respecto de hijos menores, tutores, curadores). 3) \textbf{Orgánica}: surge del estatuto de una persona jurídica (presidente de una sociedad). \\

\midrule

\textbf{¿Qué capacidad se requiere en la representación voluntaria?} &
Arts.~362 a 381 CCCN. Para el representante basta con el discernimiento, que se tiene desde los trece años para los actos lícitos. Para otorgar el poder se requiere capacidad en el representado, que debe ser una persona capaz. \\

\midrule

\textbf{¿Cuándo se requieren facultades expresas y qué se necesita en esos casos?} &
Art.~375 CCCN: las facultades para ciertos actos deben ser expresas, y para ellos se necesita un poder especial del representado. Entre esos actos: reconocer hijos, aceptar herencias, contraer matrimonio, hacer donaciones, avalar, constituir hipotecas, otorgar testamentos, transar, comprometer en árbitros y otros actos de especial trascendencia. \\

\midrule

\multicolumn{2}{p{\dimexpr\textwidth-2\tabcolsep\relax}}{
\textbf{Resumen:}
La representación puede ser voluntaria (nace de un acto jurídico), legal (nace de la ley) u orgánica (nace del estatuto de una persona jurídica). En la representación voluntaria basta el discernimiento para el representante, mientras que el representado debe ser capaz para otorgar el poder. Para los actos de mayor relevancia enumerados en el art.~375 CCCN se requieren facultades expresas, es decir, un poder especial.
} \\

\bottomrule
\end{tabularx}
\end{table}

% ============================================================
% CAPÍTULO 4
% ============================================================
\chapter{Sucesores}

A un acto concurren las partes con sus representantes. Después se encuentran los \textbf{sucesores}, que son aquellos que reciben un derecho. El sucesor puede ser \textbf{universal} o \textbf{singular}.

\section{Sucesor universal}

\begin{definition}{Sucesor universal}
Es quien recibe \textbf{todo o una parte indivisa del patrimonio} de otro.
\end{definition}

Quien sucede a su padre porque este falleció se convierte en su sucesor universal. Y aunque sean dos herederos, ambos son universales porque tienen \textbf{vocación al todo}.

\begin{example}{sucesión entre dos hermanos}
Supongamos que el padre tiene un departamento, una cochera, un auto y dinero en efectivo. Los dos hermanos tienen vocación al todo: sin testamento, cada uno al 50\,\%. No es que cada uno tenga el 50\,\% del dinero más el 50\,\% del auto: cada uno tiene vocación al todo.

\smallskip
Después, en la partición, si uno se queda con el departamento y el otro con todos los demás bienes, en ese momento cada uno se convierte en sucesor singular de lo que recibió.
\end{example}

\begin{formula}{Art.~1024 CCCN --- Efectos respecto de los sucesores universales}
Los efectos del contrato se extienden activa y pasivamente a los sucesores universales, a no ser que las obligaciones que de él nacen sean inherentes a la persona, o que la transmisión sea incompatible con la naturaleza de la obligación, o esté prohibida por una cláusula del contrato o la ley.
\end{formula}

Los límites de la transmisión a los sucesores universales se ilustran con el siguiente ejemplo.

\begin{example}{obligación inherente a la persona (pintor famoso)}
Supongamos que se contrata a un pintor famoso —con nombre propio, por sus condiciones artísticas— para que realice un retrato. Si el pintor fallece, su sucesor no puede venir a cumplir el contrato, porque se trata de una obligación inherente a la persona del artista.
\end{example}

Tampoco se transmiten las obligaciones de derecho de familia, ni se transmite el usufructo o el contrato de uso y habitación.

\section{Sucesor singular}

\begin{definition}{Sucesor singular}
Es quien recibe un \textbf{derecho determinado} (por ejemplo, un inmueble o un auto); es decir, quien sigue un bien determinado.
\end{definition}

\begin{important}
Los sucesores singulares, por regla, \textbf{no se ven afectados} por los actos jurídicos de la parte de quien reciben el bien, y se asimilan a los terceros completamente extraños. La excepción está dada por:
\begin{itemize}
    \item los actos jurídicos que son \textbf{antecedentes del derecho transmitido};
    \item las \textbf{obligaciones que pesan sobre la cosa};
    \item los \textbf{derechos accesorios del objeto}.
\end{itemize}
\end{important}

\begin{example}{inmueble con contrato de locación}
Quien hereda un inmueble sujeto a un contrato de locación recibe el inmueble con el contrato de locación, aunque sea sucesor singular, porque sigue un bien determinado.
\end{example}

\begin{example}{auto con prenda}
Quien hereda un auto sujeto a una prenda hereda el auto con la prenda: la prenda lo afecta. Esa es la consecuencia de ser sucesor singular.
\end{example}

\section{Cuadro Cornell: sucesores}

\begin{table}[H]
\centering
\small
\renewcommand{\arraystretch}{1.3}
\begin{tabularx}{\textwidth}{
    >{\raggedright\arraybackslash}p{0.28\textwidth}
    >{\raggedright\arraybackslash}X
}
\toprule
\textbf{Preguntas / palabras clave} &
\textbf{Notas / respuestas} \\
\midrule

\textbf{¿Qué son los sucesores universales?} &
Quienes reciben la totalidad o una parte indivisa del patrimonio de otra persona (art.~1024 CCCN): herederos \textit{ab intestato} o testamentarios. Aunque sean varios herederos, todos son universales porque tienen vocación al todo. \\

\midrule

\textbf{¿Se extienden los efectos del contrato a los sucesores universales y cuáles son los límites?} &
Sí: los efectos del contrato se extienden activa y pasivamente a los sucesores universales (art.~1024 CCCN), salvo que las obligaciones sean inherentes a la persona, que la transmisión sea incompatible con la naturaleza de la obligación, o que esté prohibida por una cláusula del contrato o la ley. \\

\midrule

\textbf{¿Qué ejemplos de obligaciones o derechos no se transmiten a los sucesores universales?} &
La obligación de un pintor famoso contratado por sus condiciones artísticas para realizar un retrato (obligación inherente a la persona: si fallece, su sucesor no puede cumplir el contrato); las obligaciones de derecho de familia; el usufructo o el contrato de uso y habitación. \\

\midrule

\textbf{¿Qué son los sucesores singulares y qué los afecta?} &
Quienes reciben un derecho determinado (un inmueble, un auto). Se asimilan a terceros respecto de los actos del causante, pero los afectan: los antecedentes del derecho transmitido, las obligaciones y cargas que pesan sobre la cosa y los derechos accesorios del objeto (por ejemplo, la prenda sobre un auto heredado, o la locación de un inmueble heredado). \\

\midrule

\textbf{¿Cuándo un sucesor universal se convierte en singular?} &
En la partición: si, por ejemplo, de los bienes del padre un hermano se queda con el departamento y el otro con los demás bienes, cada uno se convierte en sucesor singular de lo que recibió. \\

\midrule

\multicolumn{2}{p{\dimexpr\textwidth-2\tabcolsep\relax}}{
\textbf{Resumen:}
Los sucesores universales reciben todo o una parte indivisa del patrimonio y ocupan el lugar del causante: los efectos del contrato se les extienden activa y pasivamente, salvo obligaciones inherentes a la persona, incompatibles con la naturaleza de la obligación o vedadas por el contrato o la ley (art.~1024 CCCN). Los sucesores singulares siguen un bien determinado: por regla no se ven afectados por los actos de la parte de quien lo reciben, salvo por los antecedentes del derecho transmitido, las obligaciones que pesan sobre la cosa y los derechos accesorios del objeto.
} \\

\bottomrule
\end{tabularx}
\end{table}

% ============================================================
% CAPÍTULO 5
% ============================================================
\chapter[Terceros y acciones]{Los terceros y las acciones a su disposición}

\section{Situación de los terceros}

Los terceros no pueden ni beneficiarse ni perjudicarse por los contratos, salvo que exista disposición legal o que así se haya convenido entre las partes. Por regla, los terceros son \textbf{ajenos} y no pueden verse afectados por las disposiciones de un acto en el que no fueron parte.

% TODO: contenido incompleto o ambiguo en las notas originales
% (la expresión "tanto en tales asistencias" del ejemplo original es imprecisa; se conserva tal como figura)
\begin{example}{estipulación a favor de un tercero}
Supongamos que se establece una estipulación a favor de un tercero: «yo le voy a dar plata para que vos le des a Fulanito de Tal tanto en tales asistencias». Esto es una estipulación a favor de un tercero y es perfectamente válido.
\end{example}

\section{Las acciones de los acreedores frente a terceros}

Cuando un acreedor debe actuar frente a quien no es parte en el vínculo que lo une con su deudor, el ordenamiento le concede, excepcionalmente, la \textbf{acción directa} (arts.~736 y 737 CCCN) contra el deudor de su deudor y la \textbf{acción subrogatoria} (arts.~739 y 740 CCCN) para ejercer los derechos de un deudor inactivo. Ambas se desarrollan y se comparan en el capítulo~\ref{ch:garantia}, junto con la acción de simulación y la acción de inoponibilidad, estudiadas en la Parte~VI.

\section{Cuadro Cornell: terceros y acciones}

\begin{table}[H]
\centering
\small
\renewcommand{\arraystretch}{1.3}
\begin{tabularx}{\textwidth}{
    >{\raggedright\arraybackslash}p{0.28\textwidth}
    >{\raggedright\arraybackslash}X
}
\toprule
\textbf{Preguntas / palabras clave} &
\textbf{Notas / respuestas} \\
\midrule

\textbf{¿Puede un contrato obligar o beneficiar a un tercero?} &
Por regla, no: los terceros son ajenos y no pueden ni beneficiarse ni perjudicarse por los contratos, salvo disposición legal o que se haya convenido entre las partes. Por ejemplo, es perfectamente válida una estipulación a favor de un tercero. \\

\midrule

\textbf{¿Qué es la acción directa y qué requisitos tiene?} &
Permite al acreedor cobrar lo que un tercero debe a su deudor, actuando directamente contra el deudor de su deudor, hasta el importe de su propio crédito (arts.~736--737 CCCN). La ejerce por derecho propio y en su exclusivo beneficio; es excepcional y de interpretación restrictiva. Requisitos: crédito exigible del acreedor contra su deudor; deuda correlativa exigible del tercero a favor del deudor; homogeneidad de ambos créditos; ningún crédito embargado ni indisponible; notificación de la demanda al deudor. \\

\midrule

\textbf{¿Qué es la acción subrogatoria y qué derechos quedan excluidos?} &
El acreedor de un crédito cierto, exigible o no, ocupa el lugar de su deudor remiso y ejerce judicialmente los derechos patrimoniales que este no ejerce, cuando esa omisión afecta el cobro de su acreencia (arts.~739--740 CCCN). El acreedor no goza de preferencia sobre los bienes obtenidos. Quedan excluidos los derechos inherentes a la persona del deudor, los no patrimoniales, los patrimoniales no susceptibles de ejecución forzada y los excluidos por disposición legal. \\

\midrule

\textbf{¿En qué se diferencian la acción directa y la subrogatoria?} &
Directa: exige créditos homogéneos; el acreedor actúa por derecho propio, en su nombre; beneficia solo al acreedor que la ejerce; se limita al importe de su crédito. Subrogatoria: exige la inacción del deudor que perjudica al acreedor; el acreedor actúa en nombre del deudor; lo recuperado ingresa al patrimonio del deudor y beneficia a todos los acreedores. \\

\midrule

\multicolumn{2}{p{\dimexpr\textwidth-2\tabcolsep\relax}}{
\textbf{Resumen:}
Los terceros, por regla, son ajenos al acto y no pueden beneficiarse ni perjudicarse por él, salvo disposición legal o convenio. Excepcionalmente, un acreedor puede actuar mediante la acción directa, que cobra contra el deudor del deudor y solo beneficia a quien la ejerce (arts.~736--737 CCCN), o mediante la acción subrogatoria, que ejerce en nombre del deudor inactivo y beneficia a todos los acreedores (arts.~739--740 CCCN).
} \\

\bottomrule
\end{tabularx}
\end{table}

% ============================================================
% CAPÍTULO 6
% ============================================================

\section{Síntesis general de los efectos del acto jurídico}

Los efectos del acto jurídico se rigen por dos principios estructurantes: el \textbf{\textit{nemo plus iuris}} (nadie puede transmitir un derecho mejor que el que tiene, Art.~399 CCCN) y el \textbf{efecto relativo} (los actos solo producen efectos entre las partes, Art.~1021 CCCN).

A partir de estos ejes, el sistema regula quiénes resultan alcanzados:

\begin{itemize}
    \item las \textbf{partes};
    \item sus \textbf{representantes} (voluntarios, legales u orgánicos);
    \item los \textbf{sucesores universales}, que ocupan el mismo lugar que el causante en todos los derechos y obligaciones transmisibles;
    \item los \textbf{sucesores singulares}, que solo reciben los antecedentes del derecho transmitido, las cargas sobre la cosa y sus accesorios.
\end{itemize}

Los \textbf{terceros}, en principio ajenos, pueden excepcionalmente actuar mediante la \textbf{acción directa} (Arts.~736--737 CCCN) o la \textbf{acción subrogatoria} (Arts.~739--740 CCCN).

Estas instituciones son herramientas cotidianas de la práctica jurídica: delimitan quiénes quedan obligados por los contratos, cómo se transmiten los bienes en una sucesión y de qué modo un acreedor puede cobrar su crédito cuando el deudor es negligente.
\partintro{Un acto jurídico puede existir y, sin embargo, presentar defectos que impiden que produzca todos sus efectos normales. Esta parte estudia los vicios de la voluntad (error, dolo y violencia), el vicio de lesión, que afecta el equilibrio de las prestaciones, y los vicios de simulación y de fraude, que afectan la sinceridad del acto y la garantía de los acreedores.}
\part{Vicios de los actos jurídicos: voluntad, equilibrio y sinceridad}
\chapter{Fundamentos y método del Código Civil y Comercial}
% ============================================================

\section{Introducción: ¿vicios de la voluntad o del consentimiento?}

La doctrina ha discutido si estos institutos debían denominarse \textbf{vicios de la voluntad} o \textbf{vicios del consentimiento}. La mayoría de los juristas argentinos, al igual que el propio Código, se inclina por la primera denominación, señalando que es conceptualmente más amplia que la segunda.

La razón es de fondo: hablar de vicios del consentimiento los reduciría a defectos propios de los contratos ---un tipo específico de acto jurídico---, mientras que los vicios de la voluntad afectan a todo acto voluntario.

\begin{articulo}{Art. 260 CCyCN --- Acto voluntario}
Son requisitos del acto voluntario: discernimiento, intención y libertad (condiciones internas) y manifestación de la voluntad (condición externa). Los vicios afectan a alguno de esos elementos.
\end{articulo}

\begin{articulo}{Art. 261 CCyCN --- Acto involuntario}
Determina cuándo falta el discernimiento como condición interna del acto voluntario.
\end{articulo}

Los actos voluntarios son aquellos que presentan tres condiciones internas: \textbf{discernimiento, intención y libertad}, y una condición externa que es la \textbf{manifestación de la voluntad}. Los vicios afectan alguno de estos elementos. El estudio de los vicios constituye una introducción al régimen de las \textbf{ineficacias}: la posibilidad de que el acto no produzca efectos jurídicos por la presencia de un defecto en su celebración.

Las preguntas centrales que guían este estudio son:
\begin{itemize}
    \item Bajo qué circunstancias un acto voluntario puede ser anulado por la presencia de un defecto en su celebración.
    \item Cómo se afecta la intención y cómo se afecta la libertad como elementos internos.
    \item Qué requisitos deben reunir los vicios para que pueda declararse la nulidad.
    \item Qué consecuencias jurídicas se siguen de la presencia de un vicio.
\end{itemize}

\section{Método del Código Civil y Comercial}

Desde el punto de vista metodológico, los capítulos 2 a 4 del Título Cuarto del CCyCN organizan los vicios comenzando en el artículo 265 con el error. El Código distingue dos grandes grupos bajo el término común de «vicios»:

\begin{table}[H]
\centering
\begin{tabularx}{\textwidth}{>{\bfseries\raggedright\arraybackslash}X >{\raggedright\arraybackslash}X}
\toprule
\textbf{Vicios de la voluntad} & \textbf{Vicios propios de los actos jurídicos} \\
\midrule
Error (Arts. 265--270) & Lesión \\
Dolo (Arts. 271--275) & Simulación \\
Violencia (Arts. 276--278) & Fraude \\
\bottomrule
\end{tabularx}
\caption{Clasificación de los vicios en el CCyCN}
\end{table}

Los primeros ---error, dolo y violencia--- afectan los elementos internos del acto voluntario. Los segundos ---lesión, simulación y fraude--- tienen un tratamiento posterior en el Código y son propios de los actos jurídicos.

% ============================================================
\chapter{El error como vicio de la voluntad (Arts.~265--270)}
\label{ch:error}
% ============================================================

\section{Concepto y clasificación del error de hecho}

El error afecta la \textbf{intención} como elemento interno del acto voluntario. La doctrina trata conjuntamente dos supuestos: la \textbf{ignorancia} ---la falta de conocimiento sobre algún elemento del acto--- y el \textbf{error propiamente dicho} ---un falso conocimiento---. Jurídicamente ambos se tratan como un único supuesto.

La primera gran distinción es entre el error de derecho y el error de hecho:

\begin{table}[H]
\centering
\begin{tabularx}{\textwidth}{>{\raggedright\arraybackslash}X >{\raggedright\arraybackslash}X}
\toprule
\textbf{Error de derecho} & \textbf{Error de hecho} \\
\midrule
Falso conocimiento de la normativa legal aplicable a una relación o situación jurídica. & Recae sobre los elementos fácticos del acto o circunstancias del mismo. \\[6pt]
Regulado en el Art. 8 CCyCN: la ignorancia de la ley no sirve de excusa para su incumplimiento. & Regulado en los Arts. 265 a 270 CCyCN. \\[6pt]
No puede invocarse como vicio. La ley se presume conocida por todos. & Es el que puede dar lugar a la nulidad del acto, con los requisitos del caso. \\
\bottomrule
\end{tabularx}
\caption{Error de derecho versus error de hecho}
\end{table}

\begin{articulo}{Art. 8 CCyCN --- Principio de inexcusabilidad}
La ignorancia de las leyes no sirve de excusa para su cumplimiento, si la excepción no está autorizada por el ordenamiento jurídico.
\end{articulo}

Dentro del error de hecho, la clasificación fundamental es:

\begin{itemize}
    \item \textbf{Error esencial}: recae sobre un elemento esencial del acto. Puede causar la nulidad si además es reconocible.
    \item \textbf{Error incidental}: recae sobre elementos secundarios que no son determinantes de la voluntad. No causa nulidad.
\end{itemize}

La fórmula es: cuando el error es \textbf{esencial Y reconocible}, causa la nulidad del acto (si es bilateral o unilateral de aceptación).

\section{El error reconocible --- Art. 266}

El CCyCN introduce un requisito que no estaba en el Código de Vélez: la \textbf{reconocibilidad}. En el Código anterior se exigía que el error fuera esencial y \textbf{excusable}. El nuevo Código reemplaza la excusabilidad por la reconocibilidad.

\begin{articulo}{Art. 266 CCyCN --- Error reconocible}
El error es reconocible cuando el destinatario de la declaración lo pudo conocer según la naturaleza del acto, las circunstancias de persona, tiempo y lugar.
\end{articulo}

\begin{table}[H]
\centering
\begin{tabularx}{\textwidth}{>{\raggedright\arraybackslash}X >{\raggedright\arraybackslash}X}
\toprule
\textbf{Código de Vélez} & \textbf{CCyCN (vigente)} \\
\midrule
Error esencial + excusable. & Error esencial + reconocible. \\[4pt]
Miraba al sujeto que erraba: ¿tenía razón para equivocarse? & Mira al destinatario: ¿pudo darse cuenta del error ajeno? \\[4pt]
Requisito subjetivo centrado en quien declara. & Requisito objetivo centrado en quien recibe la declaración. \\
\bottomrule
\end{tabularx}
\caption{Cambio de requisito: Código de Vélez versus CCyCN}
\end{table}

La excusabilidad del Código de Vélez miraba al sujeto que emitía la declaración: que hubiera tenido alguna razón objetiva para errar. Ese criterio fue abandonado por el nuevo Código, siguiendo el modelo del derecho italiano.

El error es reconocible cuando el destinatario de la declaración no pudo \textit{no conocer} el error, según la naturaleza del acto o las circunstancias de persona, tiempo y lugar. Existe así una expectativa de que la declaración viciada por un error esencial debió haber sido advertida por la parte destinataria.

Si el destinatario no advierte el error siendo que cualquiera lo hubiera advertido ---pero tampoco actúa con malicia (porque si actúa con malicia incurriría en dolo)---, debe admitir que el acto sea anulado. La consecuencia de que el error sea reconocible es la nulidad relativa del acto.

\section{El error esencial --- supuestos del Art. 267}

\begin{articulo}{Art. 267 CCyCN --- Supuestos de error esencial}
El error de hecho es esencial cuando recae sobre: (a) la naturaleza del acto; (b) el objeto sobre un bien o un hecho diverso o de distinta especie, o una calidad, extensión o suma diversa; (c) la cualidad sustancial del bien que haya sido determinante de la voluntad; (d) los motivos personales relevantes que hayan sido incorporados expresa o tácitamente al acto; (e) la persona con la cual se celebró o a la cual se refiere el acto, si ella fue determinante para su celebración.
\end{articulo}

\subsection{Inciso a) --- Error sobre la naturaleza del acto}

Es el error sobre el \textbf{tipo de acto} que, conforme al ordenamiento jurídico, cada parte cree celebrar.

\begin{example}{Error sobre la naturaleza del acto}
A le dice a B: «Me das ese cuadro». A interpreta que B está haciendo una donación; B interpreta que está realizando un comodato (préstamo gratuito). Hay un error sobre la naturaleza del acto: uno pensó que era una donación y el otro que era un préstamo. Consecuencia: nulidad relativa, si el error es reconocible y esencial.
\end{example}

\subsection{Inciso b) --- Error sobre el objeto}

Es el error sobre un \textbf{bien o un hecho diverso o de distinta especie}, o sobre una calidad, extensión o suma diversa.

\begin{example}{Error sobre el objeto --- distintas especies}
A cree estar comprando un caballo pura sangre para competencias; B interpreta que A está comprando un caballo de carga para trabajo en el campo. El objeto del acto ---el caballo--- no es el mismo en la mente de ambas partes. Si el error es reconocible, da lugar a la nulidad.
\end{example}

\begin{example}{Error sobre el objeto --- calidad}
A cree estar comprando un vino de excelente calidad y en realidad es un vino de calidad inferior. Si la otra parte debería haberse dado cuenta del error (reconocibilidad), el acto puede ser anulado.
\end{example}

\subsection{Inciso c) --- Error sobre la cualidad sustancial del bien}

Este supuesto presenta un factor \textbf{objetivo} y uno \textbf{subjetivo}: por un lado, la cualidad debe ser sustancial al bien (factor objetivo); por otro lado, debe haber sido determinante de la voluntad según la apreciación común a las circunstancias del caso (factor subjetivo).

\begin{table}[H]
\centering
\begin{tabularx}{\textwidth}{>{\raggedright\arraybackslash}X >{\raggedright\arraybackslash}X}
\toprule
\textbf{Código de Vélez} & \textbf{CCyCN (vigente)} \\
\midrule
Utilizaba el criterio de la cualidad sustancial sin la palabra «sustancial» en ciertos pasajes, generando debate. & Incorpora la palabra «sustancial» y la conjuga con el requisito de que haya sido «determinante de la voluntad». \\[4pt]
Discusión entre criterio objetivo y subjetivo de la cualidad. & Combina ambos factores: objetivo (cualidad sustancial del bien) + subjetivo (determinante de la voluntad según apreciación común). \\
\bottomrule
\end{tabularx}
\caption{Debate histórico sobre la cualidad sustancial}
\end{table}

\begin{example}{Error sobre la cualidad sustancial}
A está comprando una notebook y necesita que tenga buena memoria RAM porque va a trabajar con edición de videos. Esto es una cualidad sustancial de la computadora (factor objetivo), y además ha sido determinante para su voluntad de contratar (factor subjetivo). Si se le entrega una computadora sin esa capacidad y el error era reconocible, puede dar lugar a la anulabilidad del acto.
\end{example}

\subsection{Inciso d) --- Error sobre los motivos}

Este inciso se vincula con la causa de los actos jurídicos (Art. 281 CCyCN). El error puede recaer sobre los \textbf{motivos personales relevantes} que hayan sido incorporados expresa o tácitamente al acto.

\begin{example}{Error sobre los motivos}
A quiere hacer una donación a una persona pensando que esa persona le salvó la vida. En realidad no es la misma persona. El motivo ---la gratitud hacia quien le salvó la vida--- fue incorporado al acto y es esencial. Si el error es reconocible, puede anularse la donación.
\end{example}

\subsection{Inciso e) --- Error sobre la persona}

El error sobre la persona solo causa nulidad cuando la identidad de la persona fue \textbf{determinante para la celebración del acto}. Si no es determinante, es un error incidental y no produce nulidad.

\begin{example}{Error sobre la persona}
A quiere contratar a un pintor famoso para que le haga un retrato de su madre. Existe un homónimo ---alguien con el mismo nombre--- y A contrata al homónimo creyendo que es el pintor famoso. Hay un error esencial sobre la persona: era determinante quién era el pintor, el homónimo debería haberse dado cuenta de que no era el artista buscado, y el acto puede anularse.
\end{example}

\section{Error de cálculo --- Art. 268}

\begin{articulo}{Art. 268 CCyCN --- Error de cálculo}
El error de cálculo no da lugar a la nulidad del acto, sino solamente a su rectificación, excepto que sea determinante del consentimiento.
\end{articulo}

El error de cálculo no da lugar a la nulidad sino únicamente a una \textbf{rectificación}, dado que se trata de un error aritmético que no fue determinante del consentimiento.

\begin{example}{Error de cálculo}
Las partes pactan vender 40 kg de mercadería a \$100 por kg. Queda claro el valor unitario y la cantidad. Al realizar la multiplicación en el contrato, escriben mal el precio total. Eso es simplemente un error aritmético que se corrige (rectifica) pero no da lugar a la nulidad.
\end{example}

\section{Subsistencia del acto y nulidad relativa --- Art. 269}

\begin{articulo}{Art. 269 CCyCN --- Subsistencia del acto}
La parte que incurre en error no puede solicitar la nulidad del acto si la otra ofrece ejecutarlo con las modalidades y el contenido que aquella entendió celebrar.
\end{articulo}

El artículo 269 le concede a la parte destinataria de la declaración la posibilidad de \textbf{ofrecer ejecutar el acto con las modalidades y el contenido que la parte en error pretendió} que se celebrara.

\begin{example}{Subsistencia del acto}
A le dice a B: «Me das tu auto». B lo interpreta como préstamo; A como donación. A alega error y busca anular el acto. B puede decir: «Bueno, celebremos el acto bajo la modalidad del préstamo (comodato): te presto el auto». Si eso le conviene a A, se salva la nulidad. Esto responde al principio general de conservación de los actos.
\end{example}

\begin{important}
El principio de conservación de los actos es un principio general del derecho civil: se tiende a interpretar que los actos deben ser conservados, preservando el valor de los negocios jurídicos.
\end{important}

\section{Error en la declaración --- Art. 270}

Las disposiciones sobre el error se aplican tanto a la declaración de voluntad en sí como a su \textbf{transmisión}. Pueden presentarse problemas cuando quien transmite la voluntad es un \textbf{mandatario} (con representación) o simplemente un \textbf{mensajero}. En general, se aplican todas las disposiciones sobre error esencial y reconocible.

\section{Consecuencias del error: nulidad relativa y prescripción}

La sanción que se sigue de la presencia del error es la \textbf{nulidad relativa} del acto, regulada en el Art. 388 CCyCN y susceptible de confirmación.

\begin{articulo}{Arts. 2562 y 2563 CCyCN --- Prescripción}
El plazo de prescripción para la acción de nulidad por vicios del consentimiento es de \textbf{dos años}. Comienza a correr desde que el error se conoció o pudo ser conocido.
\end{articulo}

% ============================================================
\chapter{El dolo como vicio de la voluntad (Arts.~271--275)}
% ============================================================

\section{Las acepciones del dolo en el derecho civil}

La palabra «dolo» tiene múltiples significados en el CCyCN. Es fundamental no confundirlos:

\begin{table}[H]
\centering
\begin{tabularx}{\textwidth}{>{\raggedright\arraybackslash}p{0.30\textwidth} >{\raggedright\arraybackslash}X}
\toprule
\textbf{Acepción} & \textbf{Definición y referencia} \\
\midrule
1) Elemento del acto ilícito & Intención de dañar o manifiesta indiferencia por los intereses ajenos (Art. 1724 CCyCN). Componente del factor subjetivo de atribución de responsabilidad. \\[4pt]
2) Inejecución deliberada de obligación & Incumplimiento voluntario de una obligación. Figuraba en el Código de Vélez; el nuevo CCyCN no lo define con la misma fórmula. \\[4pt]
3) Vicio de la voluntad \emph{(el que se estudia aquí)} & Maniobra engañosa realizada para determinar a una persona a celebrar un acto que no tenía intención de celebrar, o a celebrarlo en condiciones más desventajosas. \\
\bottomrule
\end{tabularx}
\caption{Las tres acepciones del dolo en el CCyCN}
\end{table}

Lo que tienen en común todas estas acepciones es una suerte de \textbf{conciencia de antijuridicidad}: el dolo siempre expresa una cierta conciencia de que se está actuando de modo contrario al derecho, intencionadamente causando un perjuicio o una situación prohibida.

\section{Concepto: acción y omisión dolosa --- Art. 271}

\begin{articulo}{Art. 271 CCyCN --- Acción y omisión dolosa}
Acción dolosa es toda aserción de lo falso, disimulación de lo verdadero, cualquier artificio, astucia o maquinación que se emplee para la celebración del acto. La omisión dolosa causa los mismos efectos que la acción dolosa, cuando el acto no se habría realizado sin la reticencia u ocultación.
\end{articulo}

La incorporación de la \textbf{omisión dolosa} es una novedad bien recibida por la doctrina. Significa que la maniobra engañosa puede realizarse tanto a través de acciones como de omisiones: ocultar información o ser reticente al proporcionarla puede inducir a la otra parte a celebrar un acto que afecta sus intereses y configura, por tanto, el concepto de engaño.

Las formas que puede tomar la maniobra engañosa son:
\begin{itemize}
    \item \textbf{Aseverar lo falso}: afirmar algo que se sabe incorrecto.
    \item \textbf{Disimular lo verdadero}: ocultar la realidad.
    \item \textbf{Cualquier artificio, astucia o maquinación}: no es una mentira casual; tiene que ser una maniobra con entidad engañosa.
    \item \textbf{Omisión dolosa}: silencio o reticencia que induce al otro a celebrar el acto en condiciones perjudiciales.
\end{itemize}

\begin{note}
En el fondo de estas conductas hay un \textbf{deber de buena fe} entre las partes: un deber de lealtad y de informar. En el derecho del consumidor estos principios son especialmente relevantes: la publicidad engañosa y la omisión de información relevante son manifestaciones colaterales del problema del dolo que pueden tener regulación específica.
\end{note}

\section{Dolo esencial --- requisitos del Art. 272}

\begin{articulo}{Art. 272 CCyCN --- Dolo esencial}
El dolo es esencial y causa la nulidad del acto si es grave, si es determinante de la voluntad, si causa un daño importante y si no ha habido dolo por ambas partes.
\end{articulo}

Para que el dolo cause la nulidad del acto deben reunirse los cuatro requisitos siguientes:

\subsection*{a) Gravedad}

La gravedad es una característica de la maniobra engañosa en sí misma. Para que el dolo sea grave debe tener \textbf{entidad para engañar}: no debe ser una maniobra burda que cualquier persona advertiría a primera vista. El parámetro es el de la persona razonablemente precavida que toma las precauciones comunes. La gravedad de la maniobra es un requisito distinto del daño importante.

\subsection*{b) Determinante de la voluntad}

El dolo debe haber sido \textbf{determinante de la voluntad}. Este requisito es fundamental para distinguirlo del dolo incidental: si sin el engaño la persona igualmente hubiera celebrado el acto, no se configura dolo esencial.

\subsection*{c) Daño importante}

Debe haberse causado un daño: una afectación de un interés o de un derecho de la parte que es víctima del acto.

\subsection*{d) No dolo de ambas partes}

No debe haber habido dolo recíproco, es decir, de ambas partes. Un principio fundamental del derecho es que \textbf{nadie puede alegar su propia torpeza}: si ambas partes se engañaban mutuamente, ninguna puede pedir la nulidad.

\section{Dolo incidental --- Art. 273}

\begin{articulo}{Art. 273 CCyCN --- Dolo incidental}
El dolo incidental no es determinante de la voluntad: la parte igualmente habría celebrado el acto, aunque en condiciones distintas. No causa la nulidad del acto, pero sí puede generar una indemnización de daños y perjuicios.
\end{articulo}

\begin{example}{Dolo incidental}
A está comprando un terreno. La otra parte realiza una maniobra para hacerle creer que el terreno tiene cierta habilitación que en realidad no tiene. Pero A de todos modos iba a comprar ese terreno. El engaño le ha provocado un perjuicio ---por ejemplo, pagó más de lo que debía---, pero no fue determinante de que celebrara el acto. En ese caso, el dolo incidental solo puede dar lugar a una acción de daños y perjuicios, pero no a la nulidad.
\end{example}

\begin{table}[H]
\centering
\begin{tabularx}{\textwidth}{>{\raggedright\arraybackslash}X >{\raggedright\arraybackslash}X}
\toprule
\textbf{Dolo esencial} & \textbf{Dolo incidental} \\
\midrule
Es determinante de la voluntad: sin el engaño, el acto no se hubiera celebrado. & No es determinante: la parte igualmente hubiera celebrado el acto. \\[4pt]
Reúne los requisitos del Art. 272: grave, determinante, daño importante, sin dolo recíproco. & El engaño existe pero no alcanza los requisitos del Art. 272 en cuanto a su incidencia en la voluntad. \\[4pt]
Consecuencia: nulidad relativa del acto + daños y perjuicios. & Consecuencia: solo daños y perjuicios (si se prueban). Sin nulidad. \\
\bottomrule
\end{tabularx}
\caption{Dolo esencial versus dolo incidental}
\end{table}

\section{Dolo de tercero --- Art. 274}

\begin{articulo}{Art. 274 CCyCN --- Dolo de tercero}
El dolo de un tercero causa la nulidad del acto si es esencial. Si el dolo de un tercero es incidental, solo da lugar a daños y perjuicios. La parte que conoció o debió conocer el dolo del tercero responde solidariamente con éste por los daños causados.
\end{articulo}

El dolo de tercero es aquel cometido por alguien distinto de la parte con la cual se está contratando. Las reglas son simétricas al dolo de la propia parte:

\begin{itemize}
    \item Si el tercero reúne los requisitos del Art. 272 (dolo esencial): causa la nulidad.
    \item Si el dolo del tercero reúne los requisitos del Art. 273 (dolo incidental): solo genera responsabilidad por daños y perjuicios.
\end{itemize}

Si la parte contratante tuvo \textbf{conocimiento del dolo del tercero} y no lo comunicó, responde \textbf{solidariamente} por todos los daños causados. La responsabilidad solidaria implica que responde por el total de la deuda, no solo por su parte proporcional.

\section{Consecuencias del dolo: nulidad y daños --- Art. 275}

\begin{articulo}{Art. 275 CCyCN --- Responsabilidad por daños}
El autor del dolo esencial o incidental debe reparar el daño causado. Cuando el dolo sea del tercero, también responde la parte que al tiempo de la celebración del acto tuvo conocimiento del dolo del tercero.
\end{articulo}

Las consecuencias del dolo esencial son:

\begin{enumerate}
    \item \textbf{Nulidad relativa del acto} (Art. 388 CCyCN). Al ser relativa, está dispuesta en beneficio de la parte afectada y es susceptible de confirmación.
    \item \textbf{Acción de daños y perjuicios} en virtud del Art. 275.
    \item \textbf{Prescripción de dos años} desde el momento en que se conoció o pudo ser conocido el dolo (Arts. 2562 y 2563 CCyCN).
\end{enumerate}

% ============================================================
\chapter{La violencia como vicio de la voluntad (Arts.~276--278)}
% ============================================================

\section{Concepto de violencia: fuerza e intimidación}

La violencia es una \textbf{coerción grave de tipo físico o moral} que se ejerce sobre una persona para determinarla a realizar un acto en contra de su voluntad. Es un vicio que afecta la \textbf{libertad} como elemento interno del acto voluntario: la posibilidad de autodeterminarse y tomar decisiones sin coerciones injustas.

La violencia adopta dos grandes formas:

\begin{table}[H]
\centering
\begin{tabularx}{\textwidth}{>{\raggedright\arraybackslash}X >{\raggedright\arraybackslash}X}
\toprule
\textbf{Fuerza (Art. 276, primera parte)} & \textbf{Intimidación (Art. 276, segunda parte)} \\
\midrule
Violencia física irresistible. & Amenazas que generan el temor de sufrir un mal grave e inminente. \\[4pt]
Ejemplo: torturas físicas para que alguien firme un contrato. & Ejemplo: «Si no me vendés tu casa, voy a hacerte daño a vos o a tu familia». \\[4pt]
La coerción es de tipo corporal y directa. & La coerción opera a través del temor generado por las amenazas. \\
\bottomrule
\end{tabularx}
\caption{Formas de la violencia}
\end{table}

\begin{example}{Fuerza física irresistible}
A es torturado físicamente para que firme la escritura de venta de su casa. Aparece todo en regla: hay una escritura firmada, inscripta en el Registro de la Propiedad. Pero ese acto está viciado por fuerza física irresistible. Se inicia un juicio de nulidad para volver la situación al estado anterior a la celebración del acto. Este hecho también puede configurar delitos penales, pero aquí se analiza la dimensión civil.
\end{example}

\section{Requisitos de la intimidación --- Art. 276}

\begin{articulo}{Art. 276 CCyCN --- Fuerza e intimidación}
La fuerza irresistible y las amenazas que generan el temor de sufrir un mal grave e inminente que no se puedan contrarrestar o evitar en la persona o bienes de la parte o de un tercero, causan la nulidad del acto. La relevancia de las amenazas debe ser juzgada teniendo en cuenta la situación del amenazado y las demás circunstancias del caso.
\end{articulo}

\subsection*{a) Las amenazas deben ser injustas}

Las amenazas que configuran el vicio son \textbf{amenazas injustas}. El ejercicio legítimo de un derecho no constituye una amenaza en este sentido. Si una persona le debe dinero a otra y le dice «si no me pagás, voy a iniciar un juicio», no le está amenazando en el sentido jurídicamente relevante: está ejerciendo su derecho. Una amenaza ilícita sería, en cambio: «si no me vendés tu casa, te voy a romper la tuya».

\subsection*{b) El temor se mide según las circunstancias del amenazado}

La relevancia de las amenazas ha de ser juzgada teniendo en cuenta la \textbf{situación del amenazado y las demás circunstancias del caso}. No hay un estándar único: la misma amenaza puede configurar intimidación para una persona pero no para otra, según sus capacidades y contexto.

\begin{example}{Medición del temor según las circunstancias}
Si alguien con millones de seguidores en redes sociales amenaza con publicar que una persona es delincuente, eso puede generar un temor fundado a sufrir un daño grave a la reputación. El contexto (quién amenaza, con qué medios, a quién) determina la relevancia de la amenaza.
\end{example}

\subsection*{c) El mal debe ser grave e inminente}

\textbf{Grave}: el mal debe tener entidad suficiente. Puede ser un mal sobre la persona (integridad física, moral o espiritual) o sobre los bienes.

\textbf{Inminente}: significa que el mal va a ocurrir y que no puede ser contrarrestado ni evitado. El término «evitar» fue agregado por el nuevo Código; no figuraba en el anterior.

\subsection*{d) El mal puede recaer sobre la parte o un tercero}

En el Código de Vélez se enumeraban taxativamente las personas sobre quienes podía recaer la amenaza: cónyuge, ascendientes o descendientes. Esta enumeración generó una disputa: ¿era taxativa o ejemplificativa? La doctrina mayoritaria la interpretaba como ejemplificativa.

El nuevo Código resuelve el problema directamente indicando: \textbf{sobre la parte o un tercero}. Puede ser un familiar, un amigo, o incluso alguien que no se conoce tanto, pero cuya situación genera el temor suficiente.

\section{Violencia de tercero y responsabilidad solidaria --- Art. 278}

Al igual que en el dolo, la violencia puede provenir de la propia parte o de un tercero. En ambos casos da lugar a la nulidad del acto. La diferencia relevante es la \textbf{responsabilidad de la parte cómplice}.

\begin{articulo}{Art. 278 CCyCN --- Responsabilidad por daños causados por violencia}
El autor de la fuerza irresistible y de las amenazas debe reparar los daños causados. La parte que al tiempo de la celebración del acto tuvo conocimiento de la fuerza o de las amenazas responde solidariamente con el autor por los daños.
\end{articulo}

\begin{example}{Violencia de tercero --- complicidad de la parte}
Un tercero amenaza a A para que le venda su casa a B. A le vende la casa a B (que es la parte contratante), pero la amenaza vino del tercero. Si B tuvo conocimiento de la fuerza o de la intimidación ejercida por el tercero, responde solidariamente por todos los daños y perjuicios: responde por el total, no solo por su parte proporcional.
\end{example}

\section{El temor reverencial}

El \textbf{temor reverencial} es el respeto o sumisión de una persona hacia sus superiores (por ejemplo, de un militar hacia un oficial superior). El temor reverencial \textbf{no es causa de nulidad}.

Para configurar el vicio de violencia deben existir amenazas concretas: la sola presencia de personas con autoridad no produce una situación de afectación de la libertad. El temor a una figura de autoridad, por sí solo, es temor reverencial y no causa nulidad. Este criterio se mantiene tanto en el Código de Vélez como en el nuevo CCyCN.

\begin{note}
Si el superior jerárquico realiza amenazas concretas que reúnen los requisitos de los Arts. 276 y siguientes, sí puede configurarse el vicio de violencia y dar lugar tanto a la nulidad como a los daños y perjuicios.
\end{note}

\section{Consecuencias de la violencia: nulidad y prescripción}

Las consecuencias de la violencia son análogas a las del dolo y el error:

\begin{enumerate}
    \item \textbf{Nulidad relativa del acto}. Al ser relativa, la parte afectada puede confirmar el acto si así lo desea.
    \item \textbf{Acción de daños y perjuicios} (Art. 278 CCyCN).
\end{enumerate}

La prescripción también es de \textbf{dos años} (Arts. 2562 y 2563 CCyCN), pero existe una diferencia crucial respecto del error y el dolo: el plazo comienza a correr desde el momento en que \textbf{cesa la violencia}, no desde que se conoció el vicio.

\begin{articulo}{Diferencia en el inicio del plazo de prescripción}
Error: desde que se conoció o pudo ser conocido el error.\\
Dolo: desde que se conoció o pudo ser conocido el dolo.\\
Violencia: desde que cesa la violencia (Arts. 2562 y 2563 CCyCN).
\end{articulo}

\begin{example}{Inicio del plazo desde el cese de la violencia}
A es secuestrado durante un año y medio. Durante ese tiempo, bajo torturas, firma un contrato vendiendo su casa. Después lo liberan. Desde el momento en que cesa la violencia ---desde que recobra la libertad--- comienza a correr el plazo de dos años para iniciar la acción de nulidad. Mientras dura la coerción, la persona no puede libremente ejercer su derecho a pedir la nulidad.
\end{example}

% ============================================================

\section{Cuadro Cornell --- Vicios de la Voluntad}
% ============================================================

\begin{table}[H]
\centering
\begin{tabularx}{\textwidth}{
    >{\raggedright\arraybackslash}p{0.30\textwidth}
    >{\raggedright\arraybackslash}X
}
\toprule
\textbf{Preguntas / palabras clave} & \textbf{Notas / respuestas} \\
\midrule

\textbf{¿Por qué se prefiere hablar de «vicios de la voluntad» y no de «vicios del consentimiento»?} &
Porque hablar de vicios del consentimiento los limitaría a los contratos (un tipo específico de acto jurídico). Los vicios de la voluntad, en cambio, afectan a todo acto voluntario. El CCyCN adopta esta denominación más amplia. \\[6pt]

\textbf{¿Cuáles son los elementos internos del acto voluntario y qué vicio afecta a cada uno?} &
Los tres elementos internos son discernimiento, intención y libertad (Art. 260 CCyCN). La \textbf{intención} es afectada por el error y el dolo; la \textbf{libertad} es afectada por la violencia (fuerza o intimidación). \\[6pt]

\textbf{¿Qué diferencia hay entre error esencial y error incidental?} &
El error esencial recae sobre un elemento esencial del acto (naturaleza, objeto, cualidad sustancial, motivos determinantes, persona). El error incidental recae sobre elementos secundarios no determinantes de la voluntad. Solo el error esencial ---y reconocible--- puede causar la nulidad. \\[6pt]

\textbf{¿Qué significa que un error sea «reconocible» (Art. 266)?} &
Que el destinatario de la declaración pudo conocer el error según la naturaleza del acto y las circunstancias de persona, tiempo y lugar. Es un criterio objetivo centrado en el destinatario, que reemplaza la «excusabilidad» del Código de Vélez (criterio subjetivo centrado en quien erraba). \\[6pt]

\textbf{¿Qué sucede si el destinatario ofrece cumplir el acto según lo que el otro quiso? (Art. 269)} &
Se evita la nulidad del acto. La parte en error no puede solicitar la nulidad si el destinatario ofrece ejecutar el acto con las modalidades y el contenido que aquella pretendió celebrar. Esto responde al principio de conservación de los actos. \\[6pt]

\textbf{¿Qué es el dolo como vicio de la voluntad? (Art. 271)} &
Una maniobra engañosa: toda aserción de lo falso, disimulación de lo verdadero, cualquier artificio, astucia o maquinación para inducir a una persona a celebrar un acto. La omisión dolosa (reticencia u ocultación) tiene los mismos efectos que la acción dolosa cuando el acto no se habría realizado sin ella. \\[6pt]

\textbf{¿Cuáles son los requisitos del dolo esencial? (Art. 272)} &
Para causar la nulidad, el dolo debe ser: (1) \textbf{grave} ---con entidad para engañar---; (2) \textbf{determinante de la voluntad} ---sin el engaño, el acto no se hubiera celebrado---; (3) causante de un \textbf{daño importante}; (4) \textbf{unilateral} ---no haber habido dolo de ambas partes---. \\[6pt]

\textbf{¿Qué diferencia al dolo incidental del dolo esencial? (Art. 273)} &
El dolo incidental no es determinante de la voluntad: la parte igualmente hubiera celebrado el acto. No causa la nulidad, pero sí puede dar lugar a una indemnización por daños y perjuicios si se prueban. \\[6pt]

\textbf{¿Puede el dolo cometido por un tercero causar la nulidad? (Art. 274)} &
Sí. El dolo de tercero causa la nulidad si es esencial, y solo daños y perjuicios si es incidental. Si la parte contratante tuvo conocimiento del dolo del tercero, responde solidariamente por todos los daños. \\[6pt]

\textbf{¿Cuáles son las dos formas de la violencia? (Art. 276)} &
\textbf{Fuerza}: coerción física irresistible. \textbf{Intimidación}: amenazas injustas que generan el temor de sufrir un mal grave e inminente que no puede contrarrestarse ni evitarse, sobre la persona o bienes de la parte o de un tercero. \\[6pt]

\textbf{¿Cómo se mide si las amenazas son relevantes como intimidación?} &
Se juzga teniendo en cuenta la situación del amenazado y las demás circunstancias del caso. No hay un estándar único: la misma amenaza puede configurar intimidación para una persona pero no para otra, según sus capacidades y contexto. \\[6pt]

\textbf{¿Qué es el temor reverencial y tiene efectos jurídicos?} &
Es el respeto o sumisión de una persona hacia su superior jerárquico (por ejemplo, un militar ante un oficial superior). El temor reverencial solo, sin amenazas concretas, no es causa de nulidad. Solo cuando se suman amenazas reales que reúnen los requisitos del Art. 276 puede configurarse el vicio de violencia. \\[6pt]

\textbf{¿Desde cuándo corre el plazo de prescripción para cada vicio?} &
En todos los casos el plazo es de dos años (Arts. 2562 y 2563 CCyCN). La diferencia está en el punto de inicio: \textbf{error y dolo}: desde que se conoció o pudo ser conocido el vicio. \textbf{Violencia}: desde que cesa la violencia. \\[6pt]

\textbf{¿Cuál es el tipo de nulidad que producen los vicios de la voluntad?} &
Todos los vicios de la voluntad (error, dolo, violencia) producen \textbf{nulidad relativa} (Art. 388 CCyCN). Eso significa que está dispuesta en beneficio de la parte afectada, puede ser confirmada por ella y prescribe a los dos años. \\

\midrule

\multicolumn{2}{p{\textwidth}}{%
\textbf{Resumen:} Los vicios de la voluntad son defectos que afectan los elementos internos del acto voluntario ---intención y libertad--- y que, cuando reúnen los requisitos que establece el CCyCN, producen la nulidad relativa del acto. El \textbf{error} (Arts. 265--270) afecta la intención por una falsa o errónea representación de la realidad: para causar la nulidad debe ser esencial ---recaer sobre la naturaleza del acto, el objeto, la cualidad sustancial, los motivos determinantes o la persona--- y reconocible por el destinatario, criterio que el nuevo Código tomó del derecho italiano en reemplazo de la excusabilidad del Código de Vélez. El \textbf{dolo} (Arts. 271--275) afecta la intención mediante una maniobra engañosa ---acción u omisión--- que para producir la nulidad debe ser grave, determinante de la voluntad, causante de un daño importante y unilateral; si no es determinante, solo genera daños y perjuicios (dolo incidental). La \textbf{violencia} (Arts. 276--278) afecta la libertad mediante fuerza física irresistible o intimidación a través de amenazas injustas que generan el temor de sufrir un mal grave e inminente; la relevancia de las amenazas se mide según las circunstancias del amenazado. En los tres vicios, la consecuencia es la nulidad relativa del acto ---confirmable por la parte afectada--- y la posibilidad de reclamar daños y perjuicios. El plazo de prescripción es de dos años en todos los casos: para el error y el dolo, desde que se conoció o pudo conocerse el vicio; para la violencia, desde que cesa la coerción.%
} \\

\bottomrule
\end{tabularx}
\caption{Cuadro Cornell --- Vicios de la Voluntad (Error, Dolo y Violencia)}
\end{table}
\chapter[Antecedentes históricos y derecho comparado]{Antecedentes históricos y derecho comparado de la lesión}
% ===========================================================================

\section{Los vicios de los actos jurídicos según el Código Civil y Comercial}

El estudio de la lesión se inserta dentro del análisis general de los vicios propios de los actos jurídicos. En el ordenamiento argentino vigente, estos vicios son tres: la \textbf{lesión}, la \textbf{simulación} y el \textbf{fraude}.

A diferencia del tratamiento que recibían los vicios de la voluntad —cada uno regulado en un capítulo propio—, en el Código Civil y Comercial (CCC) la lesión, la simulación y el fraude están tratados dentro del capítulo sexto del título cuarto, distribuidos en tres secciones.

\begin{note}
La lesión no estaba incluida originariamente en el Código de Vélez Sarsfield. Se trata de un instituto que resulta central para comprender el funcionamiento concreto de la justicia conmutativa en el derecho privado patrimonial, razón por la cual merece un desarrollo extenso y diferenciado, con antecedentes históricos propios.
\end{note}

\section{Concepto de lesión: una distinción necesaria}

\begin{definition}{Lesión como vicio de los actos jurídicos}
Cuando se habla de lesión como vicio de los actos jurídicos, \textbf{no} se hace referencia a la lesión física —como en el derecho penal, donde el delito de lesiones se clasifica en graves, levísimas y leves—. Se trata de un concepto jurídico distinto, vinculado a la existencia de una \textbf{desproporción entre las prestaciones} de un acto jurídico.
\end{definition}

En todo acto jurídico —especialmente en la compraventa— existe la expectativa de que las prestaciones sean equivalentes. Este principio se vincula con la denominada \textbf{justicia conmutativa}, entendida como la justicia entre particulares que busca una igualdad aritmética: que lo que una parte da sea igual a lo que recibe.

La pregunta central que atraviesa la construcción histórica del instituto es la siguiente: en principio, se supone que las partes son libres y pactan voluntariamente los términos de su acuerdo, de modo que no habría razón para invalidarlo si no existió error, dolo ni violencia. Sin embargo, cuando aparece una desproporción notable entre las prestaciones, surge el interrogante de en qué punto esa desproporción se vuelve tan grave como para dar lugar a la lesión.

A este elemento objetivo —la desproporción— se suma un segundo elemento, de naturaleza subjetiva: la explotación, por una de las partes, de la inferioridad de la otra, con el fin de obtener una ventaja patrimonial evidentemente desproporcionada.

\section{El derecho romano y el Código Justinianeo}

En el derecho romano existen dos rescriptos atribuidos a los emperadores Diocleciano y Maximino que habrían instituido la figura, aunque se ha discutido si se trata de textos insertados con posterioridad. Lo indudable es que el instituto de la lesión aparece recogido en el Código Justinianeo, la recopilación del ordenamiento jurídico romano elaborada bajo el Emperador Justiniano.

\begin{formula}{Código Justinianeo (Corpus Iuris Civilis)}
``Si tú o tu padre hubiera vendido por menor precio una cosa, es humano que, restituyendo todo el precio a los compradores, recobres el fondo vendido mediando la autoridad del juez, aunque si el comprador lo prefiere, reciban lo que les falta al justo precio. Pero se considera que el precio es menor si no se hubiera recibido ni la mitad del verdadero precio.''
\end{formula}

Este instituto recibió el nombre de \textbf{lesión enorme}. Era concedida al vendedor en la compraventa de inmuebles, y exigía que la desproporción alcanzara una medida determinada: más de la mitad del valor real de la cosa.

La lesión enorme daba lugar a dos acciones posibles:

\begin{itemize}
    \item una acción para recuperar el inmueble (el ``fundo''); o bien
    \item una acción para que el comprador tuviera la opción de pagar el complemento del precio faltante, es decir, un reajuste.
\end{itemize}

\begin{example}{Reajuste del precio en el derecho romano}
Si la cosa vale 100 y se pagó 40, el comprador —a quien le interesa conservar la cosa— puede optar por pagar los 60 que faltan, en lugar de restituir el fondo.
\end{example}

\section{La Edad Media y el precio justo}

Los glosadores, que continuaron comentando y ampliando el derecho romano, extendieron la aplicación de la lesión a la mayoría de los contratos durante la Edad Media. Los canonistas, por su parte, pusieron especial énfasis en la inmoralidad del desequilibrio entre las prestaciones.

Este factor moral resulta históricamente relevante, porque en torno a la lesión siempre ha oscilado la tensión entre una mirada objetiva —centrada en la desproporción del objeto— y un reproche dirigido a quien explota una situación ajena.

En este contexto adquiere especial relevancia el aporte de \textbf{Tomás de Aquino} y su idea del \textbf{precio justo}, un concepto de difícil determinación pero de gran influencia doctrinal: cuando el precio resulta manifiestamente injusto, se configura una posible lesión.

\begin{formula}{Tomás de Aquino — El precio justo}
Aun la ley humana —en referencia al derecho romano— considera ilícito que el vendedor venda una cosa por más de lo que vale, o que el comprador la adquiera por menos de su valor, aunque no lo tiene por ilícito si la diferencia no es excesiva. Se reputa ilícito cuando el engaño supera la mitad del valor de la cosa, en clara referencia a la métrica del Código Justinianeo.
\end{formula}

\section{La legislación española como antecedente del derecho argentino}

La evolución de la figura en las fuentes españolas —antecedente directo del derecho argentino— puede sintetizarse del siguiente modo:

\begin{table}[H]
\centering
\caption{Tratamiento de la lesión en la legislación española}
\begin{tabularx}{\textwidth}{>{\bfseries\raggedright\arraybackslash}p{0.28\textwidth} X}
\toprule
\textbf{Fuente española} & \textbf{Tratamiento de la lesión} \\
\midrule
Fuero Juzgo & Rigió durante la dominación visigoda. No admite la lesión, pero combate seriamente la usura, instituto muy cercano. \\
Fuero Real & Concede la acción al vendedor cuando lo engañan en más de la mitad del valor, si el comprador no acepta pagar la diferencia. \\
Las Leyes de Partidas & Consolidan la acción tanto para el vendedor como para el comprador. \\
Novísima Recopilación & Acepta la acción en compraventa, permuta, rentas y otros contratos semejantes. \\
\bottomrule
\end{tabularx}
\end{table}

\section{El derecho comparado moderno}

\subsection{El Código Civil Francés (1804): la lesión objetiva}

El Código Napoleónico constituye el ejemplo por excelencia de lo que se ha denominado \textbf{lesión objetiva}: aquella que contempla únicamente los elementos objetivos —es decir, la desproporción— sin tomar en consideración, en la configuración de la norma, el elemento subjetivo de la explotación de la inferioridad de la otra parte.

\begin{formula}{Código Civil Francés — art. 1.674 (1804)}
Cuando el vendedor resulte lesionado en más de siete doceavos del precio del inmueble, tendrá derecho a solicitar la rescisión de la venta, aunque en el contrato haya renunciado expresamente a solicitar dicha rescisión y aunque haya declarado donar las plusvalías.
\end{formula}

La medida establecida —más de 7/12— representa un poco más que la mitad (que serían 6/12). El código francés sigue así la métrica del derecho romano, pero la precisa mediante un porcentaje específico, tomando el valor al momento de la venta.

\begin{example}{Complemento del precio en el código francés}
Se concede al comprador —lesionante— la posibilidad de ofrecer el complemento del precio con la deducción de la décima parte. Si la cosa vale 100 y se pagaron 40, el comprador puede ofrecer el complemento hasta 90, pagando los 50 restantes, y de esa manera conservar la propiedad.
\end{example}

En el derecho argentino, esta acción de complemento recibe el nombre de \textbf{reajuste equitativo} y es fijada por el juez. En el código francés, la prescripción de la acción era de dos años.

\subsection{El Código Civil Alemán — BGB (1900): la lesión subjetiva}

En 1900, la reforma del código civil alemán incorpora por primera vez la fórmula de la denominada \textbf{lesión subjetiva}, estrechamente vinculada al combate de la usura. La doctrina coincide en señalar este momento como el surgimiento de dicha categoría.

\begin{formula}{Código Civil Alemán — BGB (síntesis, 1900)}
Un acto jurídico contrario a las buenas costumbres por el cual alguien, explotando la necesidad, la ligereza o la inexperiencia de otro, obtiene para sí o para un tercero que a cambio de una prestación le prometa o entregue ventajas patrimoniales que exceden de tal forma el valor de la prestación que, teniendo en cuenta las circunstancias, existe una desproporción chocante con ella.
\end{formula}

Aparecen aquí los elementos subjetivos del lesionante —que explota la necesidad, la ligereza o la inexperiencia— y se hace prometer o entregar ventajas patrimoniales. Existe, entonces, una acción activa: el lesionante toma la iniciativa para explotar una inferioridad. Del lado del lesionado, la inferioridad se define como necesidad, ligereza o inexperiencia, términos que luego reproducirá casi textualmente el art.\ 954 argentino y, con alguna variante, el código vigente.

\begin{important}
El código alemán declaraba el acto nulo en su totalidad, sin admitir nulidades parciales ni reajuste de precios, por tratarse de un acto contrario a las buenas costumbres. Existía un deber de restituir; la acción no era confirmable, era imprescriptible y podía ser declarada de oficio.
\end{important}

\subsection{El Código Suizo (1912) y el Código Italiano (1942)}

En 1912, el Código Suizo de las Obligaciones sigue al modelo alemán con variantes propias: coloca la prescripción en un año, torna la acción anulable y permite la confirmación del acto, acercándose así a la idea de un vicio que da lugar a una nulidad relativa susceptible de confirmación.

En 1942 se sanciona el Código Civil Italiano, que ejercerá una influencia decisiva sobre la reforma del art.\ 954 argentino en 1968.

\begin{formula}{Código Civil Italiano — art. 1.448 (1942, síntesis)}
Si hubiera desproporción en las prestaciones de una de las partes, y esa desproporción dependiese del estado de una de ellas, del cual se aprovechara la otra para obtener ventajas, la parte damnificada podrá demandar la rescisión del contrato. La acción no es admisible si la lesión es inferior a la mitad del valor que la prestación ejecutada o prometida por la parte perjudicada tenía en el momento del contrato. La lesión debe perdurar hasta el momento en que se proponga la demanda.
\end{formula}

\begin{note}
La exigencia de que la lesión perdure hasta el momento de la demanda no estaba presente en los códigos anteriores, y será recogida tanto por la ley 17.711 como por el CCC vigente. En el código italiano, además, los contratos aleatorios no pueden ser rescindidos por esta vía.
\end{note}

% ---------------------------------------------------------
% CORNELL — CAPÍTULO 1
% ---------------------------------------------------------
\section{Cuadro de repaso — Método Cornell}
\begin{table}[H]
\centering
\begin{tabularx}{\textwidth}{
    >{\raggedright\arraybackslash}p{0.30\textwidth}
    >{\raggedright\arraybackslash}X
}
\toprule
\textbf{Preguntas / palabras clave} & \textbf{Notas / respuestas} \\
\midrule

\textbf{¿Qué son los vicios propios de los actos jurídicos en el CCC?} &
Son tres: lesión, simulación y fraude. Están regulados en el capítulo sexto del título cuarto del Código Civil y Comercial, distribuidos en tres secciones. \\

\textbf{¿Qué es la lesión como vicio jurídico?} &
No es la lesión física, sino una desproporción entre las prestaciones de un acto jurídico, combinada con la explotación, por una de las partes, de la inferioridad de la otra. Se vincula con la idea de justicia conmutativa (igualdad aritmética entre lo que se da y lo que se recibe). \\

\textbf{¿Cuál es el origen histórico del instituto?} &
El Código Justinianeo (siglo VI) consagró la ``lesión enorme'' en la compraventa de inmuebles: si el precio era inferior a la mitad del valor real, el vendedor podía pedir la nulidad o el complemento del precio. \\

\textbf{¿Qué distinción clave introduce el BGB alemán de 1900?} &
Introduce la lesión subjetiva: ya no basta la desproporción objetiva, sino que se requiere además la explotación de la necesidad, ligereza o inexperiencia de la otra parte. Esta fórmula fue luego adoptada por Argentina en 1968 y rige, con variantes, en el art.\ 332 del CCC. \\

\midrule

\multicolumn{2}{p{\dimexpr\textwidth-2\tabcolsep\relax}}{
\textbf{Resumen:} La lesión, desde la lesión enorme del Código Justinianeo (más de la mitad del valor en la compraventa de inmuebles), pasando por el énfasis moral de la Edad Media y el precio justo de Tomás de Aquino, evolucionó hacia dos grandes modelos de derecho comparado: la \textbf{lesión objetiva} del código francés (que exige únicamente la desproporción) y la \textbf{lesión subjetiva} del BGB alemán (que agrega la explotación de la inferioridad de la otra parte), modelo este último que —junto con el código italiano de 1942— influirá decisivamente en la incorporación de la figura al derecho argentino.
} \\

\bottomrule
\end{tabularx}
\end{table}

% ===========================================================================
\chapter{La lesión en el derecho argentino}
% ===========================================================================

\section{El Código de Vélez (1869): la exclusión de la lesión}

Vélez Sarsfield no incluyó la lesión entre los vicios de los actos jurídicos. Al tratar los actos jurídicos en el art.\ 944, y luego de referirse a la simulación, el codificador dedicó una extensa nota —ubicada luego del art.\ 943— para explicar por qué no adoptaba la figura.

\begin{note}
Técnicamente, la nota no se refiere al contenido del art.\ 943 —que trataba sobre la violencia—, sino que Vélez la ubicó allí, antes del comienzo del tratamiento de los actos jurídicos, con el propósito de explicar en algún lugar del código por qué no incorporaba la lesión. Años más tarde, Borda señaló en un fallo de 1958 que una nota al pie no constituye ley.
\end{note}

Vélez ofrece dos grandes razones para no reconocer la lesión:

\begin{formula}{Nota de Vélez al art. 943 — Primer argumento}
``Para sostener nosotros que la lesión enorme, menor o mínima vicia los actos y obtenernos por tanto de proyectar disposiciones de la materia, basta comparar las diversas legislaciones y de las diferencias entre ellas resultados que no han tenido un principio uniforme al establecer esta ley.''
\end{formula}

El codificador observa que la legislación comparada presenta numerosas diferencias: unas la conceden al vendedor, otras al comprador; unas la aplican a la compraventa, otras a distintos actos; unas imponen una medida de más de la mitad, otras no. Ante esta disparidad de criterios, Vélez concluye que no encuentra un principio uniforme y que, por lo tanto, no se siente obligado a incorporar la figura.

\begin{formula}{Nota de Vélez al art. 943 — Segundo argumento}
``Dejaríamos de ser responsables de nuestras acciones si la ley nos permitiera enmendar todos nuestros errores o todas las consecuencias de ellos. El consentimiento libre prestado sin error, dolo y violencia y con la solemnidad requerida por las leyes debe hacer irrevocables los contratos.''
\end{formula}

Este segundo argumento expresa que lo pactado sin error, dolo ni violencia es válido, y que quien se equivoca debe hacerse cargo de las consecuencias de su propio error. Se advierte aquí una visión marcadamente individualista y liberal, coincidente con otras opciones adoptadas por Vélez a lo largo de todo el código.

\section{La jurisprudencia anterior a la ley 17.711}

Durante el siglo XX, distintos fallos judiciales fueron invalidando actos manifiestamente desproporcionados, invocando generalmente el art.\ 953 —referido al objeto del acto—, que disponía que el objeto no podía consistir en cosas o hechos contrarios a la moral. Cuando se pactaba algo contrario a la moral, ello podía dar lugar a la nulidad, de modo que el vicio de lesión se introducía por vía del art.\ 953.

La nota de Vélez al art.\ 943 seguía vigente, pero algunos jueces —entre ellos, Borda, en un fallo de la Cámara Nacional de Apelaciones de 1958 conocido como el caso ``Orlando''— sostuvieron que las notas no constituyen texto legal y que la lesión resulta incompatible con las buenas costumbres.

\begin{example}{Casos de aplicación jurisprudencial de la lesión (antes de 1968)}
Préstamos usurarios, honorarios exorbitantes del administrador de una sucesión, venta a precio vil (muy bajo) de un terreno, y moderación de cláusulas penales excesivas.
\end{example}

En 1961, el Tercer Congreso de Derecho Civil tuvo una importancia decisiva para la reforma de la ley 17.711, al formular recomendaciones que constituyeron, en gran medida, la base de la reforma de 1968.

\section{La ley 17.711 y la reforma al art. 954 (1968)}

La ley 17.711 incorporó diversas figuras al Código Civil —entre ellas, el abuso del derecho y modificaciones en los efectos de la declaración de voluntad en el tiempo— y también la lesión, incorporada en el art.\ 954.

\begin{formula}{Art. 954 — Ley 17.711 (1968, texto incorporado)}
``También podrá demandarse la nulidad o modificación de los actos jurídicos cuando una de las partes, explotando la necesidad, ligereza o inexperiencia de la otra, obtuviera por medio de ellos una ventaja patrimonial evidentemente desproporcionada y sin justificación. Se presume, salvo prueba en contrario, que existe tal explotación en caso de notable desproporción de las prestaciones. Los cálculos deben hacerse según valores al tiempo del acto y la desproporción deberá subsistir al momento de la demanda. Solo el lesionado o sus herederos podrán ejercer la acción cuya prescripción operará a los cinco años. El accionante tiene opción para demandar la nulidad o el reajuste equitativo del convenio, pero la primera de estas acciones se transformará en acción de reajuste si éste fuera ofrecido por el demandado.''
\end{formula}

En este texto se incorporan claramente un elemento objetivo —la desproporción— y un elemento subjetivo —la explotación de la necesidad, ligereza o inexperiencia—. El nuevo código sustituirá luego la palabra ``ligereza'' por ``debilidad psíquica''. Cabe destacar que el texto del art.\ 954 no establece una métrica fija para medir el elemento objetivo.

\section{El art. 332 del CCC: texto y diferencias con el art. 954}

El Código Civil y Comercial vigente sigue, en lo sustancial, al art.\ 954 para regular la lesión, en un único artículo ubicado en la sección primera del capítulo sexto del título cuarto.

\begin{important}
Aunque se trate de un único artículo, el art.\ 332 contiene un desarrollo normativo extenso, que exige un análisis detenido de cada uno de sus elementos.
\end{important}

\begin{formula}{Art. 332 — Código Civil y Comercial de la Nación (vigente)}
``Puede demandarse la nulidad o modificación de los actos jurídicos cuando una de las partes, explotando la necesidad, debilidad psíquica o inexperiencia de la otra, obtuviera por medio de ellos una ventaja patrimonial evidentemente desproporcionada y sin justificación. Se presume, salvo prueba en contrario, que existe tal explotación en casos de notable desproporción de las prestaciones. Los cálculos deben hacerse según valores al tiempo del acto y la desproporción debe subsistir al momento de la demanda. El afectado tiene opción para demandar la nulidad o un reajuste equitativo del convenio, pero la primera de estas acciones se transforma en acción de reajuste si éste es ofrecido por el demandado al contestar la demanda. Solo el lesionado o sus herederos podrán ejercer la acción.''
\end{formula}

Las diferencias del CCC respecto del art.\ 954 pueden sintetizarse así:

\begin{table}[H]
\centering
\caption{Diferencias entre el art. 954 (ley 17.711) y el art. 332 (CCC)}
\begin{tabularx}{\textwidth}{>{\bfseries\raggedright\arraybackslash}p{0.20\textwidth} X}
\toprule
\textbf{Modificación} & \textbf{Explicación} \\
\midrule
Cambio 1 & La palabra ``ligereza'' es reemplazada por ``debilidad psíquica''. El nuevo código clarifica así un punto que había dado lugar a un problema interpretativo, entre quienes ubicaban la ligereza como un obrar irreflexivo sin mayores requisitos y quienes la entendían como expresión de una cierta debilidad psíquica; el CCC se pronuncia por esta última postura. \\
Cambio 2 & El plazo de prescripción se reduce de cinco a dos años. Metodológicamente resulta correcto tratar el plazo en el libro dedicado a las prescripciones, y no dentro del propio artículo. \\
\bottomrule
\end{tabularx}
\end{table}

\section{Ámbito de aplicación del art. 332}

El art.\ 332 recoge, en principio, todos los actos jurídicos en general. Sin embargo, el acto debe ser bilateral y no a título gratuito, puesto que en los actos a título gratuito siempre existe una desproporción entre las prestaciones que resulta ajena a la lógica de la lesión.

\section{El elemento objetivo: la desproporción de las prestaciones}

El primer elemento exigido por la norma es el elemento objetivo: la desproporción de las prestaciones. A diferencia del código francés, que establece un criterio fijo (7/12 partes), el CCC habla de ``ventaja evidentemente desproporcionada'' y de ``notable desproporción'', sin fijar un criterio numérico. No obstante, se exige que la diferencia sea muy notable: no bastaría, por ejemplo, una diferencia del veinte por ciento.

Un requisito adicional es que la desproporción debe ser \textbf{sin justificación}.

\begin{example}{Desproporción con justificación}
Un tío que va a asistir al casamiento de un sobrino le vende un departamento que vale 100.000 dólares en 40.000, como regalo de casamiento. Tras su muerte, otros herederos —primos del beneficiado— demandan por lesión, invocando la diferencia entre el valor real y el precio pagado. El demandado se defenderá alegando que existió una justificación: la intención de realizar una liberalidad. Podrá discutirse si ese regalo excede o no la porción disponible —cuestión de legítima—, pero no podrá hablarse de lesión en sentido propio, porque no hubo explotación sino una justificación para la desproporción.
\end{example}

Los cálculos deben realizarse tomando en cuenta los valores al momento del acto, puesto que todo vicio afecta al acto desde su origen. Sin embargo, la desproporción debe además subsistir al momento de la demanda, exigencia que proviene del código italiano.

\begin{example}{Desproporción que no subsiste}
Una persona compra muy barato un inmueble en un barrio cotizado. La desproporción existe al momento del acto, pero luego el valor de la propiedad desciende notablemente porque el barrio se deteriora y se vuelve inhabitable. Como la desproporción no subsiste al momento de la demanda, no existe acción por lesión, aun cuando el vicio se hubiera configurado al momento del acto.
\end{example}

\section{Los elementos subjetivos}

El art.\ 332 exige, además del elemento objetivo, la \textbf{explotación} por parte del lesionante. Se trata de algo más intenso que un simple ``aprovechar'' —entendido como que la otra parte no advierte la desproporción y el lesionante simplemente continúa adelante—: existe aquí una acción más activa, consistente en ``hacerse dar o prometer'' aprovechando la inferioridad de la otra parte. Este elemento resulta, en la práctica, de difícil prueba.

Del lado del lesionado, el CCC exige que se configure alguno de los siguientes tres estados de inferioridad:

\begin{table}[H]
\centering
\caption{Estados de inferioridad del lesionado (art. 332 CCC)}
\begin{tabularx}{\textwidth}{>{\bfseries\raggedright\arraybackslash}p{0.22\textwidth} X}
\toprule
\textbf{Estado de inferioridad} & \textbf{Definición y aclaración} \\
\midrule
Necesidad & Falta de lo indispensable para la vida. No basta con alegar, por ejemplo, que se tenía necesidad de asistir a la final de un mundial y por ello se vendió una propiedad muy barata. La necesidad exige una situación urgente, vital e indispensable —como cubrir una operación médica de un familiar—, no un deseo o conveniencia. \\
Debilidad psíquica & Presencia de alguna alteración mental. Reemplaza a la palabra ``ligereza'' del código anterior, que había generado dos posturas doctrinarias: obrar irreflexivo sin mayores requisitos, o expresión de una cierta debilidad psíquica. El nuevo código se pronuncia por esta última. \\
Inexperiencia & Falta del conocimiento que otorgan la vida y las relaciones. Quien se dedica profesionalmente al mercado inmobiliario durante treinta años no puede alegar inexperiencia en esa materia; distinta es la situación de una persona que se traslada de otro lugar, desconoce el ambiente y carece de formación en la materia del acto. \\
\bottomrule
\end{tabularx}
\end{table}

\begin{important}
En general se interpreta que estos tres supuestos de inferioridad constituyen una enumeración taxativa: no puede invocarse otro distinto. El proyecto de reforma del CCC del año 2018 proponía ampliar estos supuestos, pero en el derecho vigente son únicamente tres: necesidad, debilidad psíquica e inexperiencia.
\end{important}

\section{La presunción legal del art. 332}

El art.\ 332, al igual que su antecedente el art.\ 954, establece que se presume —salvo prueba en contrario— la existencia de explotación en caso de notable desproporción de las prestaciones.

\begin{important}
Esta presunción constituye una inversión de la carga de la prueba: si el lesionado prueba la notable desproporción de las prestaciones, se presume que hubo explotación, y corresponderá entonces al lesionante —demandado— probar que no existió tal explotación. Existe así una preponderancia práctica del elemento objetivo dentro de una norma que, en su configuración, es de lesión subjetiva.
\end{important}

Parte de la doctrina discute si el término ``notable'' equivale a ``evidentemente desproporcionada''. En general, se considera que lo notable es aquello que salta a la vista sin necesidad de una investigación mayor.

Existe además una discusión acerca del alcance de la presunción: una postura sostiene que solo alcanza al elemento subjetivo del lesionante, debiendo probarse por separado la necesidad, debilidad psíquica o inexperiencia del lesionado; otra postura entiende que, probada la desproporción, ambos elementos subjetivos se presumen.

\begin{note}
Esta presunción legal constituye una novedad propia del derecho argentino, sin equivalente en el código italiano, el código suizo ni el código alemán, y es la que otorga al elemento objetivo su especial preponderancia práctica.
\end{note}

\section{Las acciones, la legitimación y la prescripción}

La lesión da lugar a dos acciones:

\begin{enumerate}
    \item \textbf{Acción de nulidad}, de carácter relativo —susceptible de confirmación—, en tanto está establecida para la protección de la parte lesionada, al igual que ocurre con los restantes vicios.
    \item \textbf{Reajuste equitativo}, mediante el cual el juez corrige la desproporción en el caso concreto, determinando lo que resulte justo. A diferencia del esquema del código francés, que deducía una décima parte para el pago del complemento del precio, en el derecho argentino el reajuste queda librado a la valoración judicial, sin que la ley fije una pauta.
\end{enumerate}

\begin{important}
El reajuste equitativo prima sobre la acción de nulidad cuando es ofrecido por el demandado al contestar la demanda, en virtud del principio de conservación de los actos jurídicos.
\end{important}

La legitimación para iniciar la acción corresponde exclusivamente al lesionado o a sus herederos, tal como establecía ya el art.\ 954. El plazo de prescripción es de \textbf{dos años}, contado desde la fecha en que la obligación debía cumplirse —plazo que la ley 17.711 fijaba en cinco años.

\section{Figuras afines: cláusula penal y usura}

El art.\ 794, relativo a la cláusula penal, guarda semejanza con la lesión: los jueces pueden reducir las penas establecidas en una cláusula penal —cláusula contractual que fija sanciones para el caso de incumplimiento— cuando resulten excesivas. Se trata de una figura próxima a la lesión, en cuanto también existe desproporción susceptible de reducción judicial, aunque con una diferencia relevante: no produce nulidad, sino únicamente reajuste.

La \textbf{usura}, en cambio, constituye un delito en el derecho penal, regulado por el art.\ 175 bis del Código Penal.

\begin{formula}{Art. 175 bis — Código Penal (síntesis) — Usura}
Quien, aprovechando la necesidad, ligereza o inexperiencia de otro, se hiciera dar o prometer, en cualquier forma, para sí o para un tercero, intereses u otras ventajas pecuniarias evidentemente desproporcionadas con su prestación, u otorgare a sabiendas un crédito usurario, será reprimido.
\end{formula}

En la figura penal se requiere dolo. Debe notarse que esta norma no utiliza el término ``explotación'' sino ``aprovechar'', y exige que el autor se haga dar o prometer las ventajas: existe una acción activa dirigida a lograr esa entrega o promesa, que constituye el núcleo del delito.

\section{Debate doctrinario: la esencia de la lesión}

Existe un debate doctrinario acerca de cuál es la esencia de la lesión. Tobías, en su Tratado de Derecho Civil, identifica tres posturas posibles:

\begin{table}[H]
\centering
\caption{Posturas doctrinarias sobre la naturaleza jurídica de la lesión}
\begin{tabularx}{\textwidth}{>{\bfseries\raggedright\arraybackslash}p{0.28\textwidth} X}
\toprule
\textbf{Postura} & \textbf{Contenido} \\
\midrule
Vicio del objeto & Enfatiza el elemento objetivo: existe un problema en el objeto del acto, que resulta desproporcionado. Para Tobías, esta postura ignora la presencia de la explotación de la inferioridad. \\
Vicio de la voluntad & Enfatiza el elemento subjetivo del lesionado: por su inferioridad —necesidad, debilidad psíquica o inexperiencia— estaría viciada su voluntad. Tobías objeta que, de seguirse este criterio, resultaría irrelevante la existencia de desproporción. \\
Comportamiento reprochable (Tobías) & Es una respuesta del ordenamiento frente a un comportamiento reprochable, vinculado a la buena fe. Se reprocha la explotación. Esta postura se encuentra más cerca del antecedente alemán. \\
\bottomrule
\end{tabularx}
\end{table}

En definitiva, los tres elementos —objeto, voluntad y conducta— están en juego a la hora de determinar la esencia de la lesión, y de fondo siempre se encuentra la justicia conmutativa. El elemento objetivo resulta indudable: el objeto del acto debe ser proporcionado, y si existe desproporción, el acto no realiza plenamente la justicia. Al mismo tiempo, existe una valoración desde las buenas costumbres, que fundamenta el reproche jurídico cuya consecuencia es la nulidad.

\begin{note}
En síntesis: el vicio de lesión en el ordenamiento argentino se recorrió desde el Código de Vélez —que no lo contempló, por las razones expuestas en su nota al art.\ 943— pasando por la jurisprudencia posterior y la reforma introducida por la ley 17.711 en el art.\ 954, cuyo texto es sustancialmente reproducido por el art.\ 332 del CCC vigente. Dentro de esta figura se analizaron su fuente, su método, su ámbito de aplicación, sus elementos, las acciones a las que da lugar, la legitimación activa, el plazo de prescripción, su comparación con otras figuras afines y el problema de su esencia jurídica.
\end{note}

% ---------------------------------------------------------
% CORNELL — CAPÍTULO 2
% ---------------------------------------------------------
\section{Cuadro de repaso — Método Cornell}
\begin{longtable}{
    >{\raggedright\arraybackslash}p{0.30\textwidth}
    >{\raggedright\arraybackslash}p{0.58\textwidth}
}
\toprule
\textbf{Preguntas / palabras clave} & \textbf{Notas / respuestas} \\
\midrule
\endhead

\textbf{¿Por qué Vélez no incluyó la lesión?} &
Por dos razones: (1) la disparidad de criterios en el derecho comparado, que le impedía identificar un principio uniforme; (2) su visión liberal-individualista, según la cual, si no hubo error, dolo ni violencia, el contrato es válido y cada parte debe hacerse cargo de sus propias decisiones. \\

\textbf{¿Cómo incorporó Argentina la lesión antes de 1968?} &
La jurisprudencia —especialmente el fallo de Borda en el caso ``Orlando'' (1958)— la aplicó por vía del art.\ 953 (objeto contrario a la moral). El Tercer Congreso de Derecho Civil (1961) recomendó luego su incorporación expresa. \\

\textbf{¿Qué exige el art. 332 del CCC para configurar la lesión?} &
Un elemento objetivo: desproporción evidente y sin justificación, que exista al momento del acto y subsista al momento de la demanda. Dos elementos subjetivos: explotación por parte del lesionante, y necesidad, debilidad psíquica o inexperiencia del lesionado. \\

\textbf{¿Cómo opera la presunción del art. 332?} &
Si se prueba la notable desproporción de las prestaciones, se presume —salvo prueba en contrario— la explotación. Esto invierte la carga probatoria: es el demandado quien debe demostrar que no explotó. \\

\textbf{¿Qué acciones otorga la lesión?} &
Nulidad relativa (confirmable) o reajuste equitativo (fijado por el juez). El reajuste prima sobre la nulidad cuando es ofrecido por el demandado al contestar la demanda, en virtud del principio de conservación del acto. \\

\textbf{¿Cuál es el plazo de prescripción?} &
Dos años, contados desde que la obligación debía ser cumplida. El CCC redujo el plazo de cinco años que establecía la ley 17.711. \\

\textbf{¿Cuál es la naturaleza jurídica de la lesión según la doctrina?} &
Existe debate entre tres posturas: (1) vicio del objeto; (2) vicio de la voluntad; (3) conducta reprochable, contraria a la buena fe y a las buenas costumbres. Tobías se inclina por esta tercera postura, más próxima al modelo alemán. \\

\midrule

\multicolumn{2}{p{0.88\textwidth}}{
\textbf{Resumen:} Argentina, luego de que Vélez la excluyera por razones liberales y por la falta de un criterio uniforme en el derecho comparado, incorporó la lesión en 1968 mediante la ley 17.711, siguiendo el modelo alemán e italiano. El CCC vigente la regula en el art.\ 332 con tres elementos: uno objetivo (desproporción evidente, sin justificación, que subsista al tiempo de la demanda), uno subjetivo del lesionante (explotación activa) y uno subjetivo del lesionado (necesidad, debilidad psíquica o inexperiencia). La presunción legal de explotación ante la notable desproporción —novedad argentina sin equivalente en los códigos extranjeros— otorga preponderancia práctica al elemento objetivo. Las acciones disponibles son la nulidad relativa y el reajuste equitativo, prevaleciendo este último por el principio de conservación del acto. El plazo de prescripción es de dos años.
} \\

\bottomrule
\end{longtable}
\chapter[La simulación]{La simulación como vicio de los actos jurídicos}
\label{ch:simulacion}

Tanto la simulación como el fraude son vicios de los actos jurídicos que hacen que un acto, pese a existir, no produzca todos sus efectos normales. La diferencia entre ambos puede describirse como la que existe entre una fachada falsa y una casa vaciada por dentro: en la simulación se construye una fachada --un acto aparente-- que oculta o no oculta nada real, y por eso se ataca con la nulidad; en el fraude la casa es real, pero el deudor la vació de bienes para no pagarles a sus acreedores, y por eso se ataca con la inoponibilidad y no con la nulidad. Este capítulo trata la simulación; el siguiente, el fraude, y al final se comparan ambas figuras.

\begin{note}
Estos dos vicios explican por qué un acto que parece perfecto en el papel puede terminar siendo anulado o simplemente ignorado por la ley frente a ciertas personas. Para cualquier abogado --litigante, asesor de empresas, escribano o juez-- identificar una simulación (por ejemplo, una venta que en realidad esconde una donación para perjudicar herederos) o un fraude (un deudor que vacía su patrimonio para no pagar) constituye una herramienta central de la práctica profesional: de estos institutos depende poder recuperar un bien familiar, cobrar una deuda, defender a un cliente acusado de haber ``prestado su nombre'', o asesorar correctamente antes de firmar un contrato para no quedar involucrado, sin saberlo, en un acto simulado o fraudulento.

También resulta útil en la vida personal: son temas que aparecen en sucesiones familiares, en préstamos entre parientes, en la compra de bienes a nombre de terceros o en cualquier situación donde alguien pueda intentar aprovecharse de un acuerdo informal. Entender la diferencia entre lo lícito y lo ilícito, y saber qué prueba se necesita para desarmar un acto aparente, protege el propio patrimonio y el de la familia.

\end{note}

\section{Concepto y noción general}

La simulación es un vicio especial de los actos jurídicos, porque supone una situación en la que las partes acuerdan una declaración de voluntad o un acto jurídico que, o bien no es real, o bien es aparente y encubre otro acto verdadero.

La simulación puede tener fines lícitos o ilícitos, según las circunstancias en las que las personas involucradas la lleven a cabo.

\begin{example}{Simulación ilícita entre herederos}
Un padre o una madre --viudo o viuda-- tiene tres hijos y quiere beneficiar en realidad a uno de ellos, su hijo favorito, simulando venderle su única casa. Los otros dos hijos no se enteran del verdadero propósito del acto.

Cuando fallece la persona y los herederos inician la sucesión, se encuentran con que la casa figura vendida, pero el hijo que aparece como comprador en realidad no tenía fondos, pues atravesaba una situación económica muy mala. Ante la pregunta de los demás hermanos, queda claro que el acto tuvo que haber sido simulado: en apariencia se le vendió la casa, pero en realidad se la estaban regalando, de modo que se trató de una donación.

Frente a esto, se ataca el acto para que se declare su nulidad y se lo considere una donación, de modo que la porción de la legítima hereditaria que corresponde a los demás herederos reingrese al acervo hereditario. Este es un ejemplo de simulación \textbf{ilícita}, porque perjudica a terceros --en este caso, a los otros coherederos.
\end{example}

\begin{example}{Simulación lícita: el auto a nombre del hermano}
Una persona que se muda al exterior por tiempo indeterminado acuerda con su hermano poner el automóvil a su nombre, simulando una venta, para que este lo utilice mientras dure la ausencia, con el compromiso de restituirlo a su regreso.

La titularidad del auto queda simulada: el hermano actúa todo el tiempo como si el auto fuera suyo, y cuando la persona regresa y lo reclama, se produce la necesidad de volver las cosas al estado anterior.
\end{example}

\section{Definición doctrinaria y ubicación metodológica}

La doctrina define la simulación como la situación en la que dos personas, o una de ellas y un tercero que aparentemente no interviene en el acto, se ponen de acuerdo para declarar que celebran actos jurídicos que en realidad no se han propuesto llevar a cabo, o para aparentar un acto distinto, ocultando el acto verdadero bajo la apariencia de otro que declaran haber realizado.

\begin{definition}{Simulación}
Habrá simulación cuando dos personas, o una de ellas y un tercero que aparentemente no interviene en el acto, se ponen de acuerdo para declarar que celebran actos jurídicos que en realidad no se han propuesto llevar a cabo, o para aparentar un acto distinto, ocultando el acto verdadero bajo la apariencia de otro que declaran haber realizado.
\end{definition}

A partir de esta definición surgen los dos grandes tipos de simulación: la \textbf{absoluta}, cuando no hay nada de real, y la \textbf{relativa}, cuando la simulación encubre otro acto.

Metodológicamente, el tema se ubica en la Sección Segunda del Capítulo 6, ``De los vicios de los actos jurídicos'', del Título Cuarto (``Hechos y actos jurídicos'') del Libro Primero, Parte General, del Código Civil y Comercial (CCyC).

\section{Definición legal de la simulación}

La definición legal que trae el Código Civil y Comercial viene prácticamente igual a la que traía el código anterior en su artículo 955.

\begin{formula}{Definición legal de la simulación}
``La simulación tiene lugar cuando se encubre el carácter jurídico de un acto con la apariencia de otro, o cuando el acto contiene cláusulas que no son sinceras, o fechas que no son verdaderas, o cuando por el acto se constituyen o transmiten derechos a personas interpuestas que no son aquellas para quienes en realidad se constituyen o transmiten.''
\end{formula}

\section{Simulación absoluta y simulación relativa}

Dentro de esta clasificación se distingue, por un lado, la simulación absoluta y, por otro, la simulación relativa.

\subsection{Simulación absoluta}

Es absoluta cuando el acto nada tiene de verdadero.

\begin{example}{Simulación absoluta}
Una madre simula venderle un bien a un hijo para eludir a sus acreedores. Lo ``vende'', pero sigue viviendo allí: no hay ninguna transferencia de dinero, no hay nada real. El bien sigue siendo de ella; lo único que hace es aparentar haberlo puesto a nombre de su hijo. Esto constituye una simulación absoluta.
\end{example}

\subsection{Simulación relativa}

Es relativa aquella en que el acto aparente encubre a otro real y verdadero.

\begin{formula}{Art. 334 CCyC}
``Si el acto simulado encubre a otro real, este es plenamente eficaz, siempre que concurran los requisitos propios de esa categoría de acto, y no sea ilícito ni perjudique a un tercero.''
\end{formula}

\begin{example}{Simulación relativa}
Se simula una compraventa, pero en realidad se trata de una donación, sin que esto perjudique a terceros ni sea ilícito. Será una donación hecha por determinados fines, que aparentemente se hizo como compraventa, pero en realidad es una donación.
\end{example}

El artículo 334 deja a salvo el acto encubierto: el acto aparente se anula y queda subsistente el acto encubierto, plenamente eficaz.

\section{Simulación lícita e ilícita}

La otra clasificación relevante distingue entre simulación lícita e ilícita: la simulación \textbf{lícita} es aquella que no perjudica a terceros ni viola la ley; la simulación \textbf{ilícita} es aquella que perjudica a terceros o viola la ley.

Retomando los ejemplos ya expuestos: el caso del deudor que simula vender un bien para que sus acreedores no puedan ejecutarlo es ilícito, porque los perjudica. En cambio, el caso de la persona que pone el auto a nombre de su hermano, sin perjudicar a nadie ni hacer nada contrario a la ley --simplemente porque resulta más práctico, ya que el hermano lo va a usar y asume todas las responsabilidades del caso, con el compromiso de devolverlo--, constituye una simulación lícita.

\section{Naturaleza de la sanción: nulidad relativa}

La simulación da lugar a la nulidad, y se trata de una nulidad de tipo \textbf{relativa}, porque está establecida en beneficio de las partes.

Entre los vicios que se tratan en el Título Cuarto del Código --error, dolo, violencia (vicios de la voluntad), vicios de los actos voluntarios, lesión y simulación-- todos dan lugar a la nulidad. El fraude es el único vicio de ese título que da lugar a la inoponibilidad.

Algunos autores sostenían que, ante una simulación, se estaba frente a un acto inexistente, pero esta postura no encuentra eco en el Código Civil y Comercial, que siempre refiere a la nulidad del acto.

\begin{important}
Si al leer algún texto o algún autor se encuentra una referencia a la inexistencia del acto simulado, en realidad esto no tiene asidero: la inexistencia es una categoría de ineficacia reconocida en otros países, pero en el derecho argentino no está reconocida, salvo en algún supuesto muy específico. Por lo tanto, al hablar de simulación siempre corresponde hablar de nulidad.
\end{important}

\section{La acción de simulación entre las partes}

La acción de simulación puede ser ejercida por las partes entre sí, o por terceros.

\begin{example}{Acción entre las partes}
Una persona regresa de un viaje y reclama a su hermano la restitución del auto que le había puesto a su nombre de manera simulada, comprometiéndose este a devolvérselo a su regreso. El hermano desconoce que el acto haya sido simulado y sostiene que el auto es suyo.

Frente a ese desconocimiento, la persona debe iniciar una acción de simulación contra su hermano ante el juez, solicitando que se declare la nulidad del acto simulado, para que las cosas vuelvan al estado anterior. Este es un ejemplo de acción entre las propias partes.
\end{example}

\subsection{Procedencia de la acción entre partes}

\begin{formula}{Art. 335 CCyC (primer párrafo)}
``Los que otorgan un acto simulado ilícito, o que perjudica a terceros, no pueden ejercer acción alguna el uno contra el otro sobre la simulación, excepto que las partes no puedan obtener beneficio alguno de las resultas del acto practicado con la simulación.''
\end{formula}

Si la simulación es ilícita, la acción entre las partes no procede, por regla, salvo que ninguna de ellas pueda obtener un beneficio de la simulación. En cambio, en la simulación lícita no está presente este requisito, y la acción procede siempre.

\begin{example}{Excepción a la regla}
Una persona simuló vender su casa para eludir a sus acreedores y, luego, gracias a una herencia de un pariente lejano, obtiene dinero suficiente para pagarles. Con esa disponibilidad, inicia una acción de simulación contra el adquirente simulado de su casa, para recuperar la propiedad. Aunque la simulación fue ilícita, al dejar sin efecto el acto simulado no obtiene ningún beneficio indebido, porque ya pagó a los acreedores: estos no perdieron nada, no fueron perjudicados. Por eso, en este caso, la acción entre partes puede proceder.
\end{example}

\subsection{La prueba entre las partes: el contradocumento}

Entre las partes, la simulación debe probarse a través del \textbf{contradocumento}. El contradocumento es el instrumento donde consta la verdadera intención de las partes: si la simulación es una declaración de voluntad aparente que encubre otro acto real, o que no tenía nada de real, en el contradocumento se plasma la voluntad real, y las partes dejan en claro que se trató de un acto simulado. Durante la simulación se redacta también un contradocumento, justamente para que, ante la necesidad de probar que el acto es simulado, la parte interesada pueda exhibirlo.

En la redacción originaria de Vélez, de 1871 (vigente desde 1869), no había posibilidad de eximirse de presentar el contradocumento. Luego, la jurisprudencia fue incorporando excepciones, y en el período 1968/1969 la Ley 17.711 incorporó un párrafo al código anterior que estableció la posibilidad de prescindir del contradocumento. Esta solución fue receptada por el nuevo Código en el artículo 335, que agrega, respecto del código anterior, la exigencia de justificar las razones por las cuales el contradocumento no existe o no puede ser presentado.

\begin{formula}{Art. 335 CCyC (segundo párrafo)}
``La simulación alegada por las partes debe probarse mediante el respectivo contradocumento. Puede prescindirse de él cuando la parte justifica las razones por las cuales no existe o no puede ser presentado y median circunstancias que hacen inequívoca la simulación.''
\end{formula}

\begin{example}{Prescindencia del contradocumento}
Dos hermanos, por la cercanía y confianza entre ellos, no consideraron necesario escribir un contradocumento, o bien este se destruyó --por ejemplo, en un incendio-- y no puede ser presentado. En ese caso deben acreditarse circunstancias que hagan inequívoca la simulación: por ejemplo, demostrar que el hermano titular registral tenía bienes que en realidad seguía usando el otro hermano --el departamento, el automóvil--, de modo de dar al juez la convicción inequívoca de que efectivamente se está ante un acto simulado.
\end{example}

Complementariamente, el artículo 1028 del código anterior establecía que el contradocumento particular que altere lo expresado en un instrumento público puede invocarse entre las partes, pero no es oponible respecto de terceros.

\begin{important}
Si el acto simulado se instrumentó por escritura pública, puede utilizarse un contradocumento simplemente particular, es decir, que no sea escritura pública. Aunque la escritura pública tiene mayor fuerza probatoria que el instrumento particular, el contradocumento particular puede ser invocado entre las partes si tiene firma, ya que esta le da mayor certeza. Frente a terceros, en cambio, es inoponible. Si no tiene firma, habrá que probar las circunstancias que hacen inequívoca la simulación.
\end{important}

\section{La acción de simulación ejercida por terceros}

\subsection{Procedencia}

La acción ejercida por terceros procede cuando la simulación es ilícita. El tercero no tendría ningún interés en atacar una simulación lícita; la razón de su acción reside, precisamente, en la existencia de una simulación ilícita.

\begin{formula}{Art. 336 CCyC}
``Los terceros cuyos derechos o intereses legítimos son afectados por el acto simulado pueden demandar su nulidad.''
\end{formula}

Esta norma habilita a los terceros a iniciar una demanda de nulidad; se trata, nuevamente, de una demanda de nulidad, y no de inexistencia ni de inoponibilidad.

\subsection{La prueba por indicios}

Según el artículo 336, el tercero puede acreditar la simulación por cualquier medio de prueba: no se le exige el contradocumento, ya que este lo tienen las partes que hicieron el acto simulado, y no el tercero.

Ante la imposibilidad de recurrir al contradocumento, el tercero recurre a indicios. La doctrina, con aportes de distintos autores especializados en el tema, reúne una enumeración de indicios:

\begin{itemize}
    \item Carencia de medios económicos del adquirente para haber podido pagar el precio.
    \item El parentesco muy próximo, o la amistad íntima, entre las partes del acto.
    \item La ausencia de ejecución material del acto o contrato (por ejemplo, el vendedor que sigue viviendo en el inmueble ``vendido'').
    \item La venta de la cosa a una sociedad anónima en la que el propio enajenante o sus parientes son accionistas.
    \item La venta de un bien que el vendedor necesita para su sustento (por ejemplo, el auto con el cual trabaja y se traslada).
    \item El precio vil, es decir, un precio muy bajo en relación con el valor real de la cosa.
    \item La ausencia de incorporación del bien en las declaraciones impositivas (no haberlo denunciado ante la AFIP).
    \item Los pagos anticipados: pagar por anticipado todo el precio antes de la escrituración resulta sospechoso.
    \item La retroventa: es poco habitual que alguien compre algo con la posibilidad de que se lo revendan a él mismo.
    \item Que la enajenación no sea útil ni necesaria: no había ninguna necesidad real de vender.
\end{itemize}

Estos son indicios surgidos de la jurisprudencia, que el tercero que alega una simulación debe reunir para convencer a un juez de que lo sucedido no es real, sino que encubre otro acto, o que directamente no tiene nada de real.

\section{Efectos de la simulación frente a terceros}

La simulación ilícita que perjudica a un tercero provoca la nulidad del acto; si el acto simulado encubre a otro real, este es plenamente eficaz si concurren los requisitos propios de su categoría, y no es ilícito ni perjudica a terceros. En las simulaciones relativas por cláusulas simuladas, si es posible distinguir y separar la cláusula falsa de las demás, esta pierde valor, pero el acto sigue siendo válido en lo restante, por el principio de conservación de los actos.

\begin{formula}{Art. 337 CCyC}
``La simulación no puede oponerse a los acreedores del adquirente simulado que de buena fe hayan ejecutado los bienes comprendidos en el acto. La acción del tercero contra los subadquirentes solo procede si adquirió por título gratuito, o si es cómplice en la simulación. El subadquirente de mala fe y quien contrató de mala fe con el deudor responden solidariamente por los daños causados al tercero que ejerció la acción, si los derechos se transmitieron a un adquirente de buena fe y a título oneroso, o de otro modo se hizo imposible su recaudo.''
\end{formula}

\begin{example}{Efectos frente a terceros y subadquirentes}
Una persona simula venderle el auto a su hermano; este, a su vez, tiene sus propios acreedores, que ejecutan el auto de buena fe, sin saber nada de la simulación. En ese caso quedan protegidos, pues la lógica siempre es proteger al de buena fe y a título oneroso.

Si el hermano, a su vez, transmitió el auto a un subadquirente, la acción contra este solo procede si adquirió a título gratuito, o si fue cómplice de la simulación. Si el subadquirente es de buena fe y a título oneroso --por ejemplo, compró el vehículo mediante un aviso, sin conocer la simulación-- queda protegido. Si, en cambio, el hermano se puso de acuerdo con esa otra persona y le vendió el auto sabiendo todo, la acción procede, aunque la complicidad resulta muy difícil de demostrar.

El subadquirente de mala fe y el deudor responden por los daños: procede una acción de daños junto con la nulidad, dirigida contra el subadquirente de mala fe. En cambio, si el subadquirente es de buena fe pero a título gratuito --por ejemplo, si el bien fue donado a una institución de beneficencia-- no responde por los daños, sino en la medida en que se haya enriquecido, debiendo restituir en esa medida.
\end{example}

La regla general que atraviesa este tema es que siempre se protege a quien actuó de buena fe y a título oneroso; quien actuó de mala fe responde por los daños, y quien actuó de buena fe pero a título gratuito responde solo en la medida de su enriquecimiento. Se trata de un principio fundamental en todo el tema de los actos jurídicos dentro del sistema de derecho civil argentino.

\section{Prescripción de la acción de simulación}

\begin{formula}{Art. 2563, inc. b, CCyC}
El plazo de prescripción de dos años se computa, en la acción de simulación ejercida por parte, ``desde que, requerida una de ellas, se negó a dejar sin efecto el acto simulado''; y en la acción de simulación que ejerce un tercero, desde que conoció o pudo conocer el vicio del acto jurídico.
\end{formula}

La acción entre partes no comienza a correr desde que se celebró el acto simulado, sino desde que una de ellas desconoce, en los hechos, la simulación.

\begin{example}{Cómputo del plazo entre partes}
Una persona regresa de España, reclama el auto y su hermano le dice que el auto es suyo. Aunque hayan pasado varios años desde su partida, el plazo de prescripción empieza a correr recién cuando regresa y el hermano se niega a devolvérselo, y no desde el momento de la partida.
\end{example}

La acción ejercida por un tercero prescribe a los dos años desde que este conoció o pudo conocer el vicio del acto jurídico.

% -----------------------------------------------------------
% CORNELL --- CAPÍTULO 1
% -----------------------------------------------------------
\section{Repaso Cornell: la simulación}
\begin{table}[H]
\centering
\begin{tabularx}{\textwidth}{
    >{\raggedright\arraybackslash}p{0.28\textwidth}
    >{\raggedright\arraybackslash}X
}
\toprule
\textbf{Preguntas / palabras clave} &
\textbf{Notas / respuestas} \\
\midrule

\textbf{¿Qué es la simulación?} &
Un vicio en el que las partes acuerdan una declaración de voluntad que no es real, o que es aparente y encubre otro acto. Puede tener fines lícitos (por ejemplo, poner un auto a nombre de un hermano) o ilícitos (por ejemplo, simular una venta para perjudicar a coherederos). \\

\textbf{¿Simulación absoluta o relativa?} &
Absoluta: el acto no tiene nada de real (por ejemplo, una madre que ``vende'' un bien pero sigue siendo suya). Relativa: el acto aparente encubre a otro real y válido (art. 334 CCyC), que subsiste plenamente eficaz si no es ilícito ni perjudica a terceros. \\

\textbf{¿Simulación lícita o ilícita?} &
Lícita: no perjudica a terceros ni viola la ley. Ilícita: perjudica a terceros o viola la ley (por ejemplo, vender bienes para eludir acreedores). \\

\textbf{¿Qué sanción tiene la simulación?} &
La nulidad, de carácter relativo (no la inexistencia, categoría que el Código Civil y Comercial no reconoce para estos casos). \\

\textbf{¿Cómo se prueba entre las partes?} &
Mediante el contradocumento (art. 335 CCyC); puede prescindirse de él si se justifican las razones de su ausencia y median circunstancias que hacen inequívoca la simulación. \\

\textbf{¿Cómo prueba un tercero?} &
Por cualquier medio de prueba (art. 336 CCyC), típicamente a través de indicios: falta de medios económicos del comprador, parentesco cercano, precio vil, ausencia de ejecución material del contrato, pagos anticipados, retroventa, entre otros. \\

\textbf{¿Quién queda protegido frente a terceros?} &
Los acreedores de buena fe del adquirente simulado que hayan ejecutado bienes (art. 337 CCyC). Los subadquirentes de mala fe o a título gratuito responden; los de buena fe y título oneroso quedan protegidos. \\

\textbf{¿Cuándo prescribe la acción de simulación?} &
A los dos años (art. 2563 inc. b CCyC): entre partes, desde que una de ellas se negó a dejar sin efecto el acto; por terceros, desde que conocieron o pudieron conocer el vicio. \\

\midrule

\multicolumn{2}{p{\dimexpr\textwidth-2\tabcolsep\relax}}{
\textbf{Resumen:} La simulación y el fraude son los dos últimos vicios de los actos jurídicos estudiados en esta parte, y ambos responden a la misma preocupación: qué hacer cuando un acto jurídico, sin ser lo que aparenta, afecta a personas ajenas a él. En la simulación, las partes se ponen de acuerdo para declarar algo que no es real o que oculta otro acto; según perjudique o no a terceros o viole la ley, será ilícita o lícita, y siempre se sanciona con la nulidad relativa del acto ostensible, probándose entre las partes mediante contradocumento y, por los terceros, mediante cualquier indicio idóneo.
} \\

\bottomrule
\end{tabularx}
\end{table}

% ===========================================================
% CAPÍTULO 2 --- EL FRAUDE
% ===========================================================
\chapter[El fraude]{El fraude como vicio de los actos jurídicos}

\section{Introducción y distinción con el fraude penal}

La palabra ``fraude'' no debe confundirse con el delito de fraude, si bien puede darse que la misma conducta configure, a la vez, un fraude civil y un delito. En este capítulo se consideran los aspectos civiles del fraude, vinculados con la garantía del patrimonio a favor de los acreedores.

En la concepción integral del derecho civil y del Código Civil y Comercial, el patrimonio opera como garantía de los acreedores. Por eso, si un deudor enajena bienes para sustraerlos de sus acreedores, comete un fraude: elude, por así decirlo, la acción de sus acreedores. El código da una respuesta a esta conducta a través de determinadas disposiciones.

\section{Definición doctrinaria}

\begin{definition}{Fraude}
El fraude tiene lugar cuando una persona, en forma intencional, se desapodera de bienes o poderes para colocarse en una situación de insolvencia y eludir de esta manera la acción de sus acreedores.
\end{definition}

La insolvencia expresa una situación en la que el pasivo supera al activo, y en la que, por lo tanto, el deudor no tiene bienes suficientes para hacer frente a sus deudas.

\begin{example}{Insolvencia y dinero en efectivo}
Un deudor vende su casa para evitar que se la ejecuten. Si conserva el dinero de la venta en efectivo --en una caja fuerte, en un bolso, o de cualquier otra forma-- resulta muy difícil de embargar, salvo que en algún mandamiento se logre identificarlo. Así, la persona aparece como sin bienes, pero los acreedores pueden descubrir que hubo un acto de enajenación y atacarlo.
\end{example}

\section{La inoponibilidad como consecuencia del fraude}

El fraude no da lugar a una nulidad, sino a la \textbf{inoponibilidad}. La acción que corresponde al fraude se denomina acción de inoponibilidad en el nuevo Código; históricamente se la conocía como acción pauliana o acción revocatoria, nombre que usaba el código de Vélez. El nuevo Código, siguiendo lo que ya aceptaba la mayoría de la doctrina antes de su entrada en vigencia, habla directamente de inoponibilidad.

El acto es válido entre las partes, pero inoponible: tiene una ineficacia relativa. Esta ineficacia consiste en que el acto no produce efectos respecto de ciertas personas determinadas, que son los acreedores defraudados que inician la acción. Frente a ellos, el acto se comporta como si no se hubiera celebrado, es decir, es inoponible; pero frente a las propias partes, y frente a cualquier otra persona que no sea el acreedor accionante, el acto sigue siendo válido.

\begin{important}
El concepto de inoponibilidad --y no de nulidad-- es fundamental para diferenciar al fraude de los demás vicios de los actos jurídicos.
\end{important}

\section{La acción de inoponibilidad}

Metodológicamente, el fraude está tratado en la Sección Tercera del Capítulo 6 (``De los vicios de los actos jurídicos''), del Título Cuarto, del Libro Primero, Parte General.

\begin{formula}{Art. 338 CCyC}
``Todo acreedor puede solicitar la declaración de inoponibilidad de los actos celebrados por su deudor en fraude de sus derechos, y de las renuncias al ejercicio de derechos o facultades con los que hubiese podido mejorar o evitado empeorar su estado de fortuna.''
\end{formula}

El acreedor puede pedir la declaración de inoponibilidad tanto de los actos hechos en fraude, como de las renuncias hechas para eludir la acción de los acreedores.

\begin{example}{Renuncia en fraude a los acreedores}
Quien renuncia a cobrar una indemnización o un crédito perjudica a sus acreedores, porque estos se ven afectados por esa renuncia. Los acreedores pueden alegar que esa renuncia fue hecha en fraude a sus derechos.
\end{example}

\begin{note}
Esta inoponibilidad, como supuesto de ineficacia relativa, debe complementarse con los artículos 396 y 397 del Código, que regulan la inoponibilidad en general y que se estudian al abordar las ineficacias.
\end{note}

\section{Requisitos de procedencia de la acción}

\subsection{Crédito de causa anterior al acto impugnado}

El primer requisito es que el crédito sea de causa anterior al acto impugnado, salvo que el deudor haya actuado con el propósito de defraudar a futuros acreedores.

\begin{example}{Crédito anterior al acto impugnado}
Un deudor tiene una deuda desde el 1 de enero de 2018; a partir de esa fecha, su acreedor tiene la expectativa de que todo su patrimonio constituye la garantía del crédito. El acreedor no podría impugnar un acto realizado en 2017, porque en ese momento no existía la deuda. Si al 1 de enero de 2018 el deudor tenía un departamento y un auto, y durante 2018 enajena el departamento y se vuelve insolvente --porque además el auto se destruye en un choque--, el acreedor puede atacar la venta del departamento, porque su crédito es anterior (enero de 2018) y la venta es posterior (mayo). En cambio, si antes tenía además una cochera que vendió durante 2017, el acreedor no puede atacar esa venta, porque su crédito es posterior a ella, salvo que demuestre que la venta se hizo a sabiendas, para defraudar a los futuros acreedores.

Este sería el caso de un deudor que tiene varios inmuebles y, mientras un acreedor tramita un crédito que tarda varios meses en salir, vende todos sus inmuebles; cuando el crédito finalmente se otorga, el deudor ya no tiene nada, pero lo hizo a sabiendas, en fraude de los futuros acreedores.
\end{example}

\subsection{Que el acto haya causado o agravado la insolvencia}

El segundo requisito es que el acto haya causado o agravado la insolvencia.

\begin{example}{Insolvencia causada o agravada}
Un deudor que tiene cinco departamentos y vende uno, conservando otros cuatro con los cuales puede responder perfectamente, no incurre en fraude: este supone una situación de incapacidad e impotencia para atender los créditos.
\end{example}

\subsection{La complicidad del adquirente a título oneroso}

El tercer requisito exige distinguir según el acto haya sido a título oneroso o a título gratuito. Si es a título oneroso, debe probarse que quien contrató con el deudor conocía, o debía conocer, que el acto provocaba o agravaba la insolvencia, es decir, que fue cómplice porque conocía la situación. Si el acto fue a título gratuito, no se exige este requisito de complicidad, y la acción procede siempre, porque no hay nadie que se haya enriquecido en perjuicio de acreedores que quedaron desprovistos del bien.

\begin{example}{Complicidad del adquirente}
Un contador que conoce perfectamente la situación contable de su cliente le compra un departamento: en ese caso corresponde analizar si conocía o debía conocer el estado de insolvencia del vendedor. La complicidad es, en la práctica, difícil de probar.
\end{example}

\section{Efectos de la acción de inoponibilidad}

El acto es válido entre las partes, pero sus efectos son inoponibles al acreedor defraudado que obtiene la sentencia a su favor. A diferencia de la nulidad, que vuelve las cosas al estado anterior, en el fraude el acto sigue siendo válido entre las partes; lo único que ocurre es que, al ejecutarse la cosa, el acreedor cobra su crédito, y el excedente queda para quien había adquirido el bien.

\begin{example}{Efectos de la inoponibilidad}
Una persona vende un departamento a un pariente cercano en fraude de sus acreedores. Se prueba el fraude; la deuda era de 40.000 dólares y el departamento vale 100.000. Se ejecuta el bien, se obtienen 100.000 dólares: los acreedores cobran su parte, que son 40.000, y los 60.000 restantes no son de ellos, sino de quien había adquirido el inmueble, que se queda con ese excedente. El adquirente tiene, a su vez, un reclamo contra el deudor por lo que pagó, aunque probablemente no tenga de dónde cobrarse.

Esta es una diferencia importante entre el fraude y la simulación: en la simulación, si se declara la nulidad, las cosas vuelven al estado anterior; en el fraude, si al ejecutar la cosa hay un excedente, este se entrega a quien había adquirido la cosa, porque el acto vale entre las partes.
\end{example}

\section{El desinterés del adquirente}

\begin{formula}{Art. 341 CCyC}
El adquirente puede desinteresar a los acreedores que promovieron la acción, ofreciendo una garantía suficiente, o pagando directamente el crédito de quienes se hubiesen presentado.
\end{formula}

\begin{example}{Desinterés del adquirente}
Un adquirente que se encuentra muy conforme con el departamento adquirido --por su cercanía al trabajo y al colegio de sus hijos-- y a quien se le inicia la acción por haber sido cómplice, ofrece una garantía o paga directamente a los acreedores. En ese caso, la acción de inoponibilidad se detiene, porque los acreedores quedan desinteresados.
\end{example}

\section{Extensión de la inoponibilidad}

\begin{formula}{Art. 342 CCyC}
La declaración de inoponibilidad beneficia solo a los acreedores que promovieron la acción, y hasta el importe de sus respectivos créditos.
\end{formula}

Solo el acreedor que promueve la acción se beneficia con ella; si existen otros acreedores que no la iniciaron, no pueden beneficiarse. Retomando el ejemplo anterior: si el acreedor accionante tiene un crédito de 40.000 dólares y el departamento vale 100.000, se beneficia solo con los 40.000; los 60.000 restantes quedan para quien hubiera adquirido la cosa.

\section{Efectos frente a subadquirentes}

Este punto resulta muy similar al analizado para la simulación. Si el adquirente tiene, a su vez, sus propios acreedores, que embargan y ejecutan de buena fe, estos quedan protegidos: en el fraude, tampoco puede oponérsele al acreedor de buena fe que haya ejecutado los bienes comprendidos en el acto.

Si existe un subadquirente, corresponde analizar si es o no cómplice del fraude, y si adquirió a título oneroso o gratuito. Si es de mala fe, o si adquirió a título gratuito, en cualquiera de los dos casos procede la acción. Si el subadquirente es de buena fe y a título oneroso, queda protegido, y la acción se dirige, en forma de daños y perjuicios, contra el cómplice y el deudor que celebró el acto en fraude.

El criterio para determinar la complicidad es si el sujeto conocía el estado de insolvencia y que el acto lo agravaba. Se aplica también la regla de la responsabilidad solidaria por daños y perjuicios para quien actuó de mala fe, mientras que quien actuó de buena fe y a título gratuito responde solo en la medida de su enriquecimiento, es decir, de lo que deba restituir al acreedor defraudado.

\section{Prescripción de la acción de inoponibilidad}

El plazo de prescripción es de dos años, y se cuenta desde que el acreedor conoció o pudo conocer el vicio del acto jurídico.

\begin{important}
Con el Código de Vélez, la lesión prescribía a los cinco años y el fraude al año; el nuevo Código bajó el plazo de la lesión a dos años y elevó el del fraude también a dos años, unificando así el plazo de prescripción de todos los vicios: error, dolo, violencia, lesión, simulación y fraude prescriben, todos, a los dos años, aunque cambia el momento desde el cual comienza a correr ese plazo en cada caso.
\end{important}

A diferencia de la simulación y de los demás vicios, en el caso del fraude no se trata de una acción de nulidad, sino de una acción de inoponibilidad.

% -----------------------------------------------------------
% CORNELL --- CAPÍTULO 2
% -----------------------------------------------------------
\section{Repaso Cornell: el fraude}
\begin{table}[H]
\centering
\begin{tabularx}{\textwidth}{
    >{\raggedright\arraybackslash}p{0.28\textwidth}
    >{\raggedright\arraybackslash}X
}
\toprule
\textbf{Preguntas / palabras clave} &
\textbf{Notas / respuestas} \\
\midrule

\textbf{¿Qué es el fraude como vicio civil?} &
El acto por el cual el deudor, en forma intencional, se desapodera de bienes para colocarse en insolvencia y eludir a sus acreedores. No debe confundirse con el delito penal de fraude. \\

\textbf{¿Qué sanción tiene el fraude?} &
La inoponibilidad (ineficacia relativa), no la nulidad: el acto es válido entre las partes, pero no produce efectos frente al acreedor defraudado que obtuvo la sentencia a su favor (art. 338 CCyC). \\

\textbf{¿Qué requisitos exige la acción de inoponibilidad?} &
1) Crédito de causa anterior al acto (salvo fraude a futuros acreedores); 2) que el acto haya causado o agravado la insolvencia; 3) si es a título oneroso, que el adquirente conociera o debiera conocer ese estado (complicidad). \\

\textbf{¿Qué ocurre si hay un excedente al ejecutar el bien?} &
El excedente queda para el adquirente, no para el acreedor: a diferencia de la nulidad, el fraude no retrotrae el acto por completo, solo lo hace inoponible en la medida del crédito (arts. 341 y 342 CCyC). \\

\textbf{¿Cómo se desinteresa el adquirente?} &
Ofreciendo garantía suficiente o pagando directamente a los acreedores que promovieron la acción (art. 341 CCyC), lo que detiene la acción. \\

\textbf{¿Cuándo prescribe la acción de inoponibilidad?} &
A los dos años desde que el acreedor conoció o pudo conocer el vicio del acto (antes, con el Código de Vélez, era un año). \\

\midrule

\multicolumn{2}{p{\dimexpr\textwidth-2\tabcolsep\relax}}{
\textbf{Resumen:} En el fraude, a diferencia de la simulación, el acto es plenamente real: el deudor efectivamente enajena un bien, pero lo hace para colocarse en insolvencia y eludir a sus acreedores; por eso la sanción no es la nulidad, sino la inoponibilidad, que deja vigente el acto entre las partes y solo lo desconoce frente al acreedor que promovió la acción, protegiendo siempre al adquirente de buena fe y a título oneroso.
} \\

\bottomrule
\end{tabularx}
\end{table}

% ===========================================================
% CAPÍTULO 3 --- SÍNTESIS COMPARATIVA
% ===========================================================

\section[Síntesis comparativa]{Comparación entre simulación, fraude y acción subrogatoria}

Para cerrar el estudio de los vicios de los actos jurídicos, corresponde comparar la acción de simulación, la acción por fraude y la acción subrogatoria (esta última no es, en rigor, un vicio, pero guarda cierto parentesco con las anteriores).

\begin{table}[H]
\centering
\caption{Simulación, fraude y acción subrogatoria}
\begin{tabularx}{\textwidth}{
    >{\raggedright\arraybackslash}p{0.20\textwidth}
    >{\raggedright\arraybackslash}X
    >{\raggedright\arraybackslash}X
    >{\raggedright\arraybackslash}X
}
\toprule
\textbf{Criterio} & \textbf{Simulación} & \textbf{Fraude} & \textbf{Acción subrogatoria} \\
\midrule
¿A quién beneficia? & A todos los acreedores (vuelve el bien al patrimonio del deudor) & Solo al acreedor que la ejerce & A todos los acreedores \\
\addlinespace
¿En nombre de quién se ejerce? & En derecho propio del acreedor & En derecho propio del acreedor & Por el derecho del deudor (por ejemplo, cobrar una herencia) \\
\addlinespace
¿Requiere insolvencia del deudor? & No la requiere & Sí la requiere & No la requiere; requiere inacción del deudor \\
\addlinespace
Acción que persigue & La nulidad del acto & La inoponibilidad del acto & El fin propio del derecho que se ejerce (por ejemplo, la herencia) \\
\addlinespace
¿Quién puede ejercerla? & Cualquier persona con interés (por ejemplo, herederos entre sí) & El acreedor con interés & El acreedor con interés \\
\bottomrule
\end{tabularx}
\end{table}

Con esta comparación se cierra el estudio de los vicios de los actos jurídicos. Lo que sigue en las próximas unidades son los supuestos de ineficacia y, sobre todo, el estudio de la nulidad como el gran supuesto de ineficacia civil.

% -----------------------------------------------------------
% CORNELL --- CAPÍTULO 3
% -----------------------------------------------------------
\subsection*{Repaso Cornell: síntesis comparativa}
\begin{table}[H]
\centering
\begin{tabularx}{\textwidth}{
    >{\raggedright\arraybackslash}p{0.28\textwidth}
    >{\raggedright\arraybackslash}X
}
\toprule
\textbf{Preguntas / palabras clave} &
\textbf{Notas / respuestas} \\
\midrule

\textbf{¿En qué se diferencian simulación, fraude y acción subrogatoria?} &
La simulación persigue la nulidad y beneficia a todos los acreedores; el fraude persigue la inoponibilidad, beneficia solo al acreedor accionante y requiere insolvencia; la acción subrogatoria no es un vicio, se ejerce por el derecho del deudor, no requiere insolvencia sino inacción, y beneficia a todos los acreedores. \\

\midrule

\multicolumn{2}{p{\dimexpr\textwidth-2\tabcolsep\relax}}{
\textbf{Resumen:} Ambas figuras, junto con la acción subrogatoria, integran el sistema de protección del patrimonio como garantía común de los acreedores, y sientan las bases conceptuales --nulidad, inoponibilidad, buena y mala fe, título oneroso y gratuito-- aplicables a la ineficacia de los actos jurídicos.
% TODO: verificar consistencia entre fuentes originales: el material anuncia el estudio posterior de la nulidad como supuesto general de ineficacia, pero ese desarrollo no figura en los archivos recibidos.

} \\

\bottomrule
\end{tabularx}
\end{table}
\partintro{Así como existen hechos y actos que dan nacimiento a una relación jurídica, existen también hechos y actos que la extinguen. Esta parte estudia el cierre del ciclo de la relación jurídica por tres vías: los hechos ajenos a la voluntad de las partes, los actos jurídicos extintivos y el transcurso del tiempo (prescripción y caducidad).}
\part{La extinción de las relaciones jurídicas}
\chapter[Introducción a la extinción de las relaciones jurídicas]{Introducción: el cierre del ciclo de la relación jurídica}
%==============================================================================

\section{Objeto del tema}

El estudio de la parte general del derecho civil comprende el sujeto, el objeto y la causa de las relaciones jurídicas. Dentro de la causa se ubican los \textbf{hechos jurídicos}, que dan inicio a las relaciones y producen consecuencias jurídicas. Del mismo modo en que existen hechos y actos que dan nacimiento a una relación jurídica, existen también hechos y actos que, en algún momento, la \textbf{extinguen}. Ese es el objeto de esta parte.

\begin{note}
El desarrollo del tema es introductorio: la extinción de las obligaciones se trata con mayor profundidad en \textbf{Obligaciones}, y la extinción de los contratos, en \textbf{Contratos}.
\end{note}

En materia de hechos jurídicos extintivos se sigue la doctrina de Antonio Aguiar Anzorena.

\section{Relevancia práctica}

La prescripción y la caducidad suelen tratarse al final del estudio de la materia y no siempre se estudian con detenimiento. Sin embargo, tienen una relevancia decisiva en el ejercicio profesional: todo abogado recibirá clientes con una deuda para cobrar, y lo primero que deberá verificar es si la acción sigue vigente o ya prescribió, porque de ello depende toda la estrategia del caso.

Conocer cuándo un contrato puede resolverse, rescindirse o revocarse permite asesorar con precisión en negociaciones, redactar cláusulas de salida seguras y anticipar riesgos. Esta lógica es útil también en la vida cotidiana: saber que las deudas y los reclamos tienen plazos permite tomar decisiones a tiempo (enviar un reclamo antes de que se pierda el derecho, comprender por qué una obligación antigua ya no puede exigirse judicialmente, o entender los efectos de un arrepentimiento en una seña).

\section{Las tres vías de extinción}

Un contrato de alquiler ilustra los tres bloques temáticos.

\begin{example}{El contrato de alquiler}
El contrato nace como relación jurídica. Mientras existe, puede extinguirse:
\begin{itemize}
    \item por un \textbf{hecho ajeno a la voluntad} de las partes, por ejemplo si fallece el inquilino y era el único con derecho a habitar el inmueble (\textbf{hechos extintivos});
    \item por un \textbf{acto de voluntad} de las partes, como cuando el inquilino paga todos los meses hasta el final o decide rescindir el contrato antes de tiempo (\textbf{actos extintivos});
    \item por el \textbf{transcurso del tiempo}: si el propietario nunca reclama una deuda de alquileres atrasados, la acción para cobrarla puede prescribir, y si no ejerce ciertos derechos puntuales dentro del plazo que fija la ley, esos derechos caducan (\textbf{prescripción y caducidad}).
\end{itemize}
\end{example}

Los tres bloques constituyen las tres formas en que el derecho explica el final de una misma historia: la de una relación jurídica que nace, vive y en algún momento se apaga.

\section{Enumeración de la doctrina seguida}

\textbf{Hechos jurídicos extintivos:}
\begin{itemize}
    \item la muerte;
    \item la imposibilidad (sobrevenida);
    \item la confusión;
    \item la compensación;
    \item la caducidad.
\end{itemize}

\textbf{Actos jurídicos extintivos:}
\begin{itemize}
    \item el pago;
    \item la dación en pago;
    \item la novación;
    \item la transacción;
    \item la renuncia a los derechos;
    \item la remisión de deuda;
    \item y, como los tres principales: la resolución, la rescisión y la revocación.
\end{itemize}
\chapter[Hechos jurídicos extintivos]{Hechos jurídicos extintivos}
%==============================================================================

%------------------------------------------------------------------------------
\section{La muerte como hecho extintivo}

Los efectos jurídicos de la muerte se desarrollan en el capítulo~\ref{ch:muerte}; aquí interesa su carácter de hecho que extingue relaciones jurídicas.

%------------------------------------------------------------------------------

La muerte produce la extinción de la persona: finaliza la existencia de la persona humana. Como consecuencia, muchas relaciones jurídicas se extinguen, especialmente los \textbf{derechos extrapatrimoniales}.

También se extinguen con la muerte algunos derechos patrimoniales, como el \textbf{usufructo}: es un derecho que se constituye de por vida y, por lo tanto, se extingue con la muerte de su titular, sin transmitirse a los herederos.

\begin{important}
\textbf{El usufructo como excepción a la regla.} En materia de derechos patrimoniales, la regla general es que, a la muerte de una persona, su patrimonio se transmite a los herederos, quienes pueden ser instituidos por testamento o por disposición de la ley, existiendo algunos herederos forzosos con derecho a una porción legítima. Sin embargo, esta regla de la transmisión no rige en todos los casos: en algunos supuestos la muerte extingue directamente la relación jurídica, sin transmitirla a nadie. Uno de ellos es el usufructo.
\end{important}

%------------------------------------------------------------------------------
\section{La imposibilidad sobrevenida}
%------------------------------------------------------------------------------

El segundo hecho extintivo es la \textbf{imposibilidad}. Debe tratarse de una imposibilidad \textbf{sobreviniente} y no inicial: si la imposibilidad existiera desde el comienzo, el acto jurídico no podría formarse válidamente, porque su objeto sería imposible.

\begin{articulo}{Artículo 279 del Código Civil y Comercial}
Los actos jurídicos no pueden tener por objeto un hecho imposible.
\end{articulo}

El fundamento de esta regla se expresa en el adagio latino según el cual \textit{a lo imposible nadie está obligado}: lo imposible no puede exigirse.

La imposibilidad sobrevenida, en cambio, se produce una vez que la relación jurídica ya comenzó. Se ubica en el título de las obligaciones del Código Civil y Comercial, en el Libro Tercero.

\begin{articulo}{Artículo 955 del Código Civil y Comercial}
La imposibilidad sobrevenida, objetiva, absoluta y definitiva de la prestación, producida por caso fortuito o fuerza mayor, extingue la obligación, sin responsabilidad.
\end{articulo}

Se trata, entonces, de una extinción \textbf{sin responsabilidad}, siempre que la imposibilidad sea absolutamente imprevisible (caso fortuito o fuerza mayor). En cambio, si la obligación se vuelve imposible de cumplir por causas imputables al deudor (por ejemplo, si el deudor no pagó por su propia culpa), el deudor debe responder igualmente.

%------------------------------------------------------------------------------
\section{La confusión}
%------------------------------------------------------------------------------

\begin{definition}{Confusión}
Hecho extintivo que consiste en que la calidad de \textbf{acreedor} y la de \textbf{deudor} se reúnen en una misma persona y en un mismo patrimonio.
\end{definition}

\begin{articulo}{Artículos 931 y 932 del Código Civil y Comercial}
Regulan la extinción de la obligación por confusión, cuando las calidades de acreedor y de deudor se reúnen en una misma persona y en un mismo patrimonio.
\end{articulo}

\begin{example}{Juana y su tía}
Juana le presta 20.000 dólares a su tía. La tía fallece y Juana es su única heredera. Al heredar, Juana ocupa el lugar de su tía y se convierte, al mismo tiempo, en acreedora y deudora de la misma deuda, porque la tía le debía a ella. La obligación queda confundida y se extingue por confusión.
\end{example}

La confusión extingue la obligación en la proporción en que efectivamente se produce.

%------------------------------------------------------------------------------
\section{La compensación}
%------------------------------------------------------------------------------

\begin{definition}{Compensación}
Tiene lugar cuando dos personas reúnen, recíprocamente, la calidad de acreedor y de deudor la una de la otra. Sus créditos se compensan hasta el monto de la menor deuda.
\end{definition}

\begin{example}{Compensación de créditos recíprocos}
Si yo te debo 10.000 y vos me debés 20.000, a partir de ese momento se compensan: vos me debés 10.000, porque 20.000 menos 10.000 es 10.000. Se extinguen ambas deudas, con fuerza de pago, hasta el monto de la menor, y desde el momento en que ambas obligaciones comenzaron a coexistir en condiciones de ser compensables.
\end{example}

Se trata de un \textit{hecho} extintivo y no de un acto porque esta situación puede producirse automáticamente, sin que medie una declaración de voluntad expresa de las partes: en ese caso, la compensación opera \textbf{de pleno derecho}.

%------------------------------------------------------------------------------
\section{La caducidad (introducción)}
%------------------------------------------------------------------------------

\begin{definition}{Caducidad}
Pérdida de un derecho por no ejercitarlo en el plazo fijado por la ley. Siempre tiene \textbf{fuente legal}.
\end{definition}

Este instituto se desarrolla en profundidad junto con la prescripción (véase el capítulo~\ref{cap:prescripcion}). Desde ahora cabe señalar que la caducidad extingue el derecho no ejercido y que, a diferencia de la prescripción, \textbf{no se suspende ni se interrumpe}. Existen algunos supuestos, establecidos por la ley, en los que no corre la caducidad.

\begin{important}
No debe confundirse la \textbf{caducidad}, que extingue el derecho no ejercido, con la \textbf{prescripción liberatoria o extintiva}, que extingue la acción.
\end{important}

La caducidad rige para algunos derechos muy puntuales respecto de los cuales el Código o las leyes especiales establecen que, transcurrido cierto tiempo (generalmente plazos cortos) sin que se ejerza el derecho, este caduca.

\begin{example}{La compensación económica en el divorcio}
El derecho a exigir la compensación económica luego de un divorcio está regulado en los artículos 441 y 442 del Código Civil y Comercial. Conforme al artículo 442, la acción para reclamarla caduca a los seis meses de haberse dictado la sentencia de divorcio.
\end{example}

\begin{articulo}{Artículo 442 del Código Civil y Comercial}
La acción para reclamar la compensación económica caduca a los seis meses de haberse dictado la sentencia de divorcio.
\end{articulo}

%------------------------------------------------------------------------------
\section{Cuadro Cornell: hechos jurídicos extintivos}
%------------------------------------------------------------------------------

{\renewcommand{\arraystretch}{1.2}
\begin{longtable}{
    >{\raggedright\arraybackslash}p{\cornellL}
    >{\raggedright\arraybackslash}p{\cornellR}
}
\toprule
\textbf{Preguntas / palabras clave} &
\textbf{Notas / respuestas} \\
\midrule
\endfirsthead

\toprule
\textbf{Preguntas / palabras clave} &
\textbf{Notas / respuestas} \\
\midrule
\endhead

\textbf{¿Qué hechos extinguen relaciones jurídicas sin mediar un acto de voluntad?} &
La muerte, la imposibilidad sobrevenida, la confusión, la compensación y la caducidad. \\
\addlinespace

\textbf{¿Por qué el usufructo se extingue con la muerte de su titular?} &
Porque es un derecho constituido de por vida y no se transmite a los herederos, aunque la regla general sea que el patrimonio sí se transmite. \\
\addlinespace

\textbf{¿Qué distingue a la imposibilidad sobrevenida de la imposibilidad inicial?} &
La inicial impide que el acto se forme válidamente, porque el objeto sería imposible (art.~279 CCyC). La sobrevenida ocurre luego de nacida la obligación; si es objetiva, absoluta y definitiva, producida por caso fortuito o fuerza mayor, extingue la obligación sin responsabilidad del deudor (art.~955 CCyC). Si se debe a causas imputables al deudor, este responde igualmente. \\
\addlinespace

\textbf{¿Qué es la confusión?} &
La reunión de las calidades de acreedor y deudor en una misma persona y en un mismo patrimonio (arts.~931 y 932 CCyC), como en el caso de Juana, que hereda a su tía acreedora. \\
\addlinespace

\textbf{¿Qué es la compensación y por qué es un hecho extintivo?} &
La extinción recíproca de créditos entre dos personas que son acreedoras y deudoras entre sí, hasta el monto de la menor deuda. Es un hecho porque puede operar de pleno derecho, sin declaración de voluntad de las partes. \\
\addlinespace

\textbf{¿Qué es la caducidad y en qué se diferencia de la prescripción?} &
La pérdida de un derecho por no ejercerlo dentro del plazo fijado por la ley; siempre tiene fuente legal. Extingue el derecho no ejercido, mientras que la prescripción extintiva extingue la acción. No se suspende ni se interrumpe. Ejemplo: art.~442 CCyC (seis meses desde la sentencia de divorcio). \\

\midrule
\multicolumn{2}{p{\dimexpr\cornellL+\cornellR+2\tabcolsep\relax}}{%
\textbf{Resumen:} los hechos extintivos operan sin necesidad de un acto de voluntad: la muerte (que extingue derechos extrapatrimoniales y algunos patrimoniales, como el usufructo), la imposibilidad sobrevenida (sin responsabilidad si obedece a caso fortuito o fuerza mayor), la confusión (acreedor y deudor en una misma persona), la compensación (créditos recíprocos hasta el monto de la menor deuda) y la caducidad (pérdida del derecho no ejercido en el plazo legal).} \\
\bottomrule
\end{longtable}}

%==============================================================================
\chapter[Actos jurídicos extintivos]{Actos jurídicos extintivos}
%==============================================================================

Los actos jurídicos extintivos son aquellos que suponen una \textbf{declaración de voluntad} dirigida a poner fin a la relación jurídica.

%------------------------------------------------------------------------------
\section{El pago}
%------------------------------------------------------------------------------

El acto extintivo de una obligación por excelencia es el pago.

\begin{definition}{Pago}
Cumplimiento de la prestación que constituye el objeto de la obligación.
\end{definition}

\begin{important}
No debe asociarse el pago sólo con la entrega de dinero en efectivo: \textbf{pagar es cumplir con la obligación}.
\end{important}

\begin{example}{El pago no es sólo entregar dinero}
Si yo te presto una cosa mueble y vos te comprometés a darme plata a cambio, cuando yo te entrego la cosa pactada estoy pagando, porque cumplo mi obligación; y cuando vos me entregás el dinero, vos también estás pagando. Del mismo modo, si a una persona la contratan para hacer una asesoría de servicios jurídicos, esa persona paga cuando cumple con la prestación de hacer esa asesoría.
\end{example}

%------------------------------------------------------------------------------
\section{La dación en pago}
%------------------------------------------------------------------------------

\begin{definition}{Dación en pago}
Acto extintivo en el cual el acreedor acepta voluntariamente, en pago, una prestación diversa de la debida.
\end{definition}

\begin{example}{Cambiar un servicio por dinero}
Me comprometo a prestar servicios jurídicos pero, por un problema, no puedo hacerlo. En lugar de esa prestación, ofrezco dinero al acreedor y este lo acepta en pago de la deuda original. Eso es una dación en pago.
\end{example}

%------------------------------------------------------------------------------
\section{La novación}
%------------------------------------------------------------------------------

\begin{definition}{Novación}
Extinción de una obligación por la creación de otra nueva, destinada a reemplazarla. Aparece una cierta modificación y la nueva obligación sustituye a la anterior.
\end{definition}

\begin{example}{Reemplazo del objeto de la obligación}
Yo debía entregar una cosa determinada y, mediante novación, pactamos que ahora debo entregar otra cosa distinta. Se extingue la primera obligación y aparece la segunda en su reemplazo.
\end{example}

%------------------------------------------------------------------------------
\section{La transacción}
%------------------------------------------------------------------------------

La transacción es un \textbf{contrato}. Con ella ya se sale de los artículos ubicados en la parte de las obligaciones (en torno al número 900) para pasar a la parte de contratos especiales, dentro del mismo Libro Tercero, en torno al artículo 1641. Su ubicación como contrato es una novedad del Código actual.

La palabra «transar» tiene en derecho civil y comercial un alcance muy específico: significa poner fin a obligaciones dudosas o litigiosas.

\begin{definition}{Transacción}
Contrato por el cual las partes, haciéndose concesiones recíprocas, extinguen obligaciones dudosas o litigiosas.
\end{definition}

\begin{example}{El accidente de tránsito}
Ante un litigio (por ejemplo, un accidente de tránsito) una de las partes reclama una indemnización de 200.000 pesos y la otra ofrece una suma sensiblemente menor. Quien le pone fin a ese litigio, habitualmente, es un juez, que determina qué es lo justo. Pero si las partes desean ponerle fin anticipadamente y llegan a un acuerdo (una de ellas dice «yo reconozco mi responsabilidad, te voy a pagar cien mil pesos» y la otra acepta), están haciendo una transacción. La transacción implica ceder en parte el propio derecho: quien reclamaba 200.000 termina recibiendo 100.000.
\end{example}

En la práctica profesional, un abogado necesita contar con \textbf{poder especial} de su cliente para hacer una transacción en su nombre, y existen ciertos contratos que no admiten transacción. Estas cuestiones están reguladas en detalle en el Código y se profundizan en Obligaciones y en Contratos.

%------------------------------------------------------------------------------
\section{La renuncia}
%------------------------------------------------------------------------------

\begin{definition}{Renuncia}
Acto jurídico extintivo que consiste en la abdicación de un derecho.
\end{definition}

\begin{articulo}{Artículo 944 del Código Civil y Comercial}
Toda persona puede renunciar a los derechos conferidos por la ley cuando la renuncia no está prohibida y sólo afecta intereses privados. No se admite la renuncia anticipada de las defensas que pueden hacerse valer en juicio.
\end{articulo}

\begin{example}{Renunciar a cobrar una deuda}
«Te renuncio a cobrarte esa deuda que tengo». Es una liberalidad: en el fondo se está regalando algo, y sigue siendo un acto a título gratuito. Por ello puede ser eventualmente impugnado si hubiera acreedores interesados en que esa deuda sí se cobre.
\end{example}

%------------------------------------------------------------------------------
\section{La remisión de deuda}
%------------------------------------------------------------------------------

\begin{definition}{Remisión de deuda}
Acto jurídico que se produce cuando el acreedor entrega voluntariamente al deudor el título en el que consta la deuda (por ejemplo, un pagaré).
\end{definition}

%------------------------------------------------------------------------------
\section{La resolución}
%------------------------------------------------------------------------------

La resolución, la rescisión y la revocación son los tres institutos principales con los que se cierra el estudio de los actos jurídicos extintivos.

\begin{definition}{Resolución}
Extinción de los efectos de la relación jurídica que se produce cuando ocurre un acontecimiento que, conforme a previsiones legales precisas o a pactos de las partes, provoca dicha extinción.
\end{definition}

\begin{important}
En derecho civil, «resolución» tiene un alcance técnico: no significa resolver un problema en el sentido de encontrarle una solución, sino \textbf{ponerle fin}. En general, la resolución pone fin de manera \textbf{retroactiva}: opera hacia atrás, hacia el momento de origen de las actuaciones. Este efecto se denomina, en latín, \textit{ex tunc}.
\end{important}

\subsection{La cláusula resolutoria (o pacto comisorio)}

Es una cláusula de un contrato por la cual las partes tienen la potestad de resolverlo, es decir, de ponerle fin, en caso de que la otra parte no cumpla. Puede ser \textbf{expresa} (pactada explícitamente) o \textbf{implícita} (presente en todos los contratos bilaterales, aunque nada se haya pactado). También se la conoce como \textbf{pacto comisorio}, expreso o implícito.

\begin{articulo}{Artículo 1087 del Código Civil y Comercial}
Todo contrato bilateral tiene implícita la facultad de resolver las obligaciones emergentes de él en caso de incumplimiento de una de las partes.
\end{articulo}

Los artículos 1086, 1088 y 1089 del Código precisan los requisitos y el modo en que debe ejercerse esta facultad.
% TODO: contenido incompleto en las fuentes originales: no se profundiza en los requisitos ni en el modo de ejercicio de esta facultad.


\begin{example}{El pintor que no cumplió el día pactado}
Se pacta que un pintor venga a pintar una casa el primero de enero. El pintor se toma diez días de vacaciones y recién el diez de enero pretende que se lo acepte igual. El dueño de la casa le dice «usted incumplió porque no vino el día que habíamos pactado» y resuelve el contrato, debiendo notificarlo conforme a los requisitos correspondientes.
\end{example}

\begin{example}{El catering que llega tarde}
Se contrata un servicio de catering para un casamiento el día viernes y el proveedor pretende entregar las cosas recién el domingo. El cliente responde: «usted incumplió, yo no le voy a pagar». El proveedor incumplió su parte y, por lo tanto, el cliente tampoco cumple la suya: se resuelve el contrato.
\end{example}

\subsection{La señal o arras}

La señal (también llamada arras) es una suma de dinero que se entrega al inicio de un contrato y que puede tener dos alcances distintos:

\begin{itemize}
    \item \textbf{Facultad de arrepentirse:} debe pactarse expresamente. Quien entregó la señal puede luego arrepentirse y poner fin al contrato (cláusula de arrepentimiento).
    \item \textbf{Función confirmatoria del acto:} es el alcance que se le atribuye si no se pactó nada distinto. Quien entregó la señal debe cumplir todo el contrato igualmente.
\end{itemize}

\begin{important}
Es muy importante ver cómo se pactó la entrega de la seña, porque de ello depende si el contrato queda firme o si existe la posibilidad de arrepentirse.
\end{important}

\begin{articulo}{Artículo 1059 del Código Civil y Comercial}
Regula la señal o arras: su función confirmatoria y, si así fue pactado expresamente, la facultad de arrepentimiento.
\end{articulo}

\subsection{La condición resolutoria}

Es una cláusula que subordina la resolución (es decir, la extinción) de una relación jurídica a un \textbf{hecho futuro e incierto}.

\subsection{La frustración de la finalidad}

El Código Civil y Comercial incorpora, como causal de resolución, la frustración definitiva de la finalidad del contrato: cuando el contrato se frustra en su finalidad, la parte perjudicada puede declarar su resolución, siempre que medie una causa extraordinaria en las circunstancias.

\begin{articulo}{Artículo 1090 del Código Civil y Comercial}
Regula la frustración definitiva de la finalidad del contrato como causal de resolución, cuando obedece a una alteración extraordinaria de las circunstancias existentes al momento de su celebración.
\end{articulo}

\begin{example}{La pandemia y el casamiento frustrado}
Alguien debía entregar algo el diez de abril porque iba a celebrarse un casamiento, y este finalmente no se realizó porque se prohibieron todas las reuniones de personas. Se frustró la finalidad del contrato y la parte perjudicada puede poner fin al acto invocando esa frustración.
\end{example}

\subsection{La teoría de la imprevisión}

La teoría de la imprevisión tiene como antecedente el artículo 1198 del Código de Vélez, que no contenía originalmente este instituto: fue incorporado por la reforma de la ley 17.711. Se aplica cuando la prestación a cargo de una de las partes se torna \textbf{excesivamente onerosa} por una alteración extraordinaria de las circunstancias existentes al momento de la celebración del contrato, producida por una causa sobreviniente.

Esta figura se compara con la lesión, que también supone una desproporción entre las prestaciones, aunque se trata de institutos distintos: en la imprevisión, la alteración de las bases del contrato produce una desproporción que torna excesivamente onerosa la posición de una de las partes.

\begin{articulo}{Artículo 1091 del Código Civil y Comercial}
Autoriza, ante la excesiva onerosidad sobreviniente de la prestación, la resolución total o parcial del contrato, o su adecuación.
\end{articulo}

\begin{note}
La figura tuvo su apogeo periódicamente en Argentina con los cambios en el dólar y la hiperinflación. Sin embargo, se advierte que la inflación en Argentina «ya no es algo previsible», por lo que el instituto va teniendo, con el tiempo, situaciones de aplicación cada vez más acotadas, porque siempre debe tratarse de algo realmente imprevisible.
\end{note}

\subsection{Efectos generales de la resolución}

La resolución opera, generalmente, en forma \textbf{retroactiva} (\textit{ex tunc}).

\begin{articulo}{Artículo 1079 del Código Civil y Comercial}
Salvo disposición legal en contrario, la resolución produce efectos retroactivos, con excepción de las prestaciones cumplidas que quedan firmes en cuanto resulten equivalentes.
\end{articulo}

Si ya se han cumplido determinadas prestaciones y quedaron firmes, la resolución no las alcanza. En esto se diferencia de las dos figuras siguientes (rescisión y revocación), que operan hacia el futuro.

%------------------------------------------------------------------------------
\section{La rescisión}
%------------------------------------------------------------------------------

\begin{definition}{Rescisión}
Forma de extinción que opera a futuro: no tiene efectos retroactivos. Puede ser bilateral o unilateral, en los casos autorizados por la ley.
\end{definition}

% TODO: contenido incompleto o ambiguo en las notas originales (el apunte cita nuevamente el artículo 1079 para los efectos de la rescisión, tras haberlo citado para los efectos de la resolución)
\begin{articulo}{Artículo 1079 del Código Civil y Comercial (efectos de la extinción por declaración de una de las partes)}
La rescisión opera a futuro. Excepto estipulación en contrario, sólo produce efectos para el futuro y no afecta el derecho de terceros.
\end{articulo}

\begin{example}{Rescisión bilateral: un contrato de servicios jurídicos}
Se trata, generalmente, de contratos de tracto sucesivo: por ejemplo, un contrato de servicios jurídicos que se viene pagando todos los meses. Un día, ambas partes deciden rescindirlo de común acuerdo. Eso es una rescisión bilateral.
\end{example}

\begin{example}{Rescisión unilateral: el contrato de alquiler}
El contrato puede extinguirse total o parcialmente por la sola declaración de una de las partes, en los casos en que el propio contrato o la ley le otorgan esa facultad. El ejemplo típico es el contrato de alquiler, que puede ser rescindido por el inquilino con cierto preaviso, a partir de los seis meses de celebrado el contrato: la ley concede unilateralmente al locatario la potestad de rescindir.
\end{example}

%------------------------------------------------------------------------------
\section{La revocación}
%------------------------------------------------------------------------------

\begin{definition}{Revocación}
Acto jurídico unilateral mediante el cual una de las partes deja sin efecto la relación por su exclusiva voluntad. Al igual que la rescisión, es unilateral y produce efectos a futuro: no es retroactiva.
\end{definition}

\begin{example}{El mandato revocado}
Se contrata a un abogado para que actúe en nombre de una persona (mandato). En un momento se produce un quiebre de la confianza, o simplemente existe la voluntad de no continuar con esa representación, y se revoca el mandato. Todo lo que el abogado hizo en nombre del mandante hasta ese momento es válido, porque tenía mandato vigente; pero desde el día en que se revoca, ya no tiene mandato para actuar en su nombre.
\end{example}

\begin{example}{El testamento revocado}
Una persona otorga un testamento válido. Mientras vive, puede revocarlo y otorgar otro nuevo, quedando revocado el anterior.
\end{example}

\subsection{Cuadro comparativo: resolución, rescisión y revocación}

\begin{table}[H]
\centering
\begin{tabularx}{\textwidth}{
    >{\raggedright\arraybackslash}p{3.0cm}
    >{\raggedright\arraybackslash}X
    >{\raggedright\arraybackslash}X
}
\toprule
\textbf{Criterio} & \textbf{Resolución} & \textbf{Rescisión / Revocación} \\
\midrule
Efectos en el tiempo & Retroactivos (\textit{ex tunc}) & Hacia el futuro (\textit{ex nunc}) \\
\addlinespace
Origen típico & Incumplimiento, condición, frustración de la finalidad, imprevisión & Acuerdo de partes o facultad legal / contractual unilateral \\
\addlinespace
Partes involucradas & Puede requerir declaración de la parte cumplidora & Rescisión: bilateral o unilateral. Revocación: unilateral \\
\addlinespace
Prestaciones ya cumplidas & En principio quedan sin efecto, salvo equivalencia & Quedan firmes; no se retrotraen \\
\bottomrule
\end{tabularx}
\end{table}

%------------------------------------------------------------------------------
\section{Cuadro Cornell: actos jurídicos extintivos}
%------------------------------------------------------------------------------

{\renewcommand{\arraystretch}{1.2}
\begin{longtable}{
    >{\raggedright\arraybackslash}p{\cornellL}
    >{\raggedright\arraybackslash}p{\cornellR}
}
\toprule
\textbf{Preguntas / palabras clave} &
\textbf{Notas / respuestas} \\
\midrule
\endfirsthead

\toprule
\textbf{Preguntas / palabras clave} &
\textbf{Notas / respuestas} \\
\midrule
\endhead

\textbf{¿Cuál es el acto extintivo por excelencia?} &
El pago: el cumplimiento de la prestación que constituye el objeto de la obligación. No se limita a la entrega de dinero: pagar es cumplir con la obligación. \\
\addlinespace

\textbf{¿Qué diferencia a la dación en pago de la novación?} &
En la dación en pago el acreedor acepta voluntariamente una prestación distinta de la debida. En la novación nace una obligación nueva que reemplaza y extingue a la anterior. \\
\addlinespace

\textbf{¿Qué es la transacción?} &
Un contrato (ubicado en torno al art.~1641 CCyC) por el cual las partes, mediante concesiones recíprocas, extinguen obligaciones dudosas o litigiosas, evitando o poniendo fin a un litigio. El abogado necesita poder especial de su cliente para transar. \\
\addlinespace

\textbf{¿Qué diferencia a la renuncia de la remisión de deuda?} &
La renuncia es la abdicación de un derecho conferido por la ley, cuando no está prohibida y sólo afecta intereses privados (art.~944 CCyC); puede ser una liberalidad a título gratuito. La remisión de deuda es la entrega voluntaria, por el acreedor, del título en el que consta la deuda. \\
\addlinespace

\textbf{¿Qué efecto tiene la resolución sobre la relación jurídica?} &
Un efecto retroactivo (\textit{ex tunc}), salvo las prestaciones ya cumplidas que resulten equivalentes y queden firmes (art.~1079 CCyC). \\
\addlinespace

\textbf{¿Qué formas o causas puede adoptar la resolución?} &
La cláusula resolutoria o pacto comisorio (expreso o implícito, art.~1087 CCyC), la señal o arras (art.~1059 CCyC), la condición resolutoria, la frustración de la finalidad (art.~1090 CCyC) y la teoría de la imprevisión (art.~1091 CCyC). \\
\addlinespace

\textbf{¿Cómo influye la forma en que se pactó la señal?} &
Si se pactó expresamente como facultad de arrepentirse, quien la entregó puede poner fin al contrato. Si no se pactó nada distinto, es confirmatoria del acto y quien la entregó debe cumplir todo el contrato. \\
\addlinespace

\textbf{¿Cuándo se aplica la teoría de la imprevisión?} &
Cuando la prestación de una de las partes se torna excesivamente onerosa por una alteración extraordinaria de las circunstancias existentes al celebrarse el contrato, producida por una causa sobreviniente e imprevisible. \\
\addlinespace

\textbf{¿Qué diferencia a la rescisión y a la revocación de la resolución?} &
Ambas operan hacia el futuro (\textit{ex nunc}), sin efecto retroactivo. La rescisión puede ser bilateral o unilateral; la revocación es unilateral (ejemplos: mandato, testamento). \\

\midrule
\multicolumn{2}{p{\dimexpr\cornellL+\cornellR+2\tabcolsep\relax}}{%
\textbf{Resumen:} los actos extintivos requieren una declaración de voluntad. El pago cumple la obligación; la dación en pago la satisface con una prestación distinta; la novación la reemplaza por otra nueva; la transacción pone fin a obligaciones dudosas o litigiosas mediante concesiones recíprocas; la renuncia abdica un derecho y la remisión entrega el título de la deuda. La resolución opera hacia atrás (\textit{ex tunc}); la rescisión y la revocación, hacia el futuro.} \\
\bottomrule
\end{longtable}}

%==============================================================================
\chapter[Prescripción y caducidad]{Prescripción y caducidad}\label{cap:prescripcion}
%==============================================================================

%------------------------------------------------------------------------------
\section{Importancia de la prescripción en la práctica}
%------------------------------------------------------------------------------

La prescripción es un instituto de enorme importancia que, por tratarse generalmente al final de la materia, no siempre se estudia con detenimiento; sin embargo, tiene una relevancia decisiva en la vida profesional del abogado.

\begin{example}{El cliente que llega con una deuda vieja}
Un cliente llega con una deuda y pregunta: «¿yo tengo derecho a cobrar esto? Hace tiempo que pasó, alguna vez le mandé un mensajito para reclamarla». El abogado pregunta cuánto tiempo pasó y a veces la respuesta es «como cinco años». Ahí hay que advertir que pueden quedar apenas tres días antes de que la deuda prescriba, y actuar con urgencia. Por eso, lo primero que hay que revisar al recibir a un cliente son las \textbf{fechas}.
\end{example}

El fundamento de la prescripción es que existe la expectativa de que los derechos se ejerzan en un tiempo oportuno: no se puede vivir con la indeterminación de que en cualquier momento alguien venga a reclamar algo pasado un tiempo indefinido. La inacción hace presumir, por ley, que decayó el interés en ejercer el derecho. Si una persona tiene una deuda a su favor y pasan veinte años sin reclamarla, ya no parece tener tanto interés en cobrarla, y esto hace que la deuda (más precisamente, la acción para reclamarla) prescriba.

%------------------------------------------------------------------------------
\section{Prescripción adquisitiva y prescripción extintiva}
%------------------------------------------------------------------------------

\begin{definition}{Prescripción}
Instituto jurídico en virtud del cual el transcurso del tiempo opera la adquisición de un derecho, o la extinción de la acción por la inacción del titular del derecho.
\end{definition}

La prescripción está tratada en el Libro Sexto, Título Primero del Código Civil y Comercial, que reúne las disposiciones comunes a los derechos reales y personales. Tiene dos grandes formas, la \textbf{adquisitiva} y la \textbf{extintiva o liberatoria}, que comparten algunos elementos comunes (sobre todo, la idea del transcurso del tiempo unida a una cierta inacción) pero presentan muchas diferencias entre sí y persiguen finalidades distintas. El énfasis de esta parte está en la prescripción extintiva o liberatoria.

\subsection{La prescripción adquisitiva (usucapión)}

\begin{definition}{Prescripción adquisitiva (usucapión)}
Modo por el cual, mediante la posesión de una cosa reunidas ciertas condiciones, se puede adquirir el dominio, ante el manifiesto desinterés del propietario en la atención, cuidado y vigilancia del bien.
\end{definition}

Está regulada en el Libro Cuarto, sobre Derechos Reales, y se profundiza en esa materia.

\begin{definition}{Posesión}
Relación real: la que se tiene directamente con la cosa en los hechos, es decir, el \textit{corpus} (la detentación material de algo), unida al \textit{animus domini} (el ánimo de comportarse como dueño de la cosa, sin reconocer en ningún otro un derecho superior).
\end{definition}

\begin{articulo}{Artículos 1898 y 1899 del Código Civil y Comercial}
Para la prescripción adquisitiva de inmuebles, el plazo común es de veinte años de posesión continua e ininterrumpida, sea de buena o de mala fe. Existe un plazo más corto de diez años cuando el poseedor cuenta con justo título y buena fe.
\end{articulo}

Para las cosas muebles rigen plazos distintos: si la cosa fue hurtada o perdida, el plazo de prescripción adquisitiva es de dos años; si se trata de una cosa mueble registrable, el plazo es de diez años.

\begin{example}{El automóvil poseído por más de diez años}
Una persona recibió un automóvil del titular registral o de un funcionario autorizado, de buena fe, y lo posee durante más de diez años sin que conste a su nombre en el registro de la propiedad. Esa persona pasa a tener la propiedad del bien; para ello debe promover un juicio de prescripción adquisitiva (usucapión), que es un modo de adquirir el dominio.
\end{example}

\subsection{La prescripción extintiva o liberatoria}

\begin{definition}{Prescripción extintiva o liberatoria}
Extingue la acción por la cual una persona puede reclamar un derecho, por el transcurso del tiempo fijado por la ley y la inacción del titular. Consiste en la falta de ejercicio, en el tiempo, de las acciones tendientes a reclamar los derechos, lo que determina la imposibilidad práctica de ejercerlos en el futuro por el transcurso de los plazos legales.
\end{definition}

\begin{important}
La prescripción extintiva extingue la \textbf{acción}, y no el derecho, porque puede ocurrir que finalmente se pague una deuda ya prescripta.
\end{important}

\begin{example}{El amigo que paga una deuda de veinte años}
Una persona le prestó dinero a un amigo hace mucho tiempo y nunca se lo reclamó; pasaron veinte años. Un día se encuentran y el amigo le cuenta que está pasando una mala situación económica. El otro le dice: «hace veinte años me prestaste esta plata, tomá, te pago lo que te debía», y le paga. Tiempo después, ese amigo fallece y sus herederos preguntan si, por haber pagado una deuda prescripta, corresponde la devolución. La respuesta es que no: no se pagó algo que no se debía. La deuda existía; lo que se había extinguido era solamente la acción para reclamarla judicialmente.
\end{example}

\begin{articulo}{Artículo 2538 del Código Civil y Comercial}
El pago espontáneo de una obligación prescripta no es repetible.
\end{articulo}

«Repetible» significa que puede pedirse el reembolso de lo pagado; por lo tanto, lo pagado espontáneamente no puede reclamarse. En el Código Civil anterior, esta situación se llamaba «obligación natural». El nuevo Código evita expresamente ese nombre y utiliza la expresión «pago espontáneo», pero en los hechos sigue siendo la misma idea: una obligación que tiene fundamento en la idea natural de justicia (dar a cada uno lo suyo). La deuda existe; lo que ocurre es que prescribió y ya no es exigible judicialmente. La acción que acompaña al derecho se extingue y el derecho queda desprovisto de acción, pero sigue siendo un derecho que existía.

Consecuencia práctica: cuando la contraparte intenta hacer valer un derecho ya prescripto, el abogado debe oponer la \textbf{excepción de prescripción}, defensa procesal por la cual se señala al juez que esa acción está prescripta y por qué motivos.

\subsection{Reglas generales sobre la prescripción}

Las normas de prescripción son \textbf{imperativas}: no pueden ser modificadas por las partes. No se puede pactar, por ejemplo, que una deuda sea imprescriptible.

\begin{articulo}{Artículo 2533 del Código Civil y Comercial}
Las normas relativas a la prescripción no pueden ser modificadas por convención de las partes.
\end{articulo}

La prescripción opera a favor y en contra de todas las personas, salvo disposición legal en contrario. En las obligaciones con sujetos múltiples (varios acreedores o varios deudores) cualquier acreedor puede oponer la prescripción, aunque el obligado no la invoque, y esta puede también renunciarse. % TODO: contenido incompleto o ambiguo en las notas originales (no se precisa quién puede oponer la prescripción en las obligaciones con sujetos múltiples)

La renuncia a la prescripción sólo puede hacerse una vez que la prescripción ya se ganó: la persona a cuyo favor corrió el plazo puede renunciar a invocarla, de modo similar a la renuncia como acto extintivo. La prescripción puede ser invocada en todos los casos, salvo en los supuestos en que la ley indique expresamente lo contrario.

%------------------------------------------------------------------------------
\section{Suspensión e interrupción del cómputo de la prescripción}
%------------------------------------------------------------------------------

Es muy importante distinguir dos categorías de hechos que pueden incidir sobre el cómputo del plazo de prescripción: los \textbf{hechos suspensivos} y los \textbf{hechos interruptivos}.

\subsection{La suspensión}

\begin{definition}{Suspensión del cómputo}
Detiene el conteo del tiempo mientras dura la causa suspensiva, pero aprovecha (conserva) el período ya transcurrido antes de que la suspensión comenzara.
\end{definition}

\begin{articulo}{Artículo 2539 del Código Civil y Comercial}
La suspensión de la prescripción detiene el cómputo del tiempo por el lapso que dura, pero aprovecha el período transcurrido hasta que ella comenzó.
\end{articulo}

\begin{example}{Cómo se cuenta un plazo suspendido}
El plazo de prescripción empieza a correr el 1 de enero de 2017 y el 30 de junio de 2017 ocurre un hecho suspensivo que dura un año: el cómputo se detiene ese día. El tramo ya corrido (del 1 de enero al 30 de junio de 2017) es válido y se conserva. Cuando, un año después, el 30 de junio de 2018, se reanuda el cómputo, se sigue contando desde el sexto mes en adelante, sin perder el tiempo ya corrido.
\end{example}

\textbf{Supuestos de suspensión}

\begin{itemize}
    \item \textbf{Interpelación fehaciente.} Es el acto típico del abogado que recibe un caso y reclama la deuda al deudor (o al poseedor) de manera extrajudicial, mediante carta documento o acta notarial, con fuerza probatoria hacia terceros. Suspende el cómputo \textbf{por una sola vez y por seis meses}.
\end{itemize}

\begin{articulo}{Artículo 2541 del Código Civil y Comercial}
La interpelación fehaciente hecha por el titular del derecho contra el deudor o el poseedor suspende la prescripción por seis meses, por una sola vez.
\end{articulo}

\begin{itemize}
    \item \textbf{Pedido de mediación.} Suspende el cómputo durante todo el tiempo que dura el trámite de mediación y hasta veinte días después de que se labra el acta de cierre, momento en el cual se reanuda el cómputo del plazo.
\end{itemize}

\begin{articulo}{Artículo 2542 del Código Civil y Comercial}
El curso de la prescripción se suspende desde la expedición por medio fehaciente de la solicitud de mediación, hasta veinte días contados desde el momento en que el acta que da por concluido el procedimiento se encuentre a disposición de las partes.
\end{articulo}

\begin{itemize}
    \item \textbf{Otros supuestos especiales.} Se trata de vínculos de confianza en los que los conflictos de interés dificultan ejercer acciones entre sí. Están suspendidas las acciones:
    \begin{itemize}
        \item entre cónyuges, mientras dura el matrimonio, para no sumar un motivo más a las disputas conyugales;
        \item entre convivientes en unión convivencial;
        \item entre las personas con capacidad restringida y sus representantes, apoyos o tutores, mientras dura la responsabilidad parental, la tutela o la curatela;
        \item entre las personas jurídicas y sus administradores;
        \item a favor y en contra de los herederos, mientras dure la respectiva situación.
    \end{itemize}
\end{itemize}

\begin{articulo}{Artículo 2543 del Código Civil y Comercial}
Enumera los supuestos especiales de suspensión de la prescripción entre cónyuges, convivientes, personas con capacidad restringida y sus representantes o apoyos, personas jurídicas y sus administradores, y a favor y en contra de los herederos.
\end{articulo}

\subsection{La interrupción}

\begin{definition}{Interrupción del cómputo}
A diferencia de la suspensión, tiene por no sucedido el lapso que la precede e inicia un nuevo plazo. Es decir, invalida todo el tiempo ya transcurrido.
\end{definition}

\begin{example}{Cómo se cuenta un plazo interrumpido}
La prescripción empezó a correr el 1 de enero de 2017 y el 30 de junio de 2017 ocurre un hecho interruptivo que dura un año. Cuando finaliza ese hecho, el cómputo empieza de cero: no se aprovecha nada del tiempo ya transcurrido. Esto es lo que distingue, en la práctica, a la interrupción de la suspensión.
\end{example}

\textbf{Supuestos de interrupción}

\begin{itemize}
    \item \textbf{Reconocimiento de deuda.} Cuando la persona obligada firma un documento en el que reconoce la deuda.
\end{itemize}

\begin{articulo}{Artículo 2545 del Código Civil y Comercial}
El curso de la prescripción se interrumpe por el reconocimiento que el deudor o poseedor efectúa del derecho de aquel contra quien prescribe.
\end{articulo}

\begin{itemize}
    \item \textbf{Petición judicial.} Se interrumpe la prescripción por toda petición del titular del derecho ante la autoridad judicial que traduzca la intención de no abandonar el derecho.
\end{itemize}

\begin{articulo}{Artículo 2546 del Código Civil y Comercial}
La petición judicial interrumpe la prescripción, aunque sea defectuosa, realizada por quien no es capaz o ante un tribunal incompetente, y aunque se presente dentro del plazo de gracia previsto en el ordenamiento procesal.
\end{articulo}

La redacción del nuevo Código es muy amplia para describir este supuesto: \textbf{no se exige necesariamente una demanda} en sentido estricto, sino cualquier petición que traduzca la intención de no abandonar el derecho; puede bastar, por ejemplo, una verificación de crédito en un proceso sucesorio. Interrumpir resulta más beneficioso para el acreedor que suspender, porque invalida todo el tiempo transcurrido y hace que se cuente todo de nuevo.

Sobre el plazo de gracia: los plazos procesales vencen a las 24 horas, pero como los tribunales dejan de funcionar antes, se habilitan las dos primeras horas hábiles del día siguiente para presentar el escrito dentro de ese plazo, lo que se conoce habitualmente como «las dos primeras».

\begin{itemize}
    \item \textbf{Solicitud de arbitraje.} Equivale, en los hechos, a interponer una demanda o a un reconocimiento de deuda, y también interrumpe el cómputo de la prescripción.
\end{itemize}

\begin{articulo}{Artículo 2548 del Código Civil y Comercial}
La solicitud de arbitraje interrumpe la prescripción.
\end{articulo}

\subsection{Cuadro comparativo: suspensión e interrupción}

\begin{table}[H]
\centering
\begin{tabularx}{\textwidth}{
    >{\raggedright\arraybackslash}p{3.3cm}
    >{\raggedright\arraybackslash}X
    >{\raggedright\arraybackslash}X
}
\toprule
\textbf{Criterio} & \textbf{Suspensión} & \textbf{Interrupción} \\
\midrule
Tiempo ya transcurrido & Se conserva (se aprovecha) & Se invalida por completo \\
\addlinespace
Reinicio del cómputo & Continúa sumando desde donde se detuvo & Empieza de cero \\
\addlinespace
Supuestos & Interpelación fehaciente, mediación, causas entre cónyuges y convivientes, entre otras & Reconocimiento de deuda, petición judicial, solicitud de arbitraje \\
\bottomrule
\end{tabularx}
\end{table}

%------------------------------------------------------------------------------
\section{La dispensa y los plazos de prescripción}
%------------------------------------------------------------------------------

\subsection{La dispensa}

\begin{definition}{Dispensa}
Tercer supuesto, distinto de la suspensión y de la interrupción. A diferencia de ambas, ocurre una vez que la prescripción ya se cumplió, cuando por alguna circunstancia extraordinaria era imposible ejercer la acción en el momento en que finalizaba el plazo.
\end{definition}

En el Código anterior este supuesto se trataba como un caso de suspensión. El nuevo Código lo clarifica y lo llama \textbf{dispensa}: la doctrina ya interpretaba que se trataba de un supuesto de dispensa, y ahora se lo ubica claramente como tal.

\begin{articulo}{Artículo 2550 del Código Civil y Comercial}
El juez puede dispensar de la prescripción ya cumplida al titular de la acción, si dificultades de hecho o maniobras dolosas le obstaculizaron temporalmente el ejercicio de la acción, y si el titular hace valer sus derechos dentro de los seis meses siguientes a la cesación de los obstáculos.
\end{articulo}

\begin{example}{La persona postrada en un hospital}
Se debe recurrir ante un juez y explicarle, por ejemplo: «yo estaba en cama, totalmente postrado, al momento en que prescribía esta acción, y me era imposible ejercerla». Si el juez entiende que existieron esas circunstancias extraordinarias, concede una dispensa que permite iniciar la acción dentro de los seis meses siguientes. Apenas la persona sale del hospital, inicia la demanda y pide la dispensa de la prescripción.
\end{example}

La dispensa se aplica también a personas incapaces sin representante: por ejemplo, un niño de cinco años, sin representantes, que tenía que cobrar ciertas deudas; en ese caso se espera a que existan representantes que puedan actuar en su nombre.

\begin{important}
La diferencia central con la suspensión y la interrupción es que la dispensa se \textbf{tramita ante un juez} y recae sobre una prescripción ya cumplida.
\end{important}

\subsection{Los plazos de prescripción}

Se han visto distintos plazos de prescripción a lo largo de las clases. Por ejemplo, la acción por vicios tiene un plazo de prescripción de dos años. El plazo genérico es de cinco años.

\begin{articulo}{Artículo 2560 del Código Civil y Comercial}
El plazo de la prescripción es de cinco años, excepto que esté previsto uno diferente en la legislación local.
\end{articulo}

A partir de los artículos siguientes, el Código regula distintos plazos especiales, desde más largos hasta más cortos, según el tipo de acción de que se trate.

En general, el cómputo del plazo corre desde el día en que la prestación es \textbf{exigible}, salvo distintas excepciones:
\begin{itemize}
    \item en el caso de la violencia, el plazo comienza a correr desde que cesa la violencia;
    \item en la simulación entre las partes, el plazo corre desde que una de ellas desconoce el carácter simulado del acto.
\end{itemize}

\begin{important}
Es un tema importantísimo, porque es la fecha a partir de la cual se cuenta todo.
\end{important}

%------------------------------------------------------------------------------
\section{La caducidad en profundidad}
%------------------------------------------------------------------------------

La caducidad es también una forma de operar del transcurso del tiempo, pero produce la extinción de un \textbf{derecho no ejercido}, no de la acción, como ocurre en la prescripción.

Los plazos de caducidad están fijados específicamente por la ley: no todos los derechos tienen un plazo de caducidad, sino solamente aquellos respecto de los cuales la ley expresamente lo determina. Un ejemplo es el del artículo 442 sobre la compensación económica en materia de divorcio, cuya acción caduca a los seis meses de dictada la sentencia.

\begin{articulo}{Artículo 2567 del Código Civil y Comercial (régimen general de caducidad)}
Los plazos de caducidad no se suspenden ni se interrumpen, excepto disposición legal en contrario.
\end{articulo}

La caducidad no admite suspensión ni interrupción y, en general, sus plazos son más cortos que los de la prescripción.

Hay actos que impiden que se produzca la caducidad: el cumplimiento del acto previsto por la ley, la realización del acto jurídico correspondiente, o el reconocimiento del derecho por parte de quien podría beneficiarse de la caducidad.

Al igual que la prescripción, las normas sobre caducidad establecidas por la ley son, en general, \textbf{imperativas}: las partes no pueden renunciarla ni alterarla por disposición contractual cuando la caducidad proviene de la ley (sí es disponible cuando surge de la voluntad de las partes). % TODO: contenido incompleto o ambiguo en las notas originales (la caducidad se presenta como de fuente siempre legal en la introducción, pero aquí se menciona una caducidad que surge de la voluntad de las partes)

Cuando concurren, sobre un mismo supuesto, normas de caducidad y de prescripción, \textbf{priman las de caducidad}, y para lo no regulado se aplican, supletoriamente, las normas relativas a la prescripción.

%------------------------------------------------------------------------------
\section{Cuadro Cornell: prescripción y caducidad}
%------------------------------------------------------------------------------

{\renewcommand{\arraystretch}{1.2}
\begin{longtable}{
    >{\raggedright\arraybackslash}p{\cornellL}
    >{\raggedright\arraybackslash}p{\cornellR}
}
\toprule
\textbf{Preguntas / palabras clave} &
\textbf{Notas / respuestas} \\
\midrule
\endfirsthead

\toprule
\textbf{Preguntas / palabras clave} &
\textbf{Notas / respuestas} \\
\midrule
\endhead

\textbf{¿Qué es la prescripción?} &
El instituto por el cual el transcurso del tiempo opera la adquisición de un derecho (prescripción adquisitiva) o la extinción de la acción por inacción del titular (prescripción extintiva o liberatoria). Se trata en el Libro Sexto, Título Primero del CCyC. \\
\addlinespace

\textbf{¿Qué diferencia hay entre prescripción adquisitiva y extintiva?} &
La adquisitiva otorga el dominio por la posesión de una cosa (para inmuebles: veinte años, o diez con justo título y buena fe). La extintiva extingue la acción para reclamar un derecho, no el derecho en sí mismo. \\
\addlinespace

\textbf{¿Puede pagarse válidamente una deuda ya prescripta?} &
Sí. El pago espontáneo de una obligación prescripta no es repetible (art.~2538 CCyC): el derecho existía y sólo se había extinguido la acción para exigirlo judicialmente. Si se intenta reclamar un derecho prescripto, debe oponerse la excepción de prescripción. \\
\addlinespace

\textbf{¿Pueden las partes modificar las normas de prescripción?} &
No. Son imperativas y no pueden modificarse por convención de las partes (art.~2533 CCyC). La renuncia a la prescripción sólo puede hacerse una vez que esta ya se ganó. \\
\addlinespace

\textbf{¿Qué diferencia a la suspensión de la interrupción del cómputo?} &
La suspensión detiene el cómputo pero conserva el tiempo ya transcurrido (art.~2539 CCyC); la interrupción invalida todo lo transcurrido y reinicia el plazo desde cero. \\
\addlinespace

\textbf{¿Qué hechos suspenden la prescripción?} &
La interpelación fehaciente (seis meses, una sola vez, art.~2541 CCyC), el pedido de mediación (hasta veinte días después del acta de cierre, art.~2542 CCyC) y ciertas situaciones entre cónyuges, convivientes, personas con capacidad restringida y sus representantes, personas jurídicas y sus administradores, y herederos (art.~2543 CCyC). \\
\addlinespace

\textbf{¿Qué hechos interrumpen la prescripción?} &
El reconocimiento de deuda (art.~2545 CCyC), la petición judicial, interpretada en sentido amplio y aun si es defectuosa (art.~2546 CCyC), y la solicitud de arbitraje (art.~2548 CCyC). \\
\addlinespace

\textbf{¿Qué es la dispensa de la prescripción?} &
Una autorización judicial para ejercer la acción dentro de los seis meses siguientes a la cesación de los obstáculos, pese a que el plazo de prescripción ya se cumplió, cuando existió una imposibilidad de hecho justificada (por ejemplo, una internación) (art.~2550 CCyC). \\
\addlinespace

\textbf{¿Cuál es el plazo genérico de prescripción y desde cuándo se cuenta?} &
Cinco años (art.~2560 CCyC), salvo que esté previsto uno diferente; existen plazos especiales, más largos y más cortos. En general se cuenta desde que la prestación es exigible, con excepciones (violencia, simulación entre las partes). \\
\addlinespace

\textbf{¿Por qué la caducidad es más rígida que la prescripción?} &
Porque no admite suspensión ni interrupción (art.~2567 CCyC), sus plazos suelen ser más breves y extingue el derecho no ejercido, no sólo la acción. Cuando concurren normas de caducidad y de prescripción, priman las de caducidad. \\

\midrule
\multicolumn{2}{p{\dimexpr\cornellL+\cornellR+2\tabcolsep\relax}}{%
\textbf{Resumen:} el transcurso del tiempo puede hacer adquirir un derecho (prescripción adquisitiva) o extinguir la acción por la inacción del titular (prescripción extintiva), que deja subsistente un derecho sin acción. El cómputo puede detenerse (suspensión, que conserva el tiempo corrido), reiniciarse (interrupción, que lo invalida) o ser habilitado judicialmente una vez cumplido el plazo (dispensa). La caducidad extingue el derecho no ejercido en el plazo legal, de manera más rígida e imperativa.} \\
\bottomrule
\end{longtable}}

%==============================================================================

\section{Síntesis de la extinción de las relaciones jurídicas}
%==============================================================================

El tema recorre el cierre del ciclo de vida de las relaciones jurídicas: cómo terminan, después de haber nacido y desarrollado sus efectos. Ese final puede producirse:

\begin{itemize}
    \item por \textbf{hechos ajenos a la voluntad} de las partes: la muerte, la imposibilidad sobrevenida, la confusión o la compensación;
    \item por \textbf{actos jurídicos} mediante los cuales las partes deciden ponerle fin a la relación: el pago y la dación en pago, que la cumplen; la novación, que la reemplaza; la transacción y la renuncia, que ceden derechos; la remisión, que perdona una deuda; y la resolución, la rescisión y la revocación, que la dejan sin efecto, hacia atrás o hacia adelante según el caso;
    \item por el simple \textbf{transcurso del tiempo}, que castiga la inacción del titular de un derecho a través de la prescripción (que extingue la acción, pero no siempre el derecho) y de la caducidad (que extingue directamente el derecho no ejercido, de manera más rígida e imperativa).
\end{itemize}

El hilo conductor es la idea de que ningún derecho puede sostenerse en una incertidumbre indefinida: el ordenamiento jurídico necesita fijar, tarde o temprano, un punto final, ya sea porque los hechos de la vida lo imponen, porque las partes lo deciden o porque el tiempo, por sí solo, se encarga de hacerlo. Comprender estas reglas, y sobre todo los plazos y las causales de suspensión, interrupción y dispensa de la prescripción, es una de las herramientas de mayor relevancia práctica en el ejercicio cotidiano de la abogacía.

\end{document}
